%%-----------------------------------------------------%%
%% This file may act as the main file of your article. %%
%% Please uncomment one of the following lines.        %%
%%-----------------------------------------------------%%

%% \documentclass[CM,GP]{degruyter-crelle}       %% Equations numbered as (1), (2) etc.
\documentclass[SecEq,CM,GP]{degruyter-crelle} %% Equations numbered as (1.1), (1.2) etc.

%% Do not alter the following line.
%\firstpage{}\volume{}\volumeyear{}\copyrightyear{}\doiyear{}\doi{}\received{XXX}\revised{XXX}

\title[]{Parabolic Simple $\sL$-invariants}

\author{Yiqin}{He}{}{Beijing}
%\author{Firstname}{Lastname}{}{City}

\contact{School of Mathematical Sciences, Peking University,\;No.5 Yiheyuan Road, Haidian District, Beijing 100871, P.R.China}{1801110008@pku.edu.cn}
%\contact{Department, university, street, zip code, city, country}{e-mail}

\researchsupported{This work is funded by the NSFC Grant NO. 8200905010 and NO. 8200800065.\\ \hspace*{25pt}Conflicts of interest statement: All authors disclosed no relevant relationships.\\ \hspace*{25pt}Data Availability statement: Data openly available in a public repository.}

%% Define theorem-like environments as usual: 

\theoremstyle{plain}
 \newtheorem{theorem}{Theorem}[section]
 \newtheorem{thm}{Theorem}[section]
 \newtheorem{cor}[thm]{Corollary}
 \newtheorem{lem}[thm]{Lemma}
 \newtheorem{pro}[thm]{Proposition}
 \newtheorem{dfn}[thm]{Definition}
 \newtheorem{hypothesis}[thm]{Hypothesis}
 \newtheorem{rmk}[thm]{Remark}
 \newtheorem{notation}[thm]{Notation}
 \newtheorem{exmp}[thm]{Example}
 \newtheorem{conjecture}[thm]{Conjecture}
 \newtheorem{condition}[thm]{Condition}

%\theoremstyle{definition} 
%\newtheorem{remark}[theorem]{Remark}
 \usepackage{fullpage, amsmath}
 \usepackage{amsthm,amsfonts,url,amssymb}
 \usepackage{mathrsfs}
 \usepackage{amssymb,amscd}
 \usepackage{color}
 \usepackage{chemarrow}
 \usepackage{pictexwd,dcpic}
 %\usepackage[textwidth=18cm,centering]{geometry}
 \usepackage{extarrows}
% \usepackage[colorlinks,linkcolor=red,anchorcolor=green,citecolor=blue]{hyperref}
 \usepackage{enumitem}
\usepackage{lipsum}
 \usepackage[all,cmtip]{xy}
 \usepackage{leftidx}
 \usepackage[version=3]{mhchem}
 \usepackage{makeidx}
 \usepackage{bm}
 \usepackage{stmaryrd}
% \usepackage[applemac]{inputenc}


\newcommand{\hooklongrightarrow}{\lhook\joinrel\longrightarrow}
\newcommand{\twoheadlongrightarrow}{\relbar\joinrel\twoheadrightarrow}

\newcommand{\Aut}{\rm Aut}
\newcommand{\ad}{{\rm ad}}
\newcommand{\alg}{{\rm alg}}
\newcommand{\an}{{\rm la}}
\newcommand{\ana}{{\rm an}}
\newcommand{\adm}{{\rm ad}}
\newcommand{\Art}{{\rm Art}}
\newcommand{\ab}{{\rm ab}}
\newcommand{\alge}{{\rm alg}}

\newcommand{\bdr}{{B^+_{\mathrm{dR}}}}
\newcommand{\bersteineigenvarpik}{{\mathcal{E}_{\omepik,\fp,\bm{\lambda}_\bh}^\infty(\overline{\rho})}}
\newcommand{\bmdel}{{ \bm{\delta}_\bh}}
%\newcommand{\bI}{{\mathbb{I}}}


\newcommand{\ccyc}{{\chi_{\mathrm{cyc}}}}
%\newcommand{\con}{{\rm con}}
\newcommand{\cris}{{\rm cris}}
\newcommand{\CHom}{{\cal H}om}
\newcommand{\cEo}{{\rm \overline{\cE}}}
\newcommand{\criss}{{\rm cris}}
\newcommand{\cten}{{\rm \widehat{\otimes}_E}}
\newcommand{\crys}{{\rm crys}}
\newcommand{\cpvar}{{\mathbf{{C}}_p}}
\newcommand{\cyc}{{\rm cyc}}
\newcommand{\ct}{{\rm ct}}
\newcommand{\complocalbersteineigenvarpik}{{\widehat{\mathcal{E}_{\omepik,\fp,\bm{\lambda}_\pi}^\infty}(\overline{\rho})_x}}
\newcommand{\complocaldefvarrho}{{ \widehat{X^\Box_{\omepik,\mathbf{{h}}}}(\overline{r})_{x_L}}}
\newcommand{\complocaldefvarrhosub}{{ \widehat{\xfpxxl}}}
\newcommand{\cokerr}{\rm coker}

\newcommand{\dR}{{\rm dR}}
%\DeclareMathOperator{\deg}{deg}
\newcommand{\dett}{{\rm det}}
\newcommand{\df}{{\mathbf{DF}}}
\newcommand{\dr}{{\rm dR}}
\newcommand{\Dpik}{{\mathbf{D}}}
\newcommand{\defvarnon}{{ U^\Box_{\Omega,\underline{\bh}}(\overline{r})}}
\newcommand{\defvarr}{{\mathfrak{X}_{\overline{r}}^\Box}}
\newcommand{\defvarring}{{R_{\overline{r}}^\Box}}
\newcommand{\defvar}{{ X_{\omepik,\mathbf{{h}}}(\overline{r})}}
\newcommand{\defvarrho}{{ X_{\omepik,\mathbf{{h}}}(\overline{r})}}
\newcommand{\defvarrhonon}{{ U^\Box_{\omepik,\mathbf{{h}}_{\Dpik}}(\overline{r})}}


\newcommand{\et}{{\rm\acute{e}t}}
\newcommand{\EndO}{{\rm End}}
\newcommand{\ext}{{\rm Ext}}
\newcommand{\EX}{{E^\times}}


\newcommand{\fil}{{\rm Fil}}
\newcommand{\fge}{{\rm fg}}
\newcommand{\fsl}{{\rm fs}}
\newcommand{\frk}{{\rm rk}}


%\DeclareMathOperator{\gcd}{gcd}
\newcommand{\gal}{{\rm Gal}}
\newcommand{\Gal}{{\rm Gal}}
\newcommand{\GLN}{{\rm GL}}
\newcommand{\GL}{{\rm GL}}
\newcommand{\gen}{{\rm gen}}
\newcommand{\gr}{{\rm gr}}


\newcommand{\homo}{{\rm Hom}}
\newcommand{\hH}{{\mathrm{H}}}
\newcommand{\hpi}{{{\mathbf{h}}}}


\newcommand{\ind}{{\rm Ind}}
\newcommand{\im}{{\rm Im}}
\newcommand{\lM}{{\overline{M}}}
\newcommand{\lL}{{\overline{L}}}
\newcommand{\lamoverx}{{L^{\lrr}({\bm\lambda}_\bh)}}

\newcommand{\jsm}{{\mathcal{J}^{\infty}}}
\newcommand{\jh}{{\rm JH}}


\newcommand{\Ker}{{\rm Ker}}
%\DeclareMathOperator{\ker}{ker}


\newcommand{\leg}{{\rm lg}}
\newcommand{\lrr}{{\langle r\rangle}}
\newcommand{\lalg}{{\rm lalg}}
\newcommand{\lra}{\longrightarrow}
\newcommand{\LX}{{L^\times}}
\newcommand{\loc}{{\rm loc}}
\newcommand{\lp}{{\rm lp}}

\newcommand{\Modo}{{\textbf{Mod}}}
\newcommand{\MSpec}{{\rm MaxSpec}}


\newcommand{\normm}{{\rm Norm}}


\newcommand{\op}{{\overline{{\mathbf{P}}}}}
\newcommand{\opp}{{\overline{{\mathbf{P}}}}}
\newcommand{\ob}{{\overline{{\mathbf{B}}}}}
\newcommand{\onn}{{\overline{{\mathbf{N}}}}}
\newcommand{\opki}{{\overline{\mathcal{P}}_{\Delta_{k,i}}}}
\newcommand{\obp}{\overline{\mathbb{Q}}_p^\times}
\newcommand{\obpo}{\overline{\mathbb{Q}}_p}
\newcommand{\omepik}{\Omega_{[1,k]}}


\newcommand{\perf}{{\rm perf}}
\newcommand{\proj}{{\rm Proj}}
\newcommand{\PGLN}{{\rm PGL}}
\newcommand{\pol}{{\rm pol}}
\newcommand{\pcr}{{\rm pcr}}
\newcommand{\padelc}{{(\Delta_{\Dpik}^z)}}
\newcommand{\padel}{{(\Delta_{\Dpik}^0)}}
\newcommand{\pr}{{\rm pr}}
\newcommand{\pst}{{\rm pst}}



\newcommand{\Rkoh}{{\rm R}}
\newcommand{\Rep}{\mathrm{Rep}}
\newcommand{\red}{{\rm red}}
\newcommand{\recc}{{\rm rec}}
\newcommand{\rig}{{\rm rig}}
\newcommand{\ra}{{\longrightarrow}}
\newcommand{\rec}{{\rm rec}_L}
\newcommand{\recp}{{{\rm rec}_p}}
\newcommand{\res}{{\rm Res}}
\newcommand{\rigch}{{\mathcal{Z}_{\bL^{\lrr},\mathcal{O}_L}}}
\newcommand{\rigchl}{{\mathcal{Z}_{\bL^{\lrr},L}}}
\newcommand{\rigerch}{{\mathcal{Z}_{\bL_{\ul{n}},\mathcal{O}_L}}}
\newcommand{\rigerchl}{{\mathcal{Z}_{\bL_{\ul{n}},L}}}


\newcommand{\Spm}{\rm Sp}
\newcommand{\Spec}{\rm Spec}
\newcommand{\Spf}{{\mathrm{Spf}}}
\newcommand{\ssE}{{\mathcal{S}}}
\newcommand{\sEo}{{\rm \overline{\ssE}}}
\newcommand{\sep}{{\rm sep}}
\newcommand{\sh}{{\rm sh}}
\newcommand{\spa}{{\rm Spa}}
\newcommand{\stt}{{\rm st}}
\newcommand{\st}{{\rm St}}
\newcommand{\soc}{{\rm soc}}
\newcommand{\sm}{{\rm sm}}
\newcommand{\spr}{{\rm Sp}}
\newcommand{\spf}{{\mathrm{Spf}}}
\newcommand{\semi}{{\rm st}}
\newcommand{\sbanpik}{{\big(\mathrm{Spec}\hspace{2pt}\mathfrak{Z}_{\omepik}\big)^{\mathrm{rig}}}}
\newcommand{\sbanpi}{{\left(\Spec\mathfrak{Z}_{\Omega_{r}}\right)^{\mathrm{rig}}}}
\newcommand{\sban}{{\left(\Spec\mathfrak{Z}_{\Omega}\right)^{\mathrm{rig}}}}
\newcommand{\stx}{{x_{\rm st}}}


\newcommand{\temp}{{\rm temp}}
\newcommand{\topol}{{\rm top}}
\newcommand{\tors}{{\rm tors}}
\newcommand{\ten}{{{\otimes}_E}}
\newcommand{\tr}{{\rm tr}}
\newcommand{\tee}{{{\otimes}_{\cR_{E,L}}}}
\newcommand{\te}{{\otimes_E}}
\newcommand{\tene}{{\widehat{\otimes}_{E}}}
\newcommand{\tenp}{{\widehat{\otimes}_{E,\pi}}}
\newcommand{\teni}{{\widehat{\otimes}_{E,\iota}}}

\newcommand{\ur}{{\rm ur}}
\newcommand{\unr}{{\rm unr}}
\newcommand{\unra}{{\rm unr}}
\newcommand{\univer}{{\rm univ}}
\newcommand{\undelram}{{\delta^0}}
\newcommand{\undel}{{\delta}}
\newcommand{\univ}{{\rm univ}}
\newcommand{\ul}{\underline}
\newcommand{\usl}{\mathfrak sl}
\newcommand{\ugl}{\mathfrak gl}

\newcommand{\val}{\mathrm val}
\newcommand{\valua}{\mathrm val}
\newcommand{\vg}{{(\varphi,\Gamma)}}
\newcommand{\ver}{{\rm ver}}

\newcommand{\wdre}{{\textbf{WD}}}
\newcommand{\wt}{{\rm wt}}
\newcommand{\wte}{{\widehat{\otimes}}}
\newcommand{\wdr}{{W^+_{\mathrm{dR}}}}
\newcommand{\WD}{{\rm WD}}

\newcommand{\xfpx}{X^{X^\fp-\mathrm{aut}}_{\omepik}(x)}
\newcommand{\xfpxxl}{X^{X^\fp-\mathrm{aut}}_{\omepik}(x)_{x_L}}

\newcommand{\zuni}{{\widehat{Z^{\lrr}_{\varpi_L}}}}
\newcommand{\zlch}{{\bZ_{\bL_{\ul{n}}}(\cO_L)}}
\newcommand{\zlchl}{{\bZ_{\bL_{\ul{n}}}}}



\newcommand{\BA}{{\mathbb{A}}}
\newcommand{\BB}{{\mathbb{B}}}
\newcommand{\BC}{{\mathbb{C}}}
\newcommand{\BD}{{\mathbb{D}}}
\newcommand{\BE}{{\mathbb{E}}}
\newcommand{\BF}{{\mathbb{F}}}
\newcommand{\BG}{{\mathbb{G}}}
\newcommand{\BH}{{\mathbb{H}}}
\newcommand{\BI}{{\mathbb{I}}}
\newcommand{\BJ}{{\mathbb{J}}}
\newcommand{\BK}{{\mathbb{K}}}
\newcommand{\BL}{{\mathbb{L}}}
\newcommand{\BM}{{\mathbb{M}}}
\newcommand{\BN}{{\mathbb{N}}}
\newcommand{\BO}{{\mathbb{O}}}
\newcommand{\BP}{{\mathbb{P}}}
\newcommand{\BQ}{{\mathbb{Q}}}
\newcommand{\BR}{{\mathbb{R}}}
\newcommand{\BS}{{\mathbb{S}}}
\newcommand{\BT}{{\mathbb{T}}}
\newcommand{\BU}{{\mathbb{U}}}
\newcommand{\BV}{{\mathbb{V}}}
\newcommand{\BW}{{\mathbb{W}}}
\newcommand{\BX}{{\mathbb{X}}}
\newcommand{\BY}{{\mathbb{Y}}}
\newcommand{\BZ}{{\mathbb{Z}}}


\newcommand{\bA}{{\mathbf{A}}}
\newcommand{\bB}{{\mathbf{B}}}
\newcommand{\bC}{{\mathbf{C}}}
\newcommand{\bD}{{\mathbf{D}}}
\newcommand{\bE}{{\mathbf{E}}}
\newcommand{\be}{{\mathbf{e}}}
\newcommand{\bF}{{\mathbf{F}}}
\newcommand{\bG}{{\mathbf{G}}}
\newcommand{\bH}{{\mathbf{H}}}
\newcommand{\bI}{{\mathbf{I}}}
\newcommand{\bJ}{{\mathbf{J}}}
\newcommand{\bK}{{\mathbf{K}}}
\newcommand{\bL}{{\mathbf{L}}}
\newcommand{\bM}{{\mathbf{M}}}
\newcommand{\bN}{{\mathbf{N}}}
\newcommand{\bO}{{\mathbf{O}}}
\newcommand{\bP}{{\mathbf{P}}}
\newcommand{\bQ}{{\mathbf{Q}}}
\newcommand{\bR}{{\mathbf{R}}}
\newcommand{\bS}{{\mathbf{S}}}
\newcommand{\bT}{{\mathbf{T}}}
\newcommand{\bU}{{\mathbf{U}}}
\newcommand{\bV}{{\mathbf{V}}}
\newcommand{\bW}{{\mathbf{W}}}
\newcommand{\bX}{{\mathbf{X}}}
\newcommand{\bY}{{\mathbf{Y}}}
\newcommand{\bZ}{{\mathbf{Z}}}
\newcommand{\bz}{{\mathbf{z}}}
\newcommand{\bc}{{\mathbf{c}}}
\newcommand{\bx}{{\mathbf{x}}}
\newcommand{\bv}{{\mathbf{v}}}
\newcommand{\bh}{{\mathbf{h}}}
\newcommand{\bk}{{\mathbf{k}}}
\newcommand{\bd}{{\mathbf{d}}}
\newcommand{\bj}{{\mathbf{j}}}
\newcommand{\bq}{{\mathbf{q}}}

\newcommand{\cN}{\mathcal N}
\newcommand{\cL}{\mathcal L}
\newcommand{\co}{\mathcal O}
\newcommand{\cA}{\mathcal A}
\newcommand{\cR}{\mathcal R}
\newcommand{\cH}{\mathcal H}
\newcommand{\cC}{\mathcal C}
\newcommand{\cS}{\mathcal S}
\newcommand{\cD}{\mathcal D}
\newcommand{\cI}{\mathcal I}
\newcommand{\cW}{\mathcal W}
\newcommand{\cT}{\mathcal T}
\newcommand{\cM}{\mathcal M}
\newcommand{\cF}{\mathcal F}
\newcommand{\cV}{\mathcal V}
\newcommand{\cE}{\mathcal E}
\newcommand{\cG}{\mathcal G}
\newcommand{\cJ}{\mathcal J}
\newcommand{\cO}{\mathcal O}
\newcommand{\cP}{\mathcal P}
\newcommand{\cQ}{\mathcal Q}
\newcommand{\cU}{\mathcal U}
\newcommand{\cZ}{\mathcal Z}
\newcommand{\cY}{\mathcal Y}



\newcommand{\FA}{{\mathfrak{A}}}
\newcommand{\FB}{{\mathfrak{B}}}
\newcommand{\FC}{{\mathfrak{C}}}
\newcommand{\FD}{{\mathfrak{D}}}
\newcommand{\FE}{{\mathfrak{E}}}
\newcommand{\FF}{{\mathfrak{F}}}
\newcommand{\FG}{{\mathfrak{G}}}
\newcommand{\FH}{{\mathfrak{H}}}
\newcommand{\FI}{{\mathfrak{I}}}
\newcommand{\FJ}{{\mathfrak{J}}}
\newcommand{\FK}{{\mathfrak{K}}}
\newcommand{\FL}{{\mathfrak{L}}}
\newcommand{\FM}{{\mathfrak{M}}}
\newcommand{\FN}{{\mathfrak{N}}}
\newcommand{\FO}{{\mathfrak{O}}}
\newcommand{\FP}{{\mathfrak{P}}}
\newcommand{\FQ}{{\mathfrak{Q}}}
\newcommand{\FR}{{\mathfrak{R}}}
\newcommand{\FS}{{\mathfrak{S}}}
\newcommand{\FT}{{\mathfrak{T}}}
\newcommand{\FU}{{\mathfrak{U}}}
\newcommand{\FV}{{\mathfrak{V}}}
\newcommand{\FW}{{\mathfrak{W}}}
\newcommand{\FX}{{\mathfrak{X}}}
\newcommand{\FY}{{\mathfrak{Y}}}
\newcommand{\FZ}{{\mathfrak{Z}}}


\newcommand{\fa}{{\mathfrak{a}}}
\newcommand{\fb}{{\mathfrak{b}}}
\newcommand{\fc}{{\mathfrak{c}}}
\newcommand{\fd}{{\mathfrak{d}}}
\newcommand{\fe}{{\mathfrak{e}}}
\newcommand{\ff}{{\mathfrak{f}}}
\newcommand{\fg}{{\mathfrak{g}}}
\newcommand{\fh}{{\mathfrak{h}}}
\newcommand{\Fi}{{\mathfrak{i}}}
\newcommand{\fj}{{\mathfrak{j}}}
\newcommand{\fk}{{\mathfrak{k}}}
\newcommand{\fl}{{\mathfrak{l}}}
\newcommand{\fm}{{\mathfrak{m}}}
\newcommand{\fn}{{\mathfrak{n}}}
\newcommand{\fo}{{\mathfrak{o}}}
\newcommand{\fp}{{\mathfrak{p}}}
\newcommand{\fq}{{\mathfrak{q}}}
\newcommand{\fr}{{\mathfrak{r}}}
\newcommand{\fs}{{\mathfrak{s}}}
\newcommand{\ft}{{\mathfrak{t}}}
\newcommand{\fu}{{\mathfrak{u}}}
\newcommand{\fv}{{\mathfrak{v}}}
\newcommand{\fw}{{\mathfrak{w}}}
\newcommand{\fx}{{\mathfrak{x}}}
\newcommand{\fy}{{\mathfrak{y}}}
\newcommand{\fz}{{\mathfrak{z}}}

\newcommand{\sE}{\mathscr E}
\newcommand{\sB}{\mathscr B}
\newcommand{\sH}{\mathscr H}
\newcommand{\sC}{\mathscr C}
\newcommand{\sF}{\mathscr F}
\newcommand{\sD}{\mathscr D}
\newcommand{\sI}{\mathscr I}
\newcommand{\sA}{\mathscr A}
\newcommand{\sW}{\mathscr W}
\newcommand{\sL}{\mathcal L}
\newcommand{\sK}{\mathscr K}
\newcommand{\sN}{\mathscr N}
\newcommand{\sX}{\mathscr X}
\newcommand{\sV}{\mathscr V}




\begin{document}

\begin{abstract}
Let $L$ be a finite extension of $\bQ_p$.\;Let $\rho_L$ be a potentially semi-stable non-crystalline  $p$-adic Galois representation such that the associated $F$-semisimple Weil-Deligne representation is absolutely indecomposable.\;In this paper,\;we study Fontaine-Mazur parabolic simple $\sL$-invariants of $\rho_L$, which was previously only known in the trianguline case.\;Based on the previous work \cite{2022ext1hyq} on Breuil's parabolic simple $\sL$-invariants,\;we
attach to $\rho_L$ a locally $\bQ_p$-analytic representation $\Pi(\rho_L)$ of $\GLN_{n}(L)$, which carries the information of parabolic simple $\sL$-invariants of $\rho_L$.\;When $\rho_L$  comes from a patched automorphic representation of  $\bG(\BA_{F^+})$ (for a define unitary group $\bG$ over a totally real field $F^+$ which is compact at infinite places and $\GLN_n$ at $p$-adic places),\;we prove under mild hypothesis that $\Pi(\rho_L)$ is a subrepresentation of the associated Hecke-isotypic subspace of the Banach spaces of (patched) $p$-adic automophic forms on $\bG(\BA_{F^+})$,\;this is equivalent to say that the Breuil's parabolic simple $\sL$-invariants are equal to Fontaine-Mazur parabolic simple $\sL$-invariants.\;
\end{abstract}

 \tableofcontents %% Just for papers exceeding 50 pages.

\section{Introduction and notation}




This paper aims to investigate the $p$-adic aspects of the Langlands program.\;The $p$-adic Langlands correspondence seeks to relate $p$-adic Galois representations to Banach space representations (or locally analytic representations) of $p$-adic reductive groups.\;

Associating admissible continuous unitary  $p$-adic Banach space representation $\Pi(\rho_p)$ of $\GLN_2(\BQ_p)$ to two-dimensional continuous representation $\rho_p:\gal_{\BQ_p}\rightarrow \GLN_2(E)$,\;was first proposed by Breuil for those $\rho_p$ that are potentially semistable with distinct Hodge-Tate weight \cite{Breuil03a},\;\cite{Breuil03b}.\;Through the work of many people,\;Breuil's conjectured correspondence or the $p$-adic Langlands correspondence for $\GLN_2(\BQ_p)$ (and its  $p$-adic local-global compatibility) have been achieved largely by the the fundamental work of Colmez \cite{PADINGL2COL} and Emerton \cite{localglobalgl2}.\;The current focus of research is on higher-dimensional representations and on base fields other than the field $\bQ_p$.\;The case $n\geq 3$ is much more difficult and only a few partial results are known.\;

If the $\rho_p$ in the above Breuil's consideration is genuinely irreducible semistabelline,\;Breuil proposed that $\Pi(\rho_p)$ should be a certain explicit completion of the locally algebraic representation $\Pi(\rho_p)_{\alge}$ depending on the $\sL$-invariant $\cL(\rho_p)$ of $\rho_p$.\;This 
$\sL$-invariant can be read off explicitly from an admissible Hodge filtration on some linear algebra data via Fontaine's theory.\;



But the terminology ``$\sL$-invariant(s)" has a long history.\;The study of  $\sL$-invariants started with the work of Mazur-Tate-Teitelbaum, in which the  $\sL$-invariant $\cL(f)$ describes the derivative of the  $p$-adic  $L$-function (for a certain modular form $f$) at its exceptional zero.\;If we write $\rho_f$ for the $\gal(\overline{\BQ}/\BQ)$-representation attached to $f$ via Langlands correspondence.\;Then the derivative $\cL(f)$ can be read off explicitly from Fontaine-Mazur $\sL$-invariant of $\rho_{f,p}:=\rho_f|_{\gal_{\BQ_p}}$.\;There were later several other definitions of $\sL$-invariants via different approaches.\;Through the work of many people, these  $\sL$-invariants turned out to be all the same in the modular form case.\;

Among many different definitions of the $\sL$-invariant(s) $\cL(f)$,\;Breuil constructs an explicit finite length locally analytic representation $\Pi_{f,p}$ of $\GLN_2({\BQ}_p)$ whose isomorphism class recovers $\cL(f)$. Furthermore,\;he shows that $\Pi_{f,p}$ can be embedded into the $f$-isotypic component of completed cohomology of modular curves (we usually use the term \textit{Breuil's $\sL$-invariants} for invariants defined in terms of Breuil's constructions).\;The second result is usually called ($p$-adic) local-global compatibility.\;Therefore,\;the construction of Breuil's $\sL$-invariants is actually one of the first instances of the $p$-adic Langlands program (for $\GLN_2(\BQ_p)$).\;

In general,\;the philosophy of the $p$-adic local Langlands program,\;which goes back to Breuil's initial ideals on $\sL$-invariants,\;gives two problems:
\begin{itemize}
	\item[(I)] To find the information of $p$-adic Galois representations in $p$-adic automorphic representations. More precisely,\;we need to find the missing information  of $\rho_L$  when passing from an $n$-dimensional de Rham $p$-adic Galois representation $\rho_L$ to its associated Weil-Deligne representation,\;on the automorphic side,\;e.g.,\;in the Banach representations or locally ${\BQ}_p$-analytic  representations of $\GLN_n(L)$.\;The lost information (besides the Hodge-Tate weights) of $p$-adic Galois representations can be concretely described via  Fontaine-Mazur $\sL$-invariants.\;
	\item[(II)] More precisely,\;to seek generalizations of Breuil's $\sL$-invariants to $\GLN_n(L)$,\;i.e.,\;to recover all Fontaine-Mazur $\sL$-invariants of $\rho_L$ from a certain locally analytic representation $\Sigma(\rho_L)$.\;
\end{itemize}
%The proofs of the equality witnessed the profound development of the $p$-adic arithmetic, including the $p$-adic comparison theorems, $p$-adic Galois deformations,\;and $p$-adic Langlands program.(we usually use the term \textit{Fontaine-Mazur $\sL$-invariant} for invariants defined in terms of the Hodge filtration).\;,\;a basic problem (and a starting point) in the $p$-adic Langlands program is to %Among  these different definitions of $\sL$-invariants,\;we focus on the  (see \cite{BCLINV2}).\;
In \cite{2019DINGSimple},\;Yiwen Ding studies this problem and extends the theory of 
$\sL$-invariants to higher-rank groups. He generalizes Breuil's approach (see \cite{BCLINV2}) to patched $p$-adic automorphic forms on certain definite unitary groups,\;and defines what he calls Breuil's simple $\sL$-invariants and shows that they are equal to Fontaine-Mazur $\sL$-invariants of two-dimensional subquotients of the associated semi-stable non-crystalline (trianguline) local Galois representation (or of the associated $(\varphi,\Gamma)$-module).\;

\textit{While the trianguline $p$-adic Galois representations are studied widely,\;there are fewer examples of results for the non-trianguline $p$-adic Galois representations.\;}In this paper,\;we extend the theory of {Fontaine-Mazur  simple $\sL$-invariants} and  {Breuil's simple $\sL$-invariants} to certain potentially semi-stable non-crystalline (not necessarily trianguline) Galois representation $\rho_L$ such that  the associated smooth representation is given by  the Zelevinsky-segment.\;We first define \textit{parabolic Fontaine-Mazur simple $\sL$-invariants} for this special $p$-adic Galois representation $\rho_L$.\;On the other hand,\;we have studied \textit{parabolic Breuil's simple $\sL$-invariants} in \cite{2022ext1hyq}.\;They are $\ext$-groups between some locally $\BQ_p$-analytic generalized \textit{parabolic} Steinberg representations of $\GLN_n(L)$ attached to a Zelevinsky-segment.\;

Finally,\;we establish some local-global compatibility results,\;i.e.,\;the correspondence between these two parabolic simple $\sL$-invariants can be realized in the $p$-adic completed cohomology  of some Shimura varieties (especially,\;in the space of $p$-adic automorphic forms on certain definite unitary group).\;Our results were previously only known in the trianguline case.\;We prove this result by using the geometry of Bernstein eigenvarieties,\;which were developed by Christophe Breuil and  Yiwen  Ding \cite{Ding2021} (since the non-trianguline $p$-adic Galois representations do not occur in the "classical" theory of eigenvarieties,\;we need the framework of Bernstein eigenvarieties).\;In \cite{Ding2021},\;the authors use the geometry of these Bernstein eigenvarieties to obtain various local-global compatibility results in the \textit{generic non-trianguline} case.\;Note that our local-global compatibility results lie in the \textit{non-generic non-trianguline} case.\;This is different from that in \cite{Ding2021}.\;

Our result coincides with the work of Ding \cite{2019DINGSimple} when our potentially semi-stable non-crystalline  Galois representation $\rho_L$ is collapsed to the trianguline case.\;This paper gives a parabolic generalization of Ding's work \cite{2019DINGSimple},\;and gives new evidence for the $p$-adic Langlands correspondence.\;We sketch the main results of this paper as follows.\;

\subsection{Statements of the main results}
Let $L$ (resp. $E$) be a finite extension of $\BQ_p$.\;Suppose that $E$ is sufficiently large containing all the embeddings $\Sigma_L:=\{\sigma: L\hookrightarrow \overline{\BQ}_p\}$ of $L$ in $\overline{\BQ}_p$.\;Put $q_L:=p^{f_L}$,\;where $f_L$ denotes the unramified degree of $L$ over $\BQ_p$.\;



Fix two integers $k$ and $r$ such that $n=kr$.\;Let $\rho_L:\gal_L\rightarrow \GLN_n(E)$ be a potentially semi-stable non-crystalline $p$-adic Galois representation.\;Let $\Dpik=D_{\rig}(\rho_L)$ be the $(\varphi,\Gamma)$-module of rank $n$ over $\cR_{E,L}$ associated to the $\rho_L$.\;Let  $\bh:=(\hpi_{\tau,1}>\hpi_{\tau,2}>\cdots>\hpi_{\tau,n} )_{\tau\in \Sigma_L}$ be the Hodge-Tate weights of $\Dpik$.\;We put $\hpi_{i}=(\hpi_{\tau,i})_{\tau\in \Sigma_L}$ for $1\leq i\leq n$.\;

Let $\wdre_0$ be an $r$-dimensional (absolutely) irreducible Weil-Deligne representation of the Weil group $W_L$ over $E$.\;Let $\Delta$ be the $p$-adic differential equation associated to $\wdre_0$.\;Recall from \cite{berger2008equations} that $\Delta$ is an irreducible $(\varphi,\Gamma)$-module of rank $r$ over $\cR_{E,L}$ which is de rham of constant Hodge-Tate weight $0$ such that $D_{\pst}(\Delta)$ (forgetting the Hodge filtration) is isomorphic to the (absolutely) irreducible Deligne-Fontaine module  associated by Fontaine to  $\wdre_0$ \cite[Proposition 4.1]{breuil2007first}.

Assume that the Galois representation $\rho_L$ admits the following non-critical special parabolization (more precisely,\;the so-called \textit{non-critical special $\omepik$-filtration},\;see Definitions \ref{weaklynoncritical} and \ref{dfnnoncriticalspecial} for more precise statements).\;This is a parabolic analogue of triangulation.\;

We say that $\rho_L$ admits a non-critical special parabolization $\cF$ if $\Dpik$ admits an increasing filtration by saturated $(\varphi,\Gamma)$-submodules
\[\cF=\fil_{\bullet}^{\cF} \Dpik: \ 0 =\fil_0^{\cF} \Dpik \subsetneq \fil_1^{\cF}\Dpik \subsetneq \cdots \subsetneq \fil_{k}^{\cF} \Dpik=\Dpik,\]
such that for $1\leq i\leq k$,\;we have an injection of $(\varphi,\Gamma)$-modules over $\cR_{E,L}$ of rank $r$
\begin{equation}\label{introDpikinjection}
	\mathbf{I}_i:\gr_{i}^{\cF}\Dpik:=\fil_i^{\cF}\Dpik/\fil_{i-1}^{\cF}\Dpik \hookrightarrow {\Delta\;}\tee \cR_{E,L}(\unr(\alpha q_L^{i-k})z^{\hpi_{ir}}),
\end{equation}
for some $\alpha\in E^\times$,\;where $z^{\hpi_{ir}}:=\prod_{\tau\in  \Sigma_L}\tau(z)^{\hpi_{\tau,ir}}$.\;Moreover,\;we assume that  the Hodge-Tate weights of $\fil_i^{\cF} \Dpik$ (resp.,\;$\gr_i^{\cF} \Dpik$) are given by (i.e.,\;the so-called non-critical assumption) $$\{\hpi_{\tau,1},\hpi_{\tau,2},\cdots,\hpi_{\tau,ir}\}_{\tau\in \Sigma_L}\;\; (\text{resp.,}\;\{\hpi_{\tau,(i-1)r+1},\hpi_{\tau,(i-1)r+2},\cdots,\hpi_{\tau,ir}\}_{\tau\in \Sigma_L}).\;$$
Suppose that the $F$-semi-simple Weil-Deligne representation $\mathrm{WD}(\Dpik)^{\mathrm{F-ss}}$ associated to $\Dpik$ is (absolutely) indecomposable,\;i.e.,\;$\mathrm{WD}(\Dpik)^{\mathrm{F-ss}}\cong(r_L,N)$ (up to some unramified twist),\;where the underlying representation $r_L$ is isomorphic to $ \oplus_{i=1}^k\wdre_0|\cdot|_L^{k-i}$ and the monodromy operator $N$ is of full rank (i.e.,\;$N^{k-1}\neq 0$).\;

Note that the parameters of our non-critical special parabolization $\cF$ are \textit{not generic},\;which is different from the generic assumption in \cite[(4.13)]{Ding2021}.\;See Remark \ref{nongenediffobstru} and  Remark \ref{localmodeldescri} for some statements on the difference between these two cases.\;


\subsection*{Parabolic Fontaine-Mazur simple $\sL$-invariants}

Keep the assumptions on $\rho_L$.\;Then we can attach to $\rho_L$ the parabolic Fontaine-Mazur simple $\sL$-invariants $\sL(\rho_L)$.\;We sketch the constructions of $\sL(\rho_L)$.\;

The parabolic simple $\sL$-invariants contain certain information on the consecutive extensions $\Dpik_i^{i+1}$ of $\gr_{i+1}^{\cF}\Dpik$ by $\gr_{i}^{\cF}\Dpik$ for $1\leq i\leq k-1$ inside $\Dpik$ (which we identify the set $\Delta_k:=\{1,\cdots,k-1\}$ with the set of simple roots $\Delta_n(k):=\{r,2r,\cdots,(k-1)r\}$ of $\GLN_n$).\;

For $1\leq i\leq k$,\;we  put $\bm{\delta}_{\bh,i}:=\unr(q_L^{i-k})z^{\hpi_{ir}}$.\;Using the injections $\mathbf{I}_{i}$ and $\mathbf{I}_{i+1}$  (see (\ref{introDpikinjection})),\;
we can construct a non-degenerate pairing (but maybe not perfect):
\begin{equation}\label{introperfectpairing}
	\begin{aligned}
		\ext^1_{(\varphi,\Gamma)}(\gr_{i+1}^{\cF} \Dpik,\gr_{i}^{\cF}\Dpik)  \;\;\times\;\;  &\homo(L^\times,E) \big(\cong \hH^1_{(\varphi,\Gamma)}(\cR_{E,L})\big) \\      &\xrightarrow{\cup}
		E\big(\cong\ext^2_{(\varphi,\Gamma)}(\cR_{E,L}(\bm{\delta}_{\bh,i}\bm{\delta}^{-1}_{\bh,i+1}))\big)
	\end{aligned}
\end{equation}
where $\homo(L^\times,E)$ denotes the $d_L+1$-dimensional $E$-vector space of $E$-valued additive characters on $L^\times$.\;The extension $\Dpik_i^{i+1}$ gives an extension class $[\Dpik_{i}^{i+1}]\in \ext^1_{(\varphi,\Gamma)}(\gr_{i+1}^{\cF} \Dpik,\gr_{i}^{\cF}\Dpik)$.\;We let
\[\sL(\rho_L)_{ir}\subset \homo(L^\times,E)\]
be the $E$-vector subspace orthogonal to $[\Dpik_{i}^{i+1}]$ via the non-degenerate pairing  (\ref{introperfectpairing}),\;which is a $d_L$-dimensional $E$-vector space.\;We put $\sL(\rho_L):=\{\sL(\rho_L)_{ir}\}_{1\leq i\leq k-1}$,\;which we call the \textit{parabolic Fontaine-Mazur simple $\sL$-invariants} of $\rho_L$.\;On the other hand,\;these invariants also characterize obstructions to certain
$1$-order deformations of $\Dpik$,\;see Theorem \ref{defor1thdefor} (that we call Colmez-Greenberg-Stevens formula) for more precise statements.\;

These simple $\sL$-invariants were previously only known in the trianguline case.\;When $r=1$ (i.e.,\;the trianguline case),\;our parabolic Fontaine-Mazur simple $\sL$-invariants coincide with the Fontaine-Mazur simple  $\sL$-invariants in \cite[Page 7994]{2019DINGSimple}.\;


\subsection*{Parabolic Breuil's simple $\sL$-invariants \cite{2022ext1hyq}}

%This part has been completed in the paper \cite{2022ext1hyq}.\;By \cite{berger2008equations} and the normalized local Langlands correspondence (see \cite{scholze2013local}),\;the $p$-adic differential e

Let $\Delta_n$ be the set of simple roots of $\GLN_n$ (with respect to the Borel subgroup $\bB$ of upper triangular matrices),\;and we identify the set $\Delta_n$ with the set $\{1,2,\cdots,n-1\}$. Let $\bT$ be the torus of diagonal matrices.\;We put  $\Delta_n(k):=\{r,2r,\cdots,(k-1)r\}\subseteq \Delta_n$ and $\Delta_n^k:=\Delta_n\backslash \Delta_n(k)$. For a subset $I\subseteq \Delta_n(k)$,\;we denote by $\bP_I^{\lrr}$ the parabolic subgroup of $\GLN_n$ containing the Borel subgroup $\bB$  such that $\Delta_n(k) \backslash  I$ are precisely the simple roots of the unipotent radical of $\mathbf{P}_I^{\lrr}$.\;Let $\bL_I^{\lrr}$ be the Levi subgroup of $\bP_I^{\lrr}$ containing the group $\bT$ such that  $I\cup \Delta_n^k$ is equal to the set of simple roots of  $\bL_I^{\lrr}$.\;Let $\op_I^{\lrr}$ be the parabolic subgroup opposite to $\bP_I^{\lrr}$.\;In particular,\;we have 
\[\bL^{\lrr}:=\bL^{\lrr}_{\emptyset}=\left(\begin{array}{cccc}
	\GLN_r & 0 & \cdots & 0 \\
	0 & \GLN_r & \cdots & 0 \\
	\vdots & \vdots & \ddots & 0 \\
	0 & 0 & 0 & \GLN_r \\
\end{array}\right)\subset \bP^{\lrr}:=\bP^{\lrr}_{\emptyset}=\left(\begin{array}{cccc}
	\GLN_r & \ast & \cdots & \ast \\
	0 & \GLN_r & \cdots & \ast \\
	\vdots & \vdots & \ddots & \ast \\
	0 & 0 & 0 & \GLN_r \\
\end{array}\right).\]
For simplicity,\;if $I=\{ir\}$ for some $1\leq i\leq k-1$,\;we put $\bL^{\lrr}_{ir}:=\bL^{\lrr}_{\{ir\}}$ and $\op^{\lrr}_{ir}:=\op^{\lrr}_{\{ir\}}$.\;For $\alpha\in E^\times$,\;we denote by $\unr(\alpha)$ the unramified character of $L^\times$ sending uniformizers to $\alpha$.\;


Recall the (absolutely) irreducible Weil-Deligne representation $\wdre_0$.\;Let $\pi$ be the associated irreducible cuspidal  representation of $\GLN_r(L)$ over $E$   via the classical local Langlands correspondence (normalized as in \cite{scholze2013local}).\;Following \cite{av1980induced2},\;we consider the Zelevinsky-segment \[\Delta_{[k-1,0]}(\pi):=\pi|\det|_L^{r-1}\otimes_E \cdots \otimes_E \pi|\det|_L \otimes_E \pi,\;\]which forms an irreducible cuspidal smooth representation of $\bL^{\lrr}(L)$ over $E$.\;

Put ${\bm\lambda}_\bh:=(\hpi_{\tau,i}+i-1)_{\tau\in \Sigma_L,1\leq i\leq n}$,\;which is a dominant weight of $(\mathrm{Res}_{L/\BQ_p}\GLN_n)\times_{\BQ_p}E$ with  respect to $(\mathrm{Res}_{L/\BQ_p}\bB)\times_{\BQ_p}E$.\;Let  $\{v^{\ana}_{\op^{\lrr}_{I}}(\pi,{\bm\lambda}_\bh)\}_{I\subseteq {\Delta_n(k)}}$ (resp.,\;$\{v^{\infty}_{\op^{\lrr}_{I}}(\pi,{\bm\lambda}_\bh)\}_{I\subseteq {\Delta_n(k)}}$) be the locally $\bQ_p$-analytic (resp.,\;locally $\bQ_p$-algebraic)
generalized parabolic Steinberg representations of $\GL_n(L)$ (see Section 3).\;In particular,\;we denote by \[\st_{(r,k)}^{\infty}(\pi,{\bm\lambda}_\bh):=v^{\infty}_{\op^{\lrr}_{\emptyset}}(\pi,{\bm\lambda}_\bh) \big(\text{resp.,\;}\st_{(r,k)}^{\ana}(\pi,{\bm\lambda}_\bh):=v^{\ana}_{\op^{\lrr}}(\pi,{\bm\lambda}_\bh)\big)\] the locally algebraic (resp.,\;locally $\bQ_p$-analytic) parabolic Steinberg representation. Recall that the smooth Steinberg representation $\st_{(r,k)}^{\infty}(\pi,\underline{0})$ is the irreducible smooth representation of $\GL_n(L)$ associated
to the $F$-semi-simple Weil-Deligne representation  $\mathrm{WD}(\rho_L)^{\mathrm{F-ss}}$ via the normalized local Langlands correspondence (see \cite{scholze2013local}),\;up to some unramified twist.\;For $\alpha\in E^\times$ and $\ast\in\{\infty,\ana\}$,\;let  $\st_{(r,k)}^{\ast}(\alpha,\pi,{\bm\lambda}_\bh):=\unr(\alpha)\circ\det\otimes_E\st_{(r,k)}^{\ast}(\pi,{\bm\lambda}_\bh)$ and $v^{\ast}_{\op^{\lrr}_{I}}(\alpha,\pi,{\bm\lambda}_\bh):=\unr(\alpha)\circ\det\otimes_Ev^{\ast}_{\op^{\lrr}_{I}}(\pi,{\bm\lambda}_\bh)$.\;

The main result  of the paper \cite[Theorem 5.19]{2022ext1hyq} is to compute the extension groups of locally $\bQ_p$-analytic generalized parabolic Steinberg representations.\;This theorem implies that we can see the counterpart of parabolic Fontaine-Mazur simple $\sL$-invariants in certain locally analytic representations of $\GL_n(L)$ (or see (\ref{thmintro}),\;where we refer to as parabolic Breuil's  simple $\sL$-invariants).\;
\begin{thm}
	For $ir\in {\Delta_n(k)}$,\;there exists an isomorphism of $E$-vector spaces 
	\begin{equation}\label{thmintro1}
		\homo(L^\times,E)\xrightarrow{\sim }\ext^1_{\GL_n(L)}\big(v_{\op^{\lrr}_{ir}}^{\infty}(\alpha,\pi,{\bm\lambda}_\bh), \st_{(r,k)}^{\ana}(\alpha,\pi,{\bm\lambda}_\bh)\big).
	\end{equation}
	In particular,\;we have $\dim_E  \ext^1_{G}\big(v_{\op^{\lrr}_{ir}}^{\infty}(\alpha,\pi,{\bm\lambda}_\bh), \st_{(r,k)}^{\ana}(\alpha,\pi,{\bm\lambda}_\bh)\big)=d_L+1$.\;Let $\widetilde{\Sigma}_i^{\lrr}(\alpha,\pi,{\bm\lambda}_\bh, \psi)$ be the image of $\psi\in \homo(L^\times,E)$ via (\ref{thmintro1}).\;
\end{thm}

For $ir\in \Delta_n(k)$, let $V_i$ be an $E$-vector subspace of $\homo(L^\times,E)$  of dimension $d_L$,\;and let $V=\prod_{i=1}^kV_i$.\;The above theorem gives the following constructions.\;

We will construct  locally $\bQ_p$-analytic representations  $\widetilde{\Sigma}_i^{\lrr}(\alpha,\pi,{\bm\lambda}_\bh, V_i)$ (resp.,\;$	\widetilde{\Sigma}^{\lrr}(\alpha,\pi,{\bm\lambda}_\bh, V)$) of $\GL_n(L)$,\;such that $\widetilde{\Sigma}_i^{\lrr}(\alpha,\pi,{\bm\lambda}_\bh, \sL(\rho_L)_{ir})$ (resp.,\;$\widetilde{\Sigma}^{\lrr}(\alpha,\pi,{\bm\lambda}_\bh, V)$) is isomorphic to an extension of $v_{\op^{\lrr}_{ir}}^{\infty}(\alpha,\pi,{\bm\lambda}_\bh)^{\oplus d_L}$ (resp.,\;$\bigoplus_{ir\in \Delta_n(k)}v_{\op^{\lrr}_{ir}}^{\infty}(\alpha,\pi,{\bm\lambda}_\bh)^{\oplus d_L}$) by $\st_{(r,k)}^{\ana}(\alpha,\pi,{\bm\lambda}_\bh)$.\;We then construct  certain subrepresentations (which have a simpler and more clear structure) $\Sigma_i^{\lrr}(\alpha,\pi,{\bm\lambda}_\bh, V_i)$ of the above $\widetilde{\Sigma}_i^{\lrr}(\alpha,\pi,{\bm\lambda}_\bh, V_i)$ (resp.,\;$\Sigma^{\lrr}(\alpha,\pi,{\bm\lambda}_\bh, V)$ of $\widetilde{\Sigma}^{\lrr}(\alpha,\pi,{\bm\lambda}_\bh, V)$),\;which is isomorphic to an  extension of $v_{\op^{\lrr}_{ir}}^{\infty}(\alpha,\pi,{\bm\lambda}_\bh)^{\oplus d_L}$ (resp.,\;$\bigoplus_{ir\in \Delta_n(k)}v_{\op^{\lrr}_{ir}}^{\infty}(\alpha,\pi,{\bm\lambda}_\bh)^{\oplus d_L}$) by some subrepresentations of $\st_{(r,k)}^{\ana}(\alpha,\pi,{\bm\lambda}_\bh)$ (see Section 3,\;(\ref{subofsigma})). When $V$ is equal to the parabolic simple $\sL$-invariants $\sL(\rho_L)$,\;the locally $\bQ_p$-analytic representations $\widetilde{\Sigma}^{\lrr}(\alpha,\pi,{\bm\lambda}_\bh, \sL(\rho_L))$ and $	\Sigma^{\lrr}(\alpha,\pi,{\bm\lambda}_\bh, \sL(\rho_L))$ carry the exact information of $\{\pi,{\bm\lambda}_\bh,\sL(\rho_L)\}$,\;i.e.,\;the information on the Weil-Deligne representation associated with $\rho_L$,\;the Hodge-Tate weights of $\rho_L$,\;and the parabolic simple $\sL$-invariants $\sL(\rho_L)$ of $\rho_L$.\;

\begin{rmk}
	The above two (parabolic) simple $\sL$-invariants are extension parameters at simple roots of $\GLN_n$.\;We also have concepts of higher $\sL$-invariants,\;which are  extension parameters at non-simple roots of $\GLN_n$.\;The next goal is to explore generalizations of Breuil's  $\sL$-invariants that conjecturally correspond to Fontaine-Mazur $\sL$-invariants (or parabolic Fontaine-Mazur $\sL$-invariants) at non-simple roots.\;See \cite{schraen2011GL3},\;\cite{HigherLinvariantsGL3Qp},\;\cite{breuil2019ext1},\;and \cite{Dilogarithm} for some partial results.\;
\end{rmk}








\subsection*{Local-global compatibility}
In this paper,\;we prove some new local-global compatibility results in the $p$-adic Langlands program by studying the geometry of the patched Bernstein eigenvarieties and Bernstein paraboline varieties \cite{Ding2021}.\;

By \cite{PATCHING2016} (hence we also assume the so-called `standard Talyor-Wiles assumptions),\;we have a certain definite unitary group $\bG$ over $F^+$,\;where $F^+$ is a maximal totally real subfield of an imaginary CM field .\;Let $\fp$ be a $p$-adic place of $F^+$.\;Let $L=F^+_\fp$.\;We then have a continuous Banach representation $\Pi_\infty$ of $G=\GLN_{n}(L)$,\;which is equipped with a continuous action of certain patched Galois deformation ring $R_\infty$ commuting with the  $G$-action.\;The module $\Pi_\infty$ patches the $p$-adic automorphic forms on  $\bG$.\;See Section \ref{BENVARPARAVAR} for a summary.\;

Fix the above smooth representation $\pi$ of $\GLN_r(L)$ over $E$ (see the beginning of the previous part).\;Let $\omepik$ be the cuspidal Bernstein component of $\bL^{\lrr}(L)$ containing the Zelevinsky-segment $\Delta_{[k-1,0]}(\pi)=[\pi\otimes |\det|_L^{k-1},\cdots,\pi]$. In Section \ref{BENVARPARAVAR},\;we adapt
the arguments in \cite[Section 3.3,\;Section 4.2]{Ding2021} to our case.\;We define a patched Bernstein eigenvariety $\bersteineigenvarpik$ concerning the cuspidal Bernstein component $\omepik$ and weight ${\bm\lambda}_\bh$,\;where $\overline{\rho}$ is a certain global residue Galois representations,\;which is a suitable globalization of a certain local residue Galois representation $\overline{r}:\gal_L\rightarrow \GLN_{n}(k_E)$, where $k_E$ is the residue field of $E$.\;

Suppose that the above $p$-adic Galois representation $\rho_L$ appears on the  patched Bernstein eigenvariety $\bersteineigenvarpik$ (see Section \ref{appearpoint} for a more precise statement).\;In this case,\;we can associate with $\rho_L$ the  parabolic simple $\sL$-invariants $\sL(\rho_L)$ and  an admissible unitary Banach representation $\widehat{\Pi}(\rho_L)$ of $G$ (see \cite{PATCHING2016},\;where the method is global).\;Moreover,\;we get an injection (by a parabolic analogue of the so-called global triangulation theory;\;see Lemma \ref{detlameda}) of $G$-representations
\[\st_{(r,k)}^{\infty}(\alpha,\pi,{\bm\lambda}_\bh)\hookrightarrow \widehat{\Pi}(\rho_L)\]
for some $\alpha\in E^\times$.\;The following theorem is the second main result of the paper (see Theorem \ref{thm: lgln-main} ).
\begin{thm}\label{introthm: lgln-main}The injection $\st_{(r,k)}^{\infty}(\alpha,\pi,{\bm\lambda}_\bh)\hookrightarrow \widehat{\Pi}(\rho_L)$ extends uniquely to an injection of $G$-representations
	\[\Sigma^{\lrr}(\alpha,\pi,{\bm\lambda}_\bh, \sL(\rho_L))\hookrightarrow \widehat{\Pi}(\rho_L).\;\]
	Furthermore,\;let $0\neq \psi\in \homo(L^{\times},E)$ and $ir\in \Delta_n(k)$, an injection
	\begin{equation*}
		f: \st_{(r,k)}^{\infty}(\alpha,\pi,{\bm\lambda}_\bh)\hooklongrightarrow \widehat{\Pi}(\rho_L)   \end{equation*}
	can extend to an injection $\Sigma_{i}^{\lrr}(\alpha, \pi,{\bm\lambda}_\bh, \psi)\hookrightarrow \widehat{\Pi}(\rho_L)$ if and only if $\psi\in \sL(\rho_L)_{ir}$.\;Therefore, $\sL(\rho_L)$ can be read out from $\widehat{\Pi}(\rho_L)$.\;
\end{thm}
Indeed,\;one can prove that the injection $\st_{(r,k)}^{\infty}(\alpha,\pi,{\bm\lambda}_\bh)\hookrightarrow \widehat{\Pi}(\rho_L)$ extends uniquely to a non-zero morphism of $G$-representations
\[\widetilde{\Sigma}^{\lrr}(\alpha,\pi,{\bm\lambda}_\bh, \sL(\rho_L))\rightarrow \widehat{\Pi}(\rho_L),\;\]
if $\soc_G(\mathrm{St}_{(r,k)}^{\ana}(\pi,\underline{\lambda}))\cong \mathrm{St}_{(r,k)}^{\infty}(\pi,\underline{\lambda})$ (this is true for $(r,k)=(2,2)$,\;the general case is not known at present),\;then we can prove that this non-zero morphism is also injective.\;

\subsection{Sketch of the proof of Theorem \ref{introthm: lgln-main}}
The proof  follows along the line of \cite[Page 7996-7997]{2019DINGSimple}.\;We only show the existence of an injection $\Sigma_{i}^{\lrr}(\alpha, \pi,{\bm\lambda}_\bh, \psi)\hookrightarrow \widehat{\Pi}(\rho_L)$ if $\psi\in \sL(\rho_L)_{ir}$.\;

We first need a result on certain deformation of type $\omepik$ of $\rho_L$.\;Let $F_{\Dpik,\cF}^0$ denote the functor on $\Art_E$ (the category of Artinian local $E$-algebras with residue field $E$) ``parameterizing deformations of $\Dpik$ which admit certain $\omepik$-filtration".\;The definition of $F_{\Dpik,\cF}^0$ gives a natural map \[d\kappa:F_{\Dpik,\cF}^0(E[\epsilon]/\epsilon^2)\rightarrow
\homo(\bZ^{\lrr}(L),E)\cong \prod_{i=1}^k\homo(L^\times,E).\;\]Consider the  following composition:
\begin{equation}\label{introkappakappaL}
	\kappa:F_{\Dpik,\cF}^0(E[\epsilon]/\epsilon^2)\xrightarrow{d\kappa}  \prod_{i=1}^k\homo(L^\times,E)\xrightarrow{\kappa_L} \prod_{ir\in \Delta_n(k)}\homo(L^\times,E),
\end{equation}
where the second map $\kappa_L$ sends $(\psi_1,\psi_2,\cdots,\psi_k)$ to $(\psi_i-\psi_{i+1})_{ir\in \Delta_n(k)}$.\;We get the following proposition (see Corollary \ref{corlinvaraint} and Proposition \ref{dimtanFD0}) by the Galois cohomologies argument.\;
%,\;which is a pure Galois result and follows by Galois cohomology calculations.\;

\begin{pro}\label{introdimeF0D}The functor $F_{\Dpik,\cF}^0$ is  pro-representable.\;The tangent space $F_{\Dpik,\cF}^0(E[\epsilon]/\epsilon^2)$ has $E$-dimension $1+d_L\Big(k+\frac{n(n-r)}{2}\Big)$.\;The map $\kappa$ (\ref{introkappakappaL})  factors through a surjective map
	\[\kappa:F_{\Dpik,\cF}^0(E[\epsilon]/\epsilon^2)\twoheadlongrightarrow \sL(\rho_L).\]
\end{pro}


Recall that the  patched Galois deformation ring $R_\infty$ admits a decomposition $R_\infty\cong R^{\fp}_{\infty}\widehat{\otimes}_{\cO_E}R^{\Box}_{\overline{r}}$, where $R^{\fp}_{\infty}$ is the ``prime-to-$\fp$" part of $R_\infty$ and $R^{\Box}_{\overline{r}}$ is the maximal reduced and $p$-torsion free quotient of the universal $\cO_E$-lifting ring of $\overline{r}$.\;Let $\mathfrak{X}_{\overline{r}}^\Box=(\Spf\;R^{\Box}_{\overline{r}})^{\rig}$ be the 
rigid analytic space associated to the formal scheme $\Spf\;R^{\Box}_{\overline{r}}$.\;Let $\mathfrak{Z}_{\omepik}$ be the Bernstein center of $\omepik$ over $E$.\;By an easy variation of  \cite[Section 3.3]{Ding2021} (see Section \ref{BENVARPARAVAR} for a brief summary),\;we can construct the so-called patched Bernstein eigenvariety 
\begin{equation}\label{intropatcheigrndfn}
	\bersteineigenvarpik \hooklongrightarrow (\Spf\;R^{\fp}_{\infty})^{\rig}\times \mathfrak{X}_{\overline{r}}^\Box \times \sbanpik\times \rigch
\end{equation}
from the Jacquet-Emerton module $J_{\bP^{\lrr}(L)}(\Pi_\infty^{R_\infty-\ana})$,\;where
$\Pi_\infty^{R_\infty-\ana}$ denotes the locally $R_\infty$-analytic vectors in $\Pi_\infty$(see \cite[Section 3.1]{breuil2017interpretation}),\;$\rigch$ denotes the rigid space over $E$ parameterizing continuous characters of the center $\bZ^{\lrr}(L)$ of $\bL^{\lrr}(L)$.\;By an easy variation of  \cite[Section 3.3,\;Section 4.4]{Ding2021},\;the morphism (\ref{intropatcheigrndfn}) factors through a natural embedding
\begin{equation}\label{intropactchembedding}
	\iota_{\omepik}:\bersteineigenvarpik\hookrightarrow (\Spf\; R^{\fp}_{\infty})^{\rig}\times \defvar,
\end{equation}
where $\defvar\subset \mathfrak{X}_{\overline{r}}^\Box \times \sbanpik\times \rigch$ is the paraboline deformation variety with respect to $(\omepik,{\bm\lambda}_\bh)$ (see Section \ref{BENVARPARAVAR} for more precise statements).\;The embedding (\ref{intropactchembedding}) induces an isomorphism between $\bersteineigenvarpik$ and a union of irreducible components (equipped with the reduced closed rigid subspace structure) of $(\Spf\; R^{\fp}_{\infty})^{\rig}\times \defvar$.\;

Assume that we can associate with $\rho_L$ a point $x:=(x^{\fp}, \rho_L, \pi_{x,\bL^{\lrr}},\chi)\in \bersteineigenvarpik$.\;The embedding (\ref{intropactchembedding}) gives a point $x_L=\pr_2\circ\iota_{\omepik}(x)\in \defvarrho$ (where $\pr_2$ is the projection to the second factor).\;We now look at the local geometry of paraboline deformation variety $\defvarrho$ at point $x_L$.\;Let $\xfpx$ be the union of  irreducible components of $\defvarrho$ containing $x_L$.\;We show that the non-critical special
$\omepik$-filtration $\cF$ on $x_L$  extends to some open affinoid neighborhood around $x_L$ (see Theorem \ref{localfiltration}),\;where we use the accumulation property
at $x_L$ (see Definition \ref{accum} and Lemma \ref{accumulation}) and the non-critical assumption.\;Using this result,\;we get an upper bound of the $E$-dimension of tangent space $T_{\xfpx,x_L}$ (see the proof of Proposition \ref{smoothnessdefvar}).\;


%Let $\cE(x)$ be the union of  irreducible component of $\bersteineigenvarpik$ containing $x$.\;Then $\cE(x)=X^{\fp} \times \BU^g \times \iota_{\omepik}^{-1}(\xfpx)\hookrightarrow X^{\fp} \times \BU^g \times \iota_{\omepik}^{-1}(\defvarrho)$,\;where $X^{\fp}$ is an irreducible component of $\FX_{\overline{\rho}^{\fp}}^\Box$,\;and 


The natural embedding (\ref{intropatcheigrndfn}) induces the ``weight" map 
\begin{equation}
	\omega:	\bersteineigenvarpik \longrightarrow \sbanpik\times \rigch.\;
\end{equation}
Consider the tangent map of $\omega$ at $x$:
\begin{equation}
	d\omega_x:	T_{\bersteineigenvarpik,X} \longrightarrow T_{\sbanpik\times \rigch,\omega(x)}.\;
\end{equation}
By the embedding  (\ref{intropactchembedding}) and the upper bound of $\dim_ET_{\xfpx,x_L}$,\;we can deduce the following corollary from Proposition \ref{introdimeF0D} and Theorem \ref{localfiltration}.\;
\begin{cor}\label{introcoreigentangentmap2}$\bersteineigenvarpik$ is smooth at $x$.\;The composition $\kappa_L\circ d\omega_x$ factors through a surjective map
	\begin{equation}
		T_{\bersteineigenvarpik,x}\twoheadlongrightarrow \sL(\rho_L).\;
	\end{equation}
\end{cor}


%The application of local model (see \cite[Section ]{Ding2021}) allow a better understanding of closed subspace of the (patched) Bernstein eigenvariety.\;We first have natural morphisms $\complocaldefvarrhosub  \rightarrow\complocaldefvarrho\rightarrow (\widehat{\mathfrak{X}_{\overline{r}}^\Box})_{\rho_L}\cong X_{\rho_L}$,\;where the first morphism is a closed immersion.\;To see the higher-order deformation,\;we use a  deep theory on the local model of the paraboline deformation variety $\defvarrho$ (see \cite[Section]{Ding2021}).\;By inverting $t$,\;the filtration $\cF$ on $\Dpik$ induces an increasing filtration $\cM_{\rho,\bullet}=\{\fil_{i}^\cF D[\frac{1}{t}]\}$ on $\cM$ by $(\varphi,\Gamma)$-submodules over $\cR_{E,L}[\frac{1}{t}]$.\;Let $X_{\rho_L,\cM_{\bullet}}$ (resp.,\;$X_{\rho_L,\cM_{\bullet}}^w$) be the groupoid over $\Art_E$ of deformation of the group morphism $\rho_L$ such the associated $(\varphi,\Gamma)$ module over $\cR_{E,L}[\frac{1}{t}]$ admits an $\omepik$-filtration (resp.,\;defined as $X_{\rho_L,\cM_{\bullet}}$ by adding the condition:\;the relative location of certain two filtration is controlled by $w$,\;see \cite[Section 6.4]{Ding2021}).\;The analogue of \cite[Proposition 6.4.5,\;Corollary 6.4.7]{Ding2021} in non-critical special case is as follows.\;\begin{pro}The canonical morphism $\complocaldefvarrho \rightarrow X_{\rho_L}$ factors thought a morphism\[\complocaldefvarrho \rightarrow X_{\rho_L,\cM_{\bullet}}\cong X_{\rho_L,\cM_{\bullet}}^{w_0}.\]The morphisms  $\complocaldefvarrho \rightarrow X_{\rho_L}$  and  $\complocaldefvarrho \rightarrow X_{\rho_L,\cM_{\bullet}}$ are closed immersions of groupoids over $\Art_E$.\;The same inclusions hold by replacing $\complocaldefvarrho$ with $\complocaldefvarrhosub$.\;\end{pro}

%It coincides with the evaluation of following composition (see  (\ref{omegadeltapi}),\;(\ref{tangentLinv}),and (\ref{tangentpatchvar})) on the $E[\epsilon]/\epsilon^2$-points :
%\begin{equation}\label{higherdef}
%\begin{aligned}
%	\widehat{\bersteineigenvarpik_x}\cong \widehat{Z_x}
%		&\twoheadlongrightarrow \complocaldefvarrhosub\rightarrow X_{\rho_L,\cM_{\bullet}}\rightarrow \widehat{(\cZ_{\bL^{\lrr},L})}_{\bm{\delta}_{\bh}},
%	\end{aligned}
%\end{equation}
%where $\widehat{(\cZ_{\bL^{\lrr},L})}_{\bm{\delta}_{\bh}}$ is the completion of $\cZ_{\bL^{\lrr},L}$ at $\bm{\delta}_{\bh}$,\;which pro-represents the functor of  deformations of $\bm{\delta}_{\bh}$.\;In particular,\;$\widehat{(\cZ_{\bL^{\lrr},L})}_{\bm{\delta}_{\bh}}(E[\epsilon]/\epsilon^2)=\homo(\bZ^{\lrr}(L),E)\cong\prod_{i=1}^k\homo(L^\times,E)$.\;From Proposition \ref{introdimeF0D} and a calculation of the dimension of tangent space of groupoid $X_{\rho_L,\cM_{\bullet}}$,\;

We now recall briefly the proof of the existence of an injection $\Sigma_{i}^{\lrr}(\alpha, \pi,{\bm\lambda}_\bh, \psi)\hookrightarrow \widehat{\Pi}(\rho_L)$ if $\psi\in \sL(\rho_L)_{ir}$ .\;By Corollary \ref{introcoreigentangentmap2},\;there exists $\Psi\in \mathrm{Im}(d\omega_x)$ with $\kappa_L(\Psi)=\psi.\;$Then the construction of $\bersteineigenvarpik$ implies an injection of $\bL^{\lrr}(L)$-representations
\[W\te(\mathrm{1}_\Psi\circ\mathrm{det}_{\bL^{\lrr}(L)})\te\lamoverx \hooklongrightarrow J_{\bP^{\lrr}(L)}\big(\Pi_{\infty}^{R_{\infty}-\ana}[\fm_x^\infty]\big),\]
where $[\fm_x^\infty]$ denotes the subspace of the vectors annihilated by  a certain power of $\fm_x$,\;and where $W$ denotes certain irreducible smooth representation of $\bL^{\lrr}(L)$,\;$\mathrm{1}_\Psi$ denotes the extension of two trivial characters attached to $\Psi$,\;and $\lamoverx$ is the irreducible $\BQ_p$-algebraic representation of $\bL^{\lrr}(L)$ with highest weight $\lambda_\bh$.\;Then the non-critical assumption on $\rho_L$ implies that the above injection is balanced in the sense of \cite[Definition 0.8]{emerton2007jacquet}.\;In this case,\;the Emerton's adjunction formula \cite[Theorem 0.13]{emerton2007jacquet} deduces a non-zero morphism
\[I^G_{\op^{\lrr}}\Big(W\te(\mathrm{1}_\Psi\circ\mathrm{det}_{\bL^{\lrr}(L)})\te\lamoverx\Big) \longrightarrow\Pi_{\infty}^{R_{\infty}-\ana}[\fm_x^\infty].\]
By some locally analytic representation theory,\;we can show that this non-zero morphism induces an injection:
\[\Sigma_{i}^{\lrr}(\alpha, \pi,{\bm\lambda}_\bh, \psi)\hookrightarrow \widehat{\Pi}(\rho_L).\]

%\begin{rmk}
%\end{rmk}

\begin{rmk}The key step in the proof of the main theorem is the existence of the character  $\Psi\in \mathrm{Im}(d\omega_x)$.\;This is a  consequence of the surjectivity of Corollary \ref{introcoreigentangentmap2}.\;%The smoothness of $\bersteineigenvarpik$ at $x$ is also an important result.\;
\end{rmk}
\begin{rmk}
The proof of smoothness (Corollary \ref{introcoreigentangentmap2}) only requires the information on tangent space (i.e.,$1$-order deformations).\;We can use a local model of paraboline deformation variety (see \cite[Section 6]{Ding2021}) to see more information on arbitrary deformations.\;The application of the local models (see \cite[Section 6]{Ding2021}) allows a better understanding of the local geometry of the paraboline deformation variety and (patched) Bernstein eigenvariety.\;It seems likely that we can restate the results in Section  \ref{smoothpoint} under the framework of the local model (see Remark \ref{localmodeldescri}).\;We point out the main differences.\;Since the parameter of our $\omepik$-filtration is non-generic (in the sense of \cite[(6.5)]{Ding2021}),\;the morphism
$$X_{\cM,\cM_{\bullet}}\rightarrow \widehat{(\cZ_{\bL^{\lrr},L})}_{\bm{\delta}_{\bh}}\times_{\widehat\fz^\lrr_{L}} X_{\bW,\cF}$$of groupoids in \cite[Theorem\;6.2.6]{Ding2021} is no longer formally smooth.\;Thus the discussions in  \cite[Section 6.4]{Ding2021}  cannot be applied to our case (for example,\;the final result \cite[Corollary 6.4.7]{Ding2021}).\;It is valuable to explore local-global compatibility and Breuil's locally analytic socle conjecture \cite{breuil2016socle} for  critical special (potentially) semistable non-crystalline $p$-adic Galois representations (non-generic case).\;
\end{rmk}



\subsection{General notation}\label{notationger}
Let $L$ (resp. $E$) be a finite extension of $\BQ_p$ with $\co_L$ (resp. $\co_E$) as its ring of integers and $\varpi_L$ (resp. $\varpi_E$) a uniformizer. Suppose $E$ is sufficiently large,\;containing all the embeddings of $L$ in $\overline{\BQ}_p$. Put $\Sigma_L:=\{\sigma: L\hookrightarrow \overline{\BQ}_p\} =\{\sigma: L\hookrightarrow E\}.$ the set of all the embeddings of $L$ in $\overline{\BQ}_p$ (equivalently,\;in $E$).\;Let $k_E$ be the residue field of $E$.

Let $\valua_L(\cdot)$ (resp. $\valua_p$) be the $p$-adic valuation on $\overline{\BQ_p}$ normalized by sending uniformizers of $\co_L$ (resp. of $\BZ_p$) to $1$. Let $d_L:=[L:\BQ_p]=|\Sigma_L|$, let $e_L:=\valua_L(\varpi_L)$ and let $f_L:=d_L/e_L$. We have $q_L:=p^{f_L}=|\co_L/\varpi_L|$.

Let $L_0$ denote the maximal unramified subextension of $\BQ_p$ in $L$.\;We fix a compatible system $\{\epsilon_n\}$ of $p$-th roots of unity.\;Let $L_n=L(\epsilon_n)$,\;$L_\infty=\cup_nL_n$,\;and $L'_0$ the maximal unramified subextension of $\BQ_p$ in $L_\infty$.\;We write $\gal_{L'}=\gal(\overline{\BQ}_p/L')$,\;$\Gamma_{L'}=\gal(L'_{\infty}/L')$,\;and $\cH_{L'}=\gal(\overline{\BQ}_p/L'_{\infty})$ for any subfield $L'\subset \overline{\BQ}_p$.\;We omit the subscript ${L'}$ if $L'=L$.\;

Let $X$ be a scheme locally of finite type over $E$ or a locally noetherian formal scheme over $\cO_E$ whose reduction is locally finite type over $k_E$.\;Let $X^{\mathrm{rig}}$ the associated rigid analytic space over $E$.\;If $X$ is a scheme locally of finite type over $E$ or a rigid analytic space over $E$,\;we denote by $X^{\red}$ the associated reduced Zariski-closed subspace.\;If $x$ is a point of $X$,\;we denote by $\cO_{X,x}$ (resp.,\;$k(x)$) the  local ring (resp.,\;residue field) at $x$.\;Let $\widehat{\cO}_{X,x}$ be the $\fm_{\cO_{X,x}}$-adic completion of $\cO_{X,x}$,\;and $\widehat{X}_{x}:=\mathrm{Spf}\;\widehat{\cO}_{X,x}$.\;If $x$ is a closed point of $X$,\;then  $\widehat{\cO}_{X,x}$ is a noetherian complete local $k(x)$-algebra of residue field $k(x)$.\;
%zccLet $X$ be a scheme locally of finite type over $E$,\;or a rigid analytic space over $E$,\;we denoted by $X^{\mathrm{red}}$ the associated reduced Zariski-closed subspace.\;

\subsubsection*{General setting on Reductive groups and Lie algebras}
For a Lie algebra $\fh$ over $L$, and $\sigma\in \Sigma_L$, let $\fh_{\sigma}:=\fh\otimes_{L,\sigma} E$ (which is  a Lie algebra over $E$). For $J\subseteq \Sigma_L$, let $\fh_J:=\prod_{\sigma\in J} \fg_{\sigma}$. In particular, we have $\fh_{\Sigma_L}\cong \fh\otimes_{\BQ_p} E$.\;The notation $U(\fg)$ refers to the universal enveloping algebra of $\fg$ over $L$.\;

Let $\mathbf{H}$ be an algebraic group over $L$,\;and $\fh$ be its Lie algebra over $L$.\;Let $\mathrm{Res}_{L/\BQ_p}\mathbf{H}$ be the scalar restriction of $\mathbf{H}$ from $L$ to $\BQ_p$.\;The Lie algebra of $\mathrm{Res}_{L/\BQ_p}\mathbf{H}$ over $\BQ_p$ can also be identified with $\fh$,\;where $\fh$ is regarded as a Lie algebra over $\BQ_p$.\;We write $\mathbf{H}_{/E}=(\mathrm{Res}_{L/\BQ_p}\mathbf{H})\times_{\BQ_p}E$,\;which is an algebraic group over $E$.\;The Lie algebras of $\mathbf{H}_{/E}$ over $E$ is $\fh_{\Sigma_L}$.\;
%Note that we have a decomposition $\mathbf{H}_{/E}=\prod_{\sigma\in \Sigma_L}\mathbf{H}\times_{L,\sigma}E$.\;
Let $\mathbf{G}$ be a split connected reductive group over $L$,\;and $\fg$ be its Lie algebra over $L$.\;We fix a maximal split torus $\bT$ and write $\bB$ for a choice of Borel subgroup containing $\bT$.\;We use $\bP$ for the parabolic subgroup of $\mathbf{G}$ containing $\bB$,\;and let $\bL_{\bP}$ be the Levi subgroup of $\bP$ containing $\bT$.\;Let $\bN_{\bP}$ be the unipotent radical of $\bP$.\;Then $\bP$ admits a Levi decomposition $\bP=\bL_{\bP}\bN_{\bP}$.\;The Lie algebras (over $L$) of subgroups $\bT,\bB,\bP,\bL_{\bP},\bN_{\bP}$  are denoted by $\ft,\fb,\fp,\fl_{\bP},\fn_{\bP}$ respectively.\;

Note that the group $\mathbf{G}_{/E}$ is also a split-connected reductive group over $E$ with maximal split torus $\bT_{/E}$,\;Borel subgroup $\bB_{/E}$,\;parabolic subgroup $\bP_{/E}$,\;and ${\bL_\bP}_{/E}$ is also the Levi subgroup of ${\bL_\bP}_{/E}$ containing $\bT_{/E}$.\;The Lie algebras (over $E$) of reductive groups $\mathbf{G}_{/E}$,\;$\bP_{/E}$,\;and ${\bL_\bP}_{/E}$ are given by $\fg_{\Sigma_L}$,\;$\fp_{\Sigma_L}$ and $\fl_{\bP,\Sigma_L}$.\;We use the Roman letters $G$,\;$P$,\;etc. for the $L$-points $\bG(L)$,\;$\bP(L)$.\;We view these as locally $\BQ_p$-analytic groups.


\subsubsection*{General linear group \texorpdfstring{$\GLN_n$}{Lg}}
Let $\GLN_n$ be the general linear group over $L$.\;Let $\Delta_n$ be the set of simple roots of $\GLN_n$ (with respect to the Borel subgroup $\bB$ of upper triangular matrices).\;We identify the set $\Delta_n$ with $\{1,\cdots, n-1\}$ such that $i\in \{1,\cdots, n-1\}$ corresponds to the simple root $\alpha_i:\;(x_1,\cdots, x_n)\in \ft \mapsto x_i-x_{i+1}$, where $\ft$ denotes the $L$-Lie algebra of the torus $\bT$ of diagonal matrices.\;

For a subset $I\subset \Delta_n$,\;let $\mathbf{P}_I$ be the parabolic subgroup of
$\GLN_n$ containing $\bB$ such that $\Delta_n \backslash  I$ are precisely the simple roots of the unipotent radical $\mathbf{N}_I$ of $\mathbf{P}_I$.\;Denote by $\mathbf{L}_I$ the unique Levi subgroup of $\mathbf{P}_I$ containing $\bT$ such that $I$ is equal to the set of simple roots of $\mathbf{L}_I$.\;In particular,\;we have $\mathbf{P}_{\Delta_n}=\GLN_n$, $\mathbf{P}_{\emptyset}=\mathbf{B}$.\;Let $\overline{\mathbf{P}}_I$ be the parabolic subgroup opposite to $\mathbf{P}_I$.\;Let $\mathbf{N}_I$ (resp.\;$\overline{\mathbf{N}}_I$) be the nilpotent radical of $\mathbf{P}_I$ (resp.\;$\overline{\mathbf{P}}_I$),\;let $\mathbf{Z}_I$ be the center of $\mathbf{L}_I$ and let $\mathbf{D}_I$ be the derived subgroup of $\mathbf{L}_I$.\;Then we have a Levi decomposition $\mathbf{P}_I=\mathbf{L}_I\mathbf{N}_I$ (resp.\;$\overline{\mathbf{P}}_I=\mathbf{L}_I\overline{\mathbf{N}}_I$).\;Let $\bZ_n$ be the center of $\GLN_n$.\;Let $\fg$, $\fp_I$, $\fn_I$, $\fl_I$, $\overline{\fl}_I$, $\overline{\fp}_I$, $\overline{\fn}_I$,\;$\fd_I$,\;$\fz_I$,\;$\overline{\fz}_I$ and $\ft$ be the $L$-Lie algebras of $\GLN_n$,\;$\mathbf{P}_I$,\;$\mathbf{N}_I$,\;$\mathbf{L}_I$,\;$\mathbf{L}_I/\bZ_n$,\;$\overline{\mathbf{P}}_I$,\;$\overline{\mathbf{N}}_I$,\;$\bD_I$,\;$\bZ_{I}$,\;$\bZ_{I}/\bZ_n$ and $\bT$ respectively.\;



%see \cite[section.\;9]{humphreysBGG} for precise definition.\;
Denote by $\sW_n$ ($\cong S_n$) the Weyl group of $\GLN_n$, and denote by $s_i$ the simple reflection corresponding to $i\in \Delta_n$.\;For any $I\subset \Delta_n$,\;define $\sW_{I}$ as the subgroup of $\sW_{n}$ generated by simple reflections $s_i$ with $i\in I$.\;The Weyl group of $\GLN_{n/E}$ is $\sW_{n,\Sigma_L}:=\Pi_{\sigma\in \Sigma_L}\sW_{n,\sigma}\cong S_n^{d_L}$,\;where $\sW_{n,\sigma}\cong \sW_n$ is the $\sigma$-th factor of $\sW_{n,\Sigma_L}$.\;For $i\in \Delta_n$ and $\sigma\in \Sigma_L$,\;let $s_{i,\sigma}\in \sW_{n,\sigma}$ be the simple reflection corresponding to $i\in \Delta_n$.\;Let $\rho$ be the half of the sum of positive roots of $\GLN_n$.\;

We list the basic notation in the theory of the Zelevinsky-segment.\;Let $k,r$ be two integers such that $n=kr$.\;We put ${\Delta_n(k)}:=\{r,2r,\cdots,(k-1)r\}\subseteq \Delta_{n}$ and  $\Delta_n^k:=\Delta_{n}\backslash{\Delta_n(k)}$.\;

For a subset $I\subset {\Delta_n(k)}$,\;we put
\begin{equation}
	\begin{aligned}
		&\bL^{\lrr}_I:=\bL_{\Delta_n^k\cup I},\;\bP^{\lrr}_I:=\bP_{\Delta_n^k\cup I},\;\op^{\lrr}_I:=\op_{\Delta_n^k\cup I},\\
		&\bN^{\lrr}_I:=\bN_{\Delta_n^k\cup I},\;\onn^{\lrr}_I:=\onn_{\Delta_n^k\cup I},\;\bZ^{\lrr}_I:=\bZ_{\Delta_n^k\cup I},\;\bD^{\lrr}_I:=\bD_{\Delta_n^k\cup I}.
	\end{aligned}
\end{equation}
Let $\fl^{\lrr}_I$,\;$\overline{\fl}^{\lrr}_I$,\;$\fp^{\lrr}_I$, $\fn^{\lrr}_I$, $\overline{\fp}^{\lrr}_I$, $\overline{\fn}^{\lrr}_I$,\;$\fd^{\lrr}_I$,\;$\fz^{\lrr}_I$,\;$\overline{\fz}^{\lrr}_I$ be the $L$-Lie algebras of $\bL^{\lrr}_I$,\;$\bL^{\lrr}_I/\bZ_n$, $\bP^{\lrr}_I$, $\bN^{\lrr}_I$,$\op^{\lrr}_I$, $\onn^{\lrr}_I$,\;$\bD^{\lrr}_I$ and $\bZ^{\lrr}_I$,\;$\bZ^{\lrr}_I/\bZ_n$ respectively.\;Similarly,\;for a subset $S\subset \Sigma_L$,\;we put $\sW^{\lrr}_{I}:=\sW_{\Delta_n^k\cup I},$ and $\sW^{\lrr}_{I,S}:=\sW_{\Delta_n^k\cup I,S},\;^{I}\sW_{n,S}^{\lrr}:=^{\Delta_n^k\cup I}\sW_{n,S}$.\;When $I=\emptyset$,\;we omit the subscripts $I$.\;For example,\;we have
\[\bL^{\lrr}:=\left(\begin{array}{cccc}
\GLN_r & 0 & \cdots & 0 \\
0 & \GLN_r & \cdots & 0 \\
\vdots & \vdots & \ddots & 0 \\
0 & 0 & 0 & \GLN_r \\
\end{array}\right)\subseteq \op^{\lrr}:=\left(\begin{array}{cccc}
\GLN_r & 0 & \cdots & 0 \\
\ast & \GLN_r & \cdots & 0 \\
\vdots & \vdots & \ddots & 0 \\
\ast & \ast & \cdots & \GLN_r \\
\end{array}\right)\]
The parabolic subgroups of $\GLN_n$ containing the parabolic subgroup $\op^{\lrr}$ are  $\{\op^{\lrr}_I\}_{I\subseteq \Delta_n(k)}$.\;




\subsubsection*{$P$-adic Hodge theory}

Let $\cR_{L}:=\bB_{\rig,L}^{\dagger}$ be the Robba ring associated to $L$.\;Let $A$ (resp.,\;$X$) be an $\BQ_p$-affinoid algebra (resp.\;a rigid analytic space), and let $\cR_{A,L}:=\cR_{L}\widehat{\otimes}_{\BQ_p} A$ (resp.,\;$\cR_{X,L}$) for the Robba ring associated to $L$ with $A$-coefficient (resp.,\;with $\cO_X$-coefficient) (see \cite[Definition 6.2.1]{XLROBBA}).\;Let $\delta_A:L^\times\rightarrow A^\times$ be a continuous character.\;We write  $\cR_{A,L}(\delta_A)$ for the $(\varphi,\Gamma)$-module of character type  over $\cR_{A,L}$ associated to the continuous character $\delta_A$.\;Recall the cohomology of $(\varphi,\Gamma)$-modules (and Euler-Poincar\'{e} characteristic formula and Tate duality) defined in \cite{liu2007cohomology}.\;

We recall the concept of $B$-pairs briefly.\;Let $A$ be a local artinian $E$-algebra with residue field $E$.\;Recall the $E$-$B$-pair (resp.,\;$A$-$B$-pair) defined in \cite{nakamura2009classification} (resp.,\;\cite[Section 1.1]{2015Ding}).\;An $E$-$B$-pair of $\gal_L$ is a couple $W=(W_e,W^+_{\dr})$ such that $W_e$ is a  finite $B_e\otimes_{\bQ_p}E$-module  with a continuous semi-linear $\gal_L$-action which is free as $B_e$-module,\;and $W^+_{\dr}\subset W_{\dr}:=B_{\dr}\otimes_{B_e}W_e$ is a $\gal_L$-stable $B_{\dr}^+\otimes_{\bQ_p}E$-lattice.\;By \cite[Theorem 1.36]{nakamura2009classification} (resp.,\;\cite[Section 1.1]{2015Ding}),\;there exists an equivalence of categories between the category of $E$-$B$-pairs (resp.,\;the category of $A$-$B$-pairs) and  the category of $(\varphi,\Gamma)$-modules over $\cR_{E,L}$ (resp.,\;the category of $(\varphi,\Gamma)$-modules over $\cR_{A,L}$).\;Let $\chi:L^\times\rightarrow E^\times$ (resp.,\;$\chi_A:L^\times\rightarrow A^\times$) be a continuous character,\;and let $B_E(\chi)$ (resp.,\;$B_A(\chi_A)$)  be the rank-one $E$-$B$-pair (resp.,\;the rank-one $A$-$B$-pair) associated to $\chi$ (resp.,\;$\chi_A$) (see \cite[Theorem 1.45]{nakamura2009classification}).\;Recall in  \cite{nakamura2009classification} the definition of the cohomology of $E$-$B$-pairs (the $A$-$B$-pairs can also be viewed as $E$-$B$-pairs).\;Recall that cohomology of a $(\varphi,\Gamma)$-module over $\cR_{E,L}$ and the cohomology of the associated $E$-$B$-pair via the equivalence of categories \cite[Theorem 1.36]{nakamura2009classification} is isomorphic.\;

We also need the following equivalence of categories in the $p$-adic Hodge theory.\;Let  $\mathrm{WD}_{L'/L,E}$ the category of representations $(r,N,V)$ of $W_L$,\;on an $E$-vector space $V$ of finite dimension such that $r$ is unramified when restricted to the $W_{L'}$.\;Let $\mathrm{DF}_{L'/L,E}$ be the category of Deligne-Fontaine modules,\;i.e.,\;the category of quadruples $(\varphi,N,\gal(L'/L),D)$ where $D$ is an $L'_0\otimes_{\BQ_p}E$-module free of finite rank,\;which is endowed with a Frobenius $\varphi:D\rightarrow D$ (resp.,\;an $L'_0\otimes_{\BQ_p}E$-linear endomorphism $N:D\rightarrow D$) such that $N\varphi=p\varphi N$ and an action of $\gal(L'/L)$ commuting with $\varphi$ and $N$ such that $g((l\otimes e)d)=(g(l)\otimes e)d$ for $g\in \gal(L'/L),\;l\in L'0,\;e\in E,\;d\in D$.\;Then the  Fontaine's equivalence of categories (\cite[Proposition 4.1]{breuil2007first}) assert that $\mathrm{WD}_{L'/L,E}$ and $\mathrm{DF}_{L'/L,E}$ are equivalent.\;

For $(\varphi,\Gamma)$-modules over $\cR_{E,L}$ (equivalently,\;$E$-$B$-pairs),\;we can also talk about the concepts of potentially semistable (i.e.,de Rham),\;potentially crystalline,\;semistable and crystalline.\;Recall that there is an equivalence of categories between the category of potentially semistable $(\varphi,\Gamma)$-modules over $\cR_{E,L}$ (equivalently,\;$E$-$B$-pairs) which are semistable  (resp.,\;crystalline) $(\varphi,\Gamma)$-modules over $\cR_{E,L'}$ (resp.,\;crystalline $E$-$B$-pairs of $\gal_{L'}$) to the category of $E$-filtered $(\varphi,N,\gal(L'/L))$-modules ($E$-filtered $(\varphi,\gal(L'/L))$-modules) over $L$ (i.e.,\;$E$-filtered Deligne-Fontaine modules).\;The same conclusions hold for $E$-$B$-pairs.\;The last is Berger's theory \cite[Theorem A]{berger2008equations}.\;

In this paper,\;we use  the theory of  $(\varphi,\Gamma)$-modules,\;$B$-pairs and the above equivalences of categories freely.\;


\subsubsection*{Classical and $p$-adic local Langlands correspondence}
Let $m\in \BZ_{\geq 1}$ be a positive integer,\;and $\pi$ be an irreducible smooth admissible representation of $\GLN_m(L)$.\;We denote by $\rec(\pi)$ the $F$-semi-simple Weil-Deligne representation associated to $\pi$ via the normalized local Langlands correspondence (see \cite{scholze2013local}).\;Denote by $v_m$ the character $\GLN_m(L)\rightarrow E^\times$,\;$g\mapsto |\det(g)|_L=\unr(q_L^{-1})\circ\det$.\;

We normalize the reciprocity isomorphism $\rec:L^\times\rightarrow W_L^{\mathrm{ab}}$ of local class theory such that the uniformizer $\varpi_{L}$ is mapped to a geometric Frobenius morphism,\;where $W_L^{\mathrm{ab}}$ is the abelization of the Weil group $W_L\subset \gal_L$.\;Let ${\ccyc}:\gal_L\rightarrow \bZ_p^\times$ be the $p$-adic cyclotomic character (i.e.,\;the character defined by the formula $g(\epsilon_n)=\epsilon_n^{{\ccyc}(g)}$ for any $n\geq 1$ and $g\in \gal_L$).\;Then we have 
\[{\ccyc}\circ\rec=\unr(q_L^{-1})\prod_{\tau\in \Sigma_L}\tau:L^\times\rightarrow E^\times\]
by local class theory.\;We define the Hodge-Tate weights of a de Rham $p$-adic Galois representation of $\gal_L$ as the opposite of the gaps of the filtration on the covariant de Rham functor so that the Hodge-Tate weights of the cyclotomic character is $1$.\;%Recall that the Lubin-Tate character $\chi_{\mathrm{LT}}:\gal_L\rightarrow \cO_L^\times\hookrightarrow E^\times$ makes that $\chi_{\mathrm{LT}}\circ\rec:L^\times\rightarrow \cO_L^\times$ satisfies $\chi_{\mathrm{LT}}\circ\rec(\varpi_{L})=1$ and $\chi_{\mathrm{LT}}\circ\rec|_{\cO_L^\times}=\mathrm{1}_{\cO_L^\times}$.\
For a group $A$ and $a\in A$,\;denote by $\unr(\alpha)$ the unramified character of $L^\times$ sending uniformizers to $\alpha$.\;For a character $\chi$ of $\cO_L^\times$,\;denote by $\chi_{\varpi_L}$ the character of  $L^\times$ such that $\chi_{\varpi_L}|_{\cO_L^\times}=\chi$ and $\chi_{\varpi_L}(\varpi_L)=1$.\;For a character $\delta$ of $L^\times$,\;denote by $\delta_0:=\delta|_{\cO_L^\times}$.\;If $\bk:=(\bk_{\tau})_{\tau\in \Sigma_L}\in \BZ^{\Sigma_L}$,\;denote by $z^{\bk}:=\prod_{\tau\in  \Sigma_L}\tau(z)^{\bk_{\tau}}$.\;A character $\delta:L^\times \rightarrow E^\times$ is called \emph{special} if $\delta:=\unr(q_L^{-1})z^{\bk}={\ccyc}\prod_{\tau\in  \Sigma_L}\tau(z)^{\bk_{\tau}-1}$ for some $\bk:=(\bk_{\tau})_{\tau\in \Sigma_L}\in \BZ_{\geq 1}^{\Sigma_L}$ (i.e.,\;$k_{\sigma}\in \BZ_{\geq 1}$ for all $\sigma\in \Sigma_L$).\;

The $p$-adic local Langlands correspondence is often stated as follows.\;Let $L$ (resp.\;$E$) be a finite extension of $\BQ_p$,\;where $E$ is a sufficiently large coefficient field.\;The conjectural $p$-adic local Langlands correspondence  seeking for a $1-1$ correspondence between $n$-dimensional continuous representations  of $\gal_L:=\gal(\overline{L}/L)$ over $E$  and certain admissible Banach (resp.,\;or locally analytic) representation of $\GLN_n(L)$ over $E$,\;and satisfying certain local-global compatibility.\;The conjectural $p$-adic local Langlands correspondence is compatible with (and refines) the classical local Langlands correspondence.\;We recall the feature as follows.\;Let $\rho_L:\gal_L\rightarrow \GLN_n(E)$ be a potentially semistable (thus de Rham) $p$-adic Galois representation and let $\mathrm{HT}(\rho_L):=(h_{\tau,1}>h_{\tau,2}>\cdots>h_{\tau,n} )_{\tau\in \Sigma_L}$ be the Hodge-Tate weights of $\rho_L$.\;Let $L'$ be a finite Galois extension of $L$ such that $\rho_L|_{\gal_{L'}}$ is semistable.\;We put ${\bm\lambda}_\bh=(h_{\tau,i}+i-1)_{\tau\in \Sigma_L,1\leq i\leq n}$,\;which is a  dominant weight of $(\mathrm{Res}_{L/\BQ_p}\GLN_n)\times_{\BQ_p}E$ respect to $(\mathrm{Res}_{L/\BQ_p}\bB)\times_{\BQ_p}E$.\;We associate to $\rho_L$ a locally $\bQ_p$-algebraic representation $\pi_{\mathrm{alg}}(\rho_L)$ of $G$ over $E$ in the following way:
\begin{equation}\label{assoWD}
\begin{aligned}
	\overbrace{\rho_L\longleftrightarrow D_{\pst}(\rho_L)\rightsquigarrow \df(\rho_L)}^{p-\text{Hodge\;theory}}\longleftrightarrow &\overbrace{\wdre(\rho_L)\longleftrightarrow\pi_\infty(\rho_L)}^{\text{classical\;local\;Langlands}}
\end{aligned}
\end{equation}
where $D_{\pst}(\rho_L)$ is the filtered $(\varphi,N,\gal(L'/L))$-module associated to $\rho_L$ by Fontaine's theory,\;$\df(\rho_L)$ is the underlying $(\varphi,N,\gal(L'/L))$-module of $D_{\pst}(\rho_L)$ (by forgetting the Hodge filtration). Let $\wdre(\rho_L)$ be the $n$-dimensional Weli-Deligne representation over $E$ associated to $\df(\rho_L)$ via the above Fontaine's equivalence of categories \cite[Proposition 4.1]{breuil2007first}.\;Let $\pi_\infty(\rho_L)$ be the smooth representation of $G$ associated
to the F-semi-simple Weil-Deligne representation $\wdre(\rho_L)^{\mathrm{F-ss}}$ via the normalized local Langlands correspondence.\;Put $\pi_{\mathrm{alg}}(\rho_L):=\pi_\infty(\rho_L)\otimes_EL({\bm\lambda}_\bh)$,\;which is a locally $\BQ_p$-algebraic representation of $G$.\;We expect that $\pi_{\mathrm{alg}}(\rho_L)$ is a subrepresentation of the locally $\bQ_p$-analytic representation of $G$ via the hypothetical $p$-adic local Langlands correspondence.\;

%\noindent Let $A$ be an affinoid $E$-algebra.\;If $\delta_A: L^{\times} \rightarrow A^{\times}$ is locally $\BQ_p$-analytic,\;then it induces a $\BQ_p$-linear map
%\begin{equation*}
%	L \rightarrow A, \ x\mapsto \frac{d}{dt}\delta(\exp(tx))|_{t=0},
%\end{equation*}
%and hence induces an $E$-linear map
%\begin{equation*}L\otimes_{\BQ_p} E\cong \prod_{\sigma\in \Sigma_L} E \rightarrow A.
%\end{equation*}
%Thus there exist $\wt(\delta)_{\sigma}\in A$ for all $\sigma\in \Sigma_L$, called the $\sigma$-weight of $\delta$, such that the above map is given by $(a_{\sigma})_{\sigma\in \Sigma_L}\mapsto \sum_{\sigma\in \Sigma_L} a_\sigma \wt(\delta)_{\sigma}$.\;We call  $\wt(\delta):=(\wt(\delta)_{\sigma})_{\sigma\in \Sigma_L}$ the weight of $\delta$.\;
%\\\noindent For a topological commutative group $M$,\;we use $\homo(M,E)$ (resp.\;$\homo_\infty(M,E)$) to denote the $E$-space of continuous (resp.\;locally constant) additive $E$-valued characters on $M$.\;If $M$ is a totally disconnected group,\;then the terminology ``locally constant" is often replaced by smooth.\;



\subsubsection*{Characters}
For a topological commutative group $M$,\;we use $\homo(M,E)$ (resp.\;$\homo_\infty(M,E)$) to denote the $E$-space of continuous (resp.\;locally constant) additive $E$-valued characters on $M$.\;If $M$ is a totally disconnected group,\;then the terminology "locally constant" is often replaced by smooth.\;If $M$ is a locally $L$-analytic group,\;denote by $\homo_\sigma(M,E)$ the $E$-vector space of locally $\sigma$-analytic characters on $M$ (i.e.,\;the continuous characters which are locally $\sigma$-analytic $E$-valued functions on $M$).\;The local class field theory implies a bijection $\homo(L^\times,E)\cong \homo(\gal_L,E)$.\;Let $\alpha\in E^\times$,\;denote by $\unr(\alpha)$ the unramified character of $L^\times$ sending uniformizers to $\alpha$.\;





\subsubsection*{Representation theory}
\noindent Let $\mathbf{G}$ be a split connected reductive group over $L$,\;and $\bP$ be a parabolic subgroups of $\bG$.\;
Let $G,P,N_P,L_P$ be the $L$-points of $\bG,\bP,\bN_\bP,\bL_\bP$\;respectively.\;Let $\overline{P}$ be the parabolic subgroup opposite to $P$.\;

If $V$ is a continuous representation of $G$ over $E$,\;we denote by $V^{\bQ_p-\ana}$ its locally $\bQ_p$-analytic vectors.\;If $V$ is
locally $\bQ_p$-analytic representations of $G$,\;we denote by $V^{\mathrm{sm}}$ (resp.\;$V^{\mathrm{lalg}}$) the smooth (resp,\;locally $\bQ_p$-algebraic) subrepresentation of $V$ consists of its smooth (locally $\bQ_p$-algebraic) vectors (see \cite{schneider2002banach} and \cite{Emerton2007summary} for details).\;


Let $\pi_P$ be a continuous representation of $P$ over $E$ (resp.,\;locally $\bQ_p$-analytic representations of $P$ on a locally convex $E$-vector space of compact type,\;resp.,\;smooth representations of $P$ over $E$),\;we denote by
\begin{equation}\label{smoothadj2}
\begin{aligned}
	&(\mathrm{Ind}_{P}^{G}\pi_P)^{\cC^0}:=\{f:G\rightarrow \pi_P \text{\;continuous},\;f(pg)=pf(g),\forall p\in P\},\\
	&\text{resp.,\;}(\mathrm{Ind}_{P}^{G}\pi_P)^{\bQ_p-\ana}:=\{f:G\rightarrow \pi_P \text{\;locally $\bQ_p$-analytic},\;f(pg)=pf(g),\forall p\in P\},\\
	&\text{resp.,\;}i^G_P\pi_P:=(\mathrm{Ind}_{P}^{G}\pi_P)^{\infty}=\{f:G\rightarrow \pi_P \text{\;smooth},\;f(pg)=pf(g),\forall p\in P\}
\end{aligned}
\end{equation}
the continuous parabolic induction (resp.,\;the locally $\bQ_p$-analytic parabolic induction,\;resp.,\;the (un-normalized) smooth parabolic induction) of $G$.\;It becomes a continuous representation (resp.,\;locally $\bQ_p$-analytic  representation) of $G$ over $E$ (resp.,\;on a locally convex $E$-vector space of compact type, resp., smooth representations of $G$ over $E$) by endowing the left action of $G$ by right translation on functions: $(gf)(g')=f(g'g)$.\;

%Let $(\pi,V)$ be a smooth representations of $G$ over $E$,\;we denote by
%\[r^G_PV=V/\langle \pi(n)v-v:n\in N_P,\;v\in V\rangle\]
%the Jacquet module of the smooth representation $(\pi,V)$ with respect to $P$.\;Then $r^G_P(V)$ becomes a smooth $M$-representations over $E$.\;Moreover,\;we can identity $i^G_P\pi_P$ with $((\mathrm{Ind}_{P}^{G}\pi_P)^{\bQ_p-\ana})^{\mathrm{sm}}$,\;i.e.,\;the smooth subrepresentation of $(\mathrm{Ind}_{P}^{G}\pi_P)^{\bQ_p-\ana}$.\;

%Let $H$ be a closed subgroup of $G$,\;and $\pi_H$  be a smooth representation of $H$ over $E$.\;We denote by
%\begin{equation}
%	\begin{aligned}
%		c\text{-}i^G_H\pi_H:&=(\mathrm{c}\text{-}\mathrm{Ind}_{H}^{G}\pi_H)^{\infty}\\
%		&=\{f:G\rightarrow H \text{\;smooth}: \text{the support of $f$ is compact modulo $H$},\;f(pg)=pf(g),\forall p\in P\}		
%	\end{aligned}
%\end{equation}
%the smooth induction with compact support.\;It becomes a smooth representation of $G$ over $E$ by endowing the left action of $G$ by right translation on functions.\;



%
%
%Then the normalized smooth induction (resp.\;normalized smooth restriction) is $i^G_P\delta_{G,P}^{1/2}\pi_P$ (resp.\;$\delta_{G,P}^{-1/2}r^G_P\pi_P$).\;We also denote by
%\[\mathrm{c-ind}^G_P\pi_P:=\{f:G\rightarrow \pi_P \text{\;smooth with compact support},\;f(pg)=pf(g)\}\]
%the (un-normalized) compact parabolic smooth induction.\;
%Denote by $\delta_{P}$ the modulus character of $P$.\;Let $V$ (resp.\;$W$) be a smooth $G$-representation (resp.\;$L_P$-representation) over $E$,\;then we have an adjunction formula
%\begin{equation}\label{smoothadj1}
%\begin{aligned}
%\homo_G(V,i^G_PW)\xrightarrow{\sim}\homo_{L_P}(r_P^GV,W).\;
%\end{aligned}
%\end{equation}
%Moreover,\;if $V$ and $W$ are admissible,\;we have
%\begin{equation}\label{smoothadj2}
%\begin{aligned}
%\homo_G(i^G_{\overline{P}}W,V)\xrightarrow{\sim}\homo_{L_P}(W\otimes_E\delta_{P},r^G_PV).\;
%\end{aligned}
%\end{equation}

%The analytic analogy of smooth Jacquet functor is Emerton-Jacquet functor,\;defined in categories of locally $\bQ_P$-analytic representation of $G$ on locally convex topological $E$-vector space of compact type.\;






\section{Parabolic simple \texorpdfstring{$\sL$}{Lg}-invariants}\label{sectionparasimlinv}

For a potentially semistable non-crystalline $p$-adic Galois representation $\rho_L$ which admits certain special ``parabolizations" (certain parabolic analogue of the so-called triangulations),\;we are going to define parabolic Fontaine-Mazur simple $\sL$-invariants of $\rho_L$,\;which encodes certain information on the subquotients of consecutive extensions  of the associated $(\varphi,\Gamma)$-module over the Robba rings.\;

In Section \ref{BPFORPAPAIINV},\;we recall some basic facts about the Bernstein component.\;For a fixed 
Bernstein component $\Omega$,\;we talk about an $\Omega$-filtration on a $(\varphi,\Gamma)$-module.\;In Section \ref{noncrispecfildfn},\;we focus on the non-critical special $\omepik$-filtration associated to a Zelevinsky-segment.\;We then define parabolic Fontaine-Mazur simple $\sL$-invariants in Section \ref{dfnFMparalINVSECTION}.\;In Section \ref{dfndeftyoeomega} and \ref{cgsformulaSection},\;we give a way to see parabolic Fontaine-Mazur simple $\sL$-invariants in $1$-order deformations.\;

\subsection{Background and Preliminaries}\label{BPFORPAPAIINV}

Let $\rho_L:\mathrm{Gal}_L \rightarrow\mathrm{GL}_n(E)$ be a $p$-adic Galois representation.\;Let $D_{\rig}(\rho_L)$ be the $(\varphi,\Gamma)$-module over the Robba ring $\cR_{E,L}$ associated to $\rho_L$.\;We recall that the $p$-adic Galois representation $\rho_L$ is called trianguline if $D_{\rig}(\rho_L)$ is a successive extension of $(\varphi,\Gamma)$-module of rank $1$.\;A flexible version of generalizations of triangulations is given by Breuil-Ding \cite[Definition 4.1.6]{Ding2021},\;that they call $\Omega$-filtrations,\;where $\Omega$ is a Bernstein component.\;We recall the definition of $\Omega$-filtration as follows.\;

Let $m$ be a positive integer.\;Let $\Omega_0$ be a cuspidal Bernstein component of $\GLN_{m}(L)$ (see \cite{bernstein1984centre}).\;Let $\mathfrak{Z}_{\Omega_0}$ be the Bernstein center of $\Omega_0$ over $E$ ($E$ is required to be sufficiently large such that the $\mathfrak{Z}_{\Omega_0}$ can be realized over $E$,\;by \cite[Section 3.13]{PATCHING2016}.\;Let 
$\pi_{\Omega_0}\in \Omega_0$ be a fixed irreducible smooth representation of $\GLN_{m}(L)$ over $E$ of type $\Omega_0$.\;We put
\begin{equation*}
\begin{aligned}
	\eta_{E}(\Omega_0)&:=\{\eta_{\pi}:\LX\rightarrow \EX :\pi_{\Omega_0}\cong \pi_{\Omega_0}\otimes_{E}\eta_{\pi}\circ\det\},\\
	\eta_{E}^{\unra}(\Omega_0)&:=\{\eta_{\pi}:\LX\rightarrow \EX \text{\;unramified}:\pi_{\Omega_0}\cong \pi_{\Omega_0}\otimes_{E}\eta_{\pi}\circ\det\}\\
	&\;=\{a\in\EX:\pi_{\Omega_0}\cong \pi_{\Omega_0}\otimes_{E}\unra(a)\circ\det\}.
\end{aligned}
\end{equation*}
We have $\eta_{E}^{\unra}(\Omega_0)\subset \eta_{E}(\Omega_0)$.\;Both $\eta_{E}^{\unra}(\Omega_0)\subset \eta_{E}(\Omega_0)$ and $\eta_{E}(\Omega_0)$ are finite groups and independent of the choice of $\pi_{\Omega_0}$ in $\Omega_0$.\;By comparing the central characters,\;there exists an integer $m_0|m$ such that $ \eta_{E}^{\unra}(\Omega_0)\cong \mu_{m_0}$.\;We further assume that $\mu_{m_0}\subset E$ where $\mu_{m_0}$ means the group of $m_0$-th roots of unity in $\obp$.\;By \cite[Pages\;242-243,\;Lemma\;3.24]{PATCHING2016} and \cite[Section 2.2]{Ding2021},\;we have 
\[\mathfrak{Z}_{\Omega_0}\cong E[z,z^{-1}]^{\mu_{m_0}}\cong E[z^{m_0},z^{-m_0}].\]
For any closed point $x\in \Spec(\FZ_{\Omega_0})$,\;we can associate it a smooth irreducible cuspidal representation $\pi_{x}$ of $\GLN_{m}(L)$ over $k(x)$.\;By normalized classical local Langlands correspondence (see \cite{scholze2013local}),\;we get an $F$-semi-simple Weil-Deligne representation $\wdre_{x}:=\rec(\pi_{x})$ of $W_L$ over $k(x)$.\;By the equivalence of categories (\cite[Proposition 4.1]{breuil2007first}),\;the Weil representation $\wdre_{x}$ corresponds to a Deligne-Fontaine module $\df_x$,\;which by Berger's theory \cite{berger2008equations} corresponds  to a $p$-adic differential equation ${\Delta_x}$,\;i.e.,\;a $(\phi,\Gamma)$-module of rank $m$ over $\cR_{E,L}$ which is de Rham of constant Hodge-Tate weight $0$ such that $D_{\pst}(\Delta_x)$ (forgetting the Hodge filtration) is isomorphic to the $\df_x$.\;Let $L'/L$ be a finite extension of $L$.\;Moreover,\;if $\wdre_{x}$ is unramified when restricted to $W_{L'}$,\;then $\df_x$ is a $(\varphi,N,\gal(L'/L))$-module.\;


Let $\ul{n}:=(n_1,\cdots,n_s)$ be an ordered partition  of $n$ for some positive integer $s$.\;We put \[\bL_{\ul{n}}:=\left(\begin{array}{cccc}
\GLN_{n_1} & 0 & \cdots & 0 \\
0 & \GLN_{n_2} & \cdots & 0 \\
\vdots & \vdots & \ddots & 0 \\
0 & 0 & 0 & \GLN_{n_s} \\
\end{array}\right)\subseteq \bP_{\ul{n}}:=\left(\begin{array}{cccc}
\GLN_{n_1} & \ast & \cdots & \ast \\
0 & \GLN_{n_r} & \cdots & \ast \\
\vdots & \vdots & \ddots & \ast \\
0 & 0 & \cdots & \GLN_{n_s} \\
\end{array}\right).\;\]
Note that $\bP_{\ul{n}}$ is a parabolic subgroup of $\GLN_{n}$ and  $\bL_{\ul{n}}$ is a Levi subgroup of $\bP_{\ul{n}}$.\;For each $1\leq i\leq s$,\;let $\Omega_i$ be a cuspidal  Bernstein component of $\GLN_{n_i}(L)$ and $\FZ_{\Omega_i}$ be the corresponding Bernstein centre over $E$.\;Let $\Omega:=\prod_{i=1}^s\Omega_i$ be the cuspidal Bernstein component of $\bL_{\ul{n}}(L)$.\;Then we have an isomorphism of commutative $E$-algebras $\FZ_{\Omega}\cong\otimes_{i=1}^s\FZ_{\Omega_i}$.\;For any closed point $\bx=(x_i)_{1\leq i\leq s}\in \Spec(\FZ_{\Omega})$,\;we can associate it irreducible smooth  cuspidal representations $\{\pi_{x_i}\}_{1\leq i\leq s}$ of $\{\GLN_{n_i}(L)\}_{1\leq i\leq s}$ over $k(x)$, respectively, the $F$-semi-simple Weil-Deligne representations $(\wdre_{x_i})_{1\leq i\leq s}$,\;and the $p$-adic differential equations $(\Delta_{x_i})_{1\leq i\leq s}$.\;This gives an irreducible smooth  cuspidal representation $\pi_{\bx}=\otimes_{i=1}^k\pi_{x_i}$ of  $\bL_{\ul{n}}(L)$ and an $n$-dimensional $F$-semi-simple Weil-Deligne representation $\wdre_{\bx}:=\oplus_{i=1}^s\wdre_{x_i}$ of $W_L$ over $k(x)$.\;

%Then $\omega=\otimes_{i=1}^s\pi_i$.\;For $1\leq i \leq s$,\;let $\Omega_i$ be the cuspidal Bernstein component  $\GLN_{n_i}(L)$ containing $\pi_i$.\;From the theory of types,\;each $\Omega_i$ contains a maximal simple type $(J_i,\sigma_{i}^0)$.\;We put $\sigma_{i}=\bigg(\ind^{\GLN_{n_i}(\cO_L)}_{J_i}\sigma_{i}^0\bigg)^{\infty}$,\;and $\sigma=\otimes_{i=1}^s\sigma_{i}$,\;this is an irreducible representation of $\bL_{\Omega}^0:=\bL_{\Omega}(\cO_L)$.\;Let $\FZ_{\Omega_i}$ (resp.,\;$\FZ_{\Omega_{\bL_{\Omega}}}$) be the  Bernstein centre of $\Omega_i$ (resp.,\;$\Omega_{\bL_{\Omega}}$)  over $E$.\;Then we have
%\[\mathfrak{Z}_{\Omega}\cong \EndO_{\bL_{\Omega}(L)}(	c\text{-}i^{\bL_{\Omega}(L)}_{\bL_{\Omega}^0}\sigma)\cong\otimes_{i=1}^s\mathfrak{Z}_{\Omega_i},\;\mathfrak{Z}_{\Omega_i}\cong\EndO_{\bL_{\Omega}(L)}(	c\text{-}i_{\GLN_{n_i}(\cO_L)}^{\GLN_{n_i}(L)}\sigma_i).\;\]

We now recall an $\Omega$-filtration on a $(\varphi,\Gamma)$-module (see \cite[Section 4.1.2]{Ding2021}).\;Let $Z_{\bL_{\ul{n}}}$ be the center of $\bL_{\ul{n}}$.\;We have a map $\BZ\hookrightarrow L^\times$,\;$i\mapsto \varpi_L$,\;which induces a map $\oplus_{i=1}^s\BZ\hookrightarrow Z_{\bL_{\ul{n}}}(L)$.\;We denote by $Z_{\varpi_L}$ its image,\;then we have $Z_{\bL_{\ul{n}}}(L)\cong Z_{\varpi_L}\times Z_{\bL_{\ul{n}}}(\cO_L)$.\;Denote by $\rigerch=\widehat{Z_{\bL_{\ul{n}}}(\cO_L)}$ (resp.,\;$\rigerchl=\widehat{Z_{\bL_{\ul{n}}}(L)}$,\;resp.,\;$\widehat{Z_{\varpi_L}}$) the rigid analytic space over $E$ which  parameters continuous characters of $Z_{\bL_{\ul{n}}}(\cO_L)$ (resp.,\;$Z_{\bL_{\ul{n}}}(L)$,\;resp.,\;$Z_{\varpi_L}$). Then we have $\rigerchl\cong \widehat{Z_{\varpi_L}}\times \rigerch$.\;
%We put \[\mathfrak{L}_{\Omega}:=\sban\times \rigerch.\]We denoted by $\fz_{\bL_{\ul{n}}}$,\;$\fl_{\ul{n}}^D$ the Lie algebras  of $Z_{\bL_{\ul{n}}}$,\;$\bL_{\ul{n}}^D$ over $E$,\;respectively.\;,\;and  $\bL_{\ul{n}}^D$ be the derived subgroup of $\bL_{\ul{n}}$.\;

%We also fix an isomorphism (i.e.,\;fix a choice of $\varpi_L$) $\BG_m^{\rig}\times\widehat{\cO_L^\times}\xrightarrow{\sim}\widehat{L^\times}$,\;$(a,\chi)\mapsto \unr(a)\chi_{\varpi_L}$.\;For each $1\leq i\leq s$,\;the smooth representations $\wdre_{x_i}$ induces a natural finite morphism $\iota_{\wdre_{x_i}}':\BG_m^\rig\rightarrow \Spec\FZ_{\Omega_i}$,\;$\alpha\mapsto \wdre_{x_i}\te\unr(\alpha)\circ\det$.\;
%Define by $\iota_{\wdre_{x_i}}$ the composition
%\[\iota_{\wdre_{x_i}}:\widehat{L^\times}\xrightarrow{\sim} \BG_m^{\rig}\times\widehat{\cO_L^\times} \xrightarrow{(\iota_{\wdre_{x_i}}',\mathrm{id})} \Spec\FZ_{\Omega_i}\times\widehat{\cO_L^\times}.\;\]
%We define the following composition:
%\begin{equation*}
%\begin{aligned}
%:\rigerchl\xrightarrow{\sim}\prod_{i=1}^s\widehat{L^\times}\xrightarrow{\sim}
%\prod_{i=1}^s(\BG_m^{\rig}\times\widehat{\cO_L^\times})\xrightarrow{(\iota_{\wdre_{x_i}})_{1\leq i\leq s}}\prod_{i=1}^s(\Spec\FZ_{\Omega_i})^\rig\times\widehat{\cO_L^\times}=\mathfrak{L}_{\Omega}.
%\end{aligned}
%\end{equation*}



%The natural map $Z_{\underline{\omega}}\subseteq Z_{\bL_{\Omega}}(L)\rightarrow  \mathfrak{Z}_{\Omega}$ gives a natural finite morphism $E[Z_{\underline{\varpi}_L}]\rightarrow \FZ_{\Omega}$ (observe that $\FZ_{\Omega_i}\cong E[z,z^{-1}]^{\mu_{\Omega_i}^{\unr}}$),\;which induce
%a finite morphism $\Spec \FZ_{\Omega}\rightarrow \Spec E[Z_{\underline{\varpi}_L}]$ send $x\in \Spec \bZ_{\Omega}$ to the character $z\mapsto \omega_{\pi_x}(z)$,\;where $\omega_{\pi_x}$ is the character of $\pi_x$.\;This induces an finite morphism
%\[\sban\times \rigerch\longrightarrow (\Spec E[Z_{\underline{\varpi}_L}])^{\ana}\times \rigerch\cong \prod_{i=1}^r\BG_m^\ana\times\widehat{\cO_L^\times}\cong \widehat{Z_{\underline{\omega}}}\times \rigerch\cong \rigerchl.\]

As in \cite[Section 4.1.2,\;Definition 4.1.6]{Ding2021},\;we consider:
\begin{dfn}\label{dfnomegafil} Let $D$ be a fixed $(\varphi,\Gamma)$-module of rank $n$ over $\cR_{E,L}$.\;We say that $D$ admits an \textbf{$\Omega$-filtration} $\cF$ if $D$ admits an increasing filtration by saturated $(\varphi,\Gamma)$-submodules
\[\cF=\fil_{\bullet}^{\cF} D: \ 0 =\fil_0^{\cF} D \subsetneq \fil_1^{\cF} D \subsetneq \cdots \subsetneq \fil_{k}^{\cF} D=D,\]
such that for $1\leq i\leq s$,\;
\begin{itemize}
	\item $\gr^{\cF}_{i}D$ is a $(\varphi,\Gamma)$-module of rank $n_i$;\;
	\item there exist an $E$-point $x_i\in \left(\Spec\mathfrak{Z}_{\Omega_i}\right)^{\mathrm{an}}$,\;and a continuous character $\delta_i:L^\times\rightarrow  E^\times$ such that
	$\gr^{\cF}_{i}D\tee \cR_{E,L}(\delta_{i}^{-1})\hookrightarrow \Delta_{x_i}$.\;
\end{itemize}
Put $\underline{x}=(x_i)$ and $\delta=\otimes_{i=1}^s\delta_i$.\;
\begin{itemize}
	\item We say $(\underline{x},\delta)\in \sban\times\rigerchl $ is a parameter of $\Omega$-filtration $\cF$ if $0$ is the minimal $\tau$-Hodge-Tate weight of $\gr^{\cF}_{i}D\tee \cR_{E,L}(\delta_{i}^{-1})$.\;
	\item We call $(\underline{x},\delta^0=\otimes_{i=1}^s\delta^0_i)\in \sban\times\rigerch $ is a parameter of the $\Omega$-filtration $\cF$ of $D$ if $(\underline{x},(\delta^0)_{\varpi_L})=\otimes_{i=1}^s(\delta^0_i)_{\varpi_L})$ is a parameter of $\Omega$-filtration $\cF$ in $\sban\times\rigerchl $.\;
\end{itemize}

%\begin{itemize}
%	\item $\gr^{\cF}_{i}D$ is a $(\varphi,\Gamma)$-module of rank $n_i$,\;and $\gr^{\cF}_{i}D\tee \cR_{E,L}((\delta_{i}^0)_{\varpi_L}^{-1})\hookrightarrow \Delta_{x_i}$,\;
%	\item for all $\tau\in \Sigma_L$,\;$0$ is the minimal $\tau$-Hodge-Tate weight of $\gr^{\cF}_{i}D\tee \cR_{E,L}((\delta_{i}^0)_{\varpi_L}^{-1}).$
%\end{itemize}
\end{dfn}
Denote by $\iota_{\Omega}$ the following morphism
\begin{equation}\label{notationbasisrigidmorphism}
\begin{aligned}
	\iota_{\Omega}:&\sban\times\rigerchl\rightarrow \sban\times\rigerch,\;
	&(\underline{x},\delta)\mapsto (\iota_{\Omega}(\underline{x}),\delta|_{\mathcal{O}_L^\times}),
\end{aligned}
\end{equation}
where $\iota_{\Omega}:\sban\xrightarrow{\sim}\sban$ is the isomorphism such that
$\pi_{\iota_{\Omega}(\underline{x})_i}=\pi_{x_i}\te \unr(\delta_i(\varpi_L))$. It is clear that $(\underline{x},\delta)\in\sban\times\rigerchl $ is a  parameter of $\Omega$-filtration if and only if $(\iota_{\Omega}(\underline{x}),\delta|_{\mathcal{O}_L^\times})$\\$\in\sban\times\rigerch$ is a parameter of $\Omega$-filtration.\;
%By \cite[Lemma 4.1.9]{Ding2021},\;the parameters of $\cF$ in $\sban\times\rigerchl $ or $\sban\times\rigerch$ are unique up to the twist of $\mu_{\Omega}$.\;

\subsection{\texorpdfstring{$(\varphi,\Gamma)$-module  over $\cR_{E,L}$ with non-critical special $\omepik$-filtration}{Lg}}\label{noncrispecfildfn}

In the sequel,\;we fix a cuspidal Bernstein component $\Omega_r$ of $\GLN_r(L)$ and an irreducible smooth cuspidal representation $\pi_{0}\in \Omega_r$ over $E$ of type $ \Omega_r$.\;We put
\[\omepik=\prod_{i=1}^k\Omega_{r,i},\;\Omega_{r,i}=\Omega_{r}\]
which is a cuspidal Bernstein component of $\bL^{\lrr}(L)$.\;Let $\FZ_{\omepik}\cong\otimes_{i=1}^k\FZ_{\Omega_{r,i}}=\FZ_{\Omega_{r}}^{\otimes k}$ be the associated 
Bernstein center over $E$.\;

Let $\pi\in \Omega_r$ be an (absolutely) irreducible smooth cuspidal representation over $E$ of type $\Omega_r$,\;which corresponds an $E$-point $x_\pi$ of $\Spec(\FZ_{\Omega_{r}})$.\;There exists a unique $\alpha_\pi \in E^\times$ such that $\pi\cong \pi_0\otimes_E\unr(\alpha_\pi)$. Via the classical local Langlands correspondence (normalized as in \cite{scholze2013local}),\;we get an $r$-dimensional absolutely irreducible $F$-semi-simple Weil-Deligne representation $\wdre_\pi$ of $W_L$ over $E$.\;Let $\df_\pi$ (resp.,\; ${\Delta_\pi}$) be the Deligne-Fontaine module (resp.,\; $p$-adic differential equation ${\Delta_\pi}$ over $\cR_{E,L}$) associated to $\wdre_\pi$ (equivalently,\;to $\pi$).\;Assume that $\wdre_\pi$ is unramified when restricted to $W_{L'}$ for some finite Galois extension $L'$ of $L$.\;Denote by $\varphi_\pi:\df_\pi\rightarrow \df_\pi$ the Frobenius semi-linear operator  on $\df_\pi$.\;Note that $\df_\pi=(\varphi_\pi,N=0,\gal(L'/L),\df_\pi)$ is an (absolutely) irreducible Deligne-Fontaine module.\;

%and let  $\Delta_{[k-1,0]}(\pi)=[\pi\otimes_Ev_r^{k-1},\cdots,\pi]$ be a Zelevinsky-segment.\;The irreducible smooth cuspidal representation $\otimes_{i=1}^{k}(\pi\otimes_E v_r^{k-i})$ lies in the cuspidal Bernstein component $\omepik$.\;



Next,\;we consider a special case of $\omepik$-filtration,\;that we call non-critical special $\omepik$-filtration.\;

Let $\Dpik$ be a potentially semistable $(\varphi,\Gamma)$-module over $\cR_{E,L}$ of rank $n$.\;Let $L'$ be a finite Galois extension of $L$ such that $\Dpik|_{L'}$ is a semistable $(\varphi,\Gamma)$-module over $\cR_{E,L'}$ of rank $n$.\;We consider the Deligne-Fontaine module associated to $\Dpik$:
$$\df(\Dpik)=(\varphi,N,\gal(L'/L),D_{\pst}(\Dpik))$$
where $D_{\pst}(\Dpik)=D_{\mathrm{st}}^{L'}(\Dpik\otimes_{\cR_{E,L}}\cR_{E,L'})$ is a finite free $L'_0\otimes_{\BQ_p}E$-module of rank $n$,\;where the $(\varphi,N)$-action on $D_{\pst}(\Dpik)$ is induced from the $(\varphi,N)$-action on $B_{\mathrm{st}}$,\;and where the 
$\gal(L'/L)$-action on  $D_{\pst}(\Dpik)$ is the residual action of $\gal_L$.\;Let  $\mathrm{WD}_{L'/L,E}$ the category of representations $(r_L,N,V)$ of $W_L$,\;on an $E$-vector space $V$ of finite dimension such that $r_L$ is unramified when restricted to the $W_{L'}$.\;By Fontaine's equivalence of categories as in \cite[Proposition 4.1]{breuil2007first},\;we can associate  to $\df(\Dpik)$ an $n$-dimensional Weli-Deligne representation $\wdre(\Dpik)\in \mathrm{WD}_{L'/L,E}$ (not necessarily $F$-semi-simple) of $W_L$ over $E$.\;

We say that $\wdre(\Dpik)$  admits an increasing special $\omepik$-filtration $\cF$ if $\wdre(\Dpik)$ admits an increasing filtration $\cF$ by Weil-Deligne subrepresentations:
\[\cF=\fil_{\bullet}^{\cF} \wdre(\Dpik): \ 0 =\fil_0^{\cF} \wdre(\Dpik) \subsetneq \fil_1^{\cF} \wdre(\Dpik) \subsetneq \cdots \subsetneq \fil_{k}^{\cF} \wdre(\Dpik)=\wdre(\Dpik),\]
and there is an irreducible smooth cuspidal representation $\pi\in \Omega_r$ over $E$ of type $ \Omega_r$ such that:
\begin{itemize}
\item For all $1\leqslant i\leqslant k$,\;we have $\gr^{\cF}_i\wdre(\Dpik)\cong \wdre_\pi\otimes_E|\rec^{-1}|^{k-i}$;
\item The monodromy operator $N$ sends $\gr^{\cF}_i\wdre(\Dpik)$ to
$\gr^{\cF}_{i-1}\wdre(\Dpik)$ via the zero or identity map on $\wdre_\pi$ for $2\leq i\leq k$,\;and sends $\gr^{\cF}_1\wdre(\Dpik)$ to zero.\;
\end{itemize}
Note that the formulation on the second assumption originated from the fact that any  $F$-semi-simple Weil-Deligne representation is isomorphic to a direct sum of (absolutely) indecomposable objects.\;In particular,\;the second  assumption implies that the $F$-semisimplification $\wdre(\Dpik)^{F\mathrm{-ss}}$ of $\wdre(\Dpik)$ is an absolutely indecomposable (but not irreducible) $F$-semi-simple Weil-Deligne representation if and only if the monodromy operator $N$ has full monodromy rank (i.e.,\;$N^{k-1}\neq 0$ and $N^k=0$).\;

By \cite[Proposition 4.1]{breuil2007first},\;the special $\omepik$-filtration $\cF$ on $\wdre(\Dpik)$ corresponds to an increasing  special  $\omepik$-filtration on $\df(\Dpik)$ (still denoted by $\cF$)  by Deligne-Fontaine submodules
\[\cF=\fil_{\bullet}^{\cF}\df(\Dpik): \ 0 =\fil_0^{\cF} \df(\Dpik) \subsetneq \fil_1^{\cF} \df(\Dpik) \subsetneq \cdots \subsetneq \fil_{k}^{\cF} \df(\Dpik)=\df(\Dpik),\]
such that $\fil_i^{\cF} \df(\Dpik)$ is associated to $\fil_i^{\cF} \wdre(\Dpik)$ via 
\cite[Proposition 4.1]{breuil2007first}.\;We then  see that \[\gr^{\cF}_i\df(\Dpik)\cong (\varphi_{\pi,i},N_{\gr^{\cF}_i\df(\Dpik)}=0,\gal(L'/L),\df_{\pi,i})\] for $1\leqslant i\leqslant k$,\;where $\df_{\pi,i}$ is isomorphic to $\df_\pi$ as a module,\;endowed with a Frobenius morphism $\varphi_{\pi,i}=p^{i-k}\varphi_\pi$ (i.e.,\;a twist of $\varphi_\pi$ by $p^{i-k}$).\;The monodromy operator $N$ on Deligne-Fontaine module $\df(\Dpik)$ is zero on  $(p^{1-k}\varphi_\pi,N=0,\gal(L'/L),\df_\pi)$,\;and sending $(p^{i-k}\varphi_\pi,N=0,\gal(L'/L),\df_\pi)$ to
$(p^{(i-1)-k}\varphi_\pi,N=0,\gal(L'/L),\df_\pi)$ via the zero or identity map on $\df_\pi$ for  $2\leq i\leq k$.\;


%As a $\phi$-module,\;we see that $\fil_j^{\cF} \df(\Dpik)=\bigoplus_{i=1}^j(p^{i-k}\varphi_\pi,N=0,\gal(L'/L),\df_\pi)$.\;The has full monodromy rank,\;


Let  ${\Delta_\Dpik}$  be the $p$-adic differential equation over $\cR_{E,L}$ associated to the Deligne-Fontaine module $\df(\Dpik)$.\;The special $\omepik$-filtration on $\df(\Dpik)$ induces a special  $\omepik$-filtration $\fil_{\bullet}^{\cF}{\Delta_\Dpik}=\{\fil_{i}^{\cF}{\Delta_\Dpik}\}$ on ${\Delta_\Dpik}$ by saturated $(\varphi,\Gamma)$-submodules over $\cR_{E,L}$,\;such that $\fil_{i}^{\cF}{\Delta_\Dpik}$ is the $p$-adic differential equation over $\cR_{E,L}$ associated to $\fil_i^{\cF} \df(\Dpik)$.\;In particular,\;we see that $\gr^{\cF}_i{\Delta_\Dpik}\cong {\Delta_\pi}\tee \cR_{E,L}(\unr(q_L^{i-k}))$ for $1\leqslant i\leqslant k$.\;Consider
\[\cM_\Dpik=\Dpik\big[\frac{1}{t}\big]\cong \Delta_\Dpik \big[\frac{1}{t}\big]\]
By inverting $t$,\;the filtration $\cF$ on $\Delta_\Dpik$ induces an increasing filtration $\cF:= \fil_{i}^{\cF}\cM_\Dpik:=\fil_{i}^{\cF}{\Delta_\Dpik}\big[\frac{1}{t}\big]$ on $\cM_\Dpik$ by $(\varphi,\Gamma)$-submodules over $\cR_{E,L}\big[\frac{1}{t}\big]$.\;Therefore,\;this filtration $\cF$ on $\cM_\Dpik=D\big[\frac{1}{t}\big]$ induces a filtration on $\Dpik$:
\[\cF=\fil_{\bullet}^{\cF}\Dpik: \ 0 =\fil_0^{\cF}\Dpik \subsetneq \fil_1^{\cF}\Dpik \subsetneq \cdots \subsetneq \fil_{k}^{\cF} \Dpik=\Dpik,\;\fil_i^{\cF}\Dpik=(\fil_{i}^{\cF}\cM_\Dpik)\cap \Dpik,\]
by saturated $(\varphi,\Gamma)$-submodules of $\Dpik$ over $\cR_{E,L}$.\;

Since $\Dpik$ is potentially semistable,\;it is de Rham.\;Hence we have \[D_{\dr}(\Dpik)\cong (D_{\pst}(\Dpik)\otimes_{L_0'}L')^{\gal(L'/L)},\;\] which  is a free $L\otimes_{\BQ_p}E$-module of rank $n$.\;The module $D_{\dr}(\Dpik)$ is equipped with a natural Hodge filtration.\;We assume that $D_{\dr}(\Dpik)$  has distinct Hodge-Tate weights $\bh:=(\hpi_{\tau,1}>\hpi_{\tau,2}>\cdots>\hpi_{\tau,n} )_{\tau\in \Sigma_L}$. Denote by $\hpi_{i}=(\hpi_{\tau,i})_{\tau\in \Sigma_L}$ for $1\leq i\leq n$.\;

%Hence,\;for each $\tau\in \Sigma_L$,\;the natural Hodge filtration can be expressed by the following complete flag:
%\[\fil_{\bullet}^{H}D_{\dr}(\Dpik)_\tau: \ 0 \subsetneq \fil_{-\hpi_{\tau,n}}^{H} D_{\dr}(\Dpik)_\tau \subsetneq \fil_{-\hpi_{\tau,n-1}}^{\cF}D_{\dr}(\Dpik)_\tau \subsetneq \cdots \subsetneq \fil_{{-\hpi_{\tau,1}}}^{H}D_{\dr}(\Dpik)_\tau=D_{\dr}(\Dpik)_\tau.\]
%For each $\tau\in \Sigma_L$,\;we now fix a a basis of $D_{\dr}(\Dpik)_\tau$.\;Then the Hodge filtration $\fil_{\bullet}^{H}$  corresponds to an $E$-point $(g_{\tau}\bB(E))_{\tau\in \Sigma_L}\in \bG_{/E}/\bB_{/E}$.\;
Now by Berger's equivalence of categories,\;we see that $\fil_i^{\cF}\Dpik$ corresponds to the filtered Delingen-Fontaine module $\fil_i^{\cF} \df(\Dpik)$ equipped with the induced filtration from the Hodge filtration on $D_{\pst}(\Dpik)$. We say that $\cF$ is \textbf{non-critical} if the Hodge-Tate weights of $\fil_i^{\cF}\Dpik$ are given by $$\{\hpi_{\tau,1},\hpi_{\tau,2},\cdots,\hpi_{\tau,ir}\}_{\tau\in \Sigma_L}.\;$$
This implies that the Hodge-Tate weights of $\gr_i^{\cF} \Dpik$ are $$\{\hpi_{\tau,(i-1)r+1},\hpi_{\tau,(i-1)r+2},\cdots,\hpi_{\tau,ir}\}_{\tau\in \Sigma_L}.\;$$
In this case,\;using Berger's equivalence of categories \cite[Theorem  A]{berger2008equations} and comparing the weight (or see \cite[(2.4)]{Ding2021}),\;we have an injection of $(\varphi,\Gamma)$-modules over $\cR_{E,L}$

\begin{align}\label{Dpikinjection}
	\mathbf{I}_{i}:	\gr_{i}^{\cF} \Dpik \hookrightarrow&\;{\Delta_\pi}\tee \cR_{E,L}(\unr(q_L^{i-k}))\tee \cR_{E,L}(z^{\hpi_{ir}})\\=&\;\gr^{\cF}_i{\Delta_\Dpik}\tee \cR_{E,L}(z^{\hpi_{ir}}),\nonumber
\end{align}
for $i=1,\cdots,k$.\;Note that this injection (\ref{Dpikinjection}) is equivalent to the existence of an isomorphism of $(\varphi,\Gamma)$-modules over $\cR_{E,L}[1/t]$:
$\gr_{i}^{\cF}[1/t]\cong\gr^{\cF}_i{\Delta_\Dpik}\tee \cR_{E,L}(z^{\hpi_{ir}})[1/t]$ or an injection (by comparing Hodge-Tate weights) $$\gr^{\cF}_i{\Delta_\Dpik}  \hookrightarrow \gr_{i}^{\cF}\Dpik\tee \cR_{E,L}(z^{\hpi_{(i-1)r+1}}).\;$$ 

This implies that the $(\varphi,\Gamma)$-module $\Dpik$ admits a non-critical $\omepik$-filtration $\cF$.\;The parameters of $\cF$ in $\sbanpik\times\rigchl $ or $\sbanpik\times\rigch$ are given as follows.\;Recall that $\pi\cong \pi_0\otimes\unr(\alpha_\pi )$ for some $\alpha_\pi \in E^\times$.\;

\begin{dfn}\label{weaklynoncritical}\textbf{(Non-critical special weakly $\omepik$-filtration)} Keep the above assumption.\;We say that the $(\varphi,\Gamma)$-module $\Dpik$ admits a non-critical special weakly $\omepik$-filtration $\cF$ with parameter
\begin{equation}
	\begin{aligned}
		&(\bx_{\pi},\bmdel)\in \sbanpik\times\rigchl,\;\hspace{110pt}\\
		&\bx_{\pi}=(\bx_{\pi,i}\cong x_{\pi})_{1\leq i\leq k},\;\\
		&\bmdel:=(\bm{\delta}_{\bh,i}=\unr(  q_L^{i-k}){z^{\hpi_{ir}}})_{1\leq i\leq k},\;
	\end{aligned}
\end{equation}
or with parameter
\begin{equation}
	\begin{aligned}
		&(\widetilde{\bx}_{\pi,\bh},\widetilde{\bm{\delta}}_\bh)\in \sbanpik\times\rigch,\;\hspace{100pt}\\
		&\widetilde{\bx}_{\pi,\bh}=(\widetilde\bx_{\pi,i})_{1\leq i\leq k},\;\pi_{\widetilde\bx_{\pi,\bh,i}}\cong \pi_0\te\unr(\alpha_\pi  q_L^{i-k}z^{\hpi_{ir}}(\varpi_L)),\;\\
		&\widetilde{\bm{\delta}}_\bh=(\widetilde{\bm{\delta}}_{\bh,i}={z^{\hpi_{ir}}}|_{\co_L^\times})_{1\leq i\leq k}.
	\end{aligned}
\end{equation}
\end{dfn}
\begin{rmk}For convenience,\;we may use these two kinds of parameters depending on the situation.\;By \cite[Lemma 4.1.9]{Ding2021},\;the parameters of $\cF$ in $\sbanpik\times\rigchl$ or $\sbanpik\times$\\$\rigch$ are in general not unique.\;All the parameters of $\cF$ in $\sbanpik\times\rigch$ are of form $(\underline{x}',\delta')$ such that,\;for $1\leq i\leq k$,\;we have $\mathbf{WD}_{x_i}\cong \mathbf{WD}_{\bx_{\pi,i}}\otimes_E\unr(\alpha_i)$ and $\delta'_i=\bm{\delta}_{\bh,i}\unr(\alpha_i^{-1})\eta_i$ for some $\alpha_i\in \overline{E}^\times$ and $\eta_i\in \mu_{\Omega_{r,i}}=\mu_{\Omega_{r}}$.\;All the parameters of $\cF$ in $\sbanpik\times\rigch$ are of the form $(\underline{x}',\delta'')$ such that,\;for $1\leq i\leq k$,\;we have $\mathbf{WD}_{x_i}\cong \mathbf{WD}_{\widetilde\bx_{\pi,i}}\otimes_E\unr(\eta_i(\varpi_{L}))$ and $\delta''_i={z^{\hpi_{ir}}}\eta_i|_{\co_L^\times}$ for some $\eta_i\in \mu_{\Omega_{r}}$.\; 
\end{rmk}

\begin{dfn}\label{dfnnoncriticalspecial}\textbf{(Non-critical special $\omepik$-filtration)} We call an $\omepik$-filtration $\cF$ on $\Dpik$ is non-critical special with parameter $$(\bx_{\pi},\bmdel)\in \sbanpik\times\rigchl ,\;$$(resp.,\;with parameter $(\widetilde{\bx}_{\pi,\bh},\widetilde{\bm{\delta}}_\bh)\in \sbanpik\times\rigch $)\;if\;$\Dpik$ admits a  non-critical special weakly $\omepik$-filtration $\cF$ (see Definition \ref{weaklynoncritical}) with parameter $(\bx_{\pi},\bmdel)\in \sbanpik\times\rigchl $ (resp.,\;with parameter $(\widetilde{\bx}_{\pi,\bh},\widetilde{\bm{\delta}}_\bh)\in \sbanpik\times\rigch $ ),\;and the subquotient $\Dpik^{i+1}_{i}$ is non-split for all $ir\in\Delta_n(k)$.\; %(equivalently,\;$({\Delta_{\Dpik}})^{i+1}_{i}$ is potentially semistable and potentially noncrystalline for all $s\in\Delta_n(k)$,\;see the Lemma \ref{criofnoncrystalline}).\;
\end{dfn}

Throughout section \ref{sectionparasimlinv},\;we fix such a $(\varphi,\Gamma)$-module $\Dpik$ (in Definition  \ref{weaklynoncritical} or Definition \ref{dfnnoncriticalspecial}).\;For $1\leq i\leq j\leq k$,\;we put $\Dpik^{j}_i:=\fil_j^{\cF}\Dpik/\fil_{i-1}^{\cF}\Dpik$. We also use $\cF$ to denote the induced non-critical special $\Omega_{[i,j]}$-filtration on $\Dpik^{j}_i$,\;where $\Omega_{[i,j]}:=\prod_{l=i}^j\Omega_{r,l}$.\;Clearly,\;a parameter of this non-critical special  $\Omega_{[i,j]}$-filtration is $((\bx_{0,s})_{i\leq s\leq j},(\bm{\delta}_{\bh,s})_{i\leq s\leq j})$.\;




\subsection{Deformation of type \texorpdfstring{$\omepik$}{Lg}}\label{defdfnomepik}\label{dfndeftyoeomega}

In this section,\;we study certain paraboline deformations of $(\varphi,\Gamma)$-modules which admit an $\omepik$-filtration, under the framework of \cite[Section 4.1]{Ding2021}.\;


%\begin{description}
%  \item [(1)] $\gr^{\cF}_{i}D$ is a $(\varphi,\Gamma)$-module of rank $r$,\;
%  \item [(2)]There exist an $E$-point $x_i\in \sban$,\;and a continuous character $\undel$ of $\zlch$ such that
%$D_{x_i}:=\gr^{\cF}_{i}D\tee \cR_{E,L}((\delta_{i}^0)_{\varpi_L}^{-1})\hookrightarrow \sD_{x_i}$,\;and  $0$ is the minimal $\tau$-Hodge-Tate weight of $\gr^{\cF}_{i}D\tee \cR_{E,L}((\delta_{i}^0)_{\varpi_L}^{-1})$,
%  \item [(3)]$\pi_{x_{i+1}}=\pi_{x_i}\otimes_E\eta_i$ for some smooth character $\eta_i:L^\times\rightarrow E$ such that
%             $\eta_i^{-1}(\delta_{i}^0)_{\varpi_L}(\delta_{i+1}^0)_{\varpi_L}^{-1}$ is special,\;
%  \item [(4)]$D^{i+1}_i:=\fil_{i+1}^{\cF}/\fil_{i-1}^{\cF}$ is nonsplit,\;
%  \item [(5)] $w_{\cF,\tau}=1$ for all $\tau\in \Sigma_L$.\;
%\end{description}
We first recall the contents of \cite[Section 4.1.1]{Ding2021}.\;Let $\Delta$ be an irreducible $(\varphi,\Gamma)$-module of rank $r$ over $\cR_{E,L}$,\;de Rham of constant Hodge-Tate weight $0$.\;Let $D$ be a $(\varphi,\Gamma)$-module of rank $r$ over $\cR_{E,L}$ such that there exists a continuous character $\delta: L^\times\rightarrow E^\times$ such that we have an injection $D\otimes_{\cR_{E,L}}\cR_{E,L}({\delta}^{-1})\hooklongrightarrow {\Delta}$ of $(\varphi,\Gamma)$-modules of rank $r$ over $\cR_{E,L}$.\;

In \cite[Section 4.1.1]{Ding2021},\;the authors consider the following functor:
\[F_{D}^0:\Art(E):=\{\text{Artinian local $E$-algebra with residue field $E$}\}\longrightarrow \{\text{sets}\}\]
sends $A$ to the set of isomorphism classes $\{(D_A,\pi_A,\delta_A)\}/\sim$,\;where
\begin{description}
\item[(1)] $D_A$ is a $(\varphi,\Gamma)$-module of rank $r$ over $\cR_{A,L}$ with $\pi_A:D_A\otimes_AE\cong D$;
\item[(2)] ${\delta}_{A}:L^\times\rightarrow A^\times$ is a continuous character such that ${\delta}_{A}\equiv {\delta}\mod \mathfrak{m}_A$;
\item[(3)] there exists an injection of $(\varphi,\Gamma)$-modules of rank $r$ over $\cR_{A,L}$:\;\;$D_A\hooklongrightarrow {\Delta}\otimes_{\cR_{E,L}}\cR_{A,L}({\delta}_{A})$.
\end{description}

\begin{rmk}We remark some easy facts about $F_{D}^0$.\;We consider such deformations for 
two reasons.\;
\begin{itemize}
	\item We consider $D_{E[\epsilon]/\epsilon^2}$ of $D$ over $E[\epsilon]/\epsilon^2$.\;The condition that $D_{E[\epsilon]/\epsilon^2}$ is a subobject of the form ${\Delta}\otimes_{\cR_{E,L}}\cR_{A,L}({\delta}_{E[\epsilon]/\epsilon^2})$ is more flexible than that $D_{E[\epsilon]/\epsilon^2}$ itself is of that form.\;In fact,\;by 
	\cite[Proposition 4.1.4]{Ding2021}) or Lemma \ref{deformationker},\;we see that the following map \[\ext^1_{(\varphi,\Gamma)}(D,D)\rightarrow\ext^1_{(\varphi,\Gamma)}(D,{\Delta}\otimes_{\cR_{E,L}}\cR_{E,L}({\delta}))\] is not injective,\;and its kernel has dimension $d_L\frac{r(r-1)}{2}$.\;
	\item We can view $F_{D}^0(E[\epsilon]/\epsilon^2)$ as a "determinant" part of $F_{D}(E[\epsilon]/\epsilon^2)$,\;which is easier to control.\;
\end{itemize}
\end{rmk}


Recall that in Definition  \ref{weaklynoncritical} we have fixed a $(\varphi,\Gamma)$-module $\Dpik$ which admits a non-critical special weakly  $\omepik$-filtration $\cF$ with parameter $(\bx_{\pi},\bmdel)\in \sbanpik\times\rigchl $.\;

Denote by $F_{\Dpik,\cF}^0$ the deformation functor
\[F_{\Dpik,\cF}^0:\Art(E):=\{\text{Artinian local $E$-algebra with residue field $E$}\}\longrightarrow \{\text{sets}\}\]
sends $A$ to the set of isomorphism classes $\{(D_A,\pi_A,\cF_A,\delta_A)\}/\sim$,\;where
\begin{itemize}
\item[(1)] $D_A$ is a $(\varphi,\Gamma)$-module of rank $n$ over $\cR_{A,L}$ with $\pi_A:D_A\otimes_AE\cong \Dpik$;
\item[(2)] $\cF_A=\fil_{\bullet}D_A$ is an increasing filtration of $(\varphi,\Gamma)$-module over $\cR_{A,L}$ on $D_A$,\;such that $\fil_{i}D_A$ are direct summand of $D_A$ as $\cR_{A,L}$-modules,\;and $\pi_A(\fil_{i}D_A)\cong \fil_{i}^\cF \Dpik$;
\item[(3)] ${\delta}_A=({\delta}_{A,i})_{1\leq i\leq k}$ where ${\delta}_{A,i}:L^\times\rightarrow A^\times$ is a continuous character such that ${\delta}_{A,i}\equiv \bm{\delta}_{\bh,i}(\mod \mathfrak{m}_A)$;
\item[(4)] there exists an injection of $(\varphi,\Gamma)$-modules of rank $r$ over $\cR_{A,L}$:
\begin{equation}\label{injectiondefF0deforma}
	\gr_iD_A\hooklongrightarrow {\Delta_\pi}\otimes_{\cR_{E,L}}\cR_{A,L}({\delta}_{A,i}).
\end{equation}
\end{itemize}
Note that the injection (\ref{injectiondefF0deforma}) is equivalent to the existence of an isomorphism of $(\varphi,\Gamma)$-modules over $\cR_{A,L}[1/t]$:
$\gr_iD_A[1/t]\cong{\Delta_\pi}\otimes_{\cR_{E,L}}\cR_{A,L}({\delta}_{A,i})[1/t]$.\;This deformation functor encodes certain paraboline deformations of $(\varphi,\Gamma)$-modules of $\Dpik$ which admit an $\omepik$-filtration.\;

Define by $F_{\Dpik,\cF}$ the deformation functor  which sends $A\in \Art(E)$ to the isomorphism classes $F_{\Dpik,\cF}(A):=\{(D_A,\pi_A,\cF_A)\}/\sim$.\;Moreover,\;the map sending the data $\{(D_A,\pi_A,\cF_A,\delta_A)\}$ (resp.,\;\\$\{(D_A,\pi_A,\cF_A)\}$) to $\{(\gr_iD_A,\pi_A|_{\gr_i^{\cF}D_A},{\delta}_{A,i})_{1\leq i\leq k}\}$ (resp.,\;$\{(\gr_iD_A,\pi_A|_{\gr_i^{\cF}D_A}\}$) induces a natural morphism
\begin{equation}\label{FmorphismGRD}
\Upsilon^0:F_{\Dpik,\cF}^0\longrightarrow \prod_{i=1}^kF_{\gr_i^{\cF}\Dpik}^0\;\;\; (\text{resp.,\;}\Upsilon:F_{\Dpik,\cF}\longrightarrow \prod_{i=1}^kF_{\gr_i^{\cF}\Dpik})
\end{equation}
By definition,\;we have a natural morphism $F_{\Dpik,\cF}^0\rightarrow F_{\Dpik,\cF}$.\;We see that
\[F_{\Dpik,\cF}^0=\sF_{D,\cF}\times_{\prod_{i=1}^k F_{\gr_iD}}\prod_{i=1}^k F_{\gr_iD}^0.\]
We have the following $E$-linear maps
\begin{equation}\label{kappakappaL}
\begin{aligned}
	\kappa:F_{\Dpik,\cF}^0(E[\epsilon]/\epsilon^2)\xrightarrow{\Upsilon^0(E[\epsilon]/\epsilon^2)} \prod_{i=1}^k&F_{\gr_i^{\cF}D}^0(E[\epsilon]/\epsilon^2)\\
	&\xrightarrow{\omega^0} \prod_{i=1}^k\homo(L^\times,E)\xrightarrow{\kappa_L} \prod_{ir\in \Delta_n(k)}\homo(L^\times,E),
\end{aligned}
\end{equation}
where the first map is obtained by evaluating the morphism $\Upsilon^0$ on $E[\epsilon]/\epsilon^2$-points,\;and the second map is given by 
$((D_{A,i},\pi_{A,i},\delta_{A,i}))_{1\leq i\leq k} 
\mapsto((
\delta_i^{-1}-1)/\epsilon)_{1\leq i\leq k}$.\;The last map $\kappa_L$ sends $(\psi_1,\psi_2,\cdots,\psi_k)$ to $(\psi_i-\psi_{i+1})_{ir\in \Delta_n(k)}$.\;By definition,\;the composition of $\Upsilon^0(E[\epsilon]/\epsilon^2)$ and  $\omega^0$ is given by
\[(\Dpik_{E[\epsilon]/\epsilon^2},\pi_{E[\epsilon]/\epsilon^2},\cF_{E[\epsilon]/\epsilon^2},\delta_{E[\epsilon]/\epsilon^2})\mapsto (({\delta}_{E[\epsilon]/\epsilon^2,i}\delta_i^{-1}-1)/\epsilon)_{1\leq i\leq k}.\]

In Section \ref{cgsformulaSection},\;we see that the functor $F_{\Dpik,\cF}^0$ is pro-representable.\;We next compute the $E$-dimension of tangent space  $F_{\Dpik,\cF}^0(E[\epsilon]/\epsilon^2)$ \;by using Colmez-Greenberg-Stevens formula.\;

\begin{rmk}\label{nongenediffobstru}
The deformation functor $F_{D,\cF'}^0$ in \cite[Section 4.1.2]{Ding2021} has generic parameters (\cite[(4.13)]{Ding2021}).\;This functor is formally smooth,\;which has a simpler structure than our  $F_{\Dpik,\cF}^0$.\;Proposition 4.1.17 in \cite{Ding2021} also computes its dimension.\;Note that the parameters of our non-critical special weakly $\omepik$-filtration $\cF$ are \textit{non-generic},\;thus the discussions in \cite[Section 4.1.2]{Ding2021} are not valid for our case.\;We study our deformation functor  $F_{\Dpik,\cF}^0$ by working a bit more carefully.\;One should note that,\;as opposed to generic case,\;there are obstructions to paraboline deformations of $\Dpik$.\;However,\;these obstructions to paraboline deformations are characterized by the parabolic Fontaine-Mazur  simple $\sL$-invariants (see Section \ref{dfnFMparalINVSECTION} and Section \ref{cgsformulaSection} below).\;This result overcomes the above difficulty.\;
\end{rmk}








%The following  lemma describes its image $\mathrm{Im}(\kappa)$.\;


%\end{proof}For part $(a)$,\;suppose we have a non-zero morphism $f:\gr^{\cF}_{j}D\longrightarrow \gr^{\cF}_{i}D\tee\cR_{E,L}(\ccyc)$,\;then for sufficiently large $N>0$,\;$f$ may induce a non-zero hence injective morphism
%\[t^N{\Delta_\pi}\hookrightarrow \gr^{\cF}_{j}D\tee \cR_{E,L}({\bm{\delta}}_{\pi,j}^{-1})\xrightarrow{f}\gr^{\cF}_{i}D\tee\cR_{E,L}(\ccyc\cdot{\bm{\delta}}_{\pi,j}^{-1})\hookrightarrow {\Delta_\pi}\tee\cR_{E,L}(\ccyc\cdot{\bm{\delta}}_{\pi,i}{\bm{\delta}}_{\pi,j}^{-1}).\]
%This induce a morphism ${\Delta_\pi}\rightarrow {\Delta_\pi}\tee\cR_{E,L}(\ccyc\cdot{\bm{\delta}}_{\pi,i}{\bm{\delta}}_{\pi,j}^{-1}z^{-\underline{N}})$.\;This implies that $$0\neq\hH^0_{(\varphi,\Gamma)}(\Delta_\pi^\vee\tee{\Delta_\pi}\tee\cR_{E,L}(\ccyc\cdot{\bm{\delta}}_{\pi,i}{\bm{\delta}}_{\pi,j}^{-1}z^{-\underline{N}}))\cong \hH^0_{(\varphi,\Gamma)}(\cR_{E,L}(\ccyc\cdot{\bm{\delta}}_{\pi,i}{\bm{\delta}}_{\pi,j}^{-1}z^{-\underline{N}})),$$
%which is impossible by \cite[Lem.\;2.1]{hellmann2016density}.\;Applying the same strategy  to part $(b)$.\;we can also get
%$$0\neq\hH^0_{(\varphi,\Gamma)}(\Delta_\pi^\vee\tee{\Delta_\pi}\tee\cR_{E,L}({\bm{\delta}}_{\pi,j}{\bm{\delta}}_{\pi,i}^{-1}z^{-\underline{N}}))\cong \hH^0_{(\varphi,\Gamma)}(\cR_{E,L}({\bm{\delta}}_{\pi,j}{\bm{\delta}}_{\pi,i}^{-1}z^{-\underline{N}})),$$
%contradicting the fact in  \cite[Lem.\;2.1]{hellmann2016density}.\;
%\end{proof}


\subsection{Parabolic Fontaine-Mazur  simple \texorpdfstring{$\sL$}{Lg}-invariants}\label{dfnFMparalINVSECTION}Suppose that $\Dpik$ admits a non-critical special weakly $\omepik$-filtration $\cF$ (see Definition \ref{weaklynoncritical}) with parameter $(\bx_{\pi},\bmdel)\in \sbanpik\times\rigchl $. In this section,\;we will attach to $(\varphi,\Gamma)$-module $\Dpik$ the parabolic Fontaine-Mazur simple $\sL$-invariants $\sL(\Dpik)$.\;

For $1\leq i\leq k$,\;we write $$\Dpik_i:=\gr_{i}^{\cF}\Dpik,\;\Delta_{\pi,i}:={\Delta_\pi}\tee \cR_{E,L}(\bm{\delta}_{\bh,i})$$ for simplicity.\;We first explain the construction of 
pairing (\ref{introperfectpairing})
\begin{equation}\label{normalperfectpairing}
\begin{aligned}
	\ext^1_{(\varphi,\Gamma)}(\Dpik_{i+1},\Dpik_{i})  \;\;\times\;\;  &\homo(L^\times,E) \big(\cong \hH^1_{(\varphi,\Gamma)}(\cR_{E,L})\big) \\      &\xrightarrow{\cup}
	E\big(\cong\ext^2_{(\varphi,\Gamma)}(\cR_{E,L}(\bm{\delta}_{\bh,i}\bm{\delta}^{-1}_{\bh,i+1}))\big)
\end{aligned}
\end{equation}
for each $1\leq i\leq k-1$.\;By \cite{liu2007cohomology},\;the cup product induces a pairing
\begin{equation}\label{usualcupproduct}
\xymatrix@C=5ex{
	\ext^1_{(\varphi,\Gamma)}(\Dpik_{i+1},\Dpik_{i})    \;\;\times\;\;  \ext^1_{(\varphi,\Gamma)}(\Dpik_{i},\Dpik_{i})  \xrightarrow{\cup}  \ext^2_{(\varphi,\Gamma)}(\Dpik_{i+1},\Dpik_{i}) .}
\end{equation}
We will show below that $\dim_E\ext^1_{(\varphi,\Gamma)}(\Dpik_{i+1},\Dpik_{i})=\dim_E\ext^1_{(\varphi,\Gamma)}(\Dpik_{i},\Dpik_{i})=1+r^2$ and \\$\ext^2_{(\varphi,\Gamma)}(\Dpik_{i+1},\Dpik_{i})\cong E$.\;If $r>1$,\;this pairing is not perfect.\;Roughly speaking,\;our pairing (\ref{normalperfectpairing}) is obtained by considering the cup product of the first term $\ext^1_{(\varphi,\Gamma)}(\Dpik_{i+1},\Dpik_{i})$ and a certain subspace $\homo(L^\times,E)$ of the second term  $\ext^1_{(\varphi,\Gamma)}(\Dpik_{i},\Dpik_{i})$ in (\ref{usualcupproduct}).\;We will prove that (\ref{normalperfectpairing}) is non-degenerate (but may not be perfect).\;Then the parabolic Fontaine-Mazur  simple $\sL$-invariants are defined in terms of this pairing.\;We need the following preliminaries.\;

The following lemma is a direct consequence of \cite[Lemma 4.1.12]{Ding2021}.\;
\begin{lem}\label{phigammacohononcrtialspecial}We have
\begin{itemize}
	\item[(a)] For any $j\neq i-1$,\;we have $$\hH^2_{(\varphi,\Gamma)}(\Dpik_{i},\Dpik_{j})\cong \hH^0_{(\varphi,\Gamma)}(\Dpik_{j},\Dpik_{i}\tee\cR_{E,L}(\ccyc))=0.$$
	\item[(b)] For any $i\neq j$,\;we have $\hH^0_{(\varphi,\Gamma)}(\Dpik_{i},\Dpik_{j})=0$.
\end{itemize}
\end{lem}
\begin{proof}Note that for the pair $(i,j)$ such that $j\neq i-1$,\;$(\bx_{\pi,i},\bm{\delta}_{\bh,i})$ and $(\bx_{\pi,j},\bm{\delta}_{\bh,j})$ satisfy the generic condition of \cite[(4.13)]{Ding2021}.\;Then $(a)$ is a direct consequence of \cite[Lemma 4.1.12]{Ding2021}.\;When $i\neq j$,\;the proof of	 $\hH^0_{(\varphi,\Gamma)}(\Dpik_{i},\Dpik_{j})=0$ in \cite[Lemma 4.1.12]{Ding2021} is also appropriate for our case (but is false for $\hH^2_{(\varphi,\Gamma)}(\Dpik_{i},\Dpik_{i-1})$).\;
\end{proof}

By (\ref{Dpikinjection}),\;the injections
\begin{equation*}
\begin{aligned}
	&\mathbf{I}_{i}:\Dpik_{i}\hookrightarrow \Delta_{\pi,i},\;\mathbf{I}_{i+1}:\Dpik_{i+1}  \hookrightarrow \Delta_{\pi,i+1}.\;
\end{aligned}
\end{equation*}
induces the following natural morphisms:
\begin{equation}\label{preliminamorphism1}
\begin{aligned}
	\ext^j_{(\varphi,\Gamma)}(\Dpik_{i+1},\Dpik_{i}) & \xrightarrow{\mathbf{I}_{i}^{2,j}}  & \ext^j_{(\varphi,\Gamma)}(\Dpik_{i+1},\Delta_{\pi,i}) & \xleftarrow{\mathbf{I}_{i+1}^{1,j}} & \ext^j_{(\varphi,\Gamma)}(\Delta_{\pi,i+1},\Delta_{\pi,i}),\\
	\ext^j_{(\varphi,\Gamma)}(\Delta_{\pi,i+1},\Delta_{\pi,i}) & \xleftarrow{\mathbf{I}_{i}^{2,j}} & \ext^j_{(\varphi,\Gamma)}(\Delta_{\pi,i+1},\Dpik_{i})& \xrightarrow{\mathbf{I}_{i+1}^{1,j}} & \ext^j_{(\varphi,\Gamma)}(\Dpik_{i+1},\Dpik_{i}).
	%\ext^1_{(\varphi,\Gamma)}(\gr_{i}^{\cF}\Dpik,\gr_{i}^{\cF}\Dpik) \hspace{10pt} \xrightarrow{\mathbf{I}_{i}} \hspace{10pt} & \ext^1_{(\varphi,\Gamma)}(\gr_{i}^{\cF}\Dpik,{\Delta_\pi}\tee \cR_{E,L}(\bm{\delta}_{\bh,i})) \\
	%	& \xleftarrow{\mathbf{I}_{i}} \hspace{10pt}  \ext^1_{(\varphi,\Gamma)}({\Delta_\pi}\tee \cR_{E,L}(\bm{\delta}_{\bh,i}),{\Delta_\pi}\tee \cR_{E,L}(\bm{\delta}_{\bh,i})).
\end{aligned}
\end{equation}
for $j=0,1,2$.\;
\begin{lem}\label{pullbackforiso}For $l\in \{1,2\}$ and $j\in\{0,1,2\}$,\;the morphisms $\mathbf{I}_{i}^{l,j}$ and $\mathbf{I}_{i+1}^{l,j}$ in (\ref{preliminamorphism1}) induce isomorphisms.\;
\end{lem}
\begin{proof}Let $D_0\hookrightarrow D_1$ be an injection of $(\phi,\Gamma)$-modules over $\cR_{E,L}$ with the same rank.\;After identifying the cohomology of $(\varphi,\Gamma)$-modules and the Galois cohomology of $E$-$B$-pairs (for example \cite[Section 3]{nakamura2009classification}),\;we deduce from the morphism $D_0\hookrightarrow D_1$ a long exact sequence 
\begin{equation}\label{isoIIIlong}
	\begin{aligned}
		0\longrightarrow   \hH^0_{(\varphi,\Gamma)}(D_0)&\longrightarrow
		\hH^0_{(\varphi,\Gamma)}(D_1)  \longrightarrow \hH^0_{(\varphi,\Gamma)}(D_1/D_0)\\
		\longrightarrow   &\hH^1_{(\varphi,\Gamma)}(D_0)\longrightarrow
		\hH^1_{(\varphi,\Gamma)}(D_1)  \longrightarrow \hH^1_{(\varphi,\Gamma)}(D_1/D_0)\\
		&\longrightarrow   \hH^2_{(\varphi,\Gamma)}(D_0)\longrightarrow
		\hH^2_{(\varphi,\Gamma)}(D_1)  \longrightarrow \hH^2_{(\varphi,\Gamma)}(D_1/D_0).
	\end{aligned}
\end{equation}
By \cite[Theorem 4.7]{liu2007cohomology},\;we have $\hH^2_{(\varphi,\Gamma)}(D_1/D_0)=0$ and 
\begin{equation}\label{equHIH2}
	\dim_E\hH^0_{(\varphi,\Gamma)}(D_1/D_0)=\dim_E\hH^1_{(\varphi,\Gamma)}(D_1/D_0).
\end{equation}
For a $(\varphi,\Gamma)$-module $D'$ over $\cR_{E,L}$,\;denoted by $\wdr(D')$ the $\bdr\otimes_{\BQ_p}E$-representation of $\gal_L$ associated to $D'$.\;For $\tau\in \Sigma_L$,\;let $\bdr_{,\tau,E}:=\bdr\otimes_{L,\tau} E$.\;By \cite[Lemma 5.1.1]{breuil2020probleme},\;we have
\begin{equation}\label{H0wdr}
	\hH^0_{(\varphi,\Gamma)}(D_1/D_0)\cong \hH^0(\gal_L,\wdr(D_1)/\wdr(D_0)),
\end{equation}
Applying the above discussion to the pairs
\begin{equation}
	\begin{aligned}
		(D_0,D_1)=&\;(\Dpik_{i+1}^\vee\tee\Dpik_{i},\Dpik_{i+1}^\vee\tee\Delta_{\pi,i}),\;\\
		\text{resp.,}&\;(\Delta_{\pi,i+1}^\vee\tee\Delta_{\pi,i},\Dpik_{i+1}^\vee\tee\Delta_{\pi,i}),\\
		\text{resp.,}&\;(\Delta_{\pi,i+1}^\vee\tee\Dpik_{i},\Delta_{\pi,i+1}^\vee\tee\Delta_{\pi,i}),\;\\
		\text{resp.,}&\;(\Delta_{\pi,i+1}^\vee\tee\Dpik_{i},\Dpik_{i+1}^\vee\tee\Dpik_{i}).\;
	\end{aligned}
\end{equation}
Note that 
\begin{equation}
	\begin{aligned}
		&\wdr(\Dpik_{i+1}^\vee\tee\Delta_{\pi,i})/\wdr(\Dpik_{i+1}^\vee\tee\Dpik_{i})\\
		&\hspace{80pt}\cong \oplus_{\tau\in \Sigma_L}\oplus_{s=1}^r\oplus_{l=1}^rt^{\hpi_{\tau,ir}-\hpi_{\tau,ir+s}}\bdr_{,\tau,E}/t^{\hpi_{\tau,(i-1)r+l}-\hpi_{\tau,ir+s}}\bdr_{,\tau,E},\\
		\text{resp.,\;}&\wdr(\Dpik_{i+1}^\vee\tee\Delta_{\pi,i})/\wdr(\Delta_{\pi,i+1}^\vee\tee\Delta_{\pi,i})\\
		&\hspace{80pt}\cong \oplus_{\tau\in \Sigma_L}\oplus_{s=1}^r\oplus_{l=1}^rt^{\hpi_{\tau,ir}-\hpi_{\tau,ir+s}}\bdr_{,\tau,E}/t^{\hpi_{\tau,ir}-\hpi_{\tau,(i+1)r}}\bdr_{,\tau,E},
		\\
		\text{resp.,\;}&		
		\wdr(\Delta_{\pi,i+1}^\vee\tee\Delta_{\pi,i})/\wdr(\Delta_{\pi,i+1}^\vee\tee\Dpik_{i})\\
		&\hspace{80pt}\cong \oplus_{\tau\in \Sigma_L}\oplus_{s=1}^r\oplus_{l=1}^rt^{\hpi_{\tau,ir}-\hpi_{\tau,(i+1)r}}\bdr_{,\tau,E}/t^{\hpi_{\tau,(i-1)r+l}-\hpi_{\tau,(i+1)r}}\bdr_{,\tau,E},\\
		\text{resp.,\;}&\wdr(\Dpik_{i+1}^\vee\tee\Dpik_{i})/\wdr(\Delta_{\pi,i+1}^\vee\tee\Dpik_{i})\\
		&\hspace{80pt}\cong \oplus_{\tau\in \Sigma_L}\oplus_{s=1}^r\oplus_{l=1}^rt^{\hpi_{\tau,(i-1)r+l}-\hpi_{\tau,ir+s}}\bdr_{,\tau,E}/t^{\hpi_{\tau,(i-1)r+l}-\hpi_{\tau,(i+1)r}}\bdr_{,\tau,E}.
	\end{aligned}
\end{equation}
All the choices of $(D_0,D_1)$ satisfy $\hH^0(\gal_L,\wdr(D_1)/\wdr(D_0))=0$.\;Then we have $$\hH^0_{(\varphi,\Gamma)}(D_1/D_0)\cong \hH^1(\gal_L,\wdr(D_1)/\wdr(D_0))=0.\;$$Now our lemma is a direct consequence of the long exact sequence (\ref{isoIIIlong}).\;
\end{proof}

The injection $\mathbf{I}_{i}:\Dpik_{i}\hookrightarrow \Delta_{\pi,i}$ also induces morphisms
\begin{equation}\label{preliminamorphism2}
\begin{aligned}
	\ext^1_{(\varphi,\Gamma)}(\Dpik_i,\Dpik_i)  \xrightarrow{\mathbf{I}_{i}'}  \ext^1_{(\varphi,\Gamma)}(\Dpik_i,\Delta_{\pi,i})  \xleftarrow{\mathbf{I}_{i}''}  \ext^1_{(\varphi,\Gamma)}(\Delta_{\pi,i},\Delta_{\pi,i}).
\end{aligned}
\end{equation}

\begin{lem}\label{deformationker}The morphism $\mathbf{I}_{i}''$ is an isomorphism,\;and $\dim_E\ker(\mathbf{I}_{i}')=d_L(1+\frac{r(r-1)}{2})$.\;
\end{lem}
\begin{proof}Proof of the injectivity of  \cite[ Proposition 4.1.4,\;(4.10)]{Ding2021} and  (\ref{isoIIIlong}),\;(\ref{equHIH2}) show that $\mathbf{I}_{i}''$ is an isomorphism.\;The dimension of the kernel $\ker(\mathbf{I}_{i}')$ is given in the last paragraph of the \cite[Proof of Proposition 4.1.4]{Ding2021}).\;
\end{proof}

Therefore for any $1\leq i\leq k-1$,\;all morphisms in (\ref{preliminamorphism1}) and (\ref{preliminamorphism2}) fit into the following commutative diagram:
\begin{equation}\label{bigcommutediag}
\xymatrix@C=5ex{
	\ext^1_{(\varphi,\Gamma)}(\Dpik_{i+1},\Dpik_{i}) \ar@{=}[d]                        & \times & \ext^1_{(\varphi,\Gamma)}(\Dpik_{i},\Dpik_{i})\hspace{9pt} \ar[d]^{\mathbf{I}_{i}'}         \ar[r]^{\cup}_{c_1}
	& \ext^2_{(\varphi,\Gamma)}(\Dpik_{i+1},\Dpik_{i}) \ar[d]_{\backsimeq}^{\mathbf{I}_{i}^{2,2}} \\
	\ext^1_{(\varphi,\Gamma)}(\Dpik_{i+1},\Dpik_{i}) \ar[d]^{\mathbf{I}_{i}^{2,1}}_{\simeq}             & \times & \ext^1_{(\varphi,\Gamma)}(\Dpik_{i},\Delta_{\pi,i})\hspace{5pt}                 \ar[r]^{\cup}_{c_2}
	& \ext^2_{(\varphi,\Gamma)}(\Dpik_{i+1},\Delta_{\pi,i}) \ar@{=}[d]           \\
	\ext^1_{(\varphi,\Gamma)}(\Dpik_{i+1},\Delta_{\pi,i})                             & \times & \ext^1_{(\varphi,\Gamma)}(\Delta_{\pi,i},\Delta_{\pi,i})\ar[u]_{\mathbf{I}_{i}''}^{\simeq}  \ar[r]^{\cup}_{c_3}
	&\ext^2_{(\varphi,\Gamma)}(\Dpik_{i+1},\Delta_{\pi,i})                       \\
	\ext^1_{(\varphi,\Gamma)}(\Delta_{\pi,i+1},\Delta_{\pi,i}) \ar[u]_{\mathbf{I}_{i+1}^{1,1}}^{\simeq}  & \times & \ext^1_{(\varphi,\Gamma)}(\Delta_{\pi,i},\Delta_{\pi,i}) \ar@{=}[u]            \ar[r]^{\cup}_{c_4}
	&\ext^2_{(\varphi,\Gamma)}(\Delta_{\pi,i+1},\Delta_{\pi,i}) \ar[u]^{\backsimeq}_{\mathbf{I}_{i+1}^{1,2}}.}
\end{equation}

Recall the $E$-$B$-pairs defined in \cite{nakamura2009classification}.\;By \cite[Theorem 1.36]{nakamura2009classification},\;there exists an equivalence of category between the category of $E$-$B$-pairs and  the category of $(\varphi,\Gamma)$-modules over $\cR_{E,L}$.\;Let $\chi:L^\times\rightarrow E^\times$ be a continuous character,\;and let $B_E(\chi)$ be the rank-one $E$-$B$-pair associated to $\chi$ (see \cite[Theorem 1.45]{nakamura2009classification}).\;

Put $\EndO(\Delta_{\pi}):=\Delta_{\pi}\tee\Delta_{\pi}^\vee$ and $\EndO^0(\Delta_{\pi})=\EndO(\Delta_{\pi})/\cR_{E,L}=(\Delta_{\pi} \tee\Delta_{\pi}^\vee)/\cR_{E,L}$.\;Then we have $\EndO(\Delta_{\pi})=\EndO^0(\Delta_{\pi})\oplus \cR_{E,L}$.\;Let $\bW=(\bW_e,\bW_{\dr}^+)$ (resp.,$\bW^0:=(\bW_e^0,\bW_{\dr}^{0,+})$) be the $E$-$B$-pair associated to the $(\varphi,\Gamma)$-module $\EndO(\Delta_{\pi})$ (resp.,\;$\EndO^0({\Delta_\pi})$).\;We put $\bW(\chi):=\bW\otimes_EB_E(\chi)$ (resp.,\;$\bW^0(\chi):=\bW^0\otimes_EB_E(\chi)$).\;Using $E$-$B$-pairs,\;we show that
\begin{lem}\label{lembigcommutediag}
The cup products $c_2$,\;$c_3$ and $c_4$ in (\ref{bigcommutediag}) are perfect pairings.\;
\end{lem}
\begin{proof}It suffices to prove that $c_4$ is perfect.\;Note that the natural morphism $$j:\bW(\bm{\delta}_{\bh,i}\bm{\delta}_{\bh,i+1}^{-1})\rightarrow \bW(\ccyc)$$ (induced by $B_E(\bm{\delta}_{\bh,i}\bm{\delta}_{\bh,i+1}^{-1})\hookrightarrow B_E(\ccyc)$) gives an exact sequence of $\gal_L$-complexes,\;
\begin{equation*}
	\begin{aligned}
		0\longrightarrow [\bW(\bm{\delta}_{\bh,i}\bm{\delta}_{\bh,i+1}^{-1})_e\oplus& \bW(\bm{\delta}_{\bh,i}\bm{\delta}_{\bh,i+1}^{-1})^+_{\dr}\rightarrow \bW(\bm{\delta}_{\bh,i}\bm{\delta}_{\bh,i+1}^{-1})_{\dr}]\\
		\longrightarrow \;& [\bW(\ccyc)_e\oplus \bW(\ccyc)^+_{\dr}\rightarrow \bW(\ccyc)_{\dr}]\\
		&\longrightarrow [\oplus_{\tau\in\Sigma_L}t\bdr_{,\tau,E}/t^{\hpi_{\tau,ir}-\hpi_{\tau,(i+1)r}}\bdr_{,\tau,E}\rightarrow0]\longrightarrow 0.\\
	\end{aligned}
\end{equation*}
Since $\hH^i(\gal_L,\oplus_{\tau\in\Sigma_L}t\bdr_{,\tau,E}/t^{\hpi_{\tau,ir}-\hpi_{\tau,(i+1)r}}\bdr_{,\tau,E})=0$ for any $i\geq 0$.\;We see that $j$ induces isomorphisms $\hH^i(\gal_L,\bW(\bm{\delta}_{\bh,i}\bm{\delta}_{\bh,i+1}^{-1}))\xrightarrow{\sim}\hH^i(\gal_L,\bW(\ccyc))$ for all $i\geq 0$.\;Furthermore,\;the following diagram commutes:
\[\xymatrix@C=3.5ex{
	\ext^1_{(\varphi,\Gamma)}(\Delta_{\pi,i+1},\Delta_{\pi,i}) \ar[d]_{\simeq}  & \times & \ext^1_{(\varphi,\Gamma)}(\Delta_{\pi,i},\Delta_{\pi,i}) \ar@{=}[d] \ar[r]^{\hspace{9pt}\cup}_{\hspace{9pt}c_4} & \ext^2_{(\varphi,\Gamma)}(\Delta_{\pi,i+1},\Delta_{\pi,i}) \ar[d]_{\simeq}\\
	\ext^1_{(\varphi,\Gamma)}(\Delta_{\pi},\Delta_{\pi}(\ccyc))   & \times & \ext^1_{(\varphi,\Gamma)}(\Delta_{\pi}(\ccyc),\Delta_{\pi}(\ccyc))\ar[r]^{\cup} & \ext^2_{(\varphi,\Gamma)}(\Delta_{\pi},\Delta_{\pi}(\ccyc)) ,}\]
where $\Delta_{\pi}(\ccyc):=\Delta_{\pi}\tee\cR_{E,L}(\ccyc)$ and $\ext^2_{(\varphi,\Gamma)}(\Delta_{\pi,i+1},\Delta_{\pi,i})\xrightarrow{\sim}\ext^2_{(\varphi,\Gamma)}(\Delta_{\pi},\Delta_{\pi}(\ccyc))\cong E$. The bottom cup product is perfect by using Tate duality.\;This completes the proof.\;
\end{proof}

The second term of (\ref{normalperfectpairing}) is given as follows.\;By the decomposition $\EndO(\Delta_{\pi})=\EndO^0(\Delta_{\pi})\oplus \cR_{E,L}$, we have decompositions
\begin{equation}\label{Linvariantdecomp}
\begin{aligned}
	&\ext^i_{(\varphi,\Gamma)}(\Delta_{\pi,i+1},\Delta_{\pi,i})\\
	=\;&\hH^i_{(\varphi,\Gamma)}(\cR_{E,L}(\bm{\delta}_{\bh,i}\bm{\delta}_{\bh,i+1}^{-1}))\oplus \hH^i_{(\varphi,\Gamma)}(\EndO^0(\Delta_{\pi})\tee\cR_{E,L}(\bm{\delta}_{\bh,i}\bm{\delta}_{\bh,i+1}^{-1}))\\
	\text{resp.,\;}& \ext^1_{(\varphi,\Gamma)}(\Delta_{\pi,i},\Delta_{\pi,i})=\hH^1_{(\varphi,\Gamma)}(\cR_{E,L})\oplus \hH^1_{(\varphi,\Gamma)}(\EndO^0(\Delta_{\pi})).
\end{aligned}
\end{equation}
for $i=0,1,2$.\;These induce the injections
\begin{equation}\label{i3p3q3}
\begin{aligned}
	&\widetilde{\mathbf{I}}_{i}^1:\hH^1_{(\varphi,\Gamma)}(\cR_{E,L}(\bm{\delta}_{\bh,i}\bm{\delta}_{\bh,i+1}^{-1}))\rightarrow \ext^1_{(\varphi,\Gamma)}(\Delta_{\pi,i+1},\Delta_{\pi,i}),\\
	&\widetilde{\mathbf{I}}_{i}^2:\hH^2_{(\varphi,\Gamma)}(\cR_{E,L}(\bm{\delta}_{\bh,i}\bm{\delta}_{\bh,i+1}^{-1}))\rightarrow \ext^2_{(\varphi,\Gamma)}(\Delta_{\pi,i+1},\Delta_{\pi,i}),\\
	\text{resp.,\;}& \widetilde{\mathbf{I}}'_{i}:\hH^1_{(\varphi,\Gamma)}(\cR_{E,L})\cong \homo(L^\times,E)\rightarrow \ext^1_{(\varphi,\Gamma)}(\Delta_{\pi,i},\Delta_{\pi,i}).\;
\end{aligned}
\end{equation}

\begin{lem}\label{isoI2i}The morphism $\widetilde{\mathbf{I}}_{i}^2$ is an isomorphism.\;
\end{lem}
\begin{proof}By using Tate duality,\;it suffices to prove that $$(\widetilde{\mathbf{I}}_{i}^2)^\vee:\ext^0_{(\varphi,\Gamma)}(\EndO(\Delta_{\pi})\tee \cR_{E,L}(\chi_{\mathrm{cyc}}\cdot\bm{\delta}_{\bh,i}^{-1}\bm{\delta}_{\bh,i+1}))
\rightarrow  \ext^0_{(\varphi,\Gamma)}(\cR_{E,L}(\chi_{\mathrm{cyc}}\cdot\bm{\delta}_{\bh,i}^{-1}\bm{\delta}_{\bh,i+1}))$$ 
is an isomorphism,\;i.e.,\;$\hH^0_{(\varphi,\Gamma)}(\EndO^0(\Delta_{\pi})\tee \cR_{E,L}(\chi_{\mathrm{cyc}}\cdot\bm{\delta}_{\bh,i}^{-1}\bm{\delta}_{\bh,i+1}))=0$.\;Indeed,\;the unique (up to scalar) nonzero element in  $\hH^0_{(\varphi,\Gamma)}(\cR_{E,L}(\chi_{\mathrm{cyc}}\cdot\bm{\delta}_{\bh,i}^{-1}\bm{\delta}_{\bh,i+1}))\cong E$ corresponds to an injection $\cR_{E,L}\hookrightarrow \cR_{E,L}(\chi_{\mathrm{cyc}}\cdot\bm{\delta}_{\bh,i}^{-1}\bm{\delta}_{\bh,i+1})$.\;Note that
\begin{eqnarray*}
	\bW^0(\chi_{\mathrm{cyc}}\cdot\bm{\delta}_{\bh,i}^{-1}\bm{\delta}_{\bh,i+1})^+_{\dr})/\bW^{0,+}_{\dr}
	\cong \oplus_{\tau\in \Sigma_L}\big(t^{\bh_{(i+1)r,\tau}-\bh_{ir,\tau}}\bdr_{,\tau,E}/\bdr_{,\tau,E}\big)^{\oplus (r^2-1)},
\end{eqnarray*}
we deduce $\hH^0_{(\varphi,\Gamma)}(\EndO^0(\Delta_{\pi})\tee \cR_{E,L}(\chi_{\mathrm{cyc}}\cdot\bm{\delta}_{\bh,i}^{-1}\bm{\delta}_{\bh,i+1})/\EndO^0(\Delta_{\pi}))=0$.\;By an easy d\'{e}vissage argument,\;we see that $\hH^0_{(\varphi,\Gamma)}(\EndO^0(\Delta_{\pi})\tee \cR_{E,L}(\chi_{\mathrm{cyc}}\cdot\bm{\delta}_{\bh,i}^{-1}\bm{\delta}_{\bh,i+1}))\xleftarrow{\sim}\hH^0_{(\varphi,\Gamma)}(\EndO^0(\Delta_{\pi}))=0$ (recall that $\Delta_{\pi}$ is irreducible).\;This completes the proof.\;
\end{proof}

We are ready to define the desired non-degenerate pairing.\;

\begin{pro}\label{pairingdef}For $1\leq i\leq k-1$,\;we have the following non-degenerate pairing:
\begin{equation}
	\xymatrix{
		\langle -,-\rangle_{\bc_i}: \ext^1_{(\varphi,\Gamma)}(\Dpik_{i+1},\Dpik_{i})  \times  \hH^1_{(\varphi,\Gamma)}(\cR_{E,L}) \ar[r]^{\hspace{20pt}\cup}
		& \ext^2_{(\varphi,\Gamma)}(\cR_{E,L}(\bm{\delta}_{\bh,i}\bm{\delta}^{-1}_{\bh,i+1}))\cong E.}
\end{equation}
\end{pro}
\begin{proof}Combining the  diagrams (\ref{bigcommutediag}) with  (\ref{i3p3q3}),\;we get the following commutative diagram:
\begin{equation}
	\xymatrix@C=2.5ex{
		\ext^1_{(\varphi,\Gamma)}(\Dpik_{i+1},\Dpik_{i}) \ar@{=}[d]                        & \times & \ext^1_{(\varphi,\Gamma)}(\Dpik_{i},\Dpik_{i})\hspace{9pt} \ar[d]^{ {\mathbf{I}}'_{i}}         \ar[r]^{\cup}_{c_1}
		& \ext^2_{(\varphi,\Gamma)}(\Dpik_{i+1},\Dpik_{i}) \ar[d]_{\backsimeq}^{{\mathbf{I}}_{i}^{2,2}} \\
		\ext^1_{(\varphi,\Gamma)}(\Dpik_{i+1},\Dpik_{i})           & \times & \ext^1_{(\varphi,\Gamma)}(\Dpik_{i},\Delta_{\pi,i})\hspace{5pt}                 \ar[r]^{\cup}_{c_2}
		& \ext^2_{(\varphi,\Gamma)}(\Dpik_{i+1},\Delta_{\pi,i})          \\
		\hH^1_{(\varphi,\Gamma)}(\cR_{E,L}(\bm{\delta}_{\bh,i}\bm{\delta}_{\bh,i+1}^{-1})) \ar@{_(->}[u]_{({\mathbf{I}}_{i}^{2,1})^{-1}\circ {\mathbf{I}}_{i+1}^{1,1}\circ \widetilde{\mathbf{I}}_{i}^1}  & \times & \hH^1_{(\varphi,\Gamma)}(\cR_{E,L})\cong \homo(L^\times,E)   \ar@{_(->}[u]_{ {\mathbf{I}}''_{i}\circ  \widetilde{\mathbf{I}}'_{i}}         \ar[r]^{\cup}_{c_5}
		&\hH^2_{(\varphi,\Gamma)}(\cR_{E,L}(\bm{\delta}_{\bh,i}\bm{\delta}_{\bh,i+1}^{-1})) \ar[u]^{\backsimeq}_{{\mathbf{I}}_{i+1}^{1,2}\circ \widetilde{\mathbf{I}}_{i}^2}.
	}
\end{equation}
Note that $\hH^2_{(\varphi,\Gamma)}(\cR_{E,L}(\bm{\delta}_{\bh,i}\bm{\delta}_{\bh,i+1}^{-1}))\cong E$.\;The desired pairing is given by the cup product of the $(2,1)$-term and $(3,2)$-term of this diagram.\;By definition,\;this pairing is non-degenerate.\;
\end{proof}

%We now get the following commutative diagram (see (\ref{i3p3q3}))
%\begin{equation}\label{bigcommutediagbelow}
%\xymatrix@C=2.5ex{
%	 \ext^1_{(\varphi,\Gamma)}(\Delta_{\pi,i+1},\Delta_{\pi,i})   & \times & \ext^1_{(\varphi,\Gamma)}(\Delta_{\pi,i},\Delta_{\pi,i})            \ar[r]^{\cup}_{c_4}
%	&\ext^2_{(\varphi,\Gamma)}(\Delta_{\pi,i+1},\Delta_{\pi,i})\cong E \\
%	\hH^1_{(\varphi,\Gamma)}(\cR_{E,L}(\bm{\delta}_{\bh,i}\bm{\delta}_{\bh,i+1}^{-1})) \ar@{_(->}[u]_{i_3}  & \times & \hH^1_{(\varphi,\Gamma)}(\cR_{E,L})\cong \homo(L^\times,E)   \ar@{_(->}[u]_{q_3}         \ar[r]^{\cup}_{c_5}
%	&\hH^2_{(\varphi,\Gamma)}(\cR_{E,L}(\bm{\delta}_{\bh,i}\bm{\delta}_{\bh,i+1}^{-1}))\cong E \ar[u]^{\simeq}_{p_3}.}
%\end{equation}

For $1\leq i\leq k-1$,\;the subobject $\Dpik^{i+1}_{i}$ gives an extension class $[\Dpik^{i+1}_{i}]\in \ext^1_{(\varphi,\Gamma)}(\Dpik_{i+1},\Dpik_{i})$. We further put $$[({\Delta_{\Dpik}})^{i+1}_{i}]=({\mathbf{I}}_{i}^{2,1})^{-1}\circ {\mathbf{I}}_{i+1}^{1,1}([\Dpik^{i+1}_{i}])\in \ext^1_{(\varphi,\Gamma)}(\Delta_{\pi,i+1},\Delta_{\pi,i}).$$
With respect to the  direct sum decomposition (\ref{Linvariantdecomp}),\;we put 
\begin{equation}\label{c0extensions}
\begin{aligned}
	&({\Delta_{\Dpik}})^{i+1}_{i}:=\padelc^{i+1}_{i}\oplus \padel^{i+1}_{i},\;
	%		&(\Dpik^{i+1}_{i})^z:=\padelc^{i+1}_{i}\big[{1}/{t}\big]\cap \Dpik^{i+1}_{i},\;(\Dpik^{i+1}_{i})^0:=\padel^{i+1}_{i}\big[{1}/{t}\big]\cap \Dpik^{i+1}_{i}.\;
\end{aligned}
\end{equation}
where $\padelc^{i+1}_{i}$ (resp.,\;$\padel^{i+1}_{i}$) is the projection of $({\Delta_{\Dpik}})^{i+1}_{i}$ to the first (resp.,\;second) factor.\;


We are ready to define the parabolic Fontaine-Mazur simple $\sL$-invariants of $\Dpik$.\;
\begin{dfn}\textbf{(Parabolic Fontaine-Mazur simple $\sL$-invariants)}Suppose that $\Dpik$ admits a
non-critical special $\omepik$-filtration with parameter \[(\bx_{\pi},\bmdel)\in \sbanpik\times\rigchl \] (resp.,\; with parameter $(\widetilde{\bx}_{\pi,\bh},\widetilde{\bm{\delta}}_\bh)\in \sbanpik\times\rigch $).\;

For $ir\in \Delta_n(k)=\{r,2r,\cdots,(k-1)r\}\}$,\;we can attach $\sL(\Dpik)_{ir}\subseteq \hH^1_{(\varphi,\Gamma)}(\cR_{E,L})$ (as a subspace of $\ext^1_{(\varphi,\Gamma)}(\Delta_{\pi,i},\Delta_{\pi,i})$) by (see Proposition \ref{pairingdef})
\begin{equation}\label{dfnl-invariant}
	\begin{aligned}
		\sL(\Dpik)_{ir}:&=\{\psi\in \hH^1_{(\varphi,\Gamma)}(\cR_{E,L})\;|\;\langle \psi,[\Dpik^{i+1}_{i}]\rangle_{\bc_i}=0\}.\;
	\end{aligned}
\end{equation}
We call $\sL(\Dpik)=\prod_{ir\in \Delta_n(k)}\sL(\Dpik)_{ir}$ the parabolic simple $\sL$-invariants of $\Dpik$.\;
\end{dfn}
\begin{rmk}By definition,\;we have $\langle \psi,[\Dpik^{i+1}_{i}]\rangle_{\bc_i}=\langle   \widetilde{\mathbf{I}}'_{i}(\psi),({\Delta_{\Dpik}})^{i+1}_{i}\rangle_{c_4}$ (see (\ref{bigcommutediag})).\;
\end{rmk}
\begin{lem}\label{criofnoncrystalline}The following are equivalent:
\begin{itemize}
	\item[(a)] For all $ir\in \Delta_n(k)$,\;$
	\Dpik^{i+1}_i$ is potentially semistable and noncrystalline,\;
	\item[(b)] For all $ir\in \Delta_n(k)$,\;$\padelc^{i+1}_{i}$ is  semistable noncrystalline,
	\item[(c)] For all $ir\in \Delta_n(k)$,\;$\homo_\infty(L^\times,E)\nsubseteq \sL(\Dpik)_{ir}$,\;
	\item[(d)]  We have $N^{k-1}\neq 0$ on $\df(\Dpik)$,\;i.e,\;$\df(\Dpik)$ is absolutely indecomposable. 
\end{itemize}
\end{lem}
\begin{proof} For any $ir\in \Delta_n(k)$,\;the cup product induces a perfect paring (see \cite[Lemma 1.13]{2015Ding})
\[\hH^1_{(\varphi,\Gamma)}(\cR_{E,L}(\bm{\delta}_{\bh,i}\bm{\delta}_{\bh,i+1}^{-1}))\times\hH^1_{(\varphi,\Gamma)}(\cR_{E,L})\xrightarrow{\cup}\hH^2_{(\varphi,\Gamma)}(\cR_{E,L}(\bm{\delta}_{\bh,s-1}\bm{\delta}_{\bh,s}^{-1}))\cong E.\]
Then 	$\hH^1_{(\varphi,\Gamma),e}(\cR_{E,L}(\bm{\delta}_{\bh,i}\bm{\delta}_{\bh,i+1}^{-1}))^\bot:=\{\psi\in \hH^1_{(\varphi,\Gamma)}(\cR_{E,L})\;|\;\langle\psi,\hH^1_{(\varphi,\Gamma),e}(\cR_{E,L}(\bm{\delta}_{\bh,i}\bm{\delta}_{\bh,i+1}^{-1}))\rangle=0\}$ is isomorphic to $\homo_\infty(L^\times,E)$ (see \cite[Proposition 1.9,\;Lemma 1.15]{2015Ding}).\;Thus we see that \[\homo_\infty(L^\times,E)\subseteq\sL(\Dpik)_{ir}\] if and only if $\padelc^{i+1}_{i}$ is crystalline.\;Therefore $(b)$ and $(c)$ are equivalent.\;

We next show that $(a)$ is equivalent to $(b)$.\;Since 	$\Dpik_{i}^{i+1}[\frac{1}{t}]=(\Delta_\Dpik)^{i+1}_{i}[\frac{1}{t}]$,\;Part (a) is equivalent to prove that $(\Delta_\Dpik)^{i+1}_{i}=\padelc^{i+1}_{i}\oplus \padel^{i+1}_{i}$ is potentially semistable noncrystalline.\;It suffices to show that $\padel^{i+1}_{i}$ is always potentially crystalline.\;Let $L'$ be a sufficiently large extension of $L$ such that the restriction of $\EndO^0({\Delta_\pi})\tee\cR_{E,L}(\bm{\delta}_{\bh,i}\bm{\delta}_{\bh,i+1}^{-1})$ to $L'$ is crystalline.\;Recall that $\bW^0$ (resp.,\; $\bW^0(\bm{\delta}_{\bh,i}\bm{\delta}_{\bh,i+1}^{-1})$) is the $E$-$B$-pair associated to $\EndO^0({\Delta_\pi})$ (resp.,\;$\EndO^0({\Delta_\pi})\tee\cR_{E,L}(\bm{\delta}_{\bh,i}\bm{\delta}_{\bh,i+1}^{-1})$, see \cite[Theorem 1.36]{nakamura2009classification}).\;We are going to show that \[\hH^1_{f}(\gal_{L'},\bW^0(\bm{\delta}_{\bh,i}\bm{\delta}_{\bh,i+1}^{-1}))=\hH^1(\gal_{L'},\bW^0(\bm{\delta}_{\bh,i}\bm{\delta}_{\bh,i+1}^{-1})),\] then \cite[Remark 2.5]{nakamura2009classification}) implies that  $\padel^{i+1}_{i}|_{L'}$ is always crystalline.\;Indeed,\;by \cite[Proposition 2.7]{nakamura2009classification}), we have
\begin{equation*}
	\begin{aligned}
		&\dim_E\hH^1_f(\gal_{L'},\bW^0(\bm{\delta}_{\bh,s-1}\bm{\delta}_{\bh,s}^{-1}))\\
		=\;&\dim_E\hH^0(\gal_{L'},\bW^0(\bm{\delta}_{\bh,s-1}\bm{\delta}_{\bh,s}^{-1}))+\dim_ED_{\dr}(\bW^0(\bm{\delta}_{\bh,s-1}\bm{\delta}_{\bh,s}^{-1}))/\fil^H_0D_{\dr}(\bW^0(\bm{\delta}_{\bh,s-1}\bm{\delta}_{\bh,s}^{-1}))\\
		=\;&\dim_E\hH^0(\gal_{L'},\bW^0(\bm{\delta}_{\bh,s-1}\bm{\delta}_{\bh,s}^{-1}))+\dim_ED_{\dr}(\bW^0(\bm{\delta}_{\bh,s-1}\bm{\delta}_{\bh,s}^{-1}))\\
		=\;&\dim_E\hH^1(\gal_{L'},\bW^0(\bm{\delta}_{\bh,s-1}\bm{\delta}_{\bh,s}^{-1})).\;
	\end{aligned}
\end{equation*}
It remains to show that $(a)$ and $(d)$ are equivalent.\;Suppose there exists $i\in \Delta_k$ such that $\Dpik^{i+1}_i|_{L'}$ is crystalline.\;$N$ is zero on $\Dpik^{i+1}_i|_{L'}$.\;Then we see that $N^{k-i-1}=0$ on $D_{\pst}(\Dpik^{k}_{i+1})$,\;$N^{i-1}=0$ on $D_{\pst}(\Dpik^{i-1}_{1})$ and $N=0$ on $D_{\pst}(\Dpik^{i+1}_{i})$.\;We deduce from an easy d\'evissage argument that the monodromy operator $N^{k-1}=N^{(k-i-1)+(i-1)+1}=0$ on $\df(\pi)$.\;The ``only if" part follows.\;Conversely,\;suppose $\Dpik^{i+1}_i|_{L'}$ is semistable noncrystalline for all $i\in \Delta_k$.\;By definition,\;we pick a basis \[\{v_{1,1},v_{1,2},\cdots,v_{1,r}\}\cup\cdots \cup\{v_{k,1},v_{k,2},\cdots,v_{k,r}\}\cup \{v_{k-1,1},v_{k-1,2},\cdots,v_{k-1,r}\}\] of $D_{\pst}(\Dpik)$ such that the matrices of  $\varphi^{f_L'}$ and $N$  under this basis is given by
\[\varphi^{f_L'}=\left(
\begin{array}{cccc}
	A_1 & \ast & \ast & \ast \\
	& q_L\cdot A_2 & \ast & \ast \\
	&  & \ddots & \ast \\
	&  &  & q_L^{(k-1)}\cdot A_k \\
\end{array}
\right),\;
N=\left(
\begin{array}{cccc}
	0 & \ast  & \ast & \ast  \\
	& 0 & \ast &  \ast \\
	&  & \ddots & \ast  \\
	&  &  & 0 \\
\end{array}
\right),
\]
where all the $A_i\cong A$ for some matrix $A\in\GLN_r(L\otimes_{\bQ_p}E)$.\;We use induction on $l$ to show that $N^{l-1}\neq 0$ on $D_{\pst}(\Dpik^{l}_{1})$ and $N^{l-1}v_{l,j}\neq 0$ for some  $1\leq j\leq r$.\;This is trivial for $l=2$.\;Suppose that $N^{l-2}\neq 0$ on $D_{\pst}(\Dpik^{l-1}_{1})$ and $N^{l-2}A_{l-1}v_{l-1,j}\neq 0$ for some  $j$.\;Then we have $\varphi^{f_L'}v_{l-1,j}-q_L^{l-1}A_{l-1}v_{l-1,j}\in \Dpik^{l-2}_{1}$, $\varphi^{f_L'}v_{l,j}-q_L^{l}A_{l}v_{l,j}\in \Dpik^{l-1}_{1}$ and 
$\varphi^{f_L'}(Nv_{l,j}-v_{l-1,j})\in D_{\pst}(\Dpik^{l-2}_{1})$.\;We show that $N^{l-1}A_{l}v_{l,j}\neq 0$.\;Indeed, we have
\begin{equation}
\begin{aligned}
	q_L^{l}N^{l-2}\varphi^{f_L'}NA_lv_{l,j}&=N^{l-2}\varphi^{f_L'}N(\varphi^{f_L'}v_{l,j}+(q_L^{l}A_{l}v_{l,j}-\varphi^{f_L'}v_{l,j}))\\
	&=N^{l-2}\varphi^{f_L'}N(\varphi^{f_L'}v_{l,j})=q_LN^{l-2}\varphi^{2f_L'}Nv_{l,j}\\
	&=q_LN^{l-2}\varphi^{2f_L'}v_{l-1,j}=q_L^2\varphi^{f_L'}N^{l-2}\varphi^{f_L'}v_{l-1,j}\\
	&=q_L^2\varphi^{f_L'}N^{l-2}(q_L^{l-1}A_{l-1}v_{l-1,j}+(\varphi^{f_L'}v_{l-1,j}-q_L^{l-1}A_{l-1}v_{l-1,j}))\\
	&=q_L^2\varphi^{f_L'}(q_L^{l-1}N^{l-2}A_{l-1}v_{l-1,j})\neq 0.
\end{aligned}
\end{equation}
Since $N^{l-2}\varphi^{f_L'}NA_lv_{l,j}=q_L^{l-2}\varphi^{f_L'}N^{l-1}A_lv_{l,j}$,\;we get that $N^{l-1}A_lv_{l,j}\neq 0$ and hence $N^{l-1}\neq 0$ on $D_{\pst}(\Dpik^{l}_{1})$.\;This completes the proof.\;
\end{proof}

For $\tau\in \Sigma_L$,\;recall that $\homo_\tau(L^\times,E)$ is a two-dimensional $E$-vector space and admits a basis $\psi_{\sigma,L}:=\tau\circ\log_p$ and $\psi_{\ur}:=\valua_L$,\;where $\log_p:L^\times\rightarrow L$ is equal to the $p$-adic logarithm when restricted to $\cO_L^\times$ and where $\log_p(p)$.\;We put $\sL(\Dpik)_{ir,\tau}:=\sL(\Dpik)_{ir}\cap \homo_\tau(L^\times,E)$.\;

\begin{rmk}\label{taulinvariantrmk}Assume that $
\Dpik^{i+1}_i$ is potentially semistable and noncrystalline for each $ir\in \Delta_n(k)$, then the Lemma  \ref{criofnoncrystalline} shows that $\homo_\infty(L^\times,E)\nsubseteq \sL(\Dpik)_{ir}$.\;Therefore,\;we have $$\dim_E\sL(\Dpik)_{ir,\tau}=1$$ for all $\tau\in \Sigma_L$ and $\sL(\Dpik)_{ir}=\bigoplus_{\tau\in \Sigma_L}\sL(\Dpik)_{ir,\tau}$.\;Moreover,\;there exists $\sL_{i,\tau}\in E $ such that $\sL(\Dpik)_{ir,\tau}$ is generated by $\psi_{i,\tau}:=\psi_{\tau,L}-\sL_{i,\tau}\psi_{\ur}$.\;
\end{rmk}
%We end this section with a remark.\;
%\begin{rmk}Similar to the philosophy in \cite{2015Ding},\;the above parabolic $\sL$-invariant have a partial version.\;More precisely,\;if we have $n$-dimensional potentially semistable Galois representation with critical $\omepik$-filtration,\;i.e,\;there is a subset $S_c(D)$ of $\Sigma_L$ such that $w_{\cF,\tau}\neq w_0$ for all $\tau\in S_c(D)$.\;Then for each $\tau\in \Sigma_L\backslash S_c(D)$ and $ir\in \Delta_n(k)$,\;we can associate parabolic $\sL$-invariant $\sL(\Dpik)_{i,\tau}$ by modifying and combining  the construction in \cite[Section 1.3.2]{2015Ding} and our construction.\;In this case,\;it is worth to research the locally analytic socle conjecture and companion points in eigenvariety in this case.\;We expert consider this problem in future work.\;
%




\subsection{Colmez-Greenberg-Stevens formula}\label{cgsformulaSection}

In \cite[Section.\;3.3]{2019DINGSimple},\;Ding established the Colmez-Greenberg-Stevens formula 
(on simple $\sL$-invariants) for a rank $n$ triangulable $(\varphi,\Gamma)$-module over $\cR_{E,L}$. In this section,\;we will establish the Colmez-Greenberg-Stevens formula (on parabolic Fontaine-Mazur simple $\sL$-invariants) for a  rank $n$ $(\varphi,\Gamma)$-module $\Dpik$,\;which
admits a non-critical special weakly $\omepik$-filtration $\cF$ with  parameter $(\bx_{\pi},\bmdel)\in \sbanpik\times\rigchl$.\;In particular,\;our Colmez-Greenberg-Stevens formula shows that we can reinterpret the parabolic Fontaine-Mazur simple $\sL$-invariants in terms of $1$-order paraboline deformations of type $\omepik$.\;The Colmez-Greenberg-Stevens formula also gives a way to compute the $E$-dimension of the tangent space  $F_{\Dpik,\cF}^0(E[\epsilon]/\epsilon^2)$,\;see Proposition \ref{dimtanFD0}.\;

Let $\widetilde{\Dpik}_1^{k-1}$ be a fixed deformation of ${\Dpik}_1^{k-1}$ over $\cR_{E[\epsilon]/\epsilon^2,L}$.\;We view it as a class $[\widetilde{\Dpik}_1^{k-1}]$ in extension group $\ext^1_{(\varphi,\Gamma)}(\Dpik_{k-1},\Dpik_{k-1})$.\;We assume that $\widetilde{\Dpik}_1^{k-1}$ admits an $\omepik$-filtration $\widetilde{\cF}_1^{k-1}$ with parameter $$((\bx_{\pi,i})_{1\leq i\leq k-1},(\bm{\delta}_{E[\epsilon]/\epsilon^2,i})_{1\leq i\leq k-1}),\;$$
i.e.,\;$\widetilde\Dpik_{i}\hookrightarrow {\Delta_\pi}\otimes_{\,\;\cR_{E,L}}\cR_{E[\epsilon]/\epsilon^2,L}(\bm{\delta}_{E[\epsilon]/\epsilon^2,i})$ for all $1\leq i\leq k-1$.\;Let $\widetilde{\Dpik}_{k-1}=\gr_{k-1}^{\widetilde{\cF}_1^{k-1}}\widetilde{\Dpik}_1^{k-1}$ be the $(k-1)$-th graded piece of $\widetilde{\Dpik}_1^{k-1}$.\;By the proof of \cite[Proposition]{Ding2021} and (\ref{bigcommutediag}),\;we see that $\mathbf{I}'_{i}([\widetilde{\Dpik}_{k-1}])\in \ext^1_{(\varphi,\Gamma)}(\Dpik_{k-1},\Delta_{\pi,k-1})$ belongs to the image of $\hH^1_{(\varphi,\Gamma)}(\cR_{E,L})$ via the injection \[ {\mathbf{I}}''_{i}\circ  \widetilde{\mathbf{I}}'_{i}: \hH^1_{(\varphi,\Gamma)}(\cR_{E,L}) \rightarrow  \ext^1_{(\varphi,\Gamma)}(\Dpik_{k-1},\Delta_{\pi,k-1}).\;\]Therefore, we can view $\mathbf{I}'_{i}([\widetilde{\Dpik}_{k-1}])$ as an element in $\hH^1_{(\varphi,\Gamma)}(\cR_{E,L})$.\;

This section is devoted to proving the following theorem.\;The basic strategy of the proof of \cite[Theorem 2.7]{HigherLinvariantsGL3Qp} is also suitable for our case,\;but the computation is much more complicated.\;The proof consists of some computations of cohomology of the $(\varphi,\Gamma)$-modules.\;We suggest skipping the proof on the first reading.\;

\begin{thm}\label{defor1thdefor}Fix the above parameters $((\bx_{\pi,i})_{1\leq i\leq k-1},(\bm{\delta}_{E[\epsilon]/\epsilon^2,i})_{1\leq i\leq k-1})$.\;Let 
$\bm{\delta}_{E[\epsilon]/\epsilon^2,k}$ be a deformation  of  $\bm{\delta}_{\bh,k}$ over $E[\epsilon]/\epsilon^2$,\;and $\widetilde\Dpik_{k}$ be a deformation of $\Dpik_{k}$ over $\cR_{E[\epsilon]/\epsilon^2,L}$ such that \[\widetilde\Dpik_{k}\hookrightarrow {\Delta_\pi}\otimes_{\,\;\cR_{E,L}}\cR_{E[\epsilon]/\epsilon^2,L}(\bm{\delta}_{E[\epsilon]/\epsilon^2,k})\]
(so that $\widetilde\Dpik_{k}\in F^0_{\Dpik_{k}}(E[\epsilon]/\epsilon^2)$).\;Then there exists a deformation $\Dpik_{E[\epsilon]/\epsilon^2}$ of $\Dpik$ over $\cR_{E[\epsilon]/\epsilon^2,L}$,\;
which admits  an $\omepik$-filtration $\widetilde{\cF}$ with parameter $((\bx_{\pi,i})_{1\leq i\leq k},(\bm{\delta}_{E[\epsilon]/\epsilon^2,i})_{1\leq i\leq k})$ and  $\gr^{\widetilde{\cF}}_{i}\Dpik_{E[\epsilon]/\epsilon^2}=\widetilde\Dpik_i$ for $1\leq i\leq k$ if and only if
\[\mathbf{I}'_{i}([\widetilde{\Dpik}_{k-1}])\otimes_{\cR_{E[\epsilon]/\epsilon^2,L}}\cR_{E[\epsilon]/\epsilon^2,L}({\delta}_{E[\epsilon]/\epsilon^2,k}^{-1}\bm{\delta}_{\bh,k})\in \sL(\Dpik)_{(k-1)r}.\]
\end{thm}
\begin{rmk}The Colmez-Greenberg-Stevens formula (on parabolic Fontaine-Mazur simple $\sL$-invariants) shows that the parabolic simple $\sL$-invariants characterize obstructions to paraboline deformations of type $\omepik$  of $\Dpik$ over $E[\epsilon]/\epsilon^2$.\;
\end{rmk}
\begin{proof}Replacing $\Dpik_{E[\epsilon]/\epsilon^2}$ and $\widetilde{\Dpik}_1^{k-1}$ by $$\Dpik_{E[\epsilon]/\epsilon^2}\otimes_{\cR_{E[\epsilon]/\epsilon^2,L}}\cR_{E[\epsilon]/\epsilon^2,L}(\bm{\delta}_{E[\epsilon]/\epsilon^2,k}^{-1}) \text{\;and\;} \widetilde{\Dpik}_1^{k-1}\otimes_{\cR_{E[\epsilon]/\epsilon^2,L}}\cR_{E[\epsilon]/\epsilon^2,L}(\bm{\delta}_{E[\epsilon]/\epsilon^2,k}^{-1})$$ respectively,\;we can assume that ${\delta}_{E[\epsilon]/\epsilon^2,k}=\bm{\delta}_{\bh,k}$.\;By twisting by $\cR_{E,L}(\bm{\delta}_{\bh,k}^{-1})$ we can assume that $\bm{\delta}_{\bh,k}=1$ and $\bm{\delta}_{\bh,k-1}=\unr(q_L^{-1})z^{\underline{h}}$ for some $\underline{h}\in \BZ_{>0}^{|\Sigma_L|}$.\;Therefore,\;we have  injections \[\widetilde\Dpik_{k}\hookrightarrow {\Delta_\pi}\otimes_{\cR_{E,L}}\cR_{E[\epsilon]/\epsilon^2,L}\] and $\Dpik_{k}\hookrightarrow {\Delta_\pi}$ (in this case,\;$\Delta_{\pi,k}=\Delta_{\pi}$).\;We put $ \Delta_{\pi,E[\epsilon]/\epsilon^2}:={\Delta_\pi}\otimes_{\cR_{E,L}}\cR_{E[\epsilon]/\epsilon^2,L}$.\;By definition,\;we have two commutative diagrams:
\begin{align}\label{extksequence}
	\xymatrix{0 \ar[r] & \Dpik_{k} \ar[d]^{}  \ar[r] & \widetilde\Dpik_{k} \ar[d]^{} \ar[r] & \Dpik_{k} \ar[d]^{} \ar[r] & 0  \\
		0 \ar[r] & \Delta_{\pi}  \ar[r] & \Delta_{\pi,E[\epsilon]/\epsilon^2} \ar[r] & \Delta_{\pi}\ar[r] & 0
	},
\end{align}
and 
\begin{equation}\label{extksequence1}
	\xymatrix{0 \ar[r] & {\Dpik}_1^{k-1} \ar[d]^{\pr}  \ar[r] & \widetilde{\Dpik}_1^{k-1}\ar[r] \ar[d]^{\pr}  &{\Dpik}_1^{k-1} \ar[r] \ar[d]^{\pr}  & 0\\
		0\ar[r] &\Dpik_{k-1}\ar[r] &\widetilde\Dpik_{k-1}  \ar[r] &\Dpik_{k-1}\ar[r] & 0  .}
\end{equation}
Combining the first row of (\ref{extksequence}) with the first row of (\ref{extksequence1}),\;we get a commutative diagram
\[\xymatrix{
	& 0          \ar[d]^{}        & 0    \ar[d]^{}    &  0           \ar[d]^{}                              &                               & \\
	0 \ar[r] & {\Dpik}_1^{k-1}\tee\Dpik_{k}^\vee \ar[d]^{}  \ar[r] &  {\Dpik}_1^{k-1}\tee \widetilde{\Dpik}_{k}^\vee \ar[d]^{}  \ar[r]   &  {\Dpik}_1^{k-1}\tee \Dpik_{k}^\vee    \ar[d]^{}  \ar[r]  & 0\\
	0 \ar[r] & \widetilde{\Dpik}_1^{k-1}\tee \Dpik_{k}^\vee \ar[d]^{}  \ar[r] &  \widetilde{\Dpik}_1^{k-1}\tee \widetilde{\Dpik}_{k}^\vee  \ar[d]^{}  \ar[r]  &  \widetilde{\Dpik}_1^{k-1}\tee \Dpik_{k}^\vee  \ar[d]^{}  \ar[r]     & 0\\
	0 \ar[r] & {\Dpik}_1^{k-1}\tee \Dpik_{k}^\vee \ar[d]^{}  \ar[r] &  {\Dpik}_1^{k-1}\tee \widetilde{\Dpik}_{k}^\vee \ar[d]^{}  \ar[r]   &  {\Dpik}_1^{k-1}\tee \Dpik_{k}^\vee  \ar[d]^{}  \ar[r]    & 0\\
	& 0                 & 0       &  0                                        &                       &
	,}\]
which induces a surjection $\widetilde{\Dpik}_1^{k-1}\tee \widetilde{\Dpik}_{k}^\vee \rightarrow {\Dpik}_1^{k-1}\tee \Dpik_{k}^\vee$.\;Let $\Dpik_{1,2}$ be its kernel.\;We define the morphisms
\begin{equation}
	\begin{aligned}
		u:\;&{\Dpik}_1^{k-1} \rightarrow ({\Dpik}_1^{k-1}\tee \widetilde{\Dpik}_{k}^\vee)\oplus(\widetilde{\Dpik}_1^{k-1}\tee \Dpik_{k}^\vee)\\
		&a\mapsto (a',b'),
	\end{aligned}
\end{equation}
where $a'$ (resp.,\;$b'$) is the image of $a$ in $({\Dpik}_1^{k-1}\tee \widetilde{\Dpik}_{k}^\vee)$ (resp.,\;of $a$ in $(\widetilde{\Dpik}_1^{k-1}\tee \Dpik_{k}^\vee)$),\;and 
\begin{equation}
	\begin{aligned}
		v:\;&({\Dpik}_1^{k-1}\tee \widetilde{\Dpik}_{k}^\vee)\oplus(\widetilde{\Dpik}_1^{k-1}\tee \Dpik_{k}^\vee) \rightarrow \Dpik_{1,2}\\
		&(a',b')\mapsto a''-b'',
	\end{aligned}
\end{equation}
where $a''$ (resp.,\;$b''$) is the image of $a'$ (resp.,\;$b''$) in $\widetilde{\Dpik}_1^{k-1}\tee \widetilde{\Dpik}_{k}^\vee$.\;It is easy to verify that with these definitions,\;the following sequences:
\begin{equation}\label{extksequence2}
	\begin{aligned}
		&0\rightarrow \Dpik_{1,2}\rightarrow \widetilde{\Dpik}_1^{k-1}\tee \widetilde{\Dpik}_{k}^\vee \rightarrow {\Dpik}_1^{k-1}\tee \Dpik_{k}^\vee\rightarrow 0,\\
		&0\rightarrow {\Dpik}_1^{k-1}\tee \Dpik_{k}^\vee \xrightarrow{u}  ({\Dpik}_1^{k-1}\tee \widetilde{\Dpik}_{k}^\vee)\oplus(\widetilde{\Dpik}_1^{k-1}\tee \Dpik_{k}^\vee) \xrightarrow{v} \Dpik_{1,2}\rightarrow 0,\;
	\end{aligned}
\end{equation}
are exact.\;

We can get similar diagrams  and short exact sequences by using the first row of (\ref{extksequence}),\;and the second row of (\ref{extksequence1}) ,\;i.e.,\;
\begin{equation}\label{extksequence34}
	\begin{aligned}
		&0\rightarrow \Dpik_{1,2}'\rightarrow \widetilde{\Dpik}_{k-1}\tee \widetilde{\Dpik}_{k}^\vee \rightarrow {\Dpik}_{k-1}\tee \Dpik_{k}^\vee\rightarrow 0,\\
		&0\rightarrow {\Dpik}_{k-1}\tee \Dpik_{k}^\vee \rightarrow  ({\Dpik}_{k-1}\tee \widetilde{\Dpik}_{k}^\vee)\oplus(\widetilde{\Dpik}_{k-1}\tee \Dpik_{k}^\vee) \rightarrow \Dpik_{1,2}'\rightarrow 0.
	\end{aligned}
\end{equation}
Using the second row of \ref{extksequence} and the second row of \ref{extksequence1}),\;we get
\begin{equation}\label{extksequence345}
	\begin{aligned}
		&0\rightarrow \Dpik_{1,2}''\rightarrow \widetilde{\Dpik}_{k-1}\tee \Delta_{\pi,E[\epsilon]/\epsilon^2}^\vee \rightarrow {\Dpik}_{k-1}\tee \Delta_\pi^\vee\rightarrow 0,\\
		&0\rightarrow {\Dpik}_{k-1}\tee \Delta_\pi^\vee \rightarrow  ({\Dpik}_{k-1}\tee \Delta_{\pi,E[\epsilon]/\epsilon^2}^\vee)\oplus(\widetilde{\Dpik}_{k-1}\tee \Delta_\pi^\vee) \rightarrow \Dpik_{1,2}''\rightarrow 0.\;
	\end{aligned}
\end{equation}


Taking cohomology of  the first row of  (\ref{extksequence2}) and (\ref{extksequence34}),\;we get a commutative diagram 
\begin{equation}\label{cgs1}
	\xymatrix{ \ar[r] & \hH^1_{(\varphi,\Gamma)}(\Dpik_{1,2}) \ar[d]^{u_1}  \ar[r] & \ext^1_{(\varphi,\Gamma)}(\widetilde{\Dpik}_{k},\widetilde{\Dpik}_1^{k-1}) \ar[d]^{v_1}  \ar[r]^{p}  & \ext^1_{(\varphi,\Gamma)}(\Dpik_{k},\Dpik_{1}^{k-1}) \ar[d]^{w_1} \ar[r]^{c} & \hH^2_{(\varphi,\Gamma)}(\Dpik_{1,2}) \ar[d]^{u_2}\\
		\ar[r] & \hH^1_{(\varphi,\Gamma)}(\Dpik_{1,2}')   \ar[r]
		& \ext^1_{(\varphi,\Gamma)}(\widetilde{\Dpik}_{k},\widetilde{\Dpik}_{k-1})  \ar[r]^{p'} & \ext^1_{(\varphi,\Gamma)}(\Dpik_{k},\Dpik_{k-1})  \ar[r]^{c'}  & \hH^2_{(\varphi,\Gamma)}(\Dpik_{1,2}')
		.}
\end{equation}
By Lemma \ref{phigammacohononcrtialspecial} and an easy d\'{e}vissage argument,\;we have $\ext^2_{(\varphi,\Gamma)}(D_k,\Dpik_{1}^{k-2})=0$.\;We deduce that
$w_1$ is surjective,\;and $u_2':\ext^2_{(\varphi,\Gamma)}(\Dpik_{k},\Dpik_{1}^{k-1})\rightarrow \ext^2_{(\varphi,\Gamma)}(\Dpik_{k},\Dpik_{k-1})$ is an isomorphism.\;

We are going to show that 
$u_1$ is surjective and $u_2$ is an isomorphism.\;Taking cohomology of the second row of (\ref{extksequence2}) and (\ref{extksequence34}),\;we get a commutative diagram:
\begin{equation}
	\begin{aligned}
		& \xymatrix@C=2.2ex{ \ar[r] & \ext^1_{(\varphi,\Gamma)}(\Dpik_{k},\Dpik_1^{k-1}) \ar[d]^{u_1'}  \ar[r] & \ext^1_{(\varphi,\Gamma)}(\widetilde{\Dpik}_{k},{\Dpik}_1^{k-1})\oplus \ext^1_{(\varphi,\Gamma)}(\Dpik_{k},\widetilde\Dpik_{1}^{k-1}) \ar[d]^{u_1''}  \ar[r]   & \hH^1_{(\varphi,\Gamma)}(\Dpik_{1,2}) \ar[d]^{u_1}\\
			\ar[r] & \ext^1_{(\varphi,\Gamma)}(\Dpik_{k},\Dpik_{k-1})  \ar[r] & \ext^1_{(\varphi,\Gamma)}(\widetilde{\Dpik}_{k},{\Dpik}_{k-1})\oplus \ext^1_{(\varphi,\Gamma)}(\Dpik_{k},\widetilde\Dpik_{k-1})   \ar[r]   & \hH^1_{(\varphi,\Gamma)}(\Dpik_{1,2}')
	}\end{aligned}
\end{equation}
\begin{equation}
	\begin{aligned}
		& \xymatrix@C=2.2ex{ \ar[r] & \ext^2_{(\varphi,\Gamma)}(\Dpik_{k},\Dpik_1^{k-1}) \ar[d]^{u_2'}  \ar[r] & \ext^2_{(\varphi,\Gamma)}(\widetilde{\Dpik}_{k},{\Dpik}_1^{k-1})\oplus \ext^2_{(\varphi,\Gamma)}(\Dpik_{k},\widetilde\Dpik_{1}^{k-1}) \ar[d]^{u_2''}  \\
			\ar[r] & \ext^2_{(\varphi,\Gamma)}(\Dpik_{k-1},\Dpik_{k-1})   \ar[r] & \ext^2_{(\varphi,\Gamma)}(\widetilde{\Dpik}_{k},{\Dpik}_{k-1})\oplus \ext^2_{(\varphi,\Gamma)}(\Dpik_{k},\widetilde\Dpik_{k-1})   
		}\\
		& \xymatrix@C=2.2ex{   \ar[r]  & \hH^2_{(\varphi,\Gamma)}(\Dpik_{1,2}) \ar[d]^{u_2}  \ar[r] & 0\\
			\ar[r]   & \hH^2_{(\varphi,\Gamma)}(\Dpik_{1,2}') \ar[r] & 0
			.}
	\end{aligned}
\end{equation}
By Lemma \ref{phigammacohononcrtialspecial} and an easy d\'{e}vissage argument,\;we see that $u'_1$ is a surjection and $u'_2$ is an isomorphism.\;Therefore,\;we deduce that  \[\ext^2_{(\varphi,\Gamma)}(\widetilde{\Dpik}_{k},{\Dpik}_1^{k-1}) \rightarrow \ext^2_{(\varphi,\Gamma)}(\widetilde{\Dpik}_{k},{\Dpik}_{k-1}),\; \ext^2_{(\varphi,\Gamma)}(\Dpik_{k},\widetilde\Dpik_{k-1}) \rightarrow \ext^2_{(\varphi,\Gamma)}(\Dpik_{k},\widetilde\Dpik_{k-1})\] are isomorphisms by $4$-lemma.\;This show that $u_2''$ is an isomorphism.\;We deduce from the above long exact sequence that $u_2$ is an isomorphism.\;At the same time,\;we see that the morphism \[\ext^1_{(\varphi,\Gamma)}(\widetilde{\Dpik}_{k},{\Dpik}_1^{k-1}) \rightarrow \ext^1_{(\varphi,\Gamma)}(\widetilde{\Dpik}_{k},{\Dpik}_{k-1}),\;\ext^1_{(\varphi,\Gamma)}(\Dpik_{k},\widetilde\Dpik_{k-1}) \rightarrow \ext^1_{(\varphi,\Gamma)}(\Dpik_{k},\widetilde\Dpik_{k-1})\] in $u_1''$ are surjective by $4$-lemma.\;Then by an  easy d\'{e}vissage argument,\;we deduce that  $u_1$ is surjective.\;We conclude that $v_1$ is a surjection by an easy diagram chasing (see (\ref{cgs1})).\;

Note that  $\Dpik_{E[\epsilon]/\epsilon^2}$ admits  a $\omepik$-filtration with parameter $((\bx_{\pi,})_{1\leq i\leq k},(\bm{\delta}_{E[\epsilon]/\epsilon^2,i})_{1\leq i\leq k})$ such that $\gr^{\widetilde{\cF}}_{i}\Dpik_{E[\epsilon]/\epsilon^2}=\widetilde\Dpik_i$ if and only if  $[\Dpik]\in \ext^1_{(\varphi,\Gamma)}(\Dpik_{k},{\Dpik}_1^{k-1})$ lies in $\mathrm{Im}(p)$ (see (\ref{cgs1})), i.e., $c([\Dpik])=0$. The above discussion implies that $\Dpik_{E[\epsilon]/\epsilon^2}$ admits the desired $\omepik$-filtration if and only if $c'([\Dpik^k_{k-1}])=0$, where 
$[\Dpik^k_{k-1}]=w_1([\Dpik])\in  \ext^1_{(\varphi,\Gamma)}(\Dpik_k,\Dpik_{k-1})$.\;

We now translate the above discussion to the side of $p$-adic differential equations.\;By the first rows of (\ref{extksequence34}) and first rows of (\ref{extksequence345}),\;we have the following commutative diagram:
\begin{equation}\label{D12D12}
	\xymatrix@C=2.2ex{ \ar[r] & \hH^1_{(\varphi,\Gamma)}(\Dpik_{1,2}'') \ar[d]^{x_1}  \ar[r] & \ext^1_{(\varphi,\Gamma)}(\Delta_{\pi,E[\epsilon]/\epsilon^2},\widetilde{\Dpik}_{k-1}) \ar[d]^{v_1'}  \ar[r]^{\;\;\;\;\;p''}  & \ext^1_{(\varphi,\Gamma)}({\Delta_\pi},\Dpik_{k-1}) \ar[d]^{w_1'(\text{i.e.,\;}\mathbf{I}_{k}^{1,1})}_{\simeq} \ar[r]^{\;\;\;\;c''} & \hH^2_{(\varphi,\Gamma)}(\Dpik_{1,2}'') \ar[d]^{x_2}\\
		\ar[r] & \hH^1_{(\varphi,\Gamma)}(\Dpik_{1,2}')   \ar[r]
		& \ext^1_{(\varphi,\Gamma)}(\widetilde{\Dpik}_{k},\widetilde{\Dpik}_{k-1})  \ar[r]^{p'}  & \ext^1_{(\varphi,\Gamma)}(\Dpik_{k},\Dpik_{k-1})  \ar[r]^{\;\;\;\;c'}  & \hH^2_{(\varphi,\Gamma)}(\Dpik_{1,2}')
		.}
\end{equation}
By Lemma \ref{pullbackforiso},\;we get that $w'_1$ is an isomorphism.\;

Replacing $\Dpik_{k}^\vee$ (resp.\;$\widetilde{\Dpik}_{k}^\vee$) in  (\ref{extksequence1}) and (\ref{extksequence2}) by $\Delta_\pi^\vee$ (resp.\;$\Delta_{\pi,E[\epsilon]/\epsilon^2}^\vee$),\;we get the following diagram,\;
\[\xymatrix{0 \ar[r] & \Delta_\pi^\vee\tee \Dpik_{k-1}\ar@{_(->}[d]  \ar[r] & \Delta_\pi^\vee\tee \widetilde{\Dpik}_{k-1}\ar[r] \ar@{_(->}[d]  & \Delta_\pi^\vee\tee \Dpik_{k-1}\ar[r] \ar@{_(->}[d]  & 0\\
	0 \ar[r] & \Dpik_{1,2}'' \ar[r] & \Delta_{\pi,E[\epsilon]/\epsilon^2}^\vee\tee \widetilde{\Dpik}_{k-1}  \ar[r] & \Delta_\pi^\vee\tee \Dpik_{k-1} \ar[r] & 0  .}\]
Combining the commutative diagram (\ref{D12D12}) with the cohomology of the above diagram,\;we get
\begin{equation}
	\begin{aligned}
		&\xymatrix{
			\ar[r] &  \ext^1_{(\varphi,\Gamma)}(\Delta_{\pi},\Dpik_{k-1})  \ar[r] \ar[d]^{x_1'}
			& \ext^1_{(\varphi,\Gamma)}(\Delta_{\pi},\widetilde{\Dpik}_{k-1})   \ar[d]^{v_1''}_{\simeq} \\
			\ar[r] &\hH^1_{(\varphi,\Gamma)}(\Dpik_{1,2}'')   \ar[r] \ar[d]^{x_1}
			& \ext^1_{(\varphi,\Gamma)}(\Delta_{\pi,E[\epsilon]/\epsilon^2},\widetilde{\Dpik}_{k-1})   \ar[d]^{v_1'} \\
			\ar[r] & \hH^1_{(\varphi,\Gamma)}(\Dpik_{1,2}')   \ar[r]
			& \ext^1_{(\varphi,\Gamma)}(\widetilde{\Dpik}_{k},\widetilde{\Dpik}_{k-1}) }\\
		&\hspace{100pt}\xymatrix{  \ar[r]^{p'''\hspace{40pt}}  & \ext^1_{(\varphi,\Gamma)}({\Delta_\pi},\Dpik_{k-1})  \ar[r]^{c'''} \ar@{=}[d]^{w_1''}_{\simeq}  &\ext^2_{(\varphi,\Gamma)}(\Delta_{\pi},\Dpik_{k-1}) \ar[d]^{x_2'}\\
			\ar[r]^{p''\hspace{40pt}}  & \ext^1_{(\varphi,\Gamma)}({\Delta_\pi},\Dpik_{k-1})  \ar[r]^{c''} \ar[d]^{w_1'}_{\simeq} &\hH^2_{(\varphi,\Gamma)}(\Dpik_{1,2}'') \ar[d]^{x_2}\\
			\ar[r]^{p'\hspace{40pt}}  & \ext^1_{(\varphi,\Gamma)}(\Dpik_{k},\Dpik_{k-1})  \ar[r]^{c'}  & \hH^2_{(\varphi,\Gamma)}(\Dpik_{1,2}').}
	\end{aligned}
\end{equation}
\text{\textbf{Claim.}} The composition $x_3=x_2\circ x_2':\ext^2_{(\varphi,\Gamma)}({\Delta_\pi},\Dpik_{k-1})\rightarrow\ext^2_{(\varphi,\Gamma)}(\Dpik_{1,2}') $ is injective.\;\\
\text{\textbf{Proof of the claim.}}
By the second row of  (\ref{extksequence34}) and second row of  (\ref{extksequence345}) and a diagram chasing,\;we get the following commutative diagram:
\begin{equation}
		\xymatrix{ \ar[r] & \ext^2_{(\varphi,\Gamma)}({\Delta_\pi},\Dpik_{k-1}) \ar@{=}[d]  \ar[r] & \ext^2_{(\varphi,\Gamma)}({\Delta_\pi},\widetilde{\Dpik}_{k-1}) \ar[d] \\
			\ar[r] & \ext^2_{(\varphi,\Gamma)}({\Delta_\pi},\Dpik_{k-1}) \ar@{=}[d]^{\simeq}_{(\ast\ast)} \ar[r]
			& \ext^2_{(\varphi,\Gamma)}({\Delta_\pi},\widetilde{\Dpik}_{k-1}) \oplus\ext^2_{(\varphi,\Gamma)}(\Delta_{\pi,E[\epsilon]/\epsilon^2},\Dpik_{k-1}) \ar[d]  \\
		\ar[r] & {\ext}^2_{(\varphi,\Gamma)}({\Dpik}_{k},\Dpik_{k-1}) \ar[r]			& {\ext}^2_{(\varphi,\Gamma)}({\Dpik}_{k},\widetilde{\Dpik}_{k-1})\oplus{\ext}^2_{(\varphi,\Gamma)}(\widetilde{\mathbf{D}}_{k},\Dpik_{k-1})}	
\end{equation}
\begin{equation}
		\hspace{250pt}	\xymatrix{  \ar[r] &\ext^2_{(\varphi,\Gamma)}({\Delta_\pi},\Dpik_{k-1}) \ar[d]^{x_2} \ar[r] &0 \\
			\ar[r]  & \hH^2_{(\varphi,\Gamma)}(\Dpik_{1,2}'') \ar[d]^{x_2'} \ar[r] &0 \\
			\ar[r]  & \hH^2_{(\varphi,\Gamma)}(\Dpik_{1,2}')\ar[r] &0 
			.}
\end{equation}		
The isomorphism in $(\ast\ast)$ follows from Lemma \ref{pullbackforiso}.\;By Lemma \ref{pullbackforiso},\;an easy d\'{e}vissage argument and $4$-Lemma,\;we see that the morphisms 
\begin{equation}
	\begin{aligned}
		&\ext^2_{(\varphi,\Gamma)}({\Delta_\pi},\widetilde{\Dpik}_{k-1}) \rightarrow \ext^2_{(\varphi,\Gamma)}(\Dpik_{k},\widetilde{\Dpik}_{k-1}),\;\\
		&\ext^2_{(\varphi,\Gamma)}(\Delta_{\pi,E[\epsilon]/\epsilon^2},\Dpik_{k-1})  \rightarrow \ext^2_{(\varphi,\Gamma)}(\Dpik_{k},\widetilde{\Dpik}_{k-1})
	\end{aligned}
	\end{equation}
are isomorphisms. This implies that $x_3$ is an injection by an easy  diagram chasing.\;This completes the proof of the claim.\;

Therefore,\;we conclude that $[\Dpik^k_{k-1}]\in \ext^1_{(\varphi,\Gamma)}(\Dpik_k,\Dpik_{k-1})$ lies in $\mathrm{Im}(p')$ if and only if the class $c'''((w_1'\circ w_1'')^{-1}([\Dpik^k_{k-1}]))=0$.\;Recall that 
$$\mathbf{I}'_{i}([\widetilde{\Dpik}_{k-1}])\in \ext^1_{(\varphi,\Gamma)}(\Dpik_{k-1},\Delta_{\pi}\tee \cR_{E,L}(\bm{\delta}_{\bh,k-1}))$$ belongs to the image of $\hH^1_{(\varphi,\Gamma)}(\cR_{E,L})$ via the injection $$ {\mathbf{I}}''_{i}\circ  \widetilde{\mathbf{I}}'_{i}: \hH^1_{(\varphi,\Gamma)}(\cR_{E,L}) \rightarrow  \ext^1_{(\varphi,\Gamma)}(\Dpik_{k-1},\Delta_{\pi}\tee \cR_{E,L}(\bm{\delta}_{\bh,k-1})).\;$$
By the naturality of cup products,\;we deduce
\[\xymatrix@C=3.2ex{
	\ext^0_{(\varphi,\Gamma)}(\Delta_{\pi},\Delta_{\pi}) \ar[d]_{c_0}  & \times \hspace{10pt}\ext^1_{(\varphi,\Gamma)}(\Delta_{\pi},\Dpik_{k-1}) \ar@{=}[d]  \ar[r]^{\cup} & \hspace{10pt}\ext^1_{(\varphi,\Gamma)}(\Delta_{\pi},\Dpik_{k-1}) \ar[d]_{c'''}\\
	\ext^1_{(\varphi,\Gamma)}(\Delta_{\pi},\Delta_{\pi})   & \times \hspace{10pt} \ext^1_{(\varphi,\Gamma)}(\Delta_{\pi},\Dpik_{k-1})  \ar[r]^{\cup}_{\hspace{5pt}\text{perfect}} & \hspace{10pt}\ext^2_{(\varphi,\Gamma)}(\Delta_{\pi},\Dpik_{k-1}) .}\]
This shows that the connection map $c'''$  is given by the cup product $\langle [\widetilde{\Dpik}_{k-1}],-\rangle$,\;i.e,\;
\[\langle-,-\rangle:\ext^1_{(\varphi,\Gamma)}(\Delta_{\pi},\Dpik_{k-1})\times \ext^1_{(\varphi,\Gamma)}({\Delta_\pi},{\Delta_\pi})\xrightarrow{\cup}\ext^2_{(\varphi,\Gamma)}(\Delta_{\pi},\Dpik_{k-1})\cong E.\;\]
So we have $c'''((w_1'\circ w_1'')^{-1}([\Dpik^k_{k-1}]))=0$ if and only if $\langle(w_1'\circ w_1'')^{-1}([\Dpik^k_{k-1}]),\mathbf{I}'_{i}([\widetilde{\Dpik}_{k-1}])\rangle=0$.\;Then the conclusion  follows from the definition of parabolic Fontaine-Mazur simple $\sL$-invariants.\;
\end{proof}

As a corollary of Theorem \ref{defor1thdefor},\;we deduce

\begin{cor}\label{corlinvaraint}The map $\kappa:F_{\Dpik,\cF}^0(E[\epsilon]/\epsilon^2)\longrightarrow \prod_{ir\in \Delta_n(k)}\homo(L^\times,E)$  (see (\ref{kappakappaL})) factors through a surjective map
$$\kappa:F_{\Dpik,\cF}^0(E[\epsilon]/\epsilon^2)\twoheadlongrightarrow \sL(\Dpik)=\prod_{sr\in \Delta_n(k)}\sL(\Dpik)_{sr}.$$
\end{cor}

We end this section with a computation of the $E$-dimension of the tangent space $F_{\Dpik,\cF}^0(E[\epsilon]/\epsilon^2)$. Unlike the generic case \cite[Proposition 4.1.16]{Ding2021},\;the corollary below will show that the morphism $\Upsilon^0$ in (\ref{FmorphismGRD}) is not formally smooth,\;so that the \cite[Proposition 4.1.17]{Ding2021} is not suitable for us.\;Instead,\;we will use Lemma \ref{fibredim}.\;We need some preliminaries.\;See (\ref{kappakappaL}) for the definition of $\Upsilon^0(E[\epsilon]/\epsilon^2)$,\;$\omega_0$ and $\kappa_L$.\;
%We put (see \ref{kappakappaL})
%\begin{equation}
%	\omega=\kappa_L\circ\omega^0:	\prod_{i=1}^kF_{\gr_i^{\cF}\Dpik}^0(E[\epsilon]/\epsilon^2)\longrightarrow \prod_{ir\in \Delta_n(k)}\homo(L^\times,E).\;
%\end{equation}
\begin{cor}\label{corlinvaraint1}The image of the map $$\Upsilon^0(E[\epsilon]/\epsilon^2):F_{\Dpik,\cF}^0(E[\epsilon]/\epsilon^2)\longrightarrow \prod_{i=1}^kF_{\gr_i^{\cF}\Dpik}^0(E[\epsilon]/\epsilon^2)$$ is given by $(\kappa_L\circ\omega^0)^{-1}(\sL(\Dpik))$.\;In particular,\;we have
$\dim_E\im\Upsilon^0(E[\epsilon]/\epsilon^2)=1+d_L(k+\frac{n(r-1)}{2})$ and  $\dim_E\mathrm{coker}(\Upsilon^0(E[\epsilon]/\epsilon^2))=k-1.$
\end{cor}
\begin{proof}By Corollary \ref{corlinvaraint},\;we see that $\im\Upsilon^0(E[\epsilon]/\epsilon^2)$ is contained in $(\kappa_L\circ\omega^0)^{-1}(\sL(\Dpik))$. Since $\kappa_L$ is a surjection with kernel $\homo(L^\times,E)$,\;it remains to show that $$\im\Upsilon^0(E[\epsilon]/\epsilon^2)=(\omega^0)^{-1}(\kappa_L^{-1}(\sL(\Dpik))).\;$$For any $(\psi_i)_{1\leq i\leq k}\in \kappa_L^{-1}(\sL(\Dpik))$,\;and $(\widetilde\Dpik_i)_{1\leq i\leq k}\in(\omega^0)^{-1}((\psi_i)_{1\leq i\leq k}))\subseteq \prod_{i=1}^kF_{\gr_i^{\cF}\Dpik}^0(E[\epsilon]/\epsilon^2)$, we need to show that there exists a deformation $\Dpik_{E[\epsilon]/\epsilon^2}$ of $\Dpik$ over $\cR_{E[\epsilon]/\epsilon^2,L}$ which admits  an $\omepik$-filtration $\widetilde{\cF}$ with parameter $((\bx_{\pi,i})_{1\leq i\leq k},(\bm{\delta}_{E[\epsilon]/\epsilon^2,i})_{1\leq i\leq k})$ and  $\gr^{\widetilde{\cF}}_{i}\Dpik_{E[\epsilon]/\epsilon^2}=\widetilde\Dpik_i$ for $1\leq i\leq k$.\;This can be done by using Theorem \ref{defor1thdefor} step by step.\;By the proof of \cite[Proposition 4.1.4]{Ding2021},\;we see that 
\[\dim_E(\omega^0)^{-1}((\psi_i)_{1\leq i\leq k}))=k\cdot d_L\frac{r(r-1)}{2}= d_L\frac{n(r-1)}{2}.\]
We see that $\dim_E\im\Upsilon^0(E[\epsilon]/\epsilon^2)= d_L\frac{n(r-1)}{2}+\dim_E\homo(L^\times,E)+\dim_E\sL(\Dpik)=1+d_L(k+\frac{n(r-1)}{2})$.
\end{proof}

In the sequel,\;we assume that $\cF$ is furthermore a non-critical special  $\omepik$-filtration on $\Dpik$ (see Definition \ref{dfnnoncriticalspecial}),\;i.e.,\;we further assume that $\Dpik^{i+1}_{i}$ is non-split for all $ir\in\Delta_n(k)$.\;We put
\[\EndO_{\cF}(\Dpik,\Dpik):=\{f\in \EndO(\Dpik,\Dpik)\;|\;f(\fil_i^{\cF}\Dpik)\subseteq \fil_i^{\cF}\Dpik\},\]
which is a saturated $(\varphi,\Gamma)$-submodule of $\EndO(D,D)$.\;For any $f\in \EndO_{\cF}(\Dpik,\Dpik)$,\;we see that $f(\Dpik_i)\subset \Dpik_i$ for all $1\leq i\leq k$.\;Thus $f$ produces an element $(f_i)_{1\leq i\leq k}\in \prod_{i=1}^k\EndO(\Dpik_i)$.\;

The following proposition computes the cohomology of $\EndO_{\cF}(\Dpik,\Dpik)$,\;and gives some properties of the functor  $F_{\Dpik,\cF}$.\;

\begin{pro}\label{dimlemmaDpik}
(1) We have $\hH^0_{(\varphi,\Gamma)}(\EndO(\Dpik,\Dpik))\cong\hH^0_{(\varphi,\Gamma)}(\EndO_{\cF}(\Dpik,\Dpik))\cong E$.\;\\
(2) We have  $\hH^2_{(\varphi,\Gamma)}(\EndO_{\cF}(\Dpik,\Dpik))=0$.\;\\
(3) The functor  $F_{\Dpik,\cF}$ is pro-representable and is formally smooth over $E$ of dimension  $1+d_L\frac{n(n+r)}{2}$.\;
\end{pro}
\begin{proof}For Part (1),\;it suffices to show that $\hH^0_{(\varphi,\Gamma)}(\EndO(\Dpik,\Dpik))\cong E$.\;We have an exact sequence of $(\varphi,\Gamma)$-modules over $\cR_{E,L}$:
\[0\rightarrow \homo_{\cR_{E,L}}(\Dpik_{k},\Dpik)\rightarrow \EndO(\Dpik,\Dpik)\rightarrow \homo_{\cR_{E,L}}(\bD^{k-1}_1,\Dpik) \rightarrow0.\]
Taking cohomology,\;we get
\[0\rightarrow \homo_{(\varphi,\Gamma)}(\Dpik_{k},\Dpik)\rightarrow  \hH^0_{(\varphi,\Gamma)}(\EndO(\Dpik,\Dpik)) \rightarrow  \homo_{(\varphi,\Gamma)}(\bD^{k-1}_1,\Dpik).\]
Since $\Dpik^k_{k-1}$ is nonsplit,\;we deduce $\homo_{(\varphi,\Gamma)}(\Dpik_{k},\Dpik^k_{k-1})=0$.\;By Lemma \ref{phigammacohononcrtialspecial} and an easy d\'{e}vissage argument,\;we see that $\homo_{(\varphi,\Gamma)}(\Dpik_{k},\Dpik^k_{k-1})\xrightarrow{\sim}\homo_{(\varphi,\Gamma)}(\Dpik_{k},\Dpik)$ and then  $\homo_{(\varphi,\Gamma)}(\Dpik_{k},\Dpik)=0$. Moreover,\;by Lemma \ref{phigammacohononcrtialspecial} and an easy d\'{e}vissage argument,\;we see that $\homo_{(\varphi,\Gamma)}(\bD^{k-1}_1,\Dpik_{k})=0$.\;We then see that $$\homo_{(\varphi,\Gamma)}(\bD^{k-1}_1,\Dpik)\cong \hH^0_{(\varphi,\Gamma)}(\EndO(\bD^{k-1}_1,\bD^{k-1}_1)).$$ Then Part (1) follows by induction on $k$.\;By \cite[Proposition 3.4]{chenevier2011infinite},\;we deduce that $F_{\Dpik,\cF}$ is pro-representable. For Part (2),\;we have a natural exact sequence of $(\varphi,\Gamma)$-modules over $\cR_{E,L}$:
\[0\rightarrow \homo_{\cR_{E,L}}(\Dpik_k,\Dpik)\rightarrow \EndO_{\cF}(\Dpik,\Dpik)\rightarrow \EndO_{\cF}(\Dpik^{k-1}_1,\Dpik^{k-1}_1) \rightarrow0,\;\]
where we also use $\cF$ to denote the induced $\Omega_{[1,k-1]}$-filtration on $\Dpik^{k-1}_1$.\;Applying the same strategy to $\Dpik_1^s$ for $1\leq s\leq k-1$,\;we conclude that the $(\varphi,\Gamma)$-module $\EndO_{\cF}(\Dpik,\Dpik)$ is isomorphic to an extension of $(\varphi,\Gamma)$-modules $\homo_{\cR_{E,L}}(\Dpik_{s},\Dpik^{s}_1)$ for $s=1,\cdots,k$.\;Therefore,\;$\EndO_{\cF}(\Dpik,\Dpik)$  is a $(\varphi,\Gamma)$-module over $\cR_{E,L}$ of rank $\frac{k(k+1)}{2}r^2$.\;By Tate duality,\;Lemma \ref{phigammacohononcrtialspecial} and an easy d\'{e}vissage argument,\;we have for  $2\leq s\leq k$ the isomorphisms:
\begin{equation*}
	\begin{aligned}
		\ext^2_{(\varphi,\Gamma)}(\Dpik_{s},\Dpik^{s}_1)&\cong \hH^0_{(\varphi,\Gamma)}(\Dpik^{s}_1,\Dpik_{s}\otimes_{\cR_{E,L}}\cR_{E,L}(\ccyc))\cong \hH^0_{(\varphi,\Gamma)}(\Dpik^{s}_{s-1},\Dpik_{s}\otimes_{\cR_{E,L}}\cR_{E,L}(\ccyc))\\
		&\cong \hH^2_{(\varphi,\Gamma)}(\Dpik_{s},\Dpik^{s}_{s-1}),
	\end{aligned}
\end{equation*}
Let $M_1:=\EndO(\Dpik_1)\otimes_{\cR_{E,L}}\cR_{E,L}(\ccyc)$ and let $M_s:=(\Dpik^{s}_{s-1})^\vee\otimes_{\cR_{E,L}}\Dpik_{s}\otimes_{\cR_{E,L}}\cR_{E,L}(\ccyc)$ for $2\leq s\leq k$,\;which is isomorphic to a nonsplit extension of $\EndO(\Dpik_{s-1},\Dpik_{s})\otimes_{\cR_{E,L}}\cR_{E,L}( \ccyc)$ by $\EndO(\Dpik_{s})\otimes_{\cR_{E,L}}\cR_{E,L}(\ccyc)$.\;We are going to show that $\hH^0_{(\varphi,\Gamma)}(M_s)=0$ for $s=1,\cdots,k$.\;By an easy variation of the proof of Lemma \ref{isoI2i},\;we get $\hH^0_{(\varphi,\Gamma)}(M_1)=0$.\;For $2\leq s\leq k$,\;if $\hH^0_{(\varphi,\Gamma)}(M_s)\neq 0$,\;then there exists a nonzero injection of $(\varphi,\Gamma)$-modules $j:\cR_{E,L}\rightarrow M_s$.\;Since 
\[\hH^0_{(\varphi,\Gamma)}(\EndO(\Dpik_s)\otimes_{\cR_{E,L}}\cR_{E,L}(\ccyc)[\frac{1}{t}])=\hH^0_{(\varphi,\Gamma)}(\EndO({\Delta_\pi})\otimes_{\cR_{E,L}}\cR_{E,L}(\ccyc)[\frac{1}{t}])=0\;\]we see that $\mathrm{Im}(j)[\frac{1}{t}]\cap M_s$ is a saturated  $(\varphi,\Gamma)$-submodule of $\EndO(\Dpik_{s-1},\Dpik_{s})\otimes_{\cR_{E,L}}\cR_{E,L}(\ccyc)$. By Tate duality and Lemma \ref{pullbackforiso},\;we have
$$\hH^0_{(\varphi,\Gamma)}(\EndO(\Dpik_{s-1},\Dpik_{s})\otimes_{\cR_{E,L}}\cR_{E,L}( \ccyc))=\hH^0_{(\varphi,\Gamma)}(\EndO({\Delta_\pi})\otimes_{\cR_{E,L}}\cR_{E,L}( \bm{\delta}_{\bh,s-1}^{-1} \bm{\delta}_{\bh,s}\ccyc)).$$
Since 
$\hH^0_{(\varphi,\Gamma)}(\EndO^0({\Delta_\pi})\otimes_{\cR_{E,L}}\cR_{E,L}( \bm{\delta}_{\bh,s-1}^{-1}\bm{\delta}_{\bh,s}\ccyc))=0$,\;by comparing Hodge-Tate weights,\;we see that 
$\mathrm{Im}(j)[1/t]\cap M_s\cong \cR_{E,L}( \bm{\delta}_{\bh,s-1}^{-1} \bm{\delta}_{\bh,s}\ccyc)$.\;This show that
$$\homo_{(\varphi,\Gamma)}(\EndO(\Dpik_{s-1},\Dpik_{s})\otimes_{\cR_{E,L}}\cR_{E,L}( \ccyc),M_s)\neq 0,\;$$ which leads a contradiction to the fact that $M_s$ is non-split.\;For Part (3),\;we deduce from Part (1) and \cite[Proposition 3.4]{chenevier2011infinite} that $F_{\Dpik,\cF}$ is pro-representable.\;By Part (2) and \cite[Proposition 4.1.15 (2)]{Ding2021},\;we deduce that $F_{\Dpik,\cF}$ is formally smooth of dimension $1+\frac{k(k+1)}{2}d_Lr^2$.\
\end{proof}

Recall that
\[F_{\Dpik,\cF}^0=\sF_{\Dpik,\cF}\times_{\prod_{i=1}^k F_{\Dpik_i}}\prod_{i=1}^k F_{\Dpik_i}^0.\]
To evaluate $\dim_EF_{\Dpik,\cF}^0(E[\epsilon]/\epsilon^2)$,\;we  need a dimension formula.\;An easy diagram chasing shows that



\begin{lem}\label{fibredim}Let $V,V'$ and $V''$ be vector spaces over $E$,\;and $u:V\rightarrow V''$ (resp.,\;$u':V'\rightarrow V''$) be $E$-linear homomorphism.\;Then we have 
\[\dim_E V\times_{V''}V':=\dim_EV-\dim_E\im u+\dim_E(u')^{-1}(\im u\cap\im u'),\]
where $V\times_{V''}V'=\{(v,v')\in V\times V':u(v)=u'(v')\in V''\}$.\;Moreover,\;\[\dim_E V\times_{V''}V'=\dim_EV+\dim_EV'-\dim_EV''\] if and only if $\dim_E\mathrm{coker}(u)=\dim_E\mathrm{coker}(V\times_{V''}V'\rightarrow V')$.\;In particular,\;if one of $u,u'$ is surjective,\;we have $\dim_E V\times_{V''}V'=\dim_EV+\dim_EV'-\dim_EV''$.\;
\end{lem}

Suggested by Lemma  \ref{fibredim} and Corollary \ref{corlinvaraint1},\;we prove (as  a Corollary of Lemma \ref{dimlemmaDpik}):
\begin{cor}\label{corlinvaraint2}The cokernel of the map $\Upsilon:F_{\Dpik,\cF}(E[\epsilon]/\epsilon^2)\longrightarrow \prod_{i=1}^kF_{\Dpik_i}(E[\epsilon]/\epsilon^2)$ has $E$-dimensional $k-1$.\;
\end{cor}
\begin{proof}By induction on $k$,\;we get the following commutative diagram of $(\varphi,\Gamma)$-modules over $\cR_{E,L}$:,\;
\[\xymatrix{0 \ar[r] & \homo_{\cR_{E,L}}(\Dpik_k,\Dpik)\ar[r] \ar[d] & \EndO_{\cF}(\Dpik)\ar[r] \ar[d] & \EndO_{\cF}(\Dpik^{k-1}_1) \ar[r]\ar[d]  &0\\
	0 \ar[r] & \EndO(\Dpik_k) \ar[r] & \prod_{i=1}^k\EndO(\Dpik_i)  \ar[r] & \prod_{i=1}^{k-1}\EndO(\Dpik_i) \ar[r] & 0  .}\]
Taking cohomology,\;we deduce from Proposition \ref{dimlemmaDpik} (and the proof of Proposition \ref{dimlemmaDpik}  (2)) the following long exact sequences
\[\xymatrix{ 
	0 \ar[r] & \ext^0_{(\varphi,\Gamma)}(\Dpik_k,\Dpik)=0 \ar[r] \ar[d] &  	\ext^0_{(\varphi,\Gamma)}(\EndO_{\cF}(\Dpik))\cong E \ar[d]   \\
	0 \ar[r] & \ext^0_{(\varphi,\Gamma)}(\Dpik_k,\Dpik_k)\cong E \ar[r] & \prod\limits_{i=1}^{k}\hH^0_{(\varphi,\Gamma)}( \EndO(\Dpik_i))\cong E^{k} }\]
\begin{equation}\label{corgamma}
	\begin{aligned}
		&\hspace{30pt}\xymatrix{ 
			\ar[r]^{\sim\hspace{60pt}} &\ext^0_{(\varphi,\Gamma)}(\EndO_{\cF}(\Dpik^{k-1}_1))\cong E \ar[r] \ar[d] &  \ext^1_{(\varphi,\Gamma)}(\Dpik_k,\Dpik) \ar[d]^{\Upsilon''}  \ar[r] &  F_{\Dpik,\cF}(E[\epsilon]/\epsilon^2) \ar[d]^{\Upsilon=\Upsilon_k}  \\
			\ar[r] &\prod\limits_{i=1}^{k-1}\hH^0_{(\varphi,\Gamma)}( \EndO(\Dpik_i))\cong E^{k-1} \ar[r] & \ext^1_{(\varphi,\Gamma)}(\Dpik_k,\Dpik_k) \ar[r] & \prod\limits_{i=1}^{k}F_{\Dpik_i}(E[\epsilon]/\epsilon^2)  }
	\end{aligned}
\end{equation}	
\[\hspace{130pt}\xymatrix{
	\ar[r]& F_{\Dpik^k_1,\cF}(E[\epsilon]/\epsilon^2) \ar[r] \ar[d]^{\Upsilon_{k-1}} &  \ext^2_{(\varphi,\Gamma)}(\Dpik_k,\Dpik)=0 \ar[d]  \\
	\ar[r]&\prod\limits_{i=1}^{k-1}F_{\Dpik_i}(E[\epsilon]/\epsilon^2) \ar[r] & \ext^2_{\Rep_{\bB_{\dr}}(\gal_L)}(\Dpik_k,\Dpik_k)=0 .}\]

By comparing dimensions,\;we see that the middle three terms of $1$-th extension groups (of two rows of (\ref{corgamma})) form two short exact sequences.\;An easy d\'{e}vissage argument shows that $\Upsilon''$ is injective and $\dim_E\mathrm{coker}\Upsilon''=1$.\;Then the snake lemma implies 
\[0\rightarrow \ker\Upsilon_k \rightarrow \ker\Upsilon_{k-1}\rightarrow \mathrm{coker}\Upsilon''\rightarrow \mathrm{coker}\Upsilon_k \rightarrow \mathrm{coker}\Upsilon_{k-1}\rightarrow 0.\]
Note that $\Upsilon_1$ is an isomorphism (so that $\ker\Upsilon_1=\mathrm{coker}\Upsilon_1=0$).\;By an induction on $k$,\;we see that $\ker\Upsilon_k\cong \ker\Upsilon_{k-1}\cong\cdots\cong \ker\Upsilon_1=0$ and $\dim_E\mathrm{coker}\Upsilon_k=1+\dim_E\mathrm{coker}\Upsilon_{k-1}=\cdots=k-1$.\;This completes the proof.\;
\end{proof}


\begin{pro}\label{dimtanFD0}	The functors $F_{\Dpik,\cF}^0$ is  pro-representable.\;Its tangent space has $E$-dimension $$\dim_EF_{\Dpik,\cF}^0(E[\epsilon]/\epsilon^2)=1+d_L\Big(k+\frac{n(n-1)}{2}\Big).$$
\end{pro}
\begin{proof}By \cite[Proposition 4.1.3]{Ding2021},\;we get that the functor $\prod_{i=1}^k F_{\gr_iD}^0$ is relatively representable over $\prod_{i=1}^k F_{\Dpik_i}$.\;By Proposition \ref{dimlemmaDpik} (3),\;the functor $F_{\Dpik,\cF}^0$ is pro-representable.\;By \cite[Proposition 4.1.4]{Ding2021},\;we see that $\prod_{i=1}^k F_{\Dpik_i}^0$ is formally smooth of dimension $k\big(1+d_L\big(1+\frac{r(r-1)}{2}\big)\big)$.\;It follows from the Lemma \ref{fibredim},\;Corollary \ref{corlinvaraint1},\;Proposition \ref{dimlemmaDpik},\;\cite[Proposition 4.1.4]{Ding2021} and Corollary \ref{corlinvaraint2} that 
\begin{equation}
	\begin{aligned}
		&\dim_EF_{\Dpik,\cF}^0(E[\epsilon]/\epsilon^2)\\
		=\;&\dim_E\sF_{\Dpik,\cF}(E[\epsilon]/\epsilon^2)-\sum_{i=1}^k\dim_EF_{\Dpik_i}(E[\epsilon]/\epsilon^2)+\sum_{i=1}^k\dim_EF^0_{\Dpik_i}(E[\epsilon]/\epsilon^2)\\
		=\;&1+d_L\frac{n(n+r)}{2}-k(1+d_Lr^2)+k\Big(1+d_L\Big(1+\frac{r(r-1)}{2}\Big)\Big),\\
		=\;&1+d_L\Big(k+\frac{n(n-1)}{2}\Big).\\
\end{aligned}\end{equation}
The result follows.\;
\end{proof}
\begin{rmk}We may ask if $F_{\Dpik,\cF}^0$ is formally smooth.\;The author does not know whether it is true.\;
\end{rmk}



\section{Summary of certain locally \texorpdfstring{$\bQ_p$}{Lg}-analytic representations}\label{BreuilLINVSUMMMANY}

Let $\ul{\lambda}:=(\lambda_{1,\sigma}, \cdots, \lambda_{n,\sigma})_{\sigma\in \Sigma_L}$ be a weight of $\ft_{\Sigma_L}$.\;For $I\subseteq \Delta_n=\{1,\cdots, n-1\}$, we call that $\ul{\lambda}$ is $I$-dominant with respect to $\bB_{/E}$ (resp. with respect to $\overline{\bB}_{/E}$) if $\lambda_{i,\sigma}\geq \lambda_{i+1,\sigma}$ \big(resp. $\lambda_{i,\sigma} \leq \lambda_{i+1,\sigma}$\big) for all $i\in I$ and $\sigma\in \Sigma_L$. We denote by $X_{I}^+$ (resp. $X_{I}^-$) the set of $I$-dominant integral weights  of $\ft_{\Sigma_L}$ with respect to $\bB_{/E}$ (resp. with respect to $\overline{\bB}_{/E}$). Note that $\ul{\lambda}\in X_{I}^+$ if and only if $-\ul{\lambda}\in X_{I}^-$. For $\ul{\lambda}\in X_I^+$, there exists a unique irreducible algebraic representation, denoted by $L(\ul{\lambda})_I$, of $(\bL_I)_{/E}$ with highest weight $\ul{\lambda}$ with respect to $(\bL_I)_{/E}\cap \bB_{/E}$.\;We put $\overline{L}(-\ul{\lambda})_I:=L(\ul{\lambda})^\vee_I$,\;which is an irreducible algebraic representation of $(\bL_I)_{/E}$ with highest weight $-\ul{\lambda}$ with respect to $(\bL_I)_{/E}\cap \ob_{/E}$.\;Denote $\chi_{\ul{\lambda}}:=L(\ul{\lambda})_{\emptyset}$. If $\ul{\lambda}\in X_{\Delta_n}^+$, let $L(\ul{\lambda}):=L(\ul{\lambda})_{\Delta_n}$.\;A $\bQ_p$-algebraic representation of $\GL_n(L)$ over $E$ is the induced action of $\GL_n(L)\subset \mathbf{G}_{/E}(E)$ on an algebraic representation of $\mathbf{G}_{/E}$.\;By abuse of notation we will use the same notations to denote $\bQ_p$-algebraic representations induced from an algebraic representation of $\mathbf{G}_{/E}$.\;

%Note that if $I\subseteq J\subseteq \Delta_n$, and $\ul{\lambda}\in X_J^+$, then we have $L(\ul{\lambda})_J^{\bN_I(L) \cap \bL_J(L)} \cong L(\ul{\lambda})_I$ (see \cite[Page 7997]{2019DINGSimple}).\;

Let $\ul{\lambda}$ be an integral weight,\;denote by $M(\ul{\lambda}):=\text{U}(\fg_{\Sigma_L})\otimes_{\text{U}(\fb_{\Sigma_L})} \ul{\lambda}$ (resp. $\overline{M}(\ul{\lambda}):=\text{U}(\fg_{\Sigma_L})\otimes_{\text{U}(\overline{\fb}_{\Sigma_L})} \ul{\lambda}$), the corresponding Verma module with respect to $\fb_{\Sigma_L}$ (resp.\;$\overline{\fb}_{\Sigma_L}$).\;Let $L(\ul{\lambda})$ (resp. $\overline{L}(\ul{\lambda})$) be the unique simple quotient of $M(\ul{\lambda})$ (resp. of $\overline{M}(\ul{\lambda})$).\;Actually,\;when $\ul{\lambda}\in X_{\Delta_n}^+$ (i.e. $-\ul{\lambda}\in X_{\Delta_n}^-$),\;$L(\ul{\lambda})$ is finite dimensional and  isomorphic to the algebraic representation $L(\ul{\lambda})$ introduced above (hence there is no conflict of notation).\;We have $\overline{L}(-\ul{\lambda})\cong L(\ul{\lambda})^{\vee}$.\;In general,\;for any subset $I$ of $\Delta_n$, and $\ul{\lambda}\in X_{I}^+$,\;we define the generalized parabolic Verma module
\begin{equation}
	\begin{aligned}
		&M_I(\ul{\lambda}):=\text{U}(\fg_{\Sigma_L})\otimes_{\text{U}(\fp_{I,\Sigma_L})} L(\ul{\lambda})_I,\;\\
	 \text{resp.,\;}&\overline{M}_I(-\ul{\lambda}):=\text{U}(\fg_{\Sigma_L})\otimes_{\text{U}(\overline{\fp}_{I,\Sigma_L})} \overline{L}(-\ul{\lambda})_I)
	\end{aligned}
\end{equation}
with respect to $\fp_{I,\Sigma_L}$ (resp.\;$\overline{\fp}_{I,\Sigma_L}$),\;see \cite[ Chapter 9]{humphreysBGG} for more precise statements.\;For $\ul{\lambda}\in X_{\Delta_n^k\cup I}^+$,\;we put (when $I=\emptyset$,\;we omit the subscripts $I$ in (\ref{lrrver}))
\begin{equation}\label{lrrver}
\begin{aligned}
	&L^{\lrr}(\ul{\lambda})_I:=L(\ul{\lambda})_{\Delta_n^k\cup I},\;\overline{L}^{\lrr}(-\ul{\lambda})_I:=\overline{L}(-\ul{\lambda})_{\Delta_n^k\cup I},\;\\
	&{M}_I^{\lrr}(\ul{\lambda}):={M}_{\Delta_n^k\cup I}(\ul{\lambda}),\;\overline{M}_I^{\lrr}(-\ul{\lambda}):=\overline{{M}}_{\Delta_n^k\cup I}(-\ul{\lambda}).
\end{aligned}
\end{equation}

%and denote by $s_i$ the simple reflection corresponding to $i\in \Delta_n$.\;For any $I\subset \Delta_n$,\;define $\sW_{I}$ be the subgroup of $\sW_{n}$ generated by simple reflections $s_i$ with $i\in I$.\;For any $I\subset \Delta_n$,\;define $\sW_{I,\sigma}$ to be the $\sigma$-copy of $\sW_{I}$,\;i.e.,\;to be the subgroup of $\sW_{n,\sigma}$ generated by simple reflections $s_i$ with $i\in I$.\;For $S\subseteq \Sigma_L$ and  $I\subset {\Delta_n}$, denote by $\sW_{I,S}:=\prod_{\sigma\in S} \sW_{I,\sigma}$.\;If $\ul{i}_{S}:=(i_{\sigma})_{\sigma\in S}\in {\Delta_n}^{|S|}$,\;denote by $s_{\ul{i}_S}:=\prod_{\sigma\in S}s_{i_\sigma,\sigma}$.\;If $H$ is a subgroup of $G$ and $(\rho,V)$ is an abstract representation of $H$ on $E$-vector space.\;We put $H^w:=w^{-1}Hw$,\;and $\rho^w(h):=\rho(whw^{-1})$.\;Then $(\rho^w,V)$ forms anabstract representation of $H^w$ on $E$-vector space.
\subsection{General constructions}
Let $\pi$ be an irreducible cuspidal representation of $\GLN_{r}(L)$ over $E$.\;For any $i\in \BZ$,\;we put $\pi(i):=\pi\otimes_E v_r^{i}$.\;Let $\Delta_{[k-1,0]}(\pi)=[\pi(k-1),\cdots,\pi(1),\pi]$ be a Zelevinsky-segment (see \cite{av1980induced2}).\;Put
\begin{equation}\label{pilrr}
\pi^{\lrr}:=\big(\otimes_{i=1}^{k}\pi(k-i)\big)\otimes_E \delta_{\op^{\lrr}(L)}^{1/2}=\otimes_{i=1}^{k}\big(\pi\otimes_E v_{r}^{-\frac{r}{2}(k-2i+1)+k-i}\big),\;
\end{equation}
where $\delta_{\op^{\lrr}_{{I}}(L)}$ is the modulus character of $\op^{\lrr}_{{I}}(L)$.\;This is an irreducible cuspidal smooth representation of $\bL^{\lrr}(L)$ over $E$.\;


%Let $\omega_{\pi}^{\lrr}$ be the central character of $\pi^{\lrr}$.\;Then the restriction of ${\pi^{\lrr}}$ on $Z_n:=\bZ_n(L)$ (recall that $\bZ_n$ is the center of $\GLN_{n}$) is given by the character $\omega_{\pi}^{\lrr}|_{Z_n}$,\;where $\omega_{\pi}^{\lrr}:=(\omega_{\pi}\otimes_Ev_r^{k-1})\otimes_E\cdots\otimes_E \omega_{\pi}$.\;Note that for any subset $I\subseteq {\Delta_n(k)}$,\;the center $Z_n$  acts on representations $\pi_I\otimes_EL^{\lrr}(\underline{\lambda})_{{I}}$,\;$i_{\op^{\lrr}_{{I}}}^{G}(\pi,\underline{\lambda})$ and $\BI_{\op^{\lrr}_{{I}}}^{G}(\pi,\underline{\lambda})$ via $\omega_{\pi}^{\lrr}\chi_{\underline{\lambda}}$.\;

In the sequel,\;we put  $G=\GL_n(L)$ for simplicity.\;Let $I$ be a subset of ${\Delta_n(k)}$.\;By \cite[Proposition 2.10]{av1980induced2}, we see that $i_{\op^{\lrr}(L)\cap\bL^{\lrr}_{I}(L)}^{\bL^{\lrr}_{I}(L)}\pi^{\lrr}$ admits a unique irreducible subrepresentation $\pi_I$.\;Let $\underline{\lambda}\in X_{\Delta_{n}}^+$, and let $I,J\subseteq {\Delta_n(k)}$.\;We put
\begin{equation*}
\begin{aligned}
	\mathbb{I}_{\op^{\lrr}_{{I}}}^{G}(\pi,\underline{\lambda})&=\Big(\mathrm{Ind}_{\op^{\lrr}_{{I}}(L)}^{G}\pi_I\otimes_EL^{\lrr}(\underline{\lambda})_{I}\Big)^{\mathbb{Q}_p-\mathrm{an}},
	i_{\op^{\lrr}_{{I}}}^{G}(\pi,\underline{\lambda})&=\left(i_{\op^{\lrr}_{{I}}(L)}^{G}\pi_{I}\right)\otimes_EL(\underline{\lambda}).\\
\end{aligned}
\end{equation*}
In the beginning of  \cite[Section 4.2]{2022ext1hyq} and \cite[Proposition 3.6]{2022ext1hyq},\;we show that if $J\supsetneq I$,\;we have an injection $\BI_{\op^{\lrr}_{{J}}}^{G}(\pi,\underline{\lambda})\hookrightarrow\BI_{\op^{\lrr}_{{I}}}^{G}(\pi,\underline{\lambda})$ and $i_{\op^{\lrr}_{{J}}}^{G}(\pi,\underline{\lambda})\hookrightarrow i_{\op^{\lrr}_{{I}}}^{G}(\pi,\underline{\lambda})$.\;Therefore,\;we put
\begin{equation*}
\begin{aligned}
	&v_{\op^{\lrr}_{{I}}}^{\ana}(\pi,\underline{\lambda})=\BI_{\op^{\lrr}_{{I}}}^{G}(\pi,\underline{\lambda})\big/\sum_{J\supsetneq I}\BI_{\op^{\lrr}_{{J}}}^{G}(\pi,\underline{\lambda}),\;\\
	&v_{\op^{\lrr}_{{I}}}^{\infty}(\pi,\underline{\lambda})=i_{\op^{\lrr}_{{I}}}^{G}(\pi,\underline{\lambda})\big/u_{\op^{\lrr}_{{I}}}^{\infty}(\pi,\underline{\lambda}),\;	u_{\op^{\lrr}_{{I}}}^{\infty}(\pi,\underline{\lambda})=\sum_{J\supsetneq I}i_{\op^{\lrr}_{{J}}}^{G}(\pi,\underline{\lambda}).\;
\end{aligned}
\end{equation*}
In particular,\;we denote by $\st_{(r,k)}^{\infty}(\pi,\underline{\lambda}):=v^{\infty}_{\op^{\lrr}_{\emptyset}}(\pi,\underline{\lambda})$ (resp.,\;$\st_{(r,k)}^{\ana}(\pi,\underline{\lambda})=v^{\ana}_{\op^{\lrr}}(\pi,\underline{\lambda})$) the locally $\bQ_p$-algebraic parabolic Steinberg representation (resp.,\;locally $\bQ_p$-analytic  parabolic Steinberg representation) with respect to the Zelevinsky-segment $\Delta_{[k-1,0]}(\pi)$ and the weight $\underline{\lambda}$.\;The main theorem of \cite{2022ext1hyq} is (see \cite[Theorem 5.19]{2022ext1hyq})
%We also put $\BI_{\op^{\lrr}_{{I}}}^G(\pi):=\BI_{\op^{\lrr}_{{I}}}^G(\pi,\underline{{0}})$,\;$v_{\op^{\lrr}_{{I}}}^{\ana}(\pi):=v_{\op^{\lrr}_{{I}}}^{\ana}(\pi,\ul{0})$, and $i_{\op^{\lrr}_{{I}}}^G(\pi):=i_{\op^{\lrr}_{{I}}}^G(\pi,\ul{0})$,\;$v_{\op^{\lrr}_{{I}}}^{\infty}(\pi):=v_{\op^{\lrr}_{{I}}}^{\infty}(\pi,\ul{0})$.\;The Jordan-Holder factors of $i_{\op^{\lrr}_{}}^G(\pi)$ are also given by $\{v^{\infty}_{\op^{\lrr}_{I}}(\pi)\}_{I\subseteq {\Delta_n(k)}}$.\;\cite[Lemma 3.7]{2022ext1hyq},\;we showed that 



\begin{thm}\label{thmintroLINV}
Let $ir\in {\Delta_n(k)}$,\;we have an $E$-vector spaces isomorphism
\begin{equation}\label{thmintro}
	\homo(L^\times,E)\xrightarrow{\sim }\ext^1_{G}\big(v_{\op^{\lrr}_{ir}}^{\infty}(\pi,\ul{\lambda}), \st_{(r,k)}^{\ana}(\pi,\ul{\lambda})\big).
\end{equation}
In particular,\;we have $\dim_E  \ext^1_{G}\big(v_{\op^{\lrr}_{ir}}^{\infty}(\pi,\ul{\lambda}), \st_{(r,k)}^{\ana}(\pi,\ul{\lambda})\big)=d_L+1$.\;
\end{thm}
For $\psi\in\homo(L^{\times},\;E)$,\;the image of $\psi$ via the above isomorphism (\ref{thmintroLINV}) can be constructed explicitly as follows.\;Recall that we have an isomorphism 
\begin{equation}
\begin{aligned}
	\iota_\upsilon: \homo(L^{\times},E) &\xrightarrow{\sim} & \homo(\bZ^{\lrr}(L),E)/\homo(\bZ^{\lrr}_{ir}(L),E),\\
	\psi& \mapsto& [(a_1I_{r},\cdots, a_{k}I_{r})\mapsto \psi(a_i/a_{i+1})].\;
\end{aligned}
\end{equation}
Therefore,\;for any $\psi\in\homo(L^{\times},\;E)$,\;we choose a lift $\Psi\in \homo(\bZ^{\lrr}(L),E)$.\;Then  $\Psi$ induces an extension $\iota(\Psi)$ of $\pi^{\lrr}$ by $\pi^{\lrr}$:
\begin{equation*}
\iota(\Psi)(a)=\pi^{\lrr}(a)\otimes_E\begin{pmatrix}
	1 & \Psi\circ\det(a) \\ 0 & 1
\end{pmatrix}, \\ \forall ~a\in \bL^{\lrr}_{J}(L).
\end{equation*}
Consider the locally $\bQ_p$-analytic parabolic induction $\big(\ind_{\op^{\lrr}(L)}^G \iota(\Psi)\otimes_EL^{\lrr}(\ul{\lambda})\big)^{\BQ_p-\ana}$,\;which lies in the following exact sequence
\begin{equation}\label{explicitexactseq}
0 \longrightarrow \BI_{\op^{\lrr}}^G(\pi,\ul{\lambda}) \longrightarrow \big(\ind_{\op^{\lrr}(L)}^G \iota(\Psi)\otimes_EL^{\lrr}(\ul{\lambda})\big)^{\BQ_p-\ana} \xrightarrow{\pr} \BI_{\op^{\lrr}}^G(\pi,\ul{\lambda}) \longrightarrow 0.
\end{equation}
Pushforward (\ref{explicitexactseq}) along surjection $\BI_{\op^{\lrr}}^G(\pi,\ul{\lambda})\twoheadrightarrow v^{\ana}_{\op^{\lrr}}(\pi,\ul{\lambda})$ and then pullback via natural injections
\begin{equation*}i_{\op^{\lrr}_{{ir}}}^G(\pi,\ul{\lambda}) \hooklongrightarrow \BI_{\op^{\lrr}_{{ir}}}^G(\pi,\ul{\lambda}) \hooklongrightarrow \BI_{\op^{\lrr}}^G(\pi,\ul{\lambda}),
\end{equation*}
we get a locally $\bQ_p$-analytic representation  $\sE_{ir}^\emptyset(\pi,\ul{\lambda},\Psi)^0$,\;which is an extension of $i_{\op^{\lrr}_{{ir}}}^G(\pi,\ul{\lambda})$ by the locally $\bQ_p$-analytic representation $\st_{(r,k)}^{\ana}(\pi,\ul{\lambda})$. It gives a cohomology class  \[[\sE_{ir}^\emptyset(\pi,\ul{\lambda},\psi)^0]\in \ext^{1}_{G}\big(i_{\op^{\lrr}_{{ir}}}^G(\pi,\ul{\lambda}), v^{\ana}_{\op^{\lrr}}(\pi,\ul{\lambda})\big)\] (we can show that this cohomology class is independent on the choice of $\Psi$).\;By \cite[Theorem 5.19]{2022ext1hyq},\;we see that the natural map
$$\ext^1_{G}\big(v_{\op^{\lrr}_{{ir}}}^{\infty}(\pi,\ul{\lambda}, \st_{(r,k)}^{\ana}(\pi,\ul{\lambda}))\hookrightarrow\ext^1_{G}\big(i_{\op^{\lrr}_{{ir}}}^G(\pi,\ul{\lambda}),\st_{(r,k)}^{\ana}(\pi,\ul{\lambda})),\;$$
induced by the surjection $i_{\op^{\lrr}_{{ir}}}^G(\pi,\ul{\lambda})\twoheadlongrightarrow v_{\op^{\lrr}_{{ir}}}^{\infty}(\pi,\ul{\lambda})$,\;is actually an isomorphism.\;Therefore,\;the pull back of $\sE_{\{i\}}^\emptyset(\pi,\ul{\lambda},\iota_\nu(\psi))^0$ via the natural injection $u_{\op^{\lrr}_{{ir}}}^{\infty}(\pi,\ul{\lambda}) \rightarrow i_{\op^{\lrr}_{{ir}}}^G(\pi,\ul{\lambda})$ is split (as an extension of $u_{\op^{\lrr}_{{ir}}}^{\infty}(\pi,\ul{\lambda})$ by $\st_{(r,k)}^{\ana}(\pi,\ul{\lambda})$).\;Quotient it by $u_{\op^{\lrr}_{{ir}}}^{\infty}(\pi,\ul{\lambda})$,\;we deduce $\widetilde{{\Sigma}}_i^{\lrr}(\pi,\ul{\lambda}, \psi)$,\;which is an extension of $v_{\op^{\lrr}_{{ir}}}^{\infty}(\pi,\ul{\lambda})$ by $\st_{(r,k)}^{\ana}(\pi,\ul{\lambda})$.\;We therefore get that $[\widetilde{{\Sigma}}_i^{\lrr}(\pi,\ul{\lambda}, \psi)]$ is the extension class associated with $\psi$ via (\ref{thmintroLINV}).\

%Define 
% \begin{equation*}
%	\begin{aligned}
	%		&\st_I^{\ana}(\pi,\underline{\lambda})=\BI_{\op^{\lrr}(L)\cap\bL^{\lrr}_{I}(L)}^{\bL^{\lrr}_{I}(L)}(\pi,\underline{\lambda})\big/\sum_{\emptyset \subsetneq K\subseteq I}\BI_{\op^{\lrr}_{K}(L)\cap\bL^{\lrr}_{I}(L)}^{\bL^{\lrr}_{I}(L)}(\pi,\underline{\lambda})\\
	%		&\bigg(\mathrm{resp}.,\;\st_I^{\infty}(\pi,\underline{\lambda})=i_{\op^{\lrr}(L)\cap \bL^{\lrr}_{I}(L)}^{\bL^{\lrr}_{I}(L)}\pi_I\big/\sum_{\emptyset \subsetneq K\subseteq I}i_{\op^{\lrr}_{{K}}(L)\cap \bL^{\lrr}_{I}(L)}^{\bL^{\lrr}_{I}(L)}{\pi^{\lrr}},\\
	%$		&\st_I^{\infty}(\pi,\underline{\lambda})=\st_I^{\infty}(\pi)\otimes_EL^{\lrr}(\underline{\lambda})_{I}\bigg).\;
	%	\end{aligned}
%\end{equation*}
%They are locally analytic (resp.,\;smooth,\;locally algebraic) generalized parabolic Steinberg representations of $\bL^{\lrr}_{I}(L)$.\;By the tranversity of parabolic inductions,\;we have:
%\begin{equation*}
%	\begin{aligned}
	%		&\left(\mathrm{Ind}_{\op^{\lrr}_{{I}}}^{G}\st_{I}^{\infty}(\pi)\right)^\infty\otimes_EL(\underline{\lambda})=i_{\op^{\lrr}}^{G}(\pi,\underline{\lambda})\big/\sum_{\emptyset \neq J\subseteq I}i_{\op^{\lrr}_{{J}}}^{G}(\pi,\underline{\lambda}),\\
	%		&\left(\mathrm{Ind}_{\op^{\lrr}_{{I}}}^{G}\st_{I}^{\ana}(\pi,\underline{\lambda})\right)^{\mathbb{Q}_p-\mathrm{an}}=\BI_{\op^{\lrr}}^{G}(\pi,\underline{\lambda})\big/\sum_{\emptyset \neq J\subseteq I}\BI_{\op^{\lrr}_{{J}}}^{G}(\pi,\underline{\lambda}).\;
	%	\end{aligned}
%$\end{equation*}
%Consider the following composition
%\begin{equation}
%	\tau_I: \left(\mathrm{Ind}_{\op^{\lrr}_{{I}}}^{G}\st_{I}^\infty(\pi,\underline{\lambda})\right)^{\mathbb{Q}_p-\mathrm{an}} \hooklongrightarrow \left(\mathrm{Ind}_{\op^{\lrr}_{{I}}}^{G}\st_{I}^{\ana}(\pi,\underline{\lambda})\right)^{\mathbb{Q}_p-\mathrm{an}} \twoheadlongrightarrow \mathrm{St}_{(r,k)}^{\ana}(\pi,\underline{\lambda}),
%\end{equation}
%and denote by $\widetilde{\Sigma}^{\lrr}_{{\Delta_n(k)}\backslash I}(\pi,\ul{\lambda})$ its image. When $I=\Delta_{k,i}$, we denote $\widetilde{\Sigma}^{\lrr}_{ir}(\pi,\ul{\lambda}):=\widetilde{\Sigma}^{\lrr}_{\{ir\}}(\pi,\ul{\lambda})$.

For $I\subset \Delta_n(k)$,\;we recall the  Orlik-Strauch functor $\cF^{G}_{\op_{{I}}}(-,-)$ (see \cite[The main theorem]{orlik2015jordan}),\;or see \cite[Section 2]{breuil2016socle}),\;which associates,\;to an object $M$ in the Bernstein-Gelfand-Gelfand (BGG) category  $\cO_{\alge}^{\overline{\fp}_{{{I}}},\Sigma_L}$,\;and an finite length smooth admissible representation $\pi_I$,\;a locally $\bQ_p$-analytic representation $\cF^{G}_{\op_{{I}}}(M,\pi_I)$.\;

For  $ir\in {\Delta_n(k)}$, $\sigma\in \Sigma_L$,\;we put $\ul{\lambda}_{\sigma}:=(\lambda_{1,\sigma}, \cdots, \lambda_{n,\sigma})$, and $\ul{\lambda}^{\sigma}:=(\lambda_{1,\sigma'}, \cdots, \lambda_{n,\sigma'})_{\sigma'\in \Sigma_L\backslash\{\sigma\}}$. We put
\begin{equation*}
\begin{aligned}
	& \widetilde{\Sigma}^{\lrr}_{i,\sigma}(\pi,\ul{\lambda}):=\Big(\big(\ind_{\op^{\lrr}_{{\Delta}_{k,i}}(L)}^G \st_{{\Delta}_{k,i}}^{\infty}(\pi,\ul{\lambda}_{\sigma}) )^{\sigma-\ana}\otimes_E L(\ul{\lambda}^{\sigma})\Big)/v_{\op^{\lrr}_{{ir}}}^{\infty}(\pi,\ul{\lambda}),\\
	&	C_{i,\sigma}:=\cF_{\op^{\lrr}_{{\Delta}_{k,i}}}^G\big(\overline{L}(-\ul{\lambda}^{\sigma})\otimes_E \overline{L}(-s_{ir,\sigma}\cdot \ul{\lambda}_{\sigma}), \st_{{\Delta}_{k,i}}^{\infty}(\pi,\ul{0})\big),
\end{aligned}
\end{equation*}
%Then we have by \cite[(6.40)]{2022ext1hyq} the injection $\widetilde{\Sigma}^{\lrr}_{ir,\sigma}(\pi,\ul{\lambda}) \hookrightarrow \widetilde{\Sigma}^{\lrr}_{ir}(\pi,\ul{\lambda})$.\;
By \cite[Proposition 4.6 and Proposition 4.9]{2022ext1hyq},\;$C_{i,\sigma}$ appears as an irreducible constituent in $\st_{(r,k)}^{\ana}(\pi,\ul{\lambda})$
with  multiplicity $1$.\;By the argument after \cite[(5.48)]{2022ext1hyq},\;$\widetilde{\Sigma}^{\lrr}_{i,\sigma}(\pi,\ul{\lambda})$ is a subrepresentation of $\st_{(r,k)}^{\ana}(\pi,\ul{\lambda})$, and $\widetilde{\Sigma}^{\lrr}_{i,\sigma}(\pi,\ul{\lambda})$ admits a subrepresentation $\Sigma_{i,\sigma}^{\lrr}(\pi,\ul{\lambda})$ which is an extension of $C_{i,\sigma}$ by $\st_{(r,k)}^{\infty}(\pi,\ul{\lambda})$.\;We put
\begin{equation}\label{subofsigma}
\begin{aligned}
	&\Sigma_i^{\lrr}(\pi,\ul{\lambda}):=\bigoplus^{\sigma\in \Sigma_L}_{\st_{(r,k)}^{\infty}(\pi,\ul{\lambda})}\Sigma_{ir,\sigma}^{\lrr}(\pi,\ul{\lambda})\hookrightarrow \widetilde{\Sigma}^{\lrr}_{i}(\pi,\ul{\lambda})=\bigoplus^{\sigma\in \Sigma_L}_{\st_{(r,k)}^{\infty}(\pi,\ul{\lambda})}\widetilde\Sigma_{i,\sigma}^{\lrr}(\pi,\ul{\lambda}),\\
	& \Sigma^{\lrr}(\pi,\ul{\lambda}):=\bigoplus^{ir\in {\Delta_n(k)}}_{\st_{(r,k)}^{\infty}(\pi,\ul{\lambda})} \Sigma_i^{\lrr}(\pi,\ul{\lambda})
	\hooklongrightarrow\widetilde{\Sigma}^{\lrr}(\pi,\ul{\lambda}):=\bigoplus^{\sigma\in \Sigma_L}_{\st_{(r,k)}^{\infty}(\pi,\ul{\lambda})}\widetilde{\Sigma}^{\lrr}_{i}(\pi,\ul{\lambda})\subseteq\st_{(r,k)}^{\ana}(\pi,\underline{\lambda}).
\end{aligned}
\end{equation}
By the definition,\;$\Sigma_i^{\lrr}(\pi,\ul{\lambda})$ is an extension of $\bigoplus_{\sigma\in \Sigma_L}C_{i,\sigma}$ by $\st_{(r,k)}^{\infty}(\pi,\ul{\lambda})$.\;Moreover,\;we have a natural isomorphism (see \cite[Proposition 5.34]{2022ext1hyq})
\begin{equation}\label{Intro:equ: BrLi}
\homo(L^\times,E)\xrightarrow{\sim}  \ext^1_G\big(v_{\op^{\lrr}_{{ir}}}^{\infty}(\pi,\ul{\lambda}), \Sigma_{i}^{\lrr}(\pi,\ul{\lambda})\big).
\end{equation}
For any $\psi \in \homo(L^{\times}, E)$,\;let  $[\Sigma_i^{\lrr}(\pi,\ul{\lambda}, \psi)]\in\ext^1_G\big(v_{\op^{\lrr}_{{ir}}}^{\infty}(\pi,\ul{\lambda}), \Sigma_{i}^{\lrr}(\pi,\ul{\lambda})\big)
$ be the 
image of $\psi$ via the isomorphism (\ref{Intro:equ: BrLi}).\;Moreover,\;we can see that $[\widetilde{{\Sigma}}_i^{\lrr}(\pi,\ul{\lambda}, \psi)]$ actually comes from $[\Sigma_i^{\lrr}(\pi,\ul{\lambda}, \psi)]$ by pushing-forward $\Sigma_{i}^{\lrr}(\pi,\ul{\lambda})\hookrightarrow \st_{(r,k)}^{\ana}(\pi,\ul{\lambda})$.\;

Let $\alpha\in E^\times$,\;we denote  $\ast(\alpha,\pi,\ul{\lambda}):=\ast(\pi,\ul{\lambda})\otimes_E\unr(\alpha)\circ\det$ for any representation $\ast(\pi,\ul{\lambda})$ of $G$ as above (for example,\;$v_{\op^{\lrr}_{ir}}^{\infty}(\alpha,\pi,\ul{\lambda}), \st_{(r,k)}^{\ana}(\alpha,\pi,\ul{\lambda})$,\;$\Sigma_{i}^{\lrr}(\alpha,\pi,\ul{\lambda})$,\;$\widetilde{{\Sigma}}_i^{\lrr}(\alpha,\pi,\ul{\lambda}, \psi)$,\;etc.).\;

For $ir\in \Delta_n(k)$ (so for $1\leq i\leq k-1$),\;let $V_i$ be an $E$-vector subspace of $\homo(L^\times,E)$ of dimension $d_L$.\;Let $\{\psi_{i,1}, \cdots, \psi_{i,d_L}\}$ be a basis of $V_i$.\;We put
\begin{eqnarray*}
\widetilde{\Sigma}_i^{\lrr}(\alpha,\ul{\lambda}, V_i)&:=&\bigoplus_{\st_{(r,k)}^{\ana}(\alpha,\pi, \ul{\lambda})}^{j=1,\cdots, d_L} \widetilde{\Sigma}_i^{\lrr}(\alpha,\pi, \ul{\lambda}, \psi_{i,j}), \\
\Sigma_i^{\lrr}(\alpha,\pi,\ul{\lambda}, V_i)&:=&\bigoplus_{\Sigma_i^{\lrr}(\alpha,\pi, \ul{\lambda})}^{j=1,\cdots, d_L} \Sigma_i^{\lrr}(\alpha,\pi,\ul{\lambda}, \psi_{i,j}).
\end{eqnarray*}
Thus $\Sigma_i^{\lrr}(\alpha,\pi,\ul{\lambda}, V_i)$ is a subrepresentation of $\widetilde{\Sigma}_i^{\lrr}(\alpha,\pi, \ul{\lambda}, V_i)$,\;both of the representations are independent of the choice of the basis of $V_i$ and determine $V_i$. Put $V:=\prod_{ir\in \Delta_n(k)}V_i$ and
\begin{eqnarray*}
\widetilde{\Sigma}^{\lrr}(\alpha,\pi,\ul{\lambda}, V)&:=&\bigoplus_{\st_{(r,k)}^{\ana}(\alpha,\pi, \ul{\lambda})}^{ir\in \Delta_n(k)} \widetilde{\Sigma}_i^{\lrr}(\alpha,\pi,\ul{\lambda}, V_i), \\
\Sigma^{\lrr}(\alpha,\pi,\ul{\lambda}, V)&:=&\bigoplus_{\st_{(r,k)}^{\infty}(\alpha,\pi, \ul{\lambda})}^{ir\in \Delta_n(k)} \Sigma_i^{\lrr}(\alpha,\pi, \ul{\lambda}, V_i).
\end{eqnarray*}
It is clear that $\Sigma^{\lrr}(\alpha,\pi,\ul{\lambda}, V)$ is a subrepresentation of $ \widetilde{\Sigma}^{\lrr}(\alpha,\pi,\ul{\lambda}, V)$.\;

\subsection{A subcandidate in the locally analytic \texorpdfstring{$\bQ_p$}{Lg}-adic local Langlands program}
The main result of this section is given as follows.\;Let $\rho_L:\gal_L\rightarrow \GLN_{n}(E)$ be a potentially semistable representation,\;and let $\Dpik=D_{\rig}(\rho_L)$ be the associated  
$(\varphi,\Gamma)$-module over $\cR_{E,L}$ of rank $n$.\;

Let  $\bh:=(\hpi_{\tau,1}>\hpi_{\tau,2}>\cdots>\hpi_{\tau,n} )_{\tau\in \Sigma_L}$ be the Hodge-Tate weights of $\rho_L$ (or $\Dpik$).\;We put $\hpi_{i}=(\hpi_{\tau,i})_{\tau\in \Sigma_L}$ for $1\leq i\leq n$.\;We put ${\bm\lambda}_\bh=(\hpi_{\tau,i}+i-1)_{\tau\in \Sigma_L,1\leq i\leq n}$,\;which is a dominant weight of $(\mathrm{Res}_{L/\BQ_p}\GLN_n)\times_{\BQ_p}E$ respect to $(\mathrm{Res}_{L/\BQ_p}\bB)\times_{\BQ_p}E$.\;

Suppose that $\Dpik$ admits a non-critical special  $\omepik$-filtration (see Definition \ref{dfnnoncriticalspecial}) with
parameter $(\bx_{\pi},\bmdel)\in \sbanpik\times\rigchl $ (or $(\widetilde{\bx}_{\pi,\bh},\widetilde{\bm{\delta}}_\bh)\in \sbanpik\times\rigch $).\;

Many information on $\rho_L$ is lost when passing from $\rho_L$ to its associated Weil-Deligne representation $\mathrm{WD}(\rho_L)$.\;We have defined the parabolic Fontaine-Mazur simple $\sL$-invariants $\sL(\rho_L)$ of  $\rho_L$.\;Then we will see that the locally $\BQ_p$-analytic representations $\Sigma^{\lrr}(\alpha_\pi ,\pi_0, {\bm\lambda}_\bh, \sL(\rho_L))$ and $\widetilde{\Sigma}^{\lrr}(\alpha_\pi ,\pi_0,{\bm\lambda}_\bh, \sL(\rho_L))$ of $G$ carry the exact information on the Weil-Deligne representation $\wdre(\rho_L)\cong\wdre(\Dpik)$ associated with $\rho_L$,\;the Hodge-Tate weights $\mathrm{HT}(\rho_L)$ of $\rho_L$,\;and the parabolic Fontaine-Mazur simple $\sL$-invariants $\sL(\rho_L)$ of $\rho_L$.\;Both of the representations  determine exactly the data $\{\wdre(\rho_L), \mathrm{HT}(\rho_L), \sL(\rho_L)\}$,\;and vice versa.




\begin{pro}We have:
\begin{description}
	\item[(1)] $\widetilde{\Sigma}^{\lrr}(\alpha,\pi,{\bm\lambda}_\bh, \sL(\rho_L))$ is isomorphic to an extension of $\bigoplus\limits_{ir\in \Delta_n(k)}v_{\op^{\lrr}_{ir}}^{\infty}(\alpha,\pi,{\bm\lambda}_\bh)^{d_L}$ by $\st_{(r,k)}^{\ana}(\alpha,\pi, {\bm\lambda}_\bh)$.\;
	\item[(2)] $\Sigma^{\lrr}(\alpha,\pi,{\bm\lambda}_\bh, \sL(\rho_L))$ has the following form:
	\begin{equation}\label{stru}
		\begindc{\commdiag}[300]
		\obj(0,3)[a]{$\st_{(r,k)}^{\infty}(\alpha,\pi, {\bm\lambda}_\bh)$}
		\obj(4,6)[b]{$C_{1,\sigma_1}$}
		\obj(8,6)[c]{$v_{\op^{\lrr}_{{r}}}^{\infty}(\alpha,\pi,{\bm\lambda}_\bh)$}
		\obj(4,5)[d]{$\vdots$}
		\obj(8,5)[e]{$\vdots$}
		\obj(4,4)[f]{$C_{1, \sigma_{d_L}}$}
		\obj(8,4)[g]{$v_{\op^{\lrr}_{{r}}}^{\infty}(\alpha,\pi, {\bm\lambda}_\bh)$}
		\obj(4,3)[h]{$\vdots$}
		\obj(8,3)[i]{$\vdots$}
		\obj(4,2)[j]{$C_{k-1,\sigma_1}$}
		\obj(8,2)[k]{$v_{\op^{\lrr}_{(k-1)r}}^{\infty}(\alpha,\pi, {\bm\lambda}_\bh)$}
		\obj(4,1)[l]{$\vdots$}
		\obj(8,1)[m]{$\vdots$}
		\obj(4,0)[n]{$C_{k-1,\sigma_{d_L}}$}
		\obj(8,0)[o]{$v_{\op^{\lrr}_{(k-1)r}}^{\infty}(\alpha,\pi, {\bm\lambda}_\bh)$}
		\mor{a}{b}{}[+1,\solidline]
		\mor{b}{c}{}[+1,\solidline]
		\mor{a}{f}{}[+2,\solidline]
		\mor{f}{g}{}[+1,\solidline]
		\mor{a}{j}{}[+1,\solidline]
		\mor{j}{k}{}[+1,\solidline]
		\mor{a}{n}{}[+1,\solidline]
		\mor{n}{o}{}[+1,\solidline]
		\enddc,
	\end{equation}
	\item[(3)] $\soc_G\Sigma_i^{\lrr}(\alpha,\pi,{\bm\lambda}_\bh, \sL(\rho_L)_{ir})\cong \soc_G \Sigma^{\lrr}(\alpha,\pi, {\bm\lambda}_\bh, \sL(\rho_L))\cong \st_{(r,k)}^{\infty}(\alpha,\pi, {\bm\lambda}_\bh)$.
	\item[(4)]The locally algebraic subrepresentation of $\widetilde{\Sigma}^{\lrr}(\alpha,\pi,{\bm\lambda}_\bh, \sL(\rho_L))$ (resp. of $\widetilde{\Sigma}_i^{\lrr}(\alpha,\pi,{\bm\lambda}_\bh, \sL(\rho_L)_{ir})$ for $ir\in \Delta_n(k)$) is isomorphic to $\st_{(r,k)}^{\infty}(\alpha,\pi, {\bm\lambda}_\bh)$.
	%if and only if $N^{k-1}\neq 0$ on $D_{\pst}(\Dpik)$ (resp. if and only if $\Dpik_i^{i+1}$ is not potentially crystalline).
\end{description}
\end{pro}
\begin{proof}Part (1) is clear.\;\;For (2),\;note that the Remark \ref{taulinvariantrmk} asserts that $\dim_E\sL(\Dpik)_{ir,\tau}=1$ for all $\tau\in \Sigma_L$,\;and $\sL(\Dpik)_{ir}=\bigoplus_{\tau\in \Sigma_L}\sL(\Dpik)_{ir,\tau}$.\;Thus,\;there exists a $\sL_{i,\tau}\in E$ such that $\sL(\Dpik)_{ir,\tau}$ is generated by $\psi_{i,\tau}:=\psi_{\tau,L}-\sL_{i,\tau}\psi_{\ur}$.\;Then $\{\psi_{i,\tau}\}_{i\in\Delta_n(k),\tau\in \Sigma_L}$ form a basis of $\sL(\Dpik)_{ir}$.\;By the same argument as in the  the discussion after \cite[Lem.\;3.3]{2019DINGSimple},\;we can obtain the decomposition
\[{\Sigma}_i^{\lrr}(\alpha,\pi, {\bm\lambda}_\bh, \psi_{i,\tau})={\Sigma}^{\lrr}_{i,\tau}(\alpha,\pi, {\bm\lambda}_\bh, \psi_{i,\tau})\oplus_{\st_{(r,k)}^{\infty}(\alpha,\pi, {\bm\lambda}_\bh)}\left(
\bigoplus_{\st_{(r,k)}^{\infty}(\alpha,\pi, {\bm\lambda}_\bh)}^{\sigma\in\Sigma_L,\sigma\neq\tau} \Sigma_{i,\sigma}^{\lrr}(\alpha,\pi, {\bm\lambda}_\bh)\right).\]
This deduces (2).\;Part (3) is a direct consequence of \cite[Proposition 5.28]{2022ext1hyq}.\;By \cite[Remark 5.21]{2022ext1hyq},\;we see that the locally algebraic subrepresentation of $\widetilde{\Sigma}_i^{\lrr}(\alpha,\pi,{\bm\lambda}_\bh, \sL(\rho_L)_{ir})$ is strictly bigger than\\ $\st_{(r,k)}^{\infty}(\alpha,\pi,{\bm\lambda}_\bh)$ if and only if $\homo_{\infty}(L^{\times}, E)\subseteq \sL(\rho_L)_{ir}$.\;Then $(4)$ follows from Lemma \ref{criofnoncrystalline} and Remark \ref{taulinvariantrmk}.\;
%which is equivalent to $D_i^{i+1}$ being crystalline (see the discussion in \S \ref{sec: lgln-LFM1}). The statement for $\widetilde{\Sigma}_i^{\lrr}(\alpha,\pi,{\bm\lambda}_\bh, \bL(D)_i)$ follows. By definition, it is easy to see $\widetilde{\Sigma}^{\lrr}(\alpha,\pi, {\bm\lambda}_\bh, \bL(D))$ has locally algebraic vectors other than $\st_n^{\infty}(\alpha,\pi, {\bm\lambda}_\bh)$ if and only if so does $\widetilde{\Sigma}_i^{\lrr}(\alpha,\pi, {\bm\lambda}_\bh, \bL(D)_i)$ for certain $i\in \Delta$. Together with Lemma \ref{monofull}, the statement for $\widetilde{\Sigma}^{\lrr}(\alpha,\pi, {\bm\lambda}_\bh, \bL(D))$ follows.
\end{proof}
\section{Local-global compatibility}

In this Chapter,\;we prove some new ($p$-adic) local-global compatibility results  for  potentially semistable non-crystalline $p$-adic Galois representation.\;

Since we want to explore the non-trianguline Galois representations,\;our local-global compatibility results are realized in the framework of patched Bernstein eigenvariety (roughly speaking,\;since we can only see finite slope $p$-adic automorphism forms or trianguline representation in the classical eigenvarieties),\;which is constructed by Christophe Breuil and Yiwen Ding (see \cite{Ding2021}).\;

We briefly describe the contents of each section.\;In Section \ref{BENVARPARAVAR},\;we first review the global patching construction (hence we also assume the so-called Talyor-Wiles hypothesis as in \cite{PATCHING2016}).\;Then we construct the patched Bernstein eigenvariety (by an easy variation of \cite[Section 3.3]{Ding2021}) and the (purely local) Bernstein paraboline varieties (see \cite[Section 4.2]{Ding2021}).\;

To state our local-global compatibility results,\;let $x$ be a point in the patched Bernstein eigenvariety such that the associated $p$-adic Galois representation $\rho_x$ admits a non-critical special  $\omepik$-filtration.\;

In Section \ref{smoothpoint},\;we show that $x$ is a smooth point of the patched Bernstein eigenvariety.\;By Section \ref{dfnFMparalINVSECTION},\;we can attach to $x$ (resp.,\;$\rho_L$) the parabolic Fontaine-Mazur simple $\sL$-invariants $\sL(\rho_L)$.\;We further show that the tangent map of the ``weight" map at point $x$ recovers the information of $\sL(\rho_L)$ (see Proposition \ref{coreigentangentmap2}).\;Via a study of parabolic Breuil's simple $\sL$-invariants (Sections \ref{BreuilLINVSUMMMANY}),\;we attach to $\rho_L$ a locally $\BQ_p$-analytic representation $\Sigma^{\lrr}(\pi, \ul{\lambda},\sL(\rho_L))$.\;The second main theorem (see Section \ref{lgcompthemainSECTION},\;our $p$-adic local-global compatibility results) of this paper asserts that  $\Sigma^{\lrr}(\pi, \ul{\lambda},\sL(\rho_L))$ is a subrepresentation of the associated Hecke-isotypic subspaces of the Banach spaces of $p$-adic automotphic form on certain (definite) unitary group (see Theorem \ref{thm: lgln-main}).\;

%On the other hand,\;for the continuous representation $\rho_L$,\;one can associate with an admissible unitary Banach representation $\widehat{\Pi}(\rho_L)$ of $\mathrm{GL}_n(L)$ (see Section 4.1 or see the paper \emph{Patching and the p-adic local langlands correspondence.},\;under the so called Talyor-Wiles hypothesis),\;which one might expect to be the right representation (up to multiplicities) corresponding to $\rho_L$ in the $p$-adic local Langlands program.
%In this case,\;we have an injection (up to some unramified twist):
%\[\mathrm{St}^{\infty}_{(r,k)}(\pi,\ul{\lambda})\hookrightarrow \widehat{\Pi}(\rho_L).\]In section 4.6,\;we show the main result of the paper.\;Keeping the assumption,\;the injection $\mathrm{St}^{\infty}_{(r,k)}(\pi,\ul{\lambda})\hookrightarrow \widehat{\Pi}(\rho_L)$ extends uniquely to a morphism 
%\[\Sigma^{\lrr}(\pi, \ul{\lambda},\sL(\rho_L))\hookrightarrow \widehat{\Pi}(\rho_L).\]
%Furthermore,\;$\sL(\rho_L)$ can be read out from $\widehat{\Pi}(\rho_L)$.\;


\subsection{Patched Bernstein eigenvarieties and Bernstein parabolic varieties}\label{BENVARPARAVAR}

Our Local-global compatibility results are realized in the space of the patched $p$-adic automorphic forms.\;More precisely,\;it is realized in the setting of \cite[Section 4.1.1]{2019DINGSimple}.\;In this section,\;we recall briefly  the patched Bernstein eigenvariety  and Bernstein paraboline variety  of Breuil-Ding (see \cite[Section 3.3,\;Section 4.2]{Ding2021}).\;Indeed,\;the patched arguments we need are slightly different from that in 
\cite[Section 3.3]{Ding2021}).\;We instead only vary all weights and levels at only one $p$-adic place.\;But the arguments in \cite[Section 3.3,\;Section 4.2]{Ding2021} can easily be adapted to our case.\;

We follow the notation of \cite[Section 4.1.1]{2019DINGSimple} and \cite[Section 2]{PATCHING2016}.\;Suppose that $p\nmid 2n$, and let $\overline{r}: \gal_L \rightarrow \GLN_n(k_E)$ be a continuous representation such that $\overline{r}$ admits a potentially crystalline lift $r_{\mathrm{pot.diag}}: \gal_L \rightarrow \GLN_n(E)$  of regular weight $\xi$ which is potentially diagonalisable.\;We can find a triple $(F,F^+, \overline{\rho})$,\;where $F$ is an imaginary CM field with maximal totally real subfield $F^+$,\;and $\overline{\rho}:\gal_{F^+} \rightarrow \cG_n(k_E)$ is a \emph{suitable globalisation} (cf. \cite[Section 2.1]{PATCHING2016}) of $\overline{r}$.\;In particular,\;for any place $v|p$ of $F^+$,\;$v$ splits in $F$, and has $F^+_v\cong L$.\;There is a place $\widetilde{v}$ of $F$ lying over $v$ with $\overline{\rho}|_{\gal_{F_{\widetilde{v}}}}\cong \overline{r}$.\;

Let $S_p$ be the set of places of $F^+$ above $p$.\;By \cite[Section 2.3]{PATCHING2016},\;we can find the following objects
\begin{equation*}
\{\widetilde{G}, v_1, \fp\in S_p, \{U_m\}_{m\in \BZ}\},
\end{equation*}
where $\widetilde{G}$ is a certain definite unitary group over $F^+$,\;$v_1$ is a certain finite place of $F^+$ prime to $p$,\;and $\{U_m=\prod_{v} U_{m,v}\}_{m\in \BZ_{\geq 0}}$ is a tower of certain compact open subgroups of $\widetilde{G}(\BA_{F^+}^{\infty})$ (see \cite[Section 4.1.1,\;Page 8028]{2019DINGSimple} or \cite[Section 2,\;Page 214]{PATCHING2016} for a precise description).\;We mention that $U_m$ has full level $\fp^m$ (resp.,\;level $\widetilde{G}(\cO_{F^+_v})$) at $\fp$ (resp.,\;$v\in S_p\setminus\{\fp\}$).\;

Let $\tau$ be the inertial type of $r_{\mathrm{pot.diag}}$.\;As in \cite[Section 4.1.1,\;Page 8029]{2019DINGSimple} or  \cite[Section 2,\;Page 215]{PATCHING2016},\;we consider the space of $p$-adic automorphic forms 	$\widehat{S}_{\xi,\tau}(U^{\fp}, \co_E)$ and $\widehat{S}_{\xi,\tau}(U^{\fp},E)$ (roughly speaking,\;the space of $p$-adic algebraic automorphic forms of fixed type $\sigma(\tau)$ (see \cite[Theorem 3.7]{PATCHING2016},\;the ``inertial local Langlands correspondence") at the place $S_p\setminus\{\fp\}$,\;full level at $\fp$,\;and whose weight is $0$ at places above $\fp$,\;and given by the regular weight $\xi$ at each of the places in $S_p\setminus\{\fp\}$).\;Note that $\widehat{S}_{\xi,\tau}(U^{\fp},E)$ is a Banach space for the supermum norm and is equipped with a continuous (unitary) action of $\GLN_n(L)$ (by right translation on functions).\;The space $\widehat{S}_{\xi,\tau}(U^{\fp},E)$  is also equipped with a faithful  action of a certain commutative global Hecke algebra $\bT^{S_p,\univ}$ over $\cO_E$ which is generated by some prime-to-$p$ Hecke operators (see \cite[Section 4.1.1,\;Page 8029]{2019DINGSimple}).\;We can associate to $\overline{\rho}$ a maximal ideal $\fm_{\overline{\rho}}$ of $\bT^{S_p,\univ}$.\;Let $\widehat{S}_{\xi,\tau}(U^{\fp},\ast)_{\fm_{\overline{\rho}}}$ be the localization of $\widehat{S}_{\xi,\tau}(U^{\fp},\ast)$ at $\fm_{\overline{\rho}}$ for $\ast\in\{\co_E,E\}$.\;
Then the action of $\bT^{S_p,\univ}$ on the localization $\widehat{S}_{\xi,\tau}(U^{\fp}, \co_E)_{\fm_{\overline{\rho}}}$ factors through certain  Hecke algebra
$\bT_{\xi,\tau}^{S_p}(U^{\fp}, \co_E)_{\fm_{\overline{\rho}}}$ (see \cite[Page 8029]{2019DINGSimple}).\;We also see that $\widehat{S}_{\xi,\tau}(U^{\fp},E)_{*}$ with $*\in \{\fm_{\overline{\rho}}, \emptyset\}$ are admissible unitary Banach representation of with invariant lattice $\widehat{S}_{\xi,\tau}(U^{\fp}, \co_E)_*$.\;

We denote by $R_{\widetilde{v}}^{\square}$ the maximal reduced and $p$-torsion free quotient of the universal $\co_E$-lifting ring of $\overline{\rho}_{\widetilde{v}}:=\overline{\rho}|_{\gal_{F_{\widetilde{v}}}}$ ($\cong \overline{r}$,\;and therefore $R_{\widetilde{v}}^{\square}\cong \defvarring$). For $v\in S_p\backslash\{\fp\}$, we denote by $R_{\widetilde{v}}^{\square, \xi,\tau}$ for the reduced and $p$-torsion free quotient of $R_{\widetilde{v}}^{\square}$ corresponding to potentially crystalline lifts of weight $\xi$ and inertial type $\tau$.\;Consider the following global deformation problem (in the terminology of \cite{clozel2008automorphy},\;see also \cite[Section 2.4]{PATCHING2016})
\begin{equation*}
\begin{aligned}
	\cS&=\bigg\{{F}/{F}^+,T^+,T,\cO_E,\overline{\rho},\chi_{\mathrm{cyc}}^{1-n}\delta_{{F}/{F}^+}^n,\{R_{\widetilde{v}_1}^{\square}\}\cup
	\{R_{\fp}^{\square}\}\cup \{R_{\widetilde{v}}^{\square, \xi,\tau}\}_{v\in S_p\backslash \{\fp\}}\bigg\}
\end{aligned}
\end{equation*}
Then by  \cite[Proposition 2.2.9]{clozel2008automorphy} (see also \cite[Section 2.4]{PATCHING2016}),\;this deformation problem is represented by a universal deformation ring $R_{\cS}^{\univ}$.\;Note that we have a natural morphism $R_{\cS}^{\univ}\rightarrow \bT_{\xi,\tau}^{S_p}(U^{\fp}, \co_E)_{\fm_{\overline{\rho}}}$.\;

Following \cite[Section 4.1.1]{2019DINGSimple} (or \cite[Section 2.8]{PATCHING2016}) we put
\begin{eqnarray*}
R^{\loc}:=R_{\widetilde{\fp}}^{\square} \widehat{\otimes} \Big(\widehat{\otimes}_{S_p\backslash\{\fp\}}R_{\widetilde{v}}^{\square, \xi,\tau}\Big)\widehat{\otimes} R_{\widetilde{v_1}}^{\square} ,
\end{eqnarray*}
where all completed tensor products are taken over $\cO_E$.\;We put $g:=q-[F^+:\BQ]\frac{n(n-1)}{2}$,\;where $q\geq [F^+:\BQ]\frac{n(n-1)}{2}$ is a certain integer as in  \cite[Section 2.6,\;Page 217]{PATCHING2016}).\;We now put 
\begin{eqnarray*}
R_{\infty}&:=&R^{\loc}\llbracket x_1,\cdots, x_g\rrbracket,\\
S_{\infty}&:=&\co_E\llbracket z_1,\cdots, z_{n^2(|S_p|+1)}, y_1,\cdots, y_q\rrbracket,
\end{eqnarray*}
where $x_i$, $y_i$, $z_i$ are formal variables.\;By the end of \cite[Section 4.1.1]{2019DINGSimple} (or \cite[Section 2.8]{PATCHING2016}),\;we get the following objects:
\begin{enumerate}
\item a continuous $R_{\infty}$-admissible unitary representation $\Pi_{\infty}$ of $G=\GLN_n(L)$ over $E$ together with a $G$-stable and $R_{\infty}$-stable unit ball $\Pi_{\infty}^o\subset \Pi_{\infty}$;
\item a morphism of local $\co_E$-algebras $S_{\infty}\ra R_{\infty}$ such that $M_{\infty}:= \homo_{\co_L}(\Pi_{\infty}^o, \co_E)$ is finite projective as $S_{\infty}\llbracket \GLN_n(\co_L)\rrbracket$-module;
\item a closed ideal $\fa$ of $R_{\infty}$, a surjection $R_{\infty}/\fa R_{\infty}\twoheadrightarrow R_{\cS}^{\univ}$ and a  $G   \times R_{\infty}/\fa R_{\infty}$-invariant isomorphism $\Pi_{\infty}[\fa]\cong \widehat{S}_{\xi,\tau}(U^{\fp},E)_{\fm_{\overline{\rho}}}$, where $R_{\infty}$ acts on $\widehat{S}_{\xi,\tau}(U^{\fp},E)_{\fm_{\overline{\rho}}}$ via $R_{\infty}/\fa R_{\infty}\twoheadrightarrow  R_{\cS}^{\univ}$.
\end{enumerate}

Let $R^{\fp}=\Big(\widehat{\otimes}_{v\in S_p\backslash\{\fp\}}R_{\widetilde{v}}^{\square, \xi,\tau}\Big)\widehat{\otimes} R_{\widetilde{v_1}}^{\square}$ and $R_{\infty}^{\fp}:=R^{\fp}\llbracket x_1,\cdots, x_g\rrbracket$.\;Then we have  $R^{\loc}=R^{\fp}\widehat{\otimes} \defvarring$ (recall that $R_{\widetilde{v}}^{\square}\cong \defvarring$ by definition) and $R_{\infty}=R_{\infty}^{\fp}\widehat{\otimes} \defvarring$.\;Let $\BU$ be the open unit ball in $\BA^1$.\;We put $\FX_{\overline{\rho}^{\fp}}^\Box:=(\Spf\;R^{\fp})^{\rig}$ and  $\mathfrak{X}_{\overline{r}}^\Box=(\Spf\;R_{\overline{r}}^\Box)^{\rig}$.\;Then we have $(\Spf\;R^{\fp}_{\infty})^{\rig}=\FX_{\overline{\rho}^{\fp}}^\Box\times \BU^g$.\;We have thus
$\FX_{\infty}:=(\Spf\;R_{\infty})^{\rig}\cong (\Spf\;R^{\fp}_{\infty})^{\rig}\times \mathfrak{X}_{\overline{r}}^\Box\cong \FX_{\overline{\rho}^{\fp}}^\Box\times \BU^g\times \mathfrak{X}_{\overline{r}}^\Box$.\;

%The patched module $\Pi_\infty$ defined in Section 4.1 is a $R_\infty$-module and 


Let $\bh:=(\hpi_{\tau,1},\hpi_{\tau,2},\cdots,\hpi_{\tau,n} )_{\tau\in \Sigma_L}$ be a strictly $\Delta_n^k$-dominant weight.\;We put $\hpi_{i}=(\hpi_{\tau,i})_{\tau\in \Sigma_L}$ for $1\leq i\leq n$ and put ${\bm\lambda}_\bh=(\hpi_{\tau,i}+i-1)_{\tau\in \Sigma_L,1\leq i\leq n}$.\;We denote by $\Pi_\infty^{R_\infty-\ana}$ the subrepresentation of $G$ of locally $R_\infty$-analytic vectors of $\Pi_\infty$ (see \cite[Section 3.1]{breuil2017interpretation}).\;Recall that $\bZ^{\lrr}(L)\cong Z^{\lrr}_{\varpi_L}\times \bZ^{\lrr}(\cO_L)$, where 
$Z^{\lrr}_{\varpi_L}$ is the image of $\oplus_{i=1}^k\BZ\hookrightarrow \bZ^{\lrr}(L),\;(m_i)\mapsto (\varpi_L^{m_i})$.\;

We recall the construction of certain  $R_{\infty}\times\FZ_{\omepik}\times Z^{\lrr}_{\varpi_L}\times \bZ^{\lrr}(\cO_L)$-module $B_{\omepik,{\bm\lambda}_\bh}(\Pi_\infty^{R_\infty-\ana})$ in \cite[Sections  3.1.1,\;3.1.2,\;3.3]{Ding2021}.\;Using Bushnell-Kutzko's theory of type,\;we can construct an absolutely irreducible smooth representation $\sigma$ of $\bL^{\lrr}(\cO_L)$ over $E$ from a maximal simple type of $\omepik$ (see the argument in \cite[Sections 3.1.1]{Ding2021}).\;We have $$\FZ_{\omepik}\cong \EndO_{\bL^{\lrr}(L)}\big(\mathrm{c-ind}_{\bL^{\lrr}(\cO_L)}^{\bL^{\lrr}(L)}\sigma\big),\;$$where ``$\mathrm{c-ind}$'' denotes the compact induction.\;We put 
\begin{equation*}
\begin{aligned}
	B_{\sigma,{\bm\lambda}_\bh}(\Pi_\infty^{R_\infty-\ana}):=&\;\homo_{\bL^{\lrr}(\cO_L)}\big(\sigma,J_{\bP^{\lrr}(L)}(\Pi_\infty^{R_\infty-\ana})_{{\bm\lambda}_\bh}\widehat{\otimes}_E\cC^{\BQ_p-\ana}(\bZ^{\lrr}(\cO_L),E)\big)\\
	\cong &\; \homo_{\bL^{\lrr}(L)}\big(c-\mathrm{ind}_{\bL^{\lrr}(\cO_L)}^{\bL^{\lrr}(L)}\sigma,J_{\bP^{\lrr}(L)}(\Pi_\infty^{R_\infty-\ana})_{{\bm\lambda}_\bh}\widehat{\otimes}_E\cC^{\BQ_p-\ana}(\bZ^{\lrr}(\cO_L),E)\big),\;
\end{aligned}
\end{equation*}
where $J_{\bP^{\lrr}(L)}(\Pi_\infty^{R_\infty-\ana})_{{\bm\lambda}_\bh}:=\homo_{\fd^{\lrr}}(L^{\lrr}({\bm\lambda}_\bh),J_{\bP^{\lrr}(L)}(\Pi_\infty^{R_\infty-\ana}))$ (recall that $\fd^{\lrr}$ is the Lie algebra of the derived subgroup $\bD^{\lrr}$ of $\bL^{\lrr}$,\;and $J_{\bP^{\lrr}(L)}$ is the Emerton-Jacquet functor \cite{emerton2006jacquet}).

As in \cite[Sections 3.1.2]{Ding2021},\;we recall various group actions on $B_{\sigma,{\bm\lambda}_\bh}(\Pi_\infty^{R_\infty-\ana})$.\;We write
\begin{equation}
\begin{aligned}
	&\;\iota_1(\bZ^{\lrr}(L))\cong \bZ^{\lrr}(L),\;\big(\text{i.e., the notation $\cZ_1$ in \cite[(3.4)]{Ding2021}})\\
	\text{resp.,}&\;\iota_0(\bZ^{\lrr}(\cO_L))\cong \bZ^{\lrr}(\cO_L),\;\big(\text{i.e., the notation $\cZ_0$ in \cite[(3.4)]{Ding2021}}\big)
\end{aligned}
\end{equation}
for  the action of  $\bZ^{\lrr}(L)$ on $B_{\sigma,{\bm\lambda}_\bh}(\Pi_\infty^{R_\infty-\ana})$ induced by the action $\bZ^{\lrr}(L)$ (resp.,\;$\bZ^{\lrr}(\cO_L)$) on \\ $J_{\bP^{\lrr}(L)}(\Pi_\infty^{R_\infty-\ana})_{{\bm\lambda}_\bh}$( resp.,\;$\cC^{\BQ_p-\ana}(\bZ^{\lrr}(\cO_L),E)$).\;

Next,\;the Bernstein centre $\FZ_{\omepik}$  acts on the module $B_{\sigma,{\bm\lambda}_\bh}(\Pi_\infty^{R_\infty-\ana})$ via the factor $c-\mathrm{ind}_{\bL^{\lrr}(\cO_L)}^{\bL^{\lrr}(L)}\sigma$, which commutes with the action of $\iota_1(\bZ^{\lrr}(L))\times \iota_0(\bZ^{\lrr}(\cO_L))$.\;Write $\cY_1$ (resp.,\;$\cY_0$) for the action $\iota_1(\bZ^{\lrr}(L))$ (resp., $\iota_0(\bZ^{\lrr}(\cO_L))$).\;

We write $\Delta_0$ for the action of  $\bZ^{\lrr}(L)$ on $B_{\sigma,{\bm\lambda}_\bh}(\Pi_\infty^{R_\infty-\ana})$ induced by the diagonal action of $\bZ^{\lrr}(L)$ on $J_{\bP^{\lrr}(L)}(\Pi_\infty^{R_\infty-\ana})_{{\bm\lambda}_\bh}\widehat{\otimes}_E\cC^{\BQ_p-\ana}(\bZ^{\lrr}(\cO_L),E)$ (see the argument before \cite[(3.4)]{Ding2021}).\;The argument before \cite[Lemma 3.1.2]{Ding2021} shows that the action of $\Delta_0$ is determined by $\Delta_0|_{Z^{\lrr}_{\varpi_L}}$ (and $\Delta_0|_{\bZ^{\lrr}(\cO_L)}$ acts via the central character of $\sigma$).

The $R_{\infty}\times\FZ_{\omepik}\times Z^{\lrr}_{\varpi_L}\times \bZ^{\lrr}(\cO_L)$-module structure of $B_{\omepik,{\bm\lambda}_\bh}(\Pi_\infty^{R_\infty-\ana})^\vee$ is given by the $R_{\infty}\times\FZ_{\omepik}\times\iota_1(Z^{\lrr}_{\varpi_L})\times \iota_0(\bZ^{\lrr}(\cO_L))$-action,\;where $R_{\infty}$ acts on $B_{\sigma,{\bm\lambda}_\bh}(\Pi_\infty^{R_\infty-\ana})$ via the factor $\Pi_\infty^{R_\infty-\ana}$.\;By \cite[Lemma 3.17]{Ding2021},\;the $\FZ_{\omepik}\times Z^{\lrr}_{\varpi_L}\times \bZ^{\lrr}(\cO_L)$-module $B_{\sigma,{\bm\lambda}_\bh}(\Pi_\infty^{R_\infty-\ana})$ does not depend on the choice $\sigma$.\;We denote hence $B_{\omepik,{\bm\lambda}_\bh}(\Pi_\infty^{R_\infty-\ana}):=B_{\sigma,{\bm\lambda}_\bh}(\Pi_\infty^{R_\infty-\ana})$.\;

We also recall how to recover the action $\cY_1$ from $\cY_0$-action and $\FZ_{\omepik}$-action.\;Denoting by $\psi_{\sigma}$ the central character of $\sigma$ (a character of $\bZ^{\lrr}(L)$).\;By \cite[(3.5)]{Ding2021},\;we see  that \[\cY_1(z_0)=\psi_{\sigma}(z_0)\cY_0(\mathrm{det}_{\bL^{\lrr}(L)}(z_0)).\;\]
for any $z_0\in \bZ^{\lrr}(\cO_L)$.\;On the other hand,\;we see that $\cY_1|_{Z^{\lrr}_{\varpi_L}}=\Delta_0|_{Z^{\lrr}_{\varpi_L}}$.\;Moreover, the $\bZ^{\lrr}(L)$-action on $B_{\omepik,{\bm\lambda}_\bh}(\widehat{S}(U^{\fp},W^{\fp})_{\overline{\rho}}^{\ana})$ induced by the map $$\bZ^{\lrr}(L)\rightarrow  \EndO_{\bL^{\lrr}(L)}\big(\mathrm{c-ind}_{\bL^{\lrr}(\cO_L)}^{\bL^{\lrr}(L)}\sigma\big)\cong\FZ_{\omepik}$$ coincides with the $\Delta_0$-action.\;Then the action $\cY_1$ can be recovered from the $\FZ_{\omepik}$-action.\;

By an easy variation of the proof of \cite[Lemma 3.1.3]{Ding2021} and the argument before \cite[(3.28)]{Ding2021},\;we see that $B_{\omepik,{\bm\lambda}_\bh}(\Pi_\infty^{R_\infty-\ana})^\vee$ is a coadmissible module over $\cO(\FX_{\infty}\times\zuni\times\rigch)$, which corresponds to a coherent sheaf $\cM_{\omepik,{\bm\lambda}_\bh}^{\infty,0}$ over $\FX_{\infty}\times \zuni\times \rigch$.\;Since the action of $Z^{\lrr}_{\varpi_L}$ factors through $\sbanpik$,\;we see that $\cM_{\omepik,{\bm\lambda}_\bh}^{\infty,0}$ gives rise to a coherent sheaf $\cM_{\omepik,{\bm\lambda}_\bh}^{\infty}$ over $\FX_{\infty}\times\sbanpik\times \rigch$ such that
\[\Gamma\Big(\FX_{\infty}\times\sbanpik\times \rigch,\cM_{\omepik,{\bm\lambda}_\bh}^{\infty}\Big)\cong B_{\omepik,{\bm\lambda}_\bh}(\Pi_\infty^{R_\infty-\ana})^\vee.\]
Let $\mathcal{E}_{\omepik,\fp,{\bm\lambda}_\bh}^\infty(\overline{\rho})\hookrightarrow\FX_{\infty}\times\sbanpik\times \rigch$ be the Zariski-closed support of $\cM_{\omepik,{\bm\lambda}_\bh}^{\infty}$.\;We call $\mathcal{E}_{\omepik,\fp,{\bm\lambda}_\bh}^\infty(\overline{\rho})$  the patched Bernstein eigenvariety.\;

By an easy variation of the proof of Proposition  3.3.2,\;Corollary \;3.3.3,\;Proposition\;3.3.4, Theorem 3.3.5 and Proposition\;3.3.6 in\cite{Ding2021}, we have
\begin{pro}\label{basicpropertyeigen}\hspace{20pt}
\begin{itemize}
	\item[(1)]  For $x=(\fm_x, \pi_{x},\chi_x)\in \FX_{\infty}\times\sbanpik\times \rigch$,\;$x\in \mathcal{E}_{\omepik,\fp,{\bm\lambda}_\bh}^\infty(\overline{\rho})$ if and only if
	\[\homo_{\bL^{\lrr}(L)}\Big(\pi_{x}\otimes_E((\chi_x)_{\varpi_{L}}\circ\mathrm{ det}_{\bL^{\lrr}(L)})\otimes_E L^{\lrr}({\bm\lambda}_\bh),J_{\bP^{\lrr}(L)}(\Pi_\infty^{R_\infty-\ana}[\fm_y] )\Big)\neq0\]
	\item[(2)] The rigid space $\mathcal{E}_{\omepik,\fp,{\bm\lambda}_\bh}^\infty(\overline{\rho})$ is reduced and equidimensional of dimension
	\[g+kd_L+n^2(|S_p|+1)+[F^+:\BQ]\frac{n(n-1)}{2}.\]
	\item[(3)] The coherent sheaf $\cM_{\omepik,{\bm\lambda}_\bh}^{\infty}$ is Cohen-Macaulay over $\mathcal{E}_{\omepik,\fp,{\bm\lambda}_\bh}^\infty(\overline{\rho})$.\;
	\item[(4)] The set of very classical non-critical generic points  (see \cite[Defintion 3.2.7,\;(3.22)]{Ding2021} for the definitions of classical,\;very classical,\;non-critical,\;and generic points respectively) is Zarisiki-dense in $\mathcal{E}_{\omepik,\fp,{\bm\lambda}_\bh}^\infty(\overline{\rho})$ and is an accumulation set.\;The set of very classical non-critical generic points  accumulates at point $x=(\fm_x, \pi_{x},\chi_x)$ with $\chi_x$ locally algebraic.\;
\end{itemize}
\end{pro}
%For $x=(\fm_x, \pi_{x},\chi_x)\in \mathcal{E}_{\omepik,\fp,{\bm\lambda}_\bh}^\infty(\overline{\rho})$,\;let $\rho_x$ be the $\gal_L$-representation associated to $x$ via $\fm_x\in (\Spf R_{\infty})^{\rig}=\FX_{\overline{\rho}^{\fp}}^\Box\times \BU^g\times \mathfrak{X}_{\overline{r}}^\Box$ in $\mathfrak{X}_{\overline{r}}^\Box$.\;Let $\mathbf{WD}_x=(\mathbf{WD}_{x,i})_{1\leq i\leq k}$ be the Weil-Deligne representation associated to $x$,\;i.e.,\;$\mathbf{WD}_{x,i}:=\rec(\pi_{x,i})(\frac{1-r}{2}-(i-1)r)$.\;
%We call a generic classical point $x$ is non-critical if $D^{\dag}_{\rig}(\rho_x)$ is non-critical in the sense of Section 2.4,\;with respect to the $\Omega$-filtration defined as in  \cite[(3.30),(3.24))]{Ding2021}.\;%Fix $\Omega$ a cuspidal Bernstein component of $\bL_{\Omega}(L)$,\;and denote by $\mathfrak{Z}_{\Omega}$ the corresponding Bernstein centre (over $E$).\;Let $V$ be an admissible locally analytic representation of $G$ over $E$.\;Then we see that $J_{\bP_\Omega}(V)$ is an essentially admissible representation of  $\bL_{\Omega}(L)$ (\cite{emerton2006jacquet}).\;Let $\underline{\lambda}$ be a integral $\bP_\Omega$-dominant weight of $G$,\;which corresponds an algebraic representation $L(\underline{\lambda})_\Omega$  of $\bL_{\Omega}(L)$ with weight $\underline{\lambda}$.\;We set
%\[J_{\bP_\Omega(L)}(V)_{\underline{\lambda}}:=\homo_{\sbanpik\times\rigch_{\Omega}^D}(L(\underline{\lambda})_\Omega,J_{\bP_\Omega(L)}(V)),\]
%which is closed $\bL_{\Omega}(L)$-subrepresentation of $J_{\bP_\Omega(L)}(V)\otimes_EL(\underline{\lambda})_\Omega^\vee$.\;Let $\sigma$ be  the irreducible smooth representation of $\bL_{\Omega}^0$ as in Section 2.3.\;We now put
%\[B_{\sigma,\ul{\lambda}}(V)\cong \big(J_{\bP_\Omega(L)}(V)_{\underline{\lambda}}\tene \cC^{\BQ_p-\an}(Z_{\bL_{\Omega}}^0,E)\ten \sigma^\vee\big)^{\bL_{\Omega}^0}.\]We explain the $\mathfrak{Z}_{\Omega}\times \iota_0(Z_{\bL_{\Omega}}(L))\times \iota_2(Z_{\bL_{\Omega}}^0)$ on $B_{\sigma,\ul{\lambda}}(V)$ as follows.\; Let $\iota_{2}(Z_{\bL_{\Omega}}(L))$ (resp.,\;$\iota_1(Z_{\bL_{\Omega}}^0)$) denotes the action of $Z_{\bL_{\Omega}}$ (resp.,\;$Z_{\bL_{\Omega}}^0$) on the$J_{\bP_\Omega(L)}(V)_{\underline{\lambda}}\tene \cC^{\BQ_p-\an}(Z_{\bL_{\Omega}}^0,E)$ by only acting on the term $J_{\bP_\Omega(L)}(V)_{\underline{\lambda}}$ (resp.,\;$\cC^{\BQ_p-\an}(Z_{\bL_{\Omega}}^0,E)$).\;The action $\iota_{2}(Z_{\bL_{\Omega}})\times \iota_1(Z_{\bL_{\Omega}}^0)$ is called standard action.\;On the other hand,\;we have\[B_{\sigma,\ul{\lambda}}(V)\cong \homo_{\bL_{\Omega}(L)}(\mathbf{c-ind}^{\bL_{\Omega}(L)}_{\bL_{\Omega}^0}\sigma,J_{\bP_\Omega(L)}(V)_{\underline{\lambda}}\tene \cC^{\BQ_p-\an}(Z_{\bL_{\Omega}}^0,E)).\]Then we deduce that $B_{\Omega,\ul{\lambda}}(V)$ is equipped with a natural action of $\mathfrak{Z}_{\Omega}\cong \EndO_{\bL_{\Omega}(L)}(\mathbf{c-ind}^{\bL_{\Omega}(L)}_{\bL_{\Omega}^0}\sigma)$.\;Meanwhile,\;we also use $\iota_0(Z_{\bL_{\Omega}}(L))$ to denotes the diagonal action of $Z_{\bL_{\Omega}}(L)$ on $B_{\sigma,\ul{\lambda}}(V)$,\;i.e.,\;with the natural action $Z_{\bL_{\Omega}}(L)$ on $J_{\bP_\Omega(L)}(V)_{\underline{\lambda}}\tene \cC^{\BQ_p-\an}(Z_{\bL_{\Omega}}^0,E)$ by acting via
%\[Z_{\bL_{\Omega}}(L)\xrightarrow{(\id,\;\pr_{\varpi_L}\circ\det_{\bL_{\Omega}}^{-1})}\iota_{2}(Z_{\bL_{\Omega}}(L))\times \iota_1(Z_{\bL_{\Omega}}^0).\]
%Then by definition,\;we see that $\iota_0(Z_{\bL_{\Omega}}(L))|_{Z_{\bL_{\Omega}}^0}=\psi_\sigma:=\sigma|_{Z_{\bL_{\Omega}}^0}$.\;This shows that
%$B_{\sigma,\ul{\lambda}}(V)$ is a closed subrepresentation of $\big(J_{\bP_\Omega(L)}(V)_{\underline{\lambda}}\tene \cC^{\BQ_p-\an}(Z_{\bL_{\Omega}}^0,E)\big)[Z_{\bL_{\Omega}}^0=\psi_\sigma]$.\;Therefore,\;we conclude that the $\iota_0(Z_{\bL_{\Omega}}(L))\times \iota_1(Z_{\bL_{\Omega}}^0)$ on $B_{\sigma,\ul{\lambda}}(V)$-action factors through the$\iota_0({Z_{\underline{\Omega}}})\times \iota_1(Z_{\bL_{\Omega}}^0)$.\;Since $\pr_{\varpi_L}\circ\det_{\bL_{\Omega}}^{-1}$ is trivial on $Z_{\varpi_L}$,\;we see that the restriction of standard action $\iota_{2}(Z_{\bL_{\Omega}})\times \iota_1(Z_{\bL_{\Omega}}^0)$ on $\iota_{2}(Z_{\varpi_L})$ coincides with the action $\iota_0(Z_{\bL_{\Omega}}(L))|_{Z_{\varpi_L}}=\iota_0(Z_{\varpi_L})$.\;This shows that the $\iota_0(Z_{\varpi_L})\times \iota_1(Z_{\bL_{\Omega}}^0)$-action coincide with standard action $\iota_{2}(Z_{\varpi_L})\times \iota_1(Z_{\bL_{\Omega}}^0)$.\;We can check that the natural action of $\mathfrak{Z}_{\Omega}$ on $B_{\sigma,\ul{\lambda}}(V)$ is commutes with the standard action of $\iota_0(Z_{\varpi_L})\times \iota_1(Z_{\bL_{\Omega}}^0)$.\;

%Furthermore,\;we have a natural map $Z_{\bL_{\Omega}}(L)\rightarrow  \EndO_{\bL_{\Omega}(L)}(\mathbf{c-ind}^{\bL_{\Omega}(L)}_{\bL_{\Omega}^0}\sigma)\cong\mathfrak{Z}_{\Omega}$ by setting $z$ to the homomorphism $v\mapsto z\cdot v$,\;for $v\in\mathbf{c-ind}^{\bL_{\Omega}(L)}_{\bL_{\Omega}^0}\sigma$.\;This also induce a $Z_{\bL_{\Omega}}(L)$-action on $B_{\sigma,\ul{\lambda}}(V)$ via the $\mathfrak{Z}_{\Omega}$-action on $B_{\sigma,\ul{\lambda}}(V)$.\;One can prove that this $Z_{\bL_{\Omega}}(L)$-action coincides with the $\iota_0(Z_{\bL_{\Omega}}(L))$-action.\;This shows that $\iota_0(Z_{\varpi_L})$-action factors through the $\mathfrak{Z}_{\Omega}$-action.\;

%Let $V$ be an admissible locally $\BQ_p$-analytic representation of $G$ over $E$,\;by the construction and \cite[Lem.\;3.1.3]{Ding2021},\;$B_{\sigma,\ul{\lambda}}(V)^\vee$ is a coadmissible module over $\cO(\zuni,E) \tene\cO(\rigch,E)$,\;and hence $B_{\sigma,\ul{\lambda}}(V)$ is an essentially admissible locally $\BQ_p$-analytic representation of $\iota_0(Z_{\varpi_L})\times \iota_1(Z_{\bL_{\Omega}}^0)$.\;The \cite[Lem.\;3.1.7]{Ding2021} shows that the $\mathfrak{Z}_{\Omega}\times \iota_0(Z_{\varpi_L})\times \iota_1(Z_{\bL_{\Omega}}^0)$-module $B_{\sigma,\ul{\lambda}}(V)$ does not depend on the choice $\sigma$.\;Therefore,\;we denote $B_{\Omega,\ul{\lambda}}(V):=B_{\sigma,\ul{\lambda}}(V)$.\;

%By the argument in \cite[Section 3.1.3]{Ding2021},\;we see that there exists a coherent sheaf $\cM_{\Omega,\ul{\lambda}}^0(V)$ over $\zuni\times \rigch$ such that $B_{\Omega,\ul{\lambda}}(V)\vee$ is isomorphic to the global sections of $\cM_{\Omega,\ul{\lambda}}^0(V)$.\;Since the action of $Z_{\underline{\omega}}$ factors  through $\sban$,\;we see that $\cM_{\Omega,\ul{\lambda}}^0(V)$ gives rise to a coherent sheaf $\cM_{\Omega,\ul{\lambda}}(V)$ over $\sban\times \rigerch$ such that \[\Gamma(\sban\times \rigerch,\cM_{\Omega,\ul{\lambda}}(V))\cong B_{\Omega,\ul{\lambda}}(V)^\vee.\]Let $\mathcal{E}_{\Omega,\ul{\lambda}}(V)$ be the Zariski-closed support of $\cM_{\Omega,\ul{\lambda}}(V)$.\;

We now recall the definition of Bernstein paraboline varieties \cite[Section 4.2]{Ding2021}.\;The Bernstein paraboline variety $\defvar$ of type {$(\omepik,\bh)$} is a subspace of $\mathfrak{X}_{\overline{r}}^\Box\times \sbanpik\times\rigch$.\;It contains a subspace $U^\Box_{\omepik,\mathbf{{h}}}(\overline{r})$  consists of the point $(\rho,\underline{x},\undelram)$ such that
\begin{itemize}\label{dfnvardef}
\item[(1)] $(\underline{x},\undelram)\in \big(\sbanpik\times\rigch\big)^{\gen}$ (the set of generic points in $\sbanpik\times\rigch$, see \cite[Section 4.2]{Ding2021}),
\item[(2)] $D_{\rig}(\rho)$ admits an $\omepik$-filtration $\cF=\fil_{\bullet}^{\cF} D_{\rig}(\rho)$ such that
\begin{equation}\label{paradefiinj}
	\begin{aligned}
		\gr_i^{\cF}D_{\rig}(\rho)\otimes_{\cR_{k(x),L}}\cR_{k(x),L}((\delta_i^0)^{-1}_{\varpi_L})\hookrightarrow \Delta_{x_i}\otimes_{\cR_{k(x),L}}\cR_{k(x),L}(z^{\bh_{ir}})
	\end{aligned}
\end{equation}
and the image has Hodge-Tate weights $(\bh_{(i-1)r+1},\cdots,\bh_{ir})$.\;By using Berger's equivalence of categories \cite[Theorem  A]{berger2008equations} and comparing the Hodge-Tate weights,\;we see that (\ref{paradefiinj}) is equivalent to 
\begin{equation}\label{dfnvardef1}
	\Delta_{x_i}\otimes_{\cR_{k(x),L}}\cR_{k(x),L}(z^{\bh_{(i-1)r+1}})  \hookrightarrow\gr_i^{\cF}D_{\rig}(\rho)\otimes_{\cR_{k(x),L}}\cR_{k(x),L}((\delta_i^0)^{-1}_{\varpi_L}).
\end{equation}
\end{itemize}
We define $\defvar$ to be the Zariski-closure of $U^\Box_{\omepik,\mathbf{{h}}}(\overline{r})$ in $\mathfrak{X}_{\overline{r}}^\Box\times \sbanpik\times\rigch$.\;By \cite[Theorem\;4.2.5,\;Corollary\;4.2.5]{Ding2021},\;we have:
\begin{pro}\label{propertyparavar}\hspace{20pt}
\begin{itemize}
	\item[(1)] The rigid space $\defvar$ is equidimensional of dimension $n^2+\left(\frac{n(n-1)}{2}+k\right)d_L$.\;
	\item[(2)] The set $U^\Box_{\omepik,\mathbf{{h}}}(\overline{r})$ is Zariski-open and  Zariski-dense in $\defvar$.
	\item[(3)] The rigid space  $U^\Box_{\omepik,\mathbf{{h}}}(\overline{r})$ is smooth over $E$,\;and the morphism $${\omega}|_{U^\Box_{\omepik,\mathbf{{h}}}(\overline{r})}:U^\Box_{\omepik,\mathbf{{h}}}(\overline{r})\rightarrow \sbanpik\times\rigch$$ is smooth.\;
	\item[(4)] Let $x=(\rho_x,\underline{x},\undelram)\in \defvar$,\;then $D_{\rig}(\rho_x)$ admits an $\omepik$-filtration $\cF=\{\fil_i^\cF D_{\rig}(\rho_x)\}$ such that,\;for all $1=1,\cdots,s$,\;
	\[\gr_{i}^{\cF}D_{\rig}(\rho_x)\otimes_{\cR_{k(x),L}}\cR_{k(x),L}((\delta_i^0)^{-1}_{\varpi_L})\bigg[\frac{1}{t}\bigg]=\Delta_{x_i}\bigg[\frac{1}{t}\bigg].\]
\end{itemize}
\end{pro}
\begin{rmk}In general,\;$(\underline{x},((\delta_i^0)_{\varpi_L}z^{\bh_{ir}}))$ is not a parameter (recall Definition \ref{dfnomegafil}) of the $\omepik$-filtration $\cF$ in (4).\;By definition,\;for $x=(\rho_x,\underline{x},\undelram)\in U^\Box_{\omepik,\mathbf{{h}}}(\overline{r})$,\;we see that $(\underline{x},((\delta_i^0)_{\varpi_L}z^{\bh_{ir}}))$ is a parameter of  $\omepik$-filtration  $\cF$.\;
\end{rmk}

The following definition is a generalization of the accumulation property in \cite[Definition 2.11]{breuil2017smoothness}.\;
\begin{dfn}\label{accum}Let $X$ be a union of irreducible components of an open subset of $\defvar$. We say that $X$ satisfies the accumulation property at $x$,\;if for any positive real number $C>0$,\;the set of potentially crystalline strictly points $x'=(\rho',\underline{x}',\delta^0_{x'})\in \defvar$ such that:
\begin{itemize}
	\item $x'$ is generic;
	\item $x'$ is non-critical;
	\item $\wt(\delta_{x,i}^0)-\wt(\delta_{x,i+1}^0)>C$ for $1\leq i\leq n-1$ and $\tau\in \Sigma_L$;
\end{itemize}
accumulate at $x$.\;
\end{dfn}


The (purely local) Bernstein paraboline variety is  closely related to the (global) patched Bernstein eigenvariety.\;Consider the composition
\begin{equation}\label{composition}
\begin{aligned}
	\bersteineigenvarpik\hooklongrightarrow \FX_{\overline{\rho}^{\fp}}^\Box&\times \mathfrak{X}_{\overline{r}}^\Box \times \sbanpik\times \rigch\\
	&\xrightarrow{\iota_{\omepik}} \FX_{\overline{\rho}^{\fp}}^\Box\times \mathfrak{X}_{\overline{r}}^\Box \times \sbanpik\times \rigch.
\end{aligned}
\end{equation}
where $\iota_{\omepik}:\sbanpik\xrightarrow{\sim}\sbanpik$ is the isomorphism such that \[\pi_{\iota_{\omepik}(x)_i}=\pi_{x_i}\otimes_E\unr(q_L^{(i-1)r+\frac{r-1}{2}})\circ\det\] for $x=(x_i)_{1\leq i\leq k}\in \sbanpik$.\;An easy variation of the proof of \cite[Theorem 3.3.9]{Ding2021} asserts that
\begin{pro}\label{linkingpatched}The composition in (\ref{composition}) factors through $\FX_{\overline{\rho}^{\fp}}^\Box\times \defvarrho$,\;i.e.,\;
\begin{equation}\label{mapbersteineigenvarpikdefvarrho}
	\Lambda:\bersteineigenvarpik\hooklongrightarrow \FX_{\overline{\rho}^{\fp}}^\Box\times \iota_{\omepik}^{-1}(\defvar).
\end{equation}	
It induces an isomorphism between $\bersteineigenvarpik$ with a union of irreducible components of the space $\FX_{\overline{\rho}^{\fp}}^\Box\times \iota_{\omepik}^{-1}(\defvar)$ equipped with the reduced closed rigid subspace structure.\;
\end{pro}
%\begin{rmk}\label{underthmlinkingpatched}In some sense,\;we can view deformation variety as purely local version of the patched Bernstein eigenvariety,\;But the constructions of  two objects are quite different.\;The first one is the support of the coherent sheaf defined by applying the ``Jacquet-Emerton'' functor to the locally analytic representation $\Pi_{\infty}^{R_{\infty}-\ana}$.\;In general,\;there is no purely local analogue of the patched module.\;\end{rmk}



%By definition,\;we have the following natural morphisms:
%\begin{equation}
%\begin{aligned}
%& {\omega}:\defvar\rightarrow \sbanpik\times\rigch,\;\\
%& \zeta:\defvar\rightarrow \mathfrak{X}_{\overline{r}}^\Box,\\
%\item ${\omega}=\varsigma_{\Omega}\circ\widetilde{\omega}:\defvar\rightarrow \sbanpik\times\rigch \rightarrow \rigerchl$,
%& \omega^1:\defvar\rightarrow \sbanpik\times\rigch\xrightarrow{\pr_1} \sban,\\
%& \omega^2:\defvar\rightarrow \sbanpik\times\rigch\xrightarrow{\pr_2} \rigerch.
%\end{aligned}
%\end{equation}
%Let $y=(\rho,\underline{x},\undelram)\in \defvar$,\;the tangent maps of the above morphisms are described as follows.\;
%\begin{itemize}
%\item $d\omega_y=d\omega^1_y\oplus d\omega^2_y:T_{\defvar,y}\rightarrow T_{\sbanpik\times\rigch,\omega(y)}= T_{\sban,\underline{x}}\oplus T_{\rigerch,\undelram}$.
%\item $d{\zeta}:T_{\defvar,x}\rightarrow T_{\mathfrak{X}_{\overline{r}}^\Box,\rho}$.\;Let $\bv\in T_{\defvar,x}$,\;seeing $d{\zeta}(\bv)$ as a element in $\ext^1(\rho,\rho)\cong \ext^1_{(\varphi,\Gamma)}(D_{\rig}(\rho),D_{\rig}(\rho))$ as an $k(x)[\epsilon]/\epsilon^2$-valued representation of $\gal_L$.\;Then the set of the Sen's $\tau$-weight of $d{\zeta}(\bv)$ is the set $\{\wt_\tau(\delta_{\bv,i}^0)+h_{j_i,\tau}|i=1,\cdots,r,j_i=s_{i-1}+1,\cdots.s_i\}$.\;        
%\item $d\widetilde{\omega}:T_{\defvar,x}\rightarrow T_{\sbanpik\times\rigch,\widetilde{\omega}(x)}$,\;
%\item $d\omega^1_y:T_{\defvar,x}\rightarrow T_{\sban,\underline{x}}$.\;
%For $d\omega_1:T_{\defvar,x}\rightarrow T_{\sban,\underline{x}}$,\;
%\item $d\omega_2:T_{\defvar,x}\rightarrow T_{\rigerch,\undelram}$.\;
%For $\bv\in T_{\defvar,x}$,\;we can see $\bv$ as a $E[\epsilon]/\epsilon^2$-valued point of $X$,\;and the composition
%\[\spr \big(E[\epsilon]/\epsilon^2\big)\xrightarrow{\bv} \defvar \rightarrow \rigerch\]gives rise to continuous character $\delta_{\bv}^0=\boxtimes_{i=1}^k\delta_{\bv,i}^0:{Z}_{L_\Omega}(\cO_L)\longrightarrow \big(E[\epsilon]/\epsilon^2\big)^\times$ such that $\delta_{\bv}^0\equiv \delta^0\mod \epsilon$.\;Then the map $d\omega_2$ sets $\bv\in T_{\defvar,x}$ to $(\wt_\tau(\delta_{\bv,i}^0(\delta_{i}^{0})^{-1}))_{1\leq i \leq k,\tau\in\Sigma_L}$.\;In \S 2.7,\;we see that $\{\psi_{\sigma,L}\}_{\sigma\in \Sigma_{L}}$ gives a basis of $\homo(\cO_L^{\times},E)$,\;then we can obtain the isomorphism $$T_{\rigerch,\undelram}\cong k(x)^{d_Lr}\xrightarrow{\sim}\homo({Z}_{L_\Omega}(\cO_L),k(x))=\prod_{i=1}^k\homo(\cO_L^{\times},E)$$ by sending $(a_{i,\tau})_{1\leq i\leq k,\tau\in\Sigma_L}$ to $\prod_{i=1}^k\big(\sum_{\tau\in\Sigma_L} a_{i,\tau}\psi_{i,\tau,L}\big)$,\;where.\;Using this isomorphism,\;the elements $d\omega_2(\bv)$ can be rewritten by $((\delta_{\bv,i}^0(\delta_{i}^{0})^{-1}-1)/\epsilon)_{1\leq i \leq k}$.\;
%\end{itemize}Let $X$ be a subspace of $\defvar$.\;The above discussion can be applied to $T_{X,x}$ by restriction.\;

%Since $\varsigma_\Omega$ is finite,\;the tangent map $d\varsigma_\Omega:T_{\sbanpik\times\rigch,\omega(x)}\xrightarrow{\sim} T_{\rigerchl,\omega(x)}$ induces an isomorphism.\;Therefore,\;we can identity $d\omega$ and $d\widetilde{\omega}$ via this isomorphism.\;

\subsection{Non-critical special point}\label{appearpoint}

Let  $\rho_L:\gal_L\rightarrow \GLN_{n}(E)$ be a potentially semistable non-crystalline representation.\;Let $\Dpik$ be the associated  
$(\varphi,\Gamma)$-module over $\cR_{E,L}$ of rank $n$.\;

Let  $\bh:=(\hpi_{\tau,1}>\hpi_{\tau,2}>\cdots>\hpi_{\tau,n} )_{\tau\in \Sigma_L}$ be the Hodge-Tate weights of $\rho_L$ (or $\Dpik$).\;We put $\hpi_{i}=(\hpi_{\tau,i})_{\tau\in \Sigma_L}$ for $1\leq i\leq n$.\;We put ${\bm\lambda}_\bh=(\hpi_{\tau,i}+i-1)_{\tau\in \Sigma_L,1\leq i\leq n}$,\;which is a dominant weight of $(\mathrm{Res}_{L/\BQ_p}\GLN_n)\times_{\BQ_p}E$ respect to $(\mathrm{Res}_{L/\BQ_p}\bB)\times_{\BQ_p}E$.\;In the sequel,\;we fix this weight $\bh$.\;

Suppose that $\Dpik$ admits a non-critical special  $\omepik$-filtration $\cF$ with parameter $$(\widetilde{\bx}_{\pi,\bh},\widetilde{\bm{\delta}}_\bh)\in \sbanpik\times\rigch $$ (or $(\bx_{\pi},\bmdel)\in \sbanpik\times\rigchl $),\;see (Definition \ref{weaklynoncritical},\;Definition \ref{dfnnoncriticalspecial}).\;In the sequel,\;we use the notations of Section \ref{noncrispecfildfn} freely.\;

Recall that $\pi\cong \pi_0\otimes\unr(\alpha_\pi )$ for some $\alpha_\pi \in E^\times$.\;We put 
\begin{equation}
\begin{aligned}
	\bullet\;	&(\breve{\mathbf{{x}}}_\pi,\underline{\mathbf{1}})\in \sbanpik\times\rigch,\\
	&\breve{\mathbf{{x}}}_\pi=(\breve{\mathbf{{x}}}_{\pi,i})_{1\leq i\leq k},\;\text{and\;}\pi_{\breve{\mathbf{{x}}}_{\pi,i}}\cong \pi_0\te\unr\big(\alpha_\pi q_L^{i-k+\frac{1-r}{2}-r(i-1)}\big)\circ\det \text{\;for\;}  1\leq i\leq k;\hspace{40pt}\\
	\bullet\;	&(\widehat{\bx}_\pi,\underline{\mathbf{1}})\in \sbanpik\times\rigch,\\
	&\widehat{\bx}_\pi=(\bx_{\pi,i})_{1\leq i\leq k},\text{and\;}\pi_{\widehat{\bx}_{\pi,i}}\cong \pi_0\te\unr(\alpha_\pi q_L^{i-k})\circ\det \text{\;for\;}  1\leq i\leq k.\;
\end{aligned}
\end{equation}

Suppose that there exists $x^{\fp}\in \FX_{\overline{\rho}^{\fp}}^\Box$ such that
$$x=(x^{\fp}, \rho_L, \breve{\mathbf{{x}}}_\pi,\underline{\mathbf{1}})\in \bersteineigenvarpik \hooklongrightarrow \FX_{\overline{\rho}^{\fp}}^\Box\times \mathfrak{X}_{\overline{r}}^\Box \times \sbanpik\times \rigch,\;$$
By definition,\;we see that $$\Lambda(x)=(x^{\fp},\iota_{\omepik}^{-1}(x_L))=(x^{\fp},\iota_{\omepik}^{-1}(\rho_L,\widehat{\bx}_\pi,\mathbf{\underline{1}}))\in \FX_{\overline{\rho}^{\fp}}^\Box\times \iota_{\omepik}^{-1}(\defvar)$$ via the injection (\ref{mapbersteineigenvarpikdefvarrho}).\;In this case,\;we call $x_L$ (resp.,\;$x$) a \textit{non-critical special point} on $\defvarrho$ (resp.,\;$\bersteineigenvarpik$).\;

%By Section 2.5,\;we attach to $\rho_L$ the parabolic simple $\sL$-invariants $\sL(\rho_L)$.

We derive several local consequences of Bernstein paraboline variety $\defvarrho$ at
the non-critical special point $x_L$.\;Recall $\complocalbersteineigenvarpik$ (resp.\;$\complocaldefvarrho$)
means the complete local ring of $\bersteineigenvarpik$ (resp.\;$\defvarrho$) at $x$ (resp.\;$x_L$).\;

\subsection{Accumulation property at non-critical special point}\label{ACC}
Let $\cE(x)$ be the union of  irreducible components of $\bersteineigenvarpik$ containing $x$,\;which is thus equidimensional of dimension \[g+nd_L+n^2(|S_p|+1)+[F^+:\BQ]\Big(\frac{n(n-1)}{2}+r\Big).\;\]Since $\bersteineigenvarpik$ is reduced at $x$ (by Proposition \ref{basicpropertyeigen} (2)),\;one has
\begin{equation}\label{OZlocal}
\widehat{\cO}_{\bersteineigenvarpik,x}\cong \widehat{\cO}_{\cE(x),x}.\;
\end{equation}
By Theorem \ref{linkingpatched},\;$\cE(x)$ has the form $\cup_{i,j} \big(X_i^{\fp}\times \BU^g \times \iota_{\omepik}^{-1}(X_{j,\fp})\big)$,\;where $X_{i,\fp}$ is an irreducible component of $\defvarrho$ containing $x_L$, and $X_i^{\fp}$ is an irreducible component of $\FX_{\overline{\rho}^{\fp}}^\Box$.\;By \cite[Theorem 3.3.8]{kisin2008potentially} and \cite[Lemma 2.5]{PATCHING2016},\;$\defvarrho$ is  smooth at $x^{\fp}$, and hence $\{X_i^{\fp}\}_i$ is a singleton $\{X^{\fp}\}$.\;Therefore,\;we have
\begin{equation}\label{OZfactor}
\cE(x)=X^{\fp} \times \BU^g \times \iota_{\omepik}^{-1}(\xfpx)\hookrightarrow X^{\fp} \times \BU^g \times \iota_{\omepik}^{-1}(\defvarrho)
\end{equation}
with $\xfpx=\cup_j X_{j,\fp}$.\;

\begin{lem}\label{accumulation} $\xfpx$ satisfies the accumulation property at $x_L$.
\end{lem}
\begin{proof}We follow the route of \cite[Proposition 3.10]{breuil2017smoothness},\;the proof of \cite[Theorem 3.9]{breuil2017interpretation} and the statement in \cite[Theorem 3.3.5]{Ding2021}.\;Recall that \cite[Theorem 3.3.5]{Ding2021} say that the set of very classical non-critical generic points (and therefore potentially crystalline) is Zariski-dense in $\bersteineigenvarpik$ and accumulate at any point $x=(\fm_x, \pi_{x},\chi_x)\in \bersteineigenvarpik$  with $\chi_x$ locally algebraic.\;We adapt the proof of this statement to our setting.\;More precisely,\;for the non-critical point $x=(x^{\fp}, \rho_L, \breve{\mathbf{{x}}}_\pi,\underline{\mathbf{1}})\in \bersteineigenvarpik $
and an irreducible component $X$ of $\bersteineigenvarpik $ containing $x$,\;
one can choose an affinoid and connected neighborhood of $x$ such that $U\rightarrow \omega_X(U)$ is a finite surjective morphism,\;and $\omega_X(U)$ is an affinoid open subset,\;by \cite[Proposition 3.3.2 (2)]{Ding2021}.\;For any sufficiently large $C$,\;we define a part $W_C$ of 
$\rigch$ consists of dominant algebraic character $\delta^0\in \rigch$ such that $\wt(\delta_{\tau,i}^0)-\wt(\delta_{\tau,i+1}^0)>C$ for $1\leq i\leq k-1$ and $\tau\in \Sigma_L$.\;Since $U$ is affinoid,\;by \cite[Proposition 3.2.14]{Ding2021} and \cite[Proposition 3.2.9 and (3.17)]{Ding2021},\;for any sufficiently large $C$,\;if a point $x'=(x'^\fp,\rho',\underline{x}',\delta^0_{x'})\in \bersteineigenvarpik$ satisfies that $\delta^0_{x'}\in W_C$\;then $x'$ is a very classical non-critical generic point.\;Then by a similar  argument of the second paragraph in the proof of \cite[Theorem 3.9]{breuil2017interpretation},\;we see that there exists an affinoid neighbourhood $U$ of $x$ in $X$ and a part $W$ of $\rigch$ such that $\omega_X^{-1}((\spf S_\infty)^\rig\times W)\cap U$ consists of very classical non-critical generic points,\;where $\omega_X$ is the restriction of the map \[\bersteineigenvarpik\xrightarrow{\kappa_z}((\Spf\;S_\infty)^\rig\times \rigch\times \BG_m^{\rig}\rightarrow (\Spf\;S_\infty)^\rig\times \rigch,\]which is defined in the \cite[discussion before Proposition 3.3.2]{Ding2021}.\;Using the isomorphism \[\cE(x)=X^{\fp} \times \BU^g \times \iota_{\omepik}^{-1}(\xfpx),\;\]this lemma is contained in the above discussion.\;
\end{proof}

Let $r^{\univer}$ be the universal framed Galois deformation of $\overline{r}$ over $\mathfrak{X}_{\overline{r}}^\Box$,\;$\delta_i^{0,\univer}$ be the universal character of $\cO_L^\times$ over $\widehat{\cO_L^\times}$,\;and $\Delta^{\univer}_{\Omega_r}$ be the universal $p$-adic differential equation over $\sbanpi$ (see \cite[Section 2.2]{Ding2021}).\;Let $X\subset \xfpx$ be a subspace,\;and let $r^{\univer}_{X}$,\;$\delta_{X,i}^{0,\univer}$ and $\Delta^{\univer}_{\Omega_r,X}$ be the pull-back of $r^{\univer}$,\;$\delta_{i}^{0,\univer}$ and $\Delta^{\univer}_{\pi}$ over $X$.\;Let $D_{\rig}\big(r^{\univer}_X\big)$ be the $(\varphi,\Gamma)$-module over $\cR_{X,L}$ associated to $r^{\univer}_{X}$.\;

For $1\leq i\leq k$,\;we set
\[\Delta_{X,i}=\Delta^{\univer}_{\Omega_r,X}\otimes_{\cR_{X,L}}\cR_{X,L}\Big(\big(\delta_{X,i}^{0,\univer}\big)_{\varpi_L}z^{{\bh}_{(i-1)r+1}}\Big).\]
By the argument after (\ref{Dpikinjection}),\;we have natural injection $\Delta_{X,x_L,i}\hookrightarrow \gr^{\cF}_i\Dpik\cong \gr_i^{\cF}D_{\rig}(r^{\univer}_{X,x_L})$ for $1\leq i\leq k$, where $\Delta_{X,x_L,i}$ is the specialization of $\Delta_{X,i}$ at $x_L$.\;We show that the  $\omepik$-filtration on $x_L$ can extend to some open affinoid neighborhood around $x_L$.\;

\begin{thm}\label{localfiltration} There exists an open affinoid neighborhood $X\subset \xfpx$ of non-critical special point $x_L$ such that the $(\varphi,\Gamma)$-module $D_{\rig}(r^{\univer}_X)[\frac{1}{t}]$ over $\cR_{X,L}$ admits a filtration $\cM_\bullet$ such that $\gr^i\cM_\bullet\cong \Delta_{X,i}[\frac{1}{t}]$.\;In particular,\;the specialization of the filtration $\cM_\bullet$ on non-critical special point $x_L$ gives an $\omepik$-filtration of $D_{\rig}(\rho_L)[\frac{1}{t}]$  with parameter $({\bx}_\pi,\bm{\delta}_{\bh})$.\;
\end{thm}
\begin{proof}
By the accumulation property at $x_L$,\;there exists a Zariski dense set $S$ of potentially crystalline and non-critical generic points such that for $z\in S$,\;we have (see the proof of \cite[Proposition 4.2.7]{Ding2021})
\[\dim_{k(z)}\ext^i_{(\varphi,\Gamma)}(\Delta_{X,z,1},D_{\rig}(r^{\univer}_{X,z}))=\left\{
\begin{array}{ll}
	1, & \hbox{$i=0$;} \\
	1+nr\cdot d_L, & \hbox{$i=1$;} \\
	0, & \hbox{$i=2$.}
\end{array}
\right.\]
Similar to the proof of  Lemma \ref{phigammacohononcrtialspecial} and Lemma \ref{pullbackforiso} (recall that $\gr_1^{\cF}D_{\rig}(r^{\univer}_{X,x_L})/\Delta_{X,x_L,1}$ is a torsion $(\varphi,\Gamma)$-module),\;we can deduce that for $s>1$,\;
\begin{description}
	\item[-] $\ext^0_{(\varphi,\Gamma)}(\Delta_{X,x_L,1},\gr_s^{\cF}D_{\rig}(r^{\univer}_{X,x_L}))=0$,\;
	\item[-] $\ext^2_{(\varphi,\Gamma)}(\Delta_{X,x_L,1},\gr_s^{\cF}D_{\rig}(r^{\univer}_{X,x_L}))=0$,\;
	\item[-] $\ext^0_{(\varphi,\Gamma)}(\Delta_{X,x_L,1},\gr_1^{\cF}D_{\rig}(r^{\univer}_{X,x_L})/\Delta_{X,x_L,1})=0$,\;
	\item[-] $\ext^2_{(\varphi,\Gamma)}(\Delta_{X,x_L,1},\gr_1^{\cF}D_{\rig}(r^{\univer}_{X,x_L})/\Delta_{X,x_L,1})=0$.\;
\end{description}
An easy computation and d\'{e}vissage argument imply
\[\dim_{k(x_L)}\ext^i_{(\varphi,\Gamma)}(\Delta_{X,x_L,1},D_{\rig}(r^{\univer}_{X,x_L}))=\left\{
\begin{array}{ll}
	1, & \hbox{$i=0$;} \\
	1+nr\cdot d_L, & \hbox{$i=1$;} \\
	0, & \hbox{$i=2$.}
\end{array}
\right.\]
By \cite[Corollary 4.7]{bergdall_2017},\;$\ext^0_{(\varphi,\Gamma)}(\Delta_{X,1},D_{\rig}(r^{\univer}_X))$ is locally free of rank one at $x_L$.\;Shrinking $X$,\;we can assume that $\ext^0_{(\varphi,\Gamma)}(\Delta_{X,1},D_{\rig}(r^{\univer}_X))$ is free of rank one at $x_L$.\;A choice of a generator allows a map
\[\iota_1:\Delta_{X,1}\longrightarrow D_{\rig}(r^{\univer}_X).\]
By \cite[Lemma 5.4 (a)]{bergdall_2017},\;$\iota_1$ is injective.\;

Now consider $Q_{X,1}:=\cokerr \iota_1$,\;which is a generalized $(\varphi,\Gamma)$-module over $\cR_{X,L}$.\;By \cite[Lemma 5.4 (b)]{bergdall_2017},\;we can assume $Q_{X,1}$ is nearly flat (in the sense of the \cite[Definition  4.1]{bergdall_2017}) by shrinking $X$.\;A difference with trianguline case is that the map $\iota_1$ is never saturated,\;and so $Q_{X,1}$ has $t$-torsions at all points.\;For any $z\in S$ (resp.,\;$x_L$),\;by definition $Q_{X,z}$ (resp.,\;$Q_{X,x_L,1}$) is an extension of
\begin{equation}
	\begin{aligned}
		&\fil_2^kD_{\rig}(r^{\univer}_{X,z}):=D_{\rig}(r^{\univer}_{X,z})/\gr_1^{\cF}D_{\rig}(r^{\univer}_{X,z})\;\\
		(\text{resp.},\;&\fil_2^kD_{\rig}(r^{\univer}_{X,x_L}):=D_{\rig}(r^{\univer}_{X,x_L})/\gr_1^{\cF}D_{\rig}(r^{\univer}_{X,x_L}))
	\end{aligned}
\end{equation}
by torsion $(\varphi,\Gamma)$-module $\gr_1^{\cF}D_{\rig}(r^{\univer}_{X,z})/\Delta_{X,z,1}$ (resp.,\;$\gr_1^{\cF}D_{\rig}(r^{\univer}_{X,x_L})/\Delta_{X,x_L,1}$).\;Moreover,\;as $z\in S$ is generic,\;we have  $$\ext^1_{(\varphi,\Gamma)}(\fil_2^kD_{\rig}(r^{\univer}_{X,z}),\gr_1^{\cF}D_{\rig}(r^{\univer}_{X,z})/\Delta_{X,z,1})=0$$ and hence $Q_{X,z,1}\cong \fil_2^kD_{\rig}(r^{\univer}_{X,z})\oplus \gr_1^{\cF}D_{\rig}(r^{\univer}_{X,z})/\Delta_{X,z,1}$.\;

For the point $x_L$,\;since $Q_{X,x_L}[\frac{1}{t}]=\fil_2^kD_{\rig}(r^{\univer}_{X,x_L})[\frac{1}{t}]$,\;then there exists a sufficiently large integer $N$,\;such that (using $E$-$B$-pairs and the same strategy as in the proof of  Lemma \ref{phigammacohononcrtialspecial} and Lemma \ref{pullbackforiso})
\[\ext^1_{(\varphi,\Gamma)}(t^N\Delta_{X,x_L,2},Q_{X,x_L,1})\xrightarrow{\sim}\ext^1_{(\varphi,\Gamma)}(t^N\Delta_{X,x_L,2},\fil_2^kD_{\rig}(r^{\univer}_{X,x_L})).\]
In this case,\;note that for $s>2$,\;
\begin{description}
	\item[-] $\ext^0_{(\varphi,\Gamma)}(t^N\Delta_{X,x_L,2},\gr_s^{\cF}D_{\rig}(r^{\univer}_{X,x_L}))=0$,\;
	\item[-] $\ext^2_{(\varphi,\Gamma)}(t^N\Delta_{X,x_L,2},\gr_s^{\cF}D_{\rig}(r^{\univer}_{X,x_L}))=0$,\;
	\item[-] $\ext^0_{(\varphi,\Gamma)}(t^N\Delta_{X,x_L,2},\gr_2^{\cF}D_{\rig}(r^{\univer}_{X,x_L})/\Delta_{X,x_L,2})=0$,\;
	\item[-] $\ext^2_{(\varphi,\Gamma)}(t^N\Delta_{X,x_L,2},\gr_2^{\cF}D_{\rig}(r^{\univer}_{X,x_L})/\Delta_{X,x_L,2})=0$.\;
\end{description}
We can obtain
\begin{equation}
	\begin{aligned}
		&\dim_E\ext^1_{(\varphi,\Gamma)}(t^N\Delta_{X,x_L,2},\fil_2^kD_{\rig}(r^{\univer}_{X,x_L}))=\left\{
		\begin{array}{ll}
			1, & \hbox{$i=0$;} \\
			1+(k-1)r^2\cdot d_L, & \hbox{$i=1$;} \\
			0, & \hbox{$i=2$.}
		\end{array}
		\right.
	\end{aligned}
\end{equation}
Fix this integer $N$,\;for any $z\in S$,\;we claim that if $\wt(\delta_{1,z}^0)-\wt(\delta_{2,z}^0)>N$,\;then we have
\[\ext^1_{(\varphi,\Gamma)}(t^N\Delta_{X,z,2},\gr_1^{\cF}D_{\rig}(r^{\univer}_{X,z})/\Delta_{X,z,1})=0.\]
Therefore,\;we can obtain
\begin{equation}
	\begin{aligned}
		&\dim_E\ext^i_{(\varphi,\Gamma)}(t^N\Delta_{X,z,2},Q_{X,z,1})=\ext^1_{(\varphi,\Gamma)}(t^N\Delta_{X,z,2},\fil_2^kD_{\rig}(r^{\univer}_{X,z}))\\
		=\;&\left\{
		\begin{array}{ll}
			1, & \hbox{$i=0$;} \\
			1+(k-1)r^2\cdot d_L, & \hbox{$i=1$;} \\
			0, & \hbox{$i=2$.}
		\end{array}
		\right.
	\end{aligned}
\end{equation}
By the  accumulation property at $x_L$,\;the potentially crystalline and non-critical generic points $z\in S$ such that $\wt(\delta_{1,z}^0)-\wt(\delta_{2,z}^0)>N$ is Zariski dense in $X$.\;By \cite[Corollary 5.4(b)]{bergdall_2017},\;we see that\\ $\ext^0_{(\varphi,\Gamma)}(t^N\Delta_{X,2},Q_{X,1})$ is locally free of rank one at $x_L$.\;By shrinking $X$,\;we can further assume that $\ext^0_{(\varphi,\Gamma)}(t^N\Delta_{X,2},Q_{X,1})$ is free of rank one at $x_L$.\;A choice of a generator allows a map
\[\iota_2:t^N\Delta_{X,2}\longrightarrow Q_{X,1}.\]
By \cite[Corollary 5.4(a)]{bergdall_2017},\;$\iota_2$ is injective.\;Now consider $Q_{X,2}:=\cokerr \iota_2$.\;Proceeding as we did for $Q_{X,1}:=\cokerr \iota_1$,\;we complete the proof step by step.\;In conclusion,\;the $(\varphi,\Gamma)$-module $D_{\rig}(r^{\univer}_X)[\frac{1}{t}]$ over $\cR_{X,L}$ admits a filtration $\cM_{\bullet}$ such that $\gr_i^{\cM_{\bullet}}D_{\rig}(r^{\univer}_X)[\frac{1}{t}]\cong \Delta_{X,i}[\frac{1}{t}]$.\;The specialization of the filtration $\cM_\bullet$ at non-critical special point $x_L$ gives a $\omepik$-filtration of $D_{\rig}(\rho_L)[\frac{1}{t}]$ such that
\begin{equation}
	\begin{aligned}
		\gr_i^{\cM_\bullet}D_{\rig}(\rho)[\frac{1}{t}]\xrightarrow{\sim}\Delta_{\widehat{\bx}_{\pi,i}}\tee \cR_{E,L}(z^{\bh_{(i-1)r+1}})[\frac{1}{t}]\cong\Delta_{\widehat{\bx}_{\pi,i}}\tee \cR_{E,L}(z^{{\bh}_{ir}})[\frac{1}{t}].
	\end{aligned}
\end{equation}
This proves the last assertion.\;
\end{proof}

\subsection{Smoothness of non-critical special point}\label{smoothpoint}
Recall that the deformation theory of Galois representations is usually studied by considering a functor whose values are the set of isomorphism classes of liftings.\;In some cases,\;it is better to consider a functor with values in groupoids,\;i.e.,\;to consider the category of liftings and the isomorphisms between them.\;If $X$ is a groupoid over $\Art_E$,\;we denote by $|X|$ the functor on  $\Art_E$ such that $|X|(A)$ is the set of isomorphim classes of the category $X(A)$.\;In this section,\;we study certain groupoids or functors whose values are  deformations of type $\omepik$ for $(\varphi,\Gamma)$-module $\Dpik$ over $\cR_{E,L}$ (and  $\cM_\Dpik:=\Dpik[\frac{1}{t}]$ over $\cR_{E,L}[\frac{1}{t}]$).\;

We also fix the $\omepik$-filtration  $\cF$ on $\Dpik$ with non-critical special parameter \[(\bx_{\pi},\bm{\delta}_{\bh})\in \sbanpik\times\rigchl.\;\]By inverting $t$,\;this filtration $\cF$ on $\Dpik$ induces an increasing filtration $\cM_{\bullet}=\{\fil_{i}^\cF \Dpik[\frac{1}{t}]\}$ on $\cM_\Dpik$ by $(\varphi,\Gamma)$-submodules over $\cR_{E,L}[\frac{1}{t}]$.\;Under the terminology in \cite[Section 6.2]{Ding2021},\;a parameter of $\omepik$-filtration $\cM_{\bullet}$ of $\cM_\Dpik$ is $(\bx_{\pi},\bm{\delta}_{\bh})\in \sbanpik\times\rigchl $.\;By \cite[Lemmma 6.2.1]{Ding2021},\;we see that all parameters of $\cM_{\bullet}$ are of the form $(\bx_{\pi}',\bm{\delta}_{\bh}')$ such that,\;for $i=1,\cdots,k$,\; $\Delta_{{\bx}_{\pi,i}'}=\Delta_{{\bx}_{\pi,i}}\otimes_{\cR_{E,L}} \cR_{E,L}(\psi_i)$ and $\bm{\delta}_{\bh,i}'=\bm{\delta}_{\bh,i}\psi_i^{-1}\eta_iz^\bk$ for some unramified character $\psi_i$ of $L^\times$,\;$\eta_i\in \mu_{\Omega_\pi}$,\;and $\bk\in \BZ^{|\Sigma_L|}$.\;

Let $X_{\cM_\Dpik}$ be the groupoid over $\Art_E$ of deformation of $\cM_\Dpik$,\;and $X_{\cM_\Dpik,\cM_{\bullet}}$ be the groupoid over $\Art_E$ of $\omepik$- deformation of $(\cM_\Dpik,\cM_{\bullet})$ (see \cite[Section 6.2]{Ding2021}).\;There is a natural morphism (by forgetting the filtration) $X_{\cM_\Dpik,\cM_{\bullet}}\rightarrow X_{\cM_\Dpik} $.\;Let $X_{\Dpik}$  be the groupoid over $\Art_E$ of deformation of $\Dpik$.\;Recall that we have natural morphism $X_{\Dpik}\rightarrow X_{\cM_\Dpik}$
by inverting $t$.\;More precisely,\;note that $|X_{\cM_\Dpik,\cM_{\bullet}}|$ is the deformation functor
\[|X_{\cM_\Dpik,\cM_{\bullet}}|:\Art_E:=\{\text{Artinian local $E$-algebra with residue field $E$}\}\longrightarrow \{\text{sets}\}\]
sends $A$ to the set of isomorphism classes $\{(\cM_A,j_A,\cM_{A,\bullet})\}/\sim$,\;where
\begin{description}
\item[(1)] $\cM_A$ is a $(\varphi,\Gamma)$-module of rank $n$ over $\cR_{A,L}$ with an isomorphism $j_A:\cM_A\otimes_AE\cong \cM_\Dpik$,
\item[(2)] $\cM_{A,\bullet}$ is an increasing  $\omepik$-filtration of $(\varphi,\Gamma)$-module over $\cR_{A,L}[\frac{1}{t}]$ on $\cM_A$,
\item[(3)] $j_A$ induces isomorphisms $j_A:\cM_{A,i}\otimes_AE\cong \cM_i$.\;
\end{description}
By \cite[Lemma 6.2.2]{Ding2021},\;there exists a unique characters $\delta_A=\otimes_{i=1}^k{\delta}_{A,i}$ such that ${\delta}_{A,i}\equiv \bm{\delta}_{\bh,i}(\mod \mathfrak{m}_A)$ and 
$(\bx_{\pi}:=(\bx_{\pi,i})_{1\leq i\leq k},\delta_A)$ is a parameter of $\cM_{A,\bullet}$,\;i.e.,\;there exists an isomorphism of $(\varphi,\Gamma)$-module of rank $r$ over $\cR_{A,L}[\frac{1}{t}]$:\;$\gr_i\cM_A\xrightarrow{\sim}{\Delta_\pi}\otimes_{\cR_{E,L}}\cR_{A,L}({\delta}_{A,i})[\frac{1}{t}]$.\;We put $X_{\Dpik,\cM_{\bullet}}:=X_{\Dpik}\times_{X_{\cM_\Dpik}}X_{\cM_\Dpik,\cM_{\bullet}}$.\;

We view the character $\bm{\delta}_{\bh}$ as a point of $\cZ_{\bL^{\lrr},L}$.\;Observe that the functor
\[A\in \Art_E\mapsto \{\delta_A:\bZ^{\lrr}(L)\rightarrow A^\times,\;{\delta}_{A,i}\equiv \bm{\delta}_{\bh,i}(\mathrm{mod}\;\mathfrak{m}_A)\}\]
is pro-representable by $\widehat{(\cZ_{\bL^{\lrr},L})}_{\bm{\delta}_{\bh}}$,\;where $\widehat{(\cZ_{\bL^{\lrr},L})}_{\bm{\delta}_{\bh}}$ is the completion of $\cZ_{\bL^{\lrr},L}$ at the point $\bm{\delta}_{\bh}$.\;Then we have a morphism  of groupoids over $\Art_E$:
\begin{equation}\label{omegadeltapinon}
\omega_{\bm{\delta}_{\bh}}:X_{\cM_\Dpik,\cM_{\bullet}}\rightarrow \widehat{(\cZ_{\bL^{\lrr},L})}_{\bm{\delta}_{\bh}},\;\;(A,\cM_A,j_A,\cM_{A,\bullet})\mapsto \delta_A,
\end{equation}
by \cite[Lemma 6.2.2]{Ding2021}.\;

The $E$-linear map $\delta_{E[\epsilon]/\epsilon^2}\mapsto (\delta_{E[\epsilon]/\epsilon^2}\bm{\delta}_{\bh}^{-1}-1)/\epsilon$ induces an isomorphism $$\widehat{(\cZ_{\bL^{\lrr},L})}_{\bm{\delta}_{\bh}}(E[\epsilon]/\epsilon^2)\xrightarrow{\sim} \homo(\bZ^{\lrr}(L),E)=\prod_{i=1}^k\homo(L^\times,E).$$Therefore,\;we
consider the following composition
\begin{equation}\label{omegadeltapi}
\omega^{\kappa}_{\bm{\delta}_{\bh}}:X_{\cM_\Dpik,\cM_{\bullet}}(E[\epsilon]/\epsilon^2)\rightarrow \widehat{(\cZ_{\bL^{\lrr},L})}_{\bm{\delta}_{\bh}}(E[\epsilon]/\epsilon^2) \xrightarrow{\kappa_L}
\prod_{ir\in \Delta_n(k)}\homo(L^\times,E),
\end{equation}
where the last map $\kappa_L$ sends $(\psi_1,\psi_2,\cdots,\psi_k)$ to $(\psi_i-\psi_{i+1})_{ir\in \Delta_n(k)}$ (see \ref{kappakappaL}).\;The composition of the morphism $\omega_{\bm{\delta}_{\bh}}$ \big(resp.,\;the map $\omega^{\kappa}_{\bm{\delta}_{\bh}}$ \big) with the natural morphism  $X_{\Dpik,\cM_{\bullet}}\rightarrow X_{\cM_\Dpik,\cM_{\bullet}}$  \big(resp.,\;the map $X_{\Dpik,\cM_{\bullet}}(E[\epsilon]/\epsilon^2)\rightarrow X_{\cM_\Dpik,\cM_{\bullet}}(E[\epsilon]/\epsilon^2)$ \big) of groupoids gives a morphism
$\omega_{\bm{\delta}_{\bh}}:X_{\Dpik,\cM_{\bullet}}\rightarrow \widehat{(\cZ_{\bL^{\lrr},L})}_{\bm{\delta}_{\bh}}$ of groupoids over $\Art_E$
\big(resp.,\;a map 
$\omega^{\kappa}_{\bm{\delta}_{\bh}}:X_{\Dpik,\cM_{\bullet}}(E[\epsilon]/\epsilon^2)\rightarrow \prod_{ir\in \Delta_n(k)}\homo(L^\times,E)$\big). Note that $\omega_{\bm{\delta}_{\bh}}$ (resp.,\; $\omega^{\kappa}_{\bm{\delta}_{\bh}}$) factors through $|X_{\Dpik,\cM_{\bullet}}|$  \big(resp.,\;$|X_{\Dpik,\cM_{\bullet}}|(E[\epsilon]/\epsilon^2)$ \big).\;

\begin{lem}We have
\begin{itemize}
	\item[(1)]	$|X_{\cM_\Dpik,\cM_{\bullet}}|$ is a subfunctor of $|X_{\cM_\Dpik}|$.\;
	\item[(2)]  $F_{\Dpik,\cF}^0$ (see Section \ref{defdfnomepik}) is a subfunctor of $|X_{\Dpik,\cM_{\bullet}}|$.\;Moreover,\;the following  diagram is commutative:
	\begin{equation}\label{LINVcompatiblemap1}
		\xymatrix{ F_{\Dpik,\cF}^0(E[\epsilon]/\epsilon^2) \ar[dr]^{\kappa,\;(\ref{kappakappaL})}  \ar@{^(->}[r] &  |X_{\Dpik,\cM_{\bullet}}|(E[\epsilon]/\epsilon^2) \ar[d]^{\omega^{\kappa}_{\bm{\delta}_{\bh}}}  \\
			& \prod\limits_{i=1}^k\homo(L^\times,E).}
	\end{equation}
\end{itemize}
\end{lem}
\begin{proof}We first to show $|X_{\cM_\Dpik,\cM_{\bullet}}|$ that is a subfunctor  $|X_{\cM_\Dpik}|$,\;i.e.,\;the $\omepik$-filtration  $\cM_{A,\bullet}$ deforming $\cM_{\bullet}$ on a deformation $\cM_A$ is unique.\;This follows from an argument analogous to   \cite[Proposition 6.2.8]{Ding2021} and \cite[Proposition 2.3.6]{AST2009324R10}.\;The proof proceeds by induction on the length of $\cM_{A,\bullet}$,\;we should show that $\cM_{A,\bullet}$ is an $\Omega$-filtration on $\cM_A$,\;then $\cM_{A,1}$ is uniquely determined as a $(\varphi,\Gamma)$-submodule of $\cM_{A}$,
${\cM}_{A,2}/{\cM}_{A,1}$ is uniquely determined as a $(\varphi,\Gamma)$-submodule of $\cM_{A}/{\cM}_{A,1}$,\;and so on.\;Now suppose that $\widetilde{\cM}_{A,1}$ is another $(\varphi,\Gamma)$-submodule of $\cM_{A}$ deforming $\cM_{1}$.\;Observe that $\widetilde{\cM}_{A,1}$ (resp.,\;$\cM_A/\cM_{A,1}$) is a successive extension of $\cM_{1}$ (resp.,\;$\cM/\cM_{1}$).\;Applying \cite[Lemma 6.2.5 (1)]{Ding2021} to the case $1=i<j$,\;we deduce  $\homo_{(\varphi,\Gamma)}(\widetilde{\cM}_{A,1},\cM_A/\cM_{A,1})=0$ by an easy d\'{e}vissage argument.\;Therefore,\;we see that $\widetilde{\cM}_{A,1}\subset \cM_{A,1}$.\;Then we are done since $\widetilde{\cM}_{A,1}$ and $\cM_{A,1}$ have the same rank.\;This completes the proof of Part (1) by proceeding as we did for $\cM_{A,1}$ step by step.\;Similar to the proof of Part (1),\;we see that $F_{\Dpik,\cF}^0$ is a subfunctor of $|X_{\Dpik}|$.\;The \cite[Lemma 6.2.5 (1)]{Ding2021} has to be replaced by Lemma \ref{phigammacohononcrtialspecial}.\;Since the morphism $F_{\Dpik,\cF}^0\rightarrow |X_{\Dpik}|$ factors through $F_{\Dpik,\cF}^0\rightarrow |X_{\Dpik,\cM_{\bullet}}|$,\;we deduce from Part (1) that $F_{\Dpik,\cF}^0$ is a subfunctor of $|X_{\Dpik,\cM_{\bullet}}|$.\;
\end{proof}

The non-critical property of $\Dpik$ implies the following isomorphism of functors.\;

\begin{pro}\label{isotwofunctor}The morphism $F_{\Dpik,\cF}^0\hookrightarrow |X_{\Dpik,\cM_{\bullet}}|$ of functors induces an isomorphism.\;
\end{pro}
\begin{proof}For any $A\in \Art_E$  and $(A,\Dpik_A,j_A,\cM_{A,\bullet})\in X_{\Dpik,\cM_{\bullet}}(A)$,\;the $\omepik$-filtration on $\cM_\Dpik$ induces an increasing filtration $\cF'_{A}:=(\fil_{i}^{\cF'_{A}}\Dpik_A)_{1\leq i\leq k}:=(\cM_{A,i}\cap \Dpik_A)_{1\leq i\leq k}$ on $\Dpik$ by $(\varphi,\Gamma)$-submodules  over $\cR_{E,L}$.\;Therefore,\;we get that 
\[\gr_{i}^{\cF_{A}}\Dpik_A\bigg[\frac{1}{t}\bigg]=\gr_{i}\cM_A\bigg[\frac{1}{t}\bigg]=\Delta_\pi\otimes_{\cR_{E,L}}\cR_{E,L}(\delta_{A,i})\bigg[\frac{1}{t}\bigg]\]
for any $1\leq i\leq k$.\;Note that $(\cM_{A,i}\cap \Dpik_A)_{1\leq i\leq k}$ may fail to be projective over $A$,\;and thus $\cF_{A}$ only gives an unsaturated filtration of $\Dpik_A$.\;We first show that $\cF'_{E}=\cF_{E}$.\;By an easy variation of \cite[Lemma 2.4.2]{AST2009324R10},\;we see that $\cF':=\cF'_{E}$ is a saturated filtration of $\Dpik$.\;Let $\{\hpi'_{\tau,1},\hpi'_{\tau,2},\cdots,\hpi'_{\tau,ir}\}_{\tau\in \Sigma_L}$ be the Hodge-Tate weights of $\fil_i^{\cF'}\Dpik$.\;By Lemma \ref{phigammacohononcrtialspecial},\;we see that $\homo_{(\varphi,\Gamma)}(\gr_{i}^{\cF'}\Dpik,\Dpik/\fil_{i}^{\cF}\Dpik)=0$ (resp.,\;$\homo_{(\varphi,\Gamma)}(\gr_{i}^{\cF}\Dpik,\Dpik/\fil_{i}^{\cF'}\Dpik)=0$) for $1\leq i\leq k$.\;Therefore,\;we see that $\gr_{i}^{\cF'}\Dpik\subset \fil_{i}^{\cF}\Dpik$ (resp.,\;$\gr_{i}^{\cF}\Dpik\subset \fil_{i}^{\cF'}\Dpik$).\;For $1\leq i\leq k$,\;we further recall that $\homo_{(\varphi,\Gamma)}(\gr_{i}^{\cF'}\Dpik,\fil_{i-1}^{\cF}\Dpik)=0$ (resp.,\;$\homo_{(\varphi,\Gamma)}(\gr_{i}^{\cF}\Dpik,\fil_{i-1}^{\cF'}\Dpik)=0$).\;Therefore.\;We have two injections of $(\varphi,\Gamma)$-modules:
\begin{equation}
	\begin{aligned}
		&\Delta_\pi\otimes_{\cR_{E,L}}\cR_{E,L}(\unr(q_L^{1-k}){z^{\hpi'_{(i-1)r+1}}})\hookrightarrow \fil_i^{\cF'}\Dpik\hookrightarrow \fil_i^{\cF}\Dpik \hookrightarrow \Delta_\pi\otimes_{\cR_{E,L}}\cR_{E,L}(\unr(q_L^{1-k}){z^{\hpi_{ir}}}),\\
		&\Delta_\pi\otimes_{\cR_{E,L}}\cR_{E,L}(\unr(q_L^{1-k}){z^{\hpi_{(i-1)r+1}}})\hookrightarrow \fil_i^{\cF}\Dpik\hookrightarrow \fil_i^{\cF'}\Dpik \hookrightarrow \Delta_\pi\otimes_{\cR_{E,L}}\cR_{E,L}(\unr(q_L^{1-k}){z^{\hpi'_{ir}}}).
	\end{aligned}
\end{equation}
This implies that $\hpi'_{(i-1)r+1}\geq \hpi_{ir}$ and $\hpi_{(i-1)r+1}\geq \hpi'_{ir}$.\;This implies $\hpi'_{\tau,i}=\hpi_{\tau,i}$ for each $1\leq i\leq k$ and $\tau\in \Sigma_{L}$ by the  non-critical assumption.\;By the uniqueness of $\omepik$-filtration with the parameter $(\bx_{\pi},\bm{\delta}_{\bh})\in \sbanpik\times\rigchl$,\;we conclude that $\cF'_{E}=\cF_{E}$.\;Furthermore,\;by the above discussion and  the uniqueness of $\omepik$-filtration with the parameter $(\bx_{\pi},\bm{\delta}_{\bh})\in \sbanpik\times\rigchl$,\;we also see that the filtration $\cF_{A}$ on $\Dpik_A$ coincides with $\cF$ on $\Dpik$ when module $\fm_A$.\;Similar to the proof of \cite[Lemma 2.2.3]{AST2009324R10},\;we get the result.\;	%The last assertion follows from Proposition \ref{dimtanFD0}.\;
\end{proof}

%Let $V$ be the associated representation of $\rho_L$.\;
% (deformations of the representation $V$).\;We can identity $X_{\rho_L}$ with the framed deformations of $V$.\;Therefore the morphisms $X_{\rho_L}\rightarrow X_{V}$ is relatively representable and formally smooth of relative dimension $n^2$.\;We also have an equivalence $X_V\xrightarrow{\sim} X_{\Dpik}$ and $X_{\rho_L} \xrightarrow{\sim} |X_{\rho_L}|$.\;

We now prove the main results of this section,\;which assert that $\xfpx$ (resp.,\;$\bersteineigenvarpik$) is smooth at the point $x_L$ (resp.,\;$x$).\;


We need to study the tangent space $T_{\defvarrho,x_L}$ (resp.,\;$T_{\bersteineigenvarpik,x}$) of  $\defvarrho$ (resp., $\bersteineigenvarpik$) at the non-critical special point $x_L$ (resp.,\;$x$).\;Recall that $ \xfpx\subset \defvarrho$ is a union of  irreducible components of $\defvarrho$ containing $x_L$.\;Consider the morphisms
\begin{equation}\label{variousmorphisms}
\begin{aligned}
	&\omega^{\Box}:=\omega|_{\xfpx}:	\xfpx\subset \defvarrho \longrightarrow \sbanpik\times\rigch,\\
	&\zeta^{\Box}:=\xfpx\subset \defvarrho \longrightarrow {\mathfrak{X}_{\overline{r}}^\Box}.\\
\end{aligned}
\end{equation}
We define $\omega^{\Box,1}:\xfpx\longrightarrow \sbanpik$ (resp.,\;$\omega^{\Box,2}:\xfpx\defvarrho \longrightarrow \rigch$) to be the composition of $\omega^{\Box}$ with projection to the $1$-th (resp.,\;$2$-th) factor.\;

We describe explicitly the tangent map of $\omega^{\Box,1}$ (resp.,\; $\omega^{\Box,2}$) at $x_L$:
\begin{equation}\label{tangentwxl}
\begin{aligned}
	d\omega_{x_L}^{\Box,1}\;(\text{resp}.,\;d\omega_{x_L}^{\Box,2}):	T_{\xfpx,X_L} \longrightarrow T_{\sbanpik,\breve{\mathbf{{x}}}_\pi}  (\text{resp}.,\;T_{\rigch,\underline{\mathbf{1}}}) .\;
\end{aligned}
\end{equation}
Then $d\omega_{x_L}^\Box=d\omega_{x_L}^{\Box,1}\oplus d\omega_{x_L}^{\Box,2}$.\;

We identify $\bv\in T_{\xfpx,x_L}$ with an $E[\epsilon]/\epsilon^2$-valued point $\Spec \big(E[\epsilon]/\epsilon^2\big)\xrightarrow{\bv} \xfpx$ of $\xfpx$. Then the composition
\[\Spec \big(E[\epsilon]/\epsilon^2\big)\xrightarrow{\bv} \xfpx \rightarrow \rigch\]
gives the point $d\omega_{x_L}^{\Box,2}(\bv)$ and thus a continuous character $\delta_{\bv}^0=\boxtimes_{i=1}^k\delta_{\bv,i}^0:\bZ^{\lrr}(\cO_L)\longrightarrow \big(E[\epsilon]/\epsilon^2\big)^\times$ such that $\delta_{\bv}^0\equiv \underline{1}\mod \epsilon$.\;Via the identification $T_{\rigch,\underline{\mathbf{1}}}\cong \homo(\bZ^{\lrr}(\cO_L),E)=\prod_{i=1}^k\homo(\cO_L^{\times},E)$, the element 
$\delta_{\bv}^0$  becomes $((\delta_{\bv,i}^0-1)/\epsilon)_{1\leq i \leq k}$.\;
%Recall that $\{\psi_{\sigma,L}\}_{\sigma\in \Sigma_{L}}$ form a basis of $\homo(\cO_L^{\times},E)$,\;we write $(\delta_{\bv,i}^0(\delta_{i}^{0})^{-1}-1)/\epsilon=\sum_{\tau\in\Sigma_L} a_{\bv,i,\tau}\psi_{i,\tau,L}$ for some $a_{\bv,i,\tau}\in E$.\;
On the other hand,\;the composition 
\[\Spec \big(E[\epsilon]/\epsilon^2\big)\xrightarrow{\bv} \xfpx \rightarrow  \sbanpik\]
corresponds to the point $d\omega_{x_L}^{\Box,1}(\bv)$.\;Recall that 
$\mathfrak{Z}_{\omepik}\cong \otimes_{i=1}^kE[z,z^{-1}]^{\mu_{\Omega_{r}}^{\unr}}$ (depending on the choice of the element $\pi_0^{\otimes k}$ in $\omepik$).\;Therefore,\;for $A\in \Art_E$,\;the $A$-valued points of $\sbanpik$ are isomorphic to the vector space of $A$-valued unramified characters on $\bZ^{\lrr}(L)$.\;In terms of the language of $p$-adic differential equations,\;
for any $\psi=(\psi_i)_{1\leq i\leq k}\in \sbanpik(\spr A)$, there exists a unique unramified character $\unr(a_{\psi_i}):L^\times\rightarrow A^\times$,\;such that $a_{\psi_i}\equiv 1\mod m_A$ and
\begin{equation}\label{defdiffequation}
\psi_i^*\Delta_{\Omega_r}\cong \Delta_{\Omega_r,\widehat{\mathbf{{x}}}_{\pi,i}}\otimes_{\cR_{k(x),L}}\cR_{A,L}(\unr(a_{\psi_i}))=\Delta_{\pi_0}\otimes_{\cR_{k(x),L}}\cR_{A,L}(\unr(a_{\psi_i}\alpha_\pi q_L^{i-k})),
\end{equation}
where $\Delta_{\Omega_r}$ is the universal $p$-adic differential equation over $\big(\mathrm{Spec}\hspace{2pt}\mathfrak{Z}_{\Omega_r}\big)^{\mathrm{rig}}$.\;For  $1\leq i\leq k$,\;we put $\mathrm{\mathbf{q}}_{k,i}:=\unr(\alpha_\pi q_L^{i-k}))\circ\det$
and $\mathrm{\mathbf{q}}_k:=\boxtimes_{i=1}^k\mathrm{\mathbf{q}}_{k,i}$.\;Therefore,\;we have
\begin{equation}\label{tangentwxl2}
\begin{aligned}
	d\omega_{x_L}^\Box:	T_{\xfpx,X_L} \longrightarrow T_{\sbanpik\times \rigch,\omega(x_L)}&=T_{\sbanpik,\breve{\mathbf{{x}}}_\pi}\oplus T_{\rigch,\underline{\mathbf{1}}}\\
	&\cong T_{\rigchl,\mathrm{\mathbf{q}}_k}\xrightarrow{-\cdot (z^{\bh_{ir}})_{1\leq i\leq k}} T_{\rigchl,\bm{\delta}_{\bh}}.\;
\end{aligned}
\end{equation}		
Therefore,\;by composting $	d\omega_{x_L}^\Box$ with $\kappa_L$,\;we can get a map
\begin{equation}\label{tangentLinveta}
\begin{aligned}
	\eta^\Box:\;T_{\xfpx,x_L}\xrightarrow{d\omega_{x_L}^\Box}T_{\rigchl,\bm{\delta}_{\bh}}&\cong\homo(\bZ^{\lrr}(L^{\times}),E)=\prod_{i=1}^k\homo(L^{\times},E) \\
	&\xrightarrow{\kappa_L} \prod_{ir\in \Delta_k}\homo(L^\times,E).
\end{aligned}
\end{equation}

Let $X_{\rho_L}$ be the groupoid over $\Art_E$ of deformations of the group morphism $\rho_L$.\;The map $\zeta^{\Box}$ (see (\ref{variousmorphisms})) induces a  natural morphism $\complocaldefvarrhosub \rightarrow (\widehat{\mathfrak{X}_{\overline{r}}^\Box})_{\rho_L}\cong X_{\rho_L}$.\;By \cite[Lemma 4.13]{breuil2017smoothness},\;there is an exact sequence of $E$-vector spaces
\begin{equation}\label{dimcaluexatsequence}
0\rightarrow K(\rho_L)\rightarrow T_{\mathfrak{X}_{\overline{r}}^\Box,\rho_L}\xrightarrow{f_{\rho_L}} \ext^1_{\gal_L}(\rho_L,\rho_L)\cong \ext^1_{(\varphi,\Gamma)}(\Dpik,\Dpik)=F_{\Dpik}(E[\epsilon]/\epsilon^2)\rightarrow 0,\;
\end{equation}
where $K(\rho_L)$ is a $k(x_L)$-vector space of $T_{\mathfrak{X}_{\overline{r}}^\Box,\rho_L}$ of dimension $n^2-\dim_E\EndO_{\gal_L}(\rho_L)$.\;

Let $\sL(\rho_L):=\sL(\Dpik)$ be the parabolic Fontaine-Mazur simple $\sL$-invariants associated to $\rho_L$ (equivalently,\;to $\Dpik$).\;

\begin{pro}\label{smoothnessdefvar}$\xfpx$ is smooth at the point $x_L$,\;and (\ref{tangentLinveta}) factors thought a surjective map
\begin{equation}\label{defortangentmap1eta}
	\eta^\Box:T_{\xfpx,x_L}\twoheadlongrightarrow \sL(\rho_L).\;
\end{equation}
\end{pro}
\begin{proof}We first show that the image of $f_{\rho_L}$ in (\ref{dimcaluexatsequence}) is contained in the $E$-vector space\\ $|X_{\Dpik,\cM_{\bullet}}|(E[\epsilon]/\epsilon^2)\cong F_{\Dpik,\cF}^0(E[\epsilon]/\epsilon^2)$.\;By Theorem \ref{localfiltration} (we use the notation of Section \ref{ACC},\;Theorem \ref{localfiltration} and its proof),\;we see that there  exists an open affinoid neighborhood $X\subset \xfpx$ of  $x_L$ such that the $(\varphi,\Gamma)$-module $D_{\rig}(r^{\univer}_X)[\frac{1}{t}]$ over $\cR_{X,L}$ admits a filtration $\cM_\bullet$ such that $\gr^i\cM_\bullet\cong \Delta_{X,i}[\frac{1}{t}]$.\;Let $A\in \Art_E$,\;and $\Spec(A)\rightarrow X$ a morphism of rigid analytic spaces sending the only point of $\Spec(A)$ to $x$.\;By pulling along $\psi:\Spec(A)\rightarrow X\rightarrow \mathfrak{X}_{\overline{r}}^\Box$,\;we obtain a deformation $\rho_A$ of $\overline{r}$ such that \[D_{\rig}(\rho_A)[\frac{1}{t}]\cong A\otimes_{\Gamma(X,\cO_X)}D_{\rig}(r^{\univer}_X)[\frac{1}{t}],\;\]and $A\otimes_{\Gamma(X,\cO_X)}\cM_\bullet $ gives a filtration $\cM_{A,\bullet}$ on $D_{\rig}(\rho_A)[\frac{1}{t}]$.\;Let $\Delta_{A,i}$ be the  pull-back of $\Delta_{X,i}$ along $\Spec(A)\rightarrow X$.\;Therefore,\;the filtration $\cM_{A,\bullet}$ on $D_{\rig}(\rho_A)[\frac{1}{t}]$ satisfies
\begin{equation}\label{Adefor}
	\begin{aligned}
		&\gr_i^{\cM_{A,\bullet}}D_{\rig}(\rho_{A})\Big[\frac{1}{t}\Big]\xrightarrow{\sim}\Delta_{A,i}\Big[\frac{1}{t}\Big]\\
		\xrightarrow{\sim}&\Delta_{\pi_0}\otimes_{\cR_{E,L}}\cR_{A,L}\big(\unr(a_{A,i}\alpha_\pi q_L^{i-k})\cdot({\delta}_{A,i})_{\varpi_L}\cdot z^{{\bh}_{(i-1)r+1}}\big)\Big[\frac{1}{t}\Big]
	\end{aligned}
\end{equation}
for $1\leq i\leq k$,\;where $\delta_A$ is the pullback of the universal character of  $\cO_L^\times$ along  $\Spec(A)\rightarrow X$,\;and $a_{A,i}$ is an element of $A$ such that $a_{A,i}\equiv 1\mod m_A$.\;Since the
%=\Delta^{\univer}_{\pi}\wte_{\cR_{X}}\cR_{X}((\delta_i^{0,\univer})_{\varpi_L}\cdot z^{\underline{h}_{s_{i-1}+1}}|_{\co_L^\times})$
point $x_L$ is non-critical,\;the uniqueness of $\omepik$-filtration on $\cM_\Dpik$ implies that $\cM_{A,\bullet}\otimes_AE\cong \cM_{\bullet}$,\;i.e.,\;the filtration $\cM_{\bullet}$ on $D_{\rig}(\rho_L)[\frac{1}{t}]$ satisfies
\[\gr_i^{\cM_\bullet}D_{\rig}(\rho_L)\Big[\frac{1}{t}\Big]\xrightarrow{\sim}\Delta_{\pi_0}\otimes_{\cR_{E,L}}\cR_{E,L}(\unr(\alpha_\pi )\bm{\delta}_{\bh,i})\Big[\frac{1}{t}\Big],\]
for $1\leq i\leq k$.\;Evaluating at $A=E[\epsilon]/\epsilon^2$-points and using Proposition \ref{isotwofunctor},\;we see that the image of $f_{\rho_L}$ in (\ref{dimcaluexatsequence}) is contained in the $E$-vector space $|X_{\Dpik,\cM_{\bullet}}|(E[\epsilon]/\epsilon^2)\cong F_{\Dpik,\cF}^0(E[\epsilon]/\epsilon^2)$.\;

The closed embedding $\xfpx\hookrightarrow {\defvarrho}\hookrightarrow \mathfrak{X}_{\overline{r}}^\Box\times \sbanpik\times\rigch$ induces an injection on tangent spaces $T_{\xfpx,x_L}\hookrightarrow T_{\mathfrak{X}_{\overline{r}}^\Box,\rho_L}\oplus T_{\sbanpik\times\rigch,(\breve{\mathbf{{x}}}_\pi,\underline{1})}$.\;Then $d\zeta^{\Box}_{x_L}$ is the composition of this injection with the projection
$$T_{\mathfrak{X}_{\overline{r}}^\Box,\rho_L}\oplus T_{\sbanpik\times\rigch,(\breve{\mathbf{{x}}}_\pi,\underline{1})}\rightarrow T_{\mathfrak{X}_{\overline{r}}^\Box,\rho_L}.\;$$We claim that the tangent map $d\zeta^{\Box}_{x_L}$ remains injective.\;Let $\bv\in T_{\xfpx,x_L}$ which maps to $0$ in $T_{\mathfrak{X}_{\overline{r}}^\Box,\rho_L}$.\;Then it maps to $0$ in $F_{\Dpik}(E[\epsilon]/\epsilon^2)$ via $f_{\rho_L}$.\;We have to show that the image of $\bv\in T_{\xfpx,x_L}$ in $T_{\sbanpik\times\rigch,(\breve{\mathbf{{x}}}_\pi,\underline{1})}$ is also zero.\;Since $F_{\Dpik,\cF}^0$ is a subfunctor of $F_{\Dpik,\cF}$,\;we see that the image $f_{\rho_L}(\bv)\in F_{\Dpik,\cF}^0(E[\epsilon]/\epsilon^2)$ is isomorphic to the trivial deformation of $\Dpik$ of type $\omepik$ over $E[\epsilon]/\epsilon^2$ (and then the parameter of trivial deformation is trivial).\;We conclude that the image of $\bv\in T_{\xfpx,x_L}$ in $T_{\sbanpik\times\rigch,(\breve{\mathbf{{x}}}_\pi,\underline{1})}$ is also zero.\;Therefore,\;we obtain from (\ref{dimcaluexatsequence}) a short exact sequence
\begin{equation}\label{dimcaluexatsequenceproof}
	\begin{aligned}
		0\rightarrow K(\rho_L)\cap T_{\xfpx,X_L}\rightarrow T_{\xfpx,X_L}\xrightarrow{f_{\rho_L}} F_{\Dpik,\cF}^0(E[\epsilon]/\epsilon^2).\;
	\end{aligned}
\end{equation}
Then Proposition \ref{dimtanFD0} gives an upper bound:
\[\dim_ET_{\xfpx,x_L}\leq n^2+d_L\Big(\frac{n(n-1)}{2}+k\Big)=\dim_E{\defvarrho}.\]
This implies that $\xfpx$ is smooth at the point $x_L$.\;At last,\;by comparing the construction,\;we see that the composition
\begin{equation}
	\begin{aligned}
		T_{\xfpx,X_L}\xrightarrow{f_{\rho_L}} F_{\Dpik,\cF}^0(E[\epsilon]/\epsilon^2)\xrightarrow{\omega_{\bm{\delta}_{\bh}}(E[\epsilon]/\epsilon^2),(\ref{omegadeltapinon})}\widehat{(\cZ_{\bL^{\lrr},L})}_{\bm{\delta}_{\bh}}(E[\epsilon]/\epsilon^2)=T_{\rigchl,\bm{\delta}_{\bh}}
	\end{aligned}
\end{equation}
coincides with $d\omega_{x_L}^\Box$ (see (\ref{tangentwxl2})).\;We deduce from Corollary \ref{corlinvaraint} the surjection (\ref{defortangentmap1eta}).\;	
\end{proof}


We now prove the main proposition of this section.\;The natural embedding $$\bersteineigenvarpik \hooklongrightarrow (\Spf\;R^{\fp}_{\infty})^{\rig}\times\mathfrak{X}_{\overline{r}}^\Box \times \sbanpik\times \rigch$$ induces the weight map 
\begin{equation}
\omega:	\bersteineigenvarpik \longrightarrow \sbanpik\times \rigch.\;
\end{equation}
Consider the tangent map of $\omega$ at point $x=(x^{\fp}, \rho_L, \breve{\mathbf{{x}}}_\pi,\underline{\mathbf{1}})$:
\begin{equation}\label{tangentpatchvar1eta}
d\omega_x:	T_{\bersteineigenvarpik,x} \longrightarrow T_{\sbanpik\times \rigch,\omega(x)}\cong  T_{\rigchl,\bm{\delta}_{\bh}},
\end{equation}
where the second isomorphism  follows by the same argument of (\ref{tangentwxl2}).\;
\begin{pro}\label{coreigentangentmap2}$\bersteineigenvarpik$ is smooth at $x$,\;and (\ref{tangentpatchvar1eta}) factors through a surjective map
\begin{equation}\label{eigentangentmap2}
	T_{\bersteineigenvarpik,x}\twoheadlongrightarrow \sL(\rho_L).\;
\end{equation}
\end{pro}
\begin{proof}Recall that $\cE(x)$ is the union of  irreducible components of $\bersteineigenvarpik$ containing $x$.\;By (\ref{OZlocal}) and (\ref{OZfactor}),\;we see that  $\kappa_L\circ d\omega_x$ is equal to the following composition
\begin{equation}\label{tangentpatchvar1}
	\begin{aligned}
		T_{\bersteineigenvarpik,x}\cong T_{\cE(x),x}
		\cong &\;T_{X^{\fp} \times \BU^g \times \iota_{\omepik}^{-1}(\xfpx),(x^{\fp},x_L)}\\
		&\twoheadlongrightarrow T_{\xfpx,x_L}\xrightarrow{\eta^\Box,\;(\ref{tangentLinveta})}\prod_{ir\in \Delta_n(k)}\homo(L^\times,E).\;
	\end{aligned}
\end{equation}
Now the result follows.\;
\end{proof}

\begin{rmk}\label{localmodeldescri}Let $X_{\rho_L,\cM_{\bullet}}$ be the groupoid over $\Art_E$ defined in \cite[Section 6.4]{Ding2021}.\;For $w\in \sW_{n,\Sigma_L}$, let $X^{w}_{\rho_L,\cM_{\bullet}}$ be the closed subgroupoid of $X_{\rho_L,\cM_{\bullet}}$ over $\Art_E$ defined in \cite[Section 6.4]{Ding2021},\;which is closely related to some varieties studied in geometric representation theory.\;Since the parameters of our $\omepik$-filtration $\cF$ are non-generic (in the sense of \cite[(6.5)]{Ding2021}),\;the morphism \[X_{\cM_{\Dpik},\cM_{\bullet}}\rightarrow \widehat{(\cZ_{\bL^{\lrr},L})}_{\bm{\delta}_{\bh}}\times_{\widehat\fz^\lrr_{L}} X_{W_{\dr}(\cM_{\Dpik}),\cF},\;\cF:=W_{\dr}(\cM_{\bullet})\] of groupoids in \cite[Theorem\;6.2.6]{Ding2021} is no longer formally smooth.\;Thus the discussions in  \cite[Section 6.4]{Ding2021}  cannot be applied to our case (for example,\;the final result \cite[Corollary 6.4.7]{Ding2021}).\;In our non-critical special case,\;it seems likely that $\complocaldefvarrho$ is also isomorphic to $X^{w_0}_{\rho_L,\cM_{\bullet}}$.\;
\end{rmk}








\subsection{Local-global compatibility and the main theorem}\label{lgcompthemainSECTION}

We are ready to establish the second main theorem.\;Suppose that $\rho_L$ appears in the patched Bernstein eigenvariety $\bersteineigenvarpik$,\;i.e.\;there exist  $x^{\fp}\in (\Spf\;R_{\infty}^{\fp})^{\rig}$,\;and  $$(\pi_{x,\bL^{\lrr}},\chi)\in \sbanpik\times\rigch $$ such that $x:=(x^{\fp}, \rho_L, \pi_{x,\bL^{\lrr}},\chi)\in \bersteineigenvarpik$.\;Let $\fm_y$ be the maximal ideal of $R_{\infty}[1/p]$ corresponding to the point $y:=(x^{\fp}, \rho_L)$ of $(\Spf\;R_{\infty})^{\rig}$ (e.g. if $\rho_L$ is attached to an automorphic representation of $\widetilde{G}$ with non-zero $U^{\fp}$-fixed vectors).\;By the argument in Section \ref{BENVARPARAVAR} (especially,\;the continuous $R_{\infty}$-admissible unitary representation $\Pi_{\infty}$ of $G$ over $E$) or  
\cite{PATCHING2016},\;we get that 
\begin{equation*}
\widehat{\Pi}(\rho_L):=\Pi_{\infty}[\fm_y]
\end{equation*}
is an admissible unitary Banach representation of $\GLN_n(L)$, which one might expect to be the right representation (up to multiplicities) corresponding to $\rho_L$ in the $p$-adic local Langlands program.\;Suppose the following holds:
\begin{itemize}
\item[(a)] The $(\varphi,\Gamma)$-module $\Dpik=D_{\rig}(\rho_L)$ over $\cR_{E,L}$ admits a non-critical special $\omepik$-filtration  with parameter $(\widetilde{\bx}_{\pi,\bh},\widetilde{\bm{\delta}}_\bh)\in \sbanpik\times\rigch $ (see Definition \ref{dfnnoncriticalspecial});
\item[(b)] $\rho_L$ is potentially semi-stable;
\item[(c)] The monodromy operator $N$ on $D_{\pst}(\rho_L)$ satisfies $N^{k-1}\neq 0$.
\end{itemize}
By the argument before \cite[Corollary 3.1.11]{Ding2021},\;$ \sbanpik\times\rigch $ has an action of $\mu_{\omepik}$.\;For $\psi\in \mu_{\omepik}$,\;we have
\[{\psi}\cdot((x_i)_{1\leq i\leq k},(\chi_i)_{1\leq i\leq k})=((x_i\te \unr(\psi_i(\varpi_L)))_{1\leq i\leq k},(\chi_i\psi_i|_{\cO_L^\times})_{1\leq i\leq k}),\;{\psi}=(\psi_i)_{1\leq i\leq k}\in \mu_{\omepik}.\]
%The non-critical assumption predicts that the non-existence for the non-trivial companion points of $x$ on the  patched Bernstein eigenvariety $\bersteineigenvarpik$.\;

By Proposition \ref{basicpropertyeigen} (1),\;$(x^{\fp}, \rho_L, \pi_{x,\bL^{\lrr}},\chi)$ appears in patched Berstein eigenvariety $\bersteineigenvarpik$ if and only if $\homo_{\bL^{\lrr}(L)}\big(\pi_{x,\bL^{\lrr}}\te\chi_{\varpi_L}\te \lamoverx, J_{\bP^{\lrr}(L)}\big(\Pi_{\infty}^{R_{\infty}-\ana}[\fm_y]\big)\big)\neq0$.\;By Proposition \ref{propertyparavar} (4) (where we use the theory of $\omepik$-filtration in families (see \cite[Appendix A.1]{Ding2021}),\;which may be viewed as a parabolic analogue of the global triangulation theory),\;\cite[Corollary 3.1.11]{Ding2021},\;and the uniqueness of $\omepik$-filtration with the parameter $(\widetilde{\bx}_{\pi,\bh},\widetilde{\bm{\delta}}_\bh)\in \sbanpik\times\rigch $ and proof of Proposition \ref{isotwofunctor},\;we deduce
\begin{lem}\label{detlameda}
For $(\pi_{x,\bL^{\lrr}},\chi)\in \sbanpik\times\rigch $,\;the eigenspace
\[\homo_{\bL^{\lrr}(L)}\Big(\pi_{x,\bL^{\lrr}}\te\chi_{\varpi_L}\te \lamoverx, J_{\bP^{\lrr}(L)}\big(\Pi_{\infty}^{R_{\infty}-\ana}[\fm_y]\big)\Big)\neq0\]
i.e.,\;$(x^{\fp}, \rho_L, \pi_{x,\bL^{\lrr}},\chi)$ appears in patched Berstein eigenvariety $\bersteineigenvarpik$ if and only if
$(\pi_{x,\bL^{\lrr}},\chi)$ belongs to the same $\mu_{\omepik}$-orbit of the point $(\breve{\mathbf{{x}}}_\pi,\underline{\mathbf{1}})\in \sbanpik\times\rigch $.\;For $\psi\in \mu_{\omepik}$,\;we put $x_{\psi}:=(x^{\fp}, \rho_L,{\psi}\cdot(\breve{\mathbf{{x}}}_\pi,\underline{\mathbf{1}}))$.\;Then the point $x$ appearing in Section \ref{appearpoint} is just $x=x_{\underline{\mathbf{1}}}$.\;
\end{lem}
For $\beta\in E^\times$,\;and $\underline{\lambda}\in X^+_{\Delta_n^k}$,\;we put $\pi_0^{\lrr}(\beta,\underline{\lambda}):=\unr(\beta)\circ\mathrm{det}_{\bL^{\lrr}(L)}\te \pi_0^{\lrr}\te L^{\lrr}(\underline{\lambda})$ and $\pi_0^{\lrr}(\beta):=\pi_0^{\lrr}(\beta,\underline{0})$ (see (\ref{pilrr}) for $\pi_0^{\lrr}$).\;Then for all $\psi\in \mu_{\omepik}$,
\begin{equation}
\begin{aligned}
	&\pi_{x_{\psi},\bL^{\lrr}}\te (\psi\cdot \underline{\mathbf{1}})_{\varpi_L}\te \lamoverx\\
	\cong &\;\unr(\alpha_\pi  q_L^{\frac{n-1}{2}})\circ\mathrm{det}_{\bL^{\lrr}(L)}\te \delta_{\bP^{\lrr}}^{1/2}\te\Delta_{[k-1,0]}(\pi_0)\te \lamoverx\\
	=&\;\pi_0^{\lrr}(\alpha,{\bm\lambda}_\bh)\te \delta_{\bP^{\lrr}},\;\alpha:=\alpha_\pi  q_L^{\frac{n-1}{2}},
\end{aligned}
\end{equation}
which is independent on the choice of $\psi$.\;Therefore,\;Lemma \ref{detlameda} implies 
\begin{equation}\label{jPinj}
\homo_{\bL^{\lrr}(L)}\big(\pi_0^{\lrr}(\alpha,{\bm\lambda}_\bh)\te \delta_{\bP^{\lrr}}, J_{\bP^{\lrr}(L)}\big(\Pi_{\infty}^{R_{\infty}-\ana}[\fm_y]\big)\big)\neq0.
\end{equation}
We define the natural map
\begin{equation}\label{adjuncitonformula}
\begin{aligned}
	\homo_{G}\big(&\BI_{\op^{\lrr}}^G(\alpha,\pi,{\bm\lambda}_\bh),\Pi_{\infty}^{R_{\infty}-\ana}[\fm_y]\big)\\
	&\longrightarrow\homo_{\bL^{\lrr}(L)}\big(J_{\bP^{\lrr}(L)}\big(\BI_{\op^{\lrr}}^G(\alpha,\pi_0,{\bm\lambda}_\bh)\big), J_{\bP^{\lrr}(L)}\big(\Pi_{\infty}^{R_{\infty}-\ana}[\fm_y]\big)\big)\\
	&\longrightarrow \homo_{\bL^{\lrr}(L)}\big(\pi_0^{\lrr}(\alpha,{\bm\lambda}_\bh)\te \delta_{\bP^{\lrr}}, J_{\bP^{\lrr}(L)}\big(\Pi_{\infty}^{R_{\infty}-\ana}[\fm_y]\big)\big),
\end{aligned}
\end{equation}
where the first map is induced by applying the Jacquet-Emerton functor,\;the second map is induced by the injection 
\[\pi_0^{\lrr}(\alpha,{\bm\lambda}_\bh)\te \delta_{\bP^{\lrr}}\hookrightarrow J_{\bP^{\lrr}(L)}\Big(i_{\op^{\lrr}}^G(\alpha,\pi_0,{\bm\lambda}_\bh\big)\Big)\hookrightarrow J_{\bP^{\lrr}(L)}\Big(\BI_{\op^{\lrr}}^G(\alpha,\pi_0,{\bm\lambda}_\bh\big)\Big)\]
(by the same argument as in the proof of \cite[Lemma 3.4]{2022ext1hyq}),\;we can prove that $J_{\bP^{\lrr}(L)}(i_{\op^{\lrr}}^G(\alpha,\pi_0,{\bm\lambda}_\bh))$ is semi-simple and $J_{\bP^{\lrr}(L)}\big(\st^{\infty}_{(r,k)}(\alpha,\pi_0,{\bm\lambda}_\bh)\big)=\pi_0^{\lrr}(\alpha,{\bm\lambda}_\bh)\te \delta_{\bP^{\lrr}}$ for the unique irreducible quotient $\st^{\infty}_{(r,k)}(\alpha,\pi_0,{\bm\lambda}_\bh)$ of $i_{\op^{\lrr}}^G(\alpha,\pi_0,{\bm\lambda}_\bh)$).

We first establish an adjunction formula,\;following the line of \cite[Proposition 4.7]{2019DINGSimple} and \cite{bergdall2018adjunction}.\;

\begin{pro}\label{proadjuncitonformularight}
The natural map constructed in (\ref{adjuncitonformula})
\begin{equation}\label{adjuncitonformularight}
	\begin{aligned}
		\homo_{G}\Big(&\BI_{\op^{\lrr}}^G(\alpha,\pi_0,{\bm\lambda}_\bh),\; \Pi_{\infty}^{R_{\infty}-\ana}[\fm_y]\Big)\\
		&\lra \homo_{\bL^{\lrr}(L)}\Big(\pi_0^{\lrr}(\alpha,{\bm\lambda}_\bh)\te \delta_{\bP^{\lrr}}, J_{\bP^{\lrr}(L)}\big(\Pi_{\infty}^{R_{\infty}-\ana}[\fm_y]\big)\Big)
	\end{aligned}
\end{equation}
is bijective.\;Moreover.\;this bijection stays true if \;$\BI_{\op^{\lrr}}^G(\alpha,\pi_0,{\bm\lambda}_\bh)$ is replaced by any subrepresentation $W$ such that $\st^{\infty}_{(r,k)}(\alpha,\pi_0,{\bm\lambda}_\bh)\subseteq W \subseteq \st^{\ana}_{(r,k)}(\alpha,\pi_0,{\bm\lambda}_\bh)$.\;In particular,\;we have an injection
\begin{equation}\label{steinberginj}
	\st_{(r,k)}^{\infty}(\alpha,\pi_0, {\bm\lambda}_\bh)\hookrightarrow\Pi_{\infty}^{R_{\infty}-\ana}[\fm_y]\subseteq\widehat{\Pi}(\rho_L).
\end{equation}
\end{pro}
\begin{proof}The proof is divided into the following steps.\;\\
\textbf{(a)} Let $f\in \homo_{G}\Big(\BI_{\op^{\lrr}}^G(\alpha,\pi_0,{\bm\lambda}_\bh), \Pi_{\infty}^{R_{\infty}-\ana}[\fm_y]\Big)$ be any non-zero map.\;By Lemma \ref{detlameda},\;\cite[Theorem 4.3]{breuil2015II},\;\cite[Corollary 3.4]{breuil2016socle},\;we see that $f$  factors through $\st^{\ana}_{(r,k)}(\alpha,\pi_0,{\bm\lambda}_\bh)$ and induces a non-zero map $\st^{\infty}_{(r,k)}(\alpha,\pi_0,{\bm\lambda}_\bh)\hookrightarrow \Pi_{\infty}^{R_{\infty}-\ana}[\fm_y]$.\;Indeed,\;suppose that $W'$ is an irreducible constituent of $\BI_{\op^{\lrr}_I}^G(\alpha,\pi_0,{\bm\lambda}_\bh)$ for some $\emptyset\neq I\subset \Delta_n(k)$.\;If $W'$ is locally algebraic (i.e.,\;$s=1$),\;then $W'$ is an irreducible constituent of $i_{\op^{\lrr}_{I}}^G(\alpha,\pi_0,\bm{\lambda}_\bh)$.\;But $$\homo_{\bL^{\lrr}(L)}\Big( J_{\bP^{\lrr}(L)}\big(i_{\op^{\lrr}_{{I}}}^G(\alpha,\pi_0,{\bm\lambda}_\bh)\big), J_{\bP^{\lrr}(L)}\big(\Pi_{\infty}^{R_{\infty}-\ana}[\fm_y]\big)\Big)=0$$
by Lemma \ref{detlameda}.\;If $W'$ is not locally algebraic,\;we deduce from  \cite[Theorem]{orlik2014jordan} that $W'$ has the form $\mathcal{F}_{\op_J^{\lrr}(L)}^G(L^{\lrr}(-s\cdot{\bm\lambda}_\bh),W'')$
with $s\in \sW_{n,\Sigma_{L}}$ satisfying $I\subseteq J$,\;$s\cdot{\bm\lambda}_\bh\in X^+_{\Delta_n^k\cup J}$,\;where $W''$ is an irreducible  constituent of $i_{\op_I^{\lrr}(L)\cap\bL^{\lrr}_J(L) }^{\bL^{\lrr}_J(L)}\pi_I$.\;We deduce from \cite[Th\'{e}or\`{e}me 4.3,\;Remarque 4.4 (i),\;Corollaire 4.5]{breuil2015II} that
\begin{equation}
	\begin{aligned}
		0\neq \homo_{G}\Big(&\mathcal{F}_{\op^{\lrr}(L)}^G((\overline{M}^{\lrr}(-s\cdot{\bm\lambda}_\bh))^\vee,\pi_0^{\lrr}),\;\Pi_{\infty}^{R_{\infty}-\ana}[\fm_y]\Big)\\
		&\cong \homo_{\bL^{\lrr}(L)}\Big(\pi_0^{\lrr}\te L^{\lrr}(s\cdot{\bm\lambda}_\bh), J_{\op^{\lrr}(L)}\big(\Pi_{\infty}^{R_{\infty}-\ana}[\fm_y]\big)\Big),
	\end{aligned}
\end{equation}
which leads a contradiction to Lemma \ref{detlameda}.\;This shows that $f$ factors through the locally $\bQ_p$-analytic representation $\st^{\ana}_{(r,k)}(\alpha,\pi_0,{\bm\lambda}_\bh)$. By the left exactness of $J_{\bP^{\lrr}(L)}$,\;we get that $f$ induces an non-zero map $\st^{\infty}_{(r,k)}(\alpha,\pi_0,{\bm\lambda}_\bh)\hookrightarrow \Pi_{\infty}^{R_{\infty}-\ana}[\fm_y]$ (so (\ref{steinberginj}) holds).\;This implies that
(\ref{adjuncitonformularight}) is injective.\\
%Therefore,\;$W'$ is an irreducible  constituent of $\mathcal{F}_{\op^{\lrr}(L)}^G(\overline{M}^{\lrr}(-s\cdot{\bm\lambda}_\bh),\pi_0^{\lrr})$.\;Recall that $\mathcal{F}_{\op^{\lrr}(L)}^G(\overline{M}^{\lrr}(-s\cdot{\bm\lambda}_\bh),\pi_0^{\lrr})$ and $\mathcal{F}_{\op^{\lrr}(L)}^G((\overline{M}^{\lrr}(-s\cdot{\bm\lambda}_\bh))^\vee,\pi_0^{\lrr})$ has the same  irreducible constituents.\;and $s\cdot{\bm\lambda}_\bh\upuparrows {{\bm\lambda}_\bh}$
\textbf{(b)}\;The proof of \cite[Theorem 4.8]{bergdall2018adjunction} and \cite[Proposition 4.7]{2019DINGSimple}
apply in our case,\;although our input is slightly different from that in \cite[Theorem 4.8]{bergdall2018adjunction} or \cite[Proposition 4.7]{2019DINGSimple}.\;We indicate below the changes.\;Recall that $M_{\infty}$ is finite projective over $S_{\infty}[[\GLN_n(\cO_L)]]$,\;we deduce from \cite[Corollary 3.9]{breuil2017interpretation} that $\Pi_{\infty}^{R_{\infty}-\ana}$ is a direct summand of $\cC^{\BQ_p-\ana}\big(\BZ_p^s \times \GLN_n(\cO_L),E\big)$ as $\GLN_n(\cO_L)$-representations where $s=n^2(|S_p|+1)+q$.\;Let $\fm \subset S_{\infty}$ be the preimage of $\fm_y$ via the morphism $S_{\infty} \ra R_{\infty}$.\;\\
\textbf{(c)}\;We put $V:=\Pi_{\infty}^{R_{\infty}-\ana}[\fm]$.\;
It is an admissible Banach representation of $\GLN_n(L)$ equipped with a continuous action of $R_{\infty}$.\;There exists a 
pro-$p$ uniform compact open subgroup $H$ of $\GLN_n(\cO_L)$ such that  $V|_{H}\cong \cC^{\BQ_p-\ana}(H,E)^{\oplus r}$ for certain $r\in \BZ_{\geq 1}$.\;By Lemma \ref{fpacycilic} below,\;we see that $V$ satisfies the first hypothesis in  \cite[Theorem 4.8]{bergdall2018adjunction}.\;On the other hand,\;by Lemma \ref{detlameda},\;$(U,\pi)=\Big(\lamoverx,\pi_0^{\lrr}(\alpha)\te \delta_{\bP^{\lrr}}\Big)$ is non-critical with respect to $\Pi_{\infty}^{R_{\infty}-\ana}[\fm_y]$ (in the terminology of \cite[Definition 4.4]{bergdall2018adjunction}).\;

We first recall that $\cC^{\lp}_c(\bN^{\lrr}(L),\pi_0^{\lrr}(\alpha,{\bm\lambda}_\bh))$ (resp.,\;$\cC^{\infty}_c(\bN^{\lrr}(L),\pi_0^{\lrr}(\alpha))$) is the locally polynomial (resp.,\;smooth) $\pi_0^{\lrr}(\alpha,{\bm\lambda}_\bh)$-valued (resp.,\;$\pi_0^{\lrr}(\alpha)$-valued) functions on $\bN^{\lrr}(L)$ with compact support,\;it has the topology defined in \cite[(2.5)]{emerton2007jacquet} (resp.\;see \cite[(2.2)]{emerton2007jacquet}).\;Both spaces $\cC^{\lp}_c(\bN^{\lrr}(L),\pi_0^{\lrr}(\alpha,{\bm\lambda}_\bh))$ and $\cC^{\infty}_c(\bN^{\lrr}(L),\pi_0^{\lrr}(\alpha))$ are convex $E$-vector space of compact type.\;We write \[A^{\lp}:=\cC^{\lp}_c(\bN^{\lrr}(L),\pi_0^{\lrr}(\alpha,{\bm\lambda}_\bh)),\; A^{\infty}:=\cC^{\infty}_c(\bN^{\lrr}(L),\pi_0^{\lrr}(\alpha))\] for simplicity (only suitable for the following commutative diagram).\;We use a similar commutative diagram to \cite[proof of Theorem 4.8]{bergdall2018adjunction} with $V$  being our $V[\fm_y]$,\;$(U,\pi)$ being our $\Big(\lamoverx,\pi_0^{\lrr}(\alpha)\te \delta_{\bP^{\lrr}}\Big)$,\;
\[\xymatrix{
	\homo_{G}\big(\BI_{\op^{\lrr}}^G(\alpha,\pi_0,{\bm\lambda}_\bh),V[\fm_y]\big) \ar[d]^{\simeq}_{(a)} \ar[r]^{(1)\hspace{40pt}} & \homo_{\bL^{\lrr}(L)}\big(\pi_0^{\lrr}(\alpha,{\bm\lambda}_\bh)\te \delta_{\bP^{\lrr}}, J_{\bP^{\lrr}(L)}\big(V[\fm_y]\big)\big) \ar[d]^{\simeq}_{(b)}  \\
	\homo^\#\big(A^{\lp},V[\fm_y]\big)  \ar[r]^{(2)\hspace{30pt}} \ar[d]^{\simeq}_{(c)} & \homo^\#\big(U(\mathfrak{g}_{\Sigma_L})\otimes_ {U(\mathfrak{p}^{\lrr}_{\Sigma_L})}A^{\lp},V[\fm_y]\big)\ar[d]^{\simeq}_{(d)}\\
	\homo^\#\big(M^{\lrr}({\bm\lambda}_\bh)^\vee\otimes_EA^{\infty},V[\fm_y]\big)  \ar[r]^{(3)} \ar[d]^{\eta_1} & \homo^\#\big(M^{\lrr}({\bm\lambda}_\bh)\otimes_EA^{\infty},V[\fm_y]\big)\\
	\homo^\#\big(L^{\lrr}({\bm\lambda}_\bh)\otimes_EA^{\infty},V[\fm_y]\big),  \ar[ur]^{\eta_2} &
}\]
where we write  $\homo^\#:=\homo_{(\mathfrak{g}_{\Sigma_L},\bP^{\lrr}(L))}$ for simplicity.\;The map $(1)$ is equal to (\ref{adjuncitonformularight}).\;We refer to 
\cite[Theorem 4.8]{bergdall2018adjunction} or \cite[Proposition 4.7]{2019DINGSimple} for the explanation for the terms and maps $(1)$,\;$(2)$,\;$(a)$,\;$(b)$,\;$(c)$ and $(d)$.\;
The maps $\eta_1$,\;$\eta_2$ and $(3)$ are induced by natural morphisms $M^{\lrr}({\bm\lambda}_\bh)\twoheadrightarrow \lamoverx\hookrightarrow M^{\lrr}({\bm\lambda}_\bh)^\vee$. By Lemma \ref{detlameda} and the same arguments as in \cite[Proposition 4.9]{bergdall2018adjunction},\;we can show that $\eta_2$ is bijective and $\eta_1$ is injective.\;

To prove that $\eta_1$ is surjective,\;we need a generalization of Step (c) of the proof of \cite[Proposition 4.7]{2019DINGSimple}.\;It suffices to prove that for any pair $(M,M')$ such that $M'/M=L(s\cdot{\bm\lambda}_\bh)$ in $\cO_{\alge}^{\overline{\fp}^{\lrr},\Sigma_L}$ with $1\neq s\cdot{\bm\lambda}_\bh\in X^+_{\Delta_n^k}$,\;the restriction map
\begin{equation*}
	\begin{aligned}
		\homo_{(\mathfrak{g}_{\Sigma_L},\bP^{\lrr}(L))}&\big(M'\otimes_E\cC^{\infty}_c(\bN^{\lrr}(L),\pi_0^{\lrr}(\alpha)),V[\fm_y]\big)\\
		&\longrightarrow \homo_{(\mathfrak{g}_{\Sigma_L},\bP^{\lrr}(L))}\big(M\otimes_E\cC^{\infty}_c(\bN^{\lrr}(L),\pi_0^{\lrr}(\alpha)),V[\fm_y]\big)
	\end{aligned}
\end{equation*}
is surjective.\;Now given a $(\mathfrak{g}_{\Sigma_L},\bP^{\lrr}(L))$-equivariant morphism $f$ in the hright side,\;we can obtain
\begin{equation}\label{proofadjformla}
	M\otimes_E\cC^{\infty}_c(\bN^{\lrr}(L),\pi_0^{\lrr}(\alpha))\rightarrow V[\fm_y]\hookrightarrow V.
\end{equation}
Observe that $V$ is equipped with a natural action of $R_{\infty}[\frac{1}{p}]$,\;we can endow with an $R_{\infty}[\frac{1}{p}]$ action on the left side of (\ref{proofadjformla}) via $R_{\infty}[\frac{1}{p}]\twoheadrightarrow R_{\infty}[\frac{1}{p}]/\fm_y\cong E$ to make that the $(\mathfrak{g}_{\Sigma_L},\bP^{\lrr}(L))$-equivariant morphism $f$ is also $R_{\infty}[\frac{1}{p}]$-equivariant.\;Let $V'$ denote the pushfoward of $M'\otimes_E\cC^{\infty}_c(\bN^{\lrr}(L),\pi_0^{\lrr}(\alpha))$ via $f$.\;Then we get an exact sequence of  $(\mathfrak{g}_{\Sigma_L},\bP^{\lrr}(L))$-modules (endowed with natural  continuous $R_{\infty}[\frac{1}{p}]$ actions)
\begin{equation}\label{adjuncitonexact1}
	0\rightarrow V \rightarrow  V' \xrightarrow{\beta_0} L(s\cdot{\bm\lambda}_\bh)\otimes_E\cC^{\infty}_c(\bN^{\lrr}(L),\pi_0^{\lrr}(\alpha))
	\rightarrow0.\;
\end{equation}
It suffices to construct  a section of $\beta_0$ (as $(\mathfrak{g}_{\Sigma_L},\bP^{\lrr}(L))\times R_{\infty}[\frac{1}{p}] $-equivariant morphism.\;Indeed, if we can obtain a section $s_0$ of $\beta_0 $,\;then $s_0\circ f'$ gives a desired lifting of $f$ to $M'\otimes_E\cC^{\infty}_c(\bN^{\lrr}(L),\pi_0^{\lrr}(\alpha))$, and gives an $(\mathfrak{g}_{\Sigma_L},\bP^{\lrr}(L))$-morphism $M'\otimes_E\cC^{\infty}_c(\bN^{\lrr}(L),\pi_0^{\lrr}(\alpha))\rightarrow V[\fm_y]$.\;
Pulling back the exact sequence (\ref{adjuncitonexact1}) along the quotient map $M^{\lrr}(s\cdot{\bm\lambda}_\bh)\twoheadrightarrow L(s\cdot{\bm\lambda}_\bh)$,\;we get an exact sequence 
\begin{equation}\label{adjuncitonexact2}
	0\rightarrow V \rightarrow  V'' \xrightarrow{\beta_1}  M^{\lrr}(s\cdot{\bm\lambda}_\bh)\otimes_E\cC^{\infty}_c(\bN^{\lrr}(L),\pi_0^{\lrr}(\alpha))
	\rightarrow0.\;
\end{equation}
By Lemma \ref{detlameda}  and the argument after \cite[(80)]{2019DINGSimple} or \cite[(15)]{bergdall2018adjunction},\;it suffices to construct a $(\mathfrak{g}_{\Sigma_L},\bP^{\lrr}(L))\times R_{\infty}[\frac{1}{p}]$-equivariant map 
\begin{equation}\label{desiredhomo}
	M^{\lrr}(s\cdot{\bm\lambda}_\bh) \otimes_E\cC^{\infty}_c(\bN^{\lrr}(L),\pi_0^{\lrr}(\alpha))\rightarrow  V''.
\end{equation}
It remains to construct a $\bP^{\lrr}(L)\times R_{\infty}[\frac{1}{p}]$-equivariant map
\begin{equation}\label{desiredhomo1}
	L^{\lrr}(s\cdot{\bm\lambda}_\bh) \otimes_E\cC^{\infty}_c(\bN^{\lrr}(L),\pi_0^{\lrr}(\alpha))\rightarrow  V''.
\end{equation}
By using Proposition \ref{emertongenerpara} below and taking the $\fm_y$-generalized eigenspaces for $R_{\infty}[\frac{1}{p}]$,\;we can get an $R_{\infty}[\frac{1}{p}]$-equivariant exact sequence of finite-dimensional $E$-vector spaces
\begin{equation}
	\begin{aligned}
		0&\rightarrow B_{\omepik,\lambda_\bh}(V)[\fz^{\lrr}=d\omega_{s\cdot{\bm\lambda}_\bh}][\fm^\infty,\fm_y^\infty][
		\fm_{\omega_{s\cdot{\bm\lambda}_\bh}}]\\
		&\rightarrow B_{\omepik,\lambda_\bh}(V'')[\fz^{\lrr}=d\omega_{s\cdot{\bm\lambda}_\bh}][\fm^\infty,\fm_y^\infty][
		\fm_{\omega_{s\cdot{\bm\lambda}_\bh}}]\\
		&\rightarrow B_{\omepik,\lambda_\bh}(L(s\cdot{\bm\lambda}_\bh)\otimes_E\cC^{\infty}_c(\bN^{\lrr}(L),\pi_0^{\lrr}(\alpha)))[\fz^{\lrr}=d\omega_{s\cdot{\bm\lambda}_\bh}][\fm^\infty,\fm_y^\infty][
		\fm_{\omega_{s\cdot{\bm\lambda}_\bh}}]\\&\rightarrow 0.
	\end{aligned}
\end{equation}
where $\omega_{s\cdot{\bm\lambda}_\bh}$ denotes the central character of  $\pi_0^{\lrr}(\alpha)\te L^{\lrr}(s\cdot{\bm\lambda}_\bh)$,\;and $\fm$ denotes the maximal ideal of $\mathfrak{Z}_{\omepik}$ associated to $\pi_0^{\lrr}(\alpha)\te \delta_{\bP^{\lrr}}$.\;By Lemma \ref{detlameda},\;we deduce that
%Using the non-criticality of $(\lamoverx,\pi_0^{\lrr}(\alpha)\te \delta_{\bP^{\lrr}})$ with respect to $\Pi_{\infty}^{R_{\infty}-\ana}[\fm_x]$,\;
\[B_{\omepik,\lambda_\bh}(V)[\fz^{\lrr}=d\omega_{s\cdot{\bm\lambda}_\bh}][\fm^\infty,\fm_y^\infty][
\fm_{\omega_{s\cdot{\bm\lambda}_\bh}}]=0.\]
Hence we obtain an isomorphism
\begin{equation}\label{526}
	\begin{aligned}
		&B_{\omepik,\lambda_\bh}(V'')[\fz^{\lrr}=d\omega_{s\cdot{\bm\lambda}_\bh}][\fm^\infty,\fm_y^\infty][
		\fm_{\omega_{s\cdot{\bm\lambda}_\bh}}]\\
		\cong\;&B_{\omepik,\lambda_\bh}(L(s\cdot{\bm\lambda}_\bh)\otimes_E\cC^{\infty}_c(\bN^{\lrr}(L),\pi_0^{\lrr}(\alpha)))[\fz^{\lrr}=d\omega_{s\cdot{\bm\lambda}_\bh}][\fm^\infty,\fm_y^\infty][
		\fm_{\omega_{s\cdot{\bm\lambda}_\bh}}]\\ \cong\;&\pi_0^{\lrr}(\alpha)\te \delta_{\bP^{\lrr}}\te L^{\lrr}(s\cdot{\bm\lambda}_\bh),
	\end{aligned}
\end{equation}
where the second isomorphism follows from \cite[Lemma 3.5.2]{emerton2006jacquet}.\;Then (\ref{desiredhomo1}) is induced by the inverse of (\ref{526}) and \cite[Theorem\;3.5.6]{emerton2006jacquet}.\;
\end{proof}

The proof of \cite[Example 4.2]{bergdall2018adjunction} gives the following Lemma.\;
\begin{lem}\label{fpacycilic}Suppose that $V$ is an admissible continuous Banach representation of $G$,\;and there exists a compact open subgroup $H$ such that $V|_H\xrightarrow{\sim}C^0(H,L)^{\oplus l}$.\;Then 
\begin{equation}
	\hH^1\big(\overline{\fp}^{\lrr}_{\Sigma_L},V^{\BQ_p-\ana}\otimes U^\vee\big)=0,
\end{equation}
for all $i\in \BZ_{\geq 1}$, and any finite-dimensional $E$-linear locally analytic representation $U$ of $\bL^{\lrr}(L)$.\;
\end{lem}
%\begin{proof}Following the route of \cite[Example 4.2]{bergdall2018adjunction},\;it is sufficient to show that $\hH^1\big(\overline{\tau}^{\lrr}_{\Sigma_L},C^{\BQ_p-\ana}(G,L)\big)=0$.\;This result follows from the same strategy as in the proof of \cite[Example 4.2]{bergdall2018adjunction} with $P$  (resp.\;$\fp_L$) being our  $\on^{\lrr}(L)\rtimes \bZ^{\lrr}(L)$ (resp.,\;$\overline{\tau}^{\lrr}_{\Sigma_L}$).\;\end{proof}
%Passing to the derivative of a character induces a natural map $\rigerchl$,\;and hence (by restriction) a map \[\bersteineigenvarpik\rightarrow\widehat{\fz^{\lrr}}.\]

The following proposition generalizes a result of \cite[Proposition 4.1]{breuil2007first}.\;Let $\fo$ be an integral weight of $\fz^{\lrr}$.\;Let $\mathcal{E}_{\omepik,\fp,{\bm\lambda}_\bh}^\infty(\overline{\rho})_{\fo}$ be the fibre of $\bersteineigenvarpik$  at $\fo$ via the morphism 
\[\bersteineigenvarpik\hookrightarrow \FX_{\infty}\times \sbanpik\times \rigch\rightarrow \rigch\rightarrow (\fz^{\lrr})^\vee,\;\]where $\rigch\rightarrow (\fz^{\lrr})^\vee$ is the differentiation map $\chi\mapsto d\chi$.\;We use $\iota_0(\fz^{\lrr})$  to emphasize the action on $B_{\omepik,{\bm\lambda}_\bh}(\Pi_\infty^{R_\infty-\ana})$ derived from $\iota_0(\bZ^{\lrr}(\cO_L))$.\;For an $E$-algebra $A$,\;$\fm\subset A$ an ideal,\;and an $A$-module,\;we denote by $M[\fm^{\infty}]$ the $A$-submodule of $M$ consisting of elements annihilated by $\fm^n$ for some $n\geq 0$.\;Then \cite[Lemma 3.1.4]{Ding2021} shows that 
\begin{equation}
\begin{aligned}
	%\Gamma(\mathcal{E}_{\omepik,\fp,{\bm\lambda}_\bh}^\infty(\overline{\rho})_{\fo},(\cM_{\omepik,{\bm\lambda}_\bh}^{\infty})_{\fo})^\vee&=
	&B_{\omepik,{\bm\lambda}_\bh}(\Pi_\infty^{R_\infty-\ana})[\iota_0(\fz^{\lrr})=\fo]\\
	=&\;\bigoplus_{\delta,\chi}B_{\omepik,{\bm\lambda}_\bh}(\Pi_\infty^{R_\infty-\ana})[\iota_0(\fz^{\lrr})=\fo][\fm_\chi][\fm_\delta^\infty]\\
	=&\;\bigoplus_{\fm\in \MSpec\mathfrak{Z}_{\omepik},\chi}B_{\omepik,{\bm\lambda}_\bh}(\Pi_\infty^{R_\infty-\ana})[\iota_0(\fz^{\lrr})=\fo][\fm_\chi][\fm^\infty],\;
\end{aligned}
\end{equation}
where $\delta$ (resp.\;$\chi$) runs through the smooth characters of $\Delta_0$ (see \cite[(3.4)]{Ding2021}) (resp.\;through the locally algebraic character of $\bZ^{\lrr}(\cO_L)$ of weight $\fo$),\;and $\fm_\delta\subset E[\Delta_0]$ (resp.\;$\fm_\chi\subset  E[\bZ^{\lrr}(\cO_L)]$) is the maximal ideal associated to $\delta$ (resp.\;$\chi$).\;Replacing character $\delta$ of $\Delta_0$ by maximal ideals $\fm$ of $\mathfrak{Z}_{\omepik}$,\;we can get the second identity by the proof of \cite[Lemma 3.1.4]{Ding2021}.\;By \cite[Lemma 3.1.4]{Ding2021},\;we see that each term in the direct sums is finite-dimensional over $E$.\;
%This is equivalent to say that $\mathcal{E}_{\omepik,\fp,{\bm\lambda}_\bh}^\infty(\overline{\rho})_{\fo}$ is discrete.\;Therefore,\;we deduce
\begin{pro}\label{emertongenerpara}Suppose that $V$ is an admissible continuous Banach representation of $G$ and there exists a compact open subgroup $H$ such that $V|_H\xrightarrow{\sim}C^0(H,L)^{\oplus l}$.\;Assume that
\begin{equation}\label{exactanaylytic}
	0\rightarrow V^{\BQ_p-\ana}\rightarrow \Pi\rightarrow \Pi_1\rightarrow 0,
\end{equation}
is an exact sequence of admissible locally $\BQ_p$-analytic representations of $G$.\;Let $\chi:\bZ^{\lrr}(\cO_L)\rightarrow E^\times$ be a continuous character of weight $\fd$,\;and let $\fm$ be a maximal ideal of $\mathfrak{Z}_{\omepik}$.\;Now we have short exact sequences
\begin{equation}\label{techexactseq1}
	\begin{aligned}
		0\rightarrow (V^{\BQ_p-\ana})^{\bN^{\lrr}(L)_0}[\fz^{\lrr}=\fo,\fd^{\lrr}=0]\rightarrow& \Pi^{\bN^{\lrr}(L)_0}[\fz^{\lrr}=\fo,\fd^{\lrr}=0] \\
		&\rightarrow \Pi_1^{\bN^{\lrr}(L)_0}[\fz^{\lrr}=\fo,\fd^{\lrr}=0]\rightarrow 0,
	\end{aligned}
\end{equation}
\begin{equation}\label{techexactseq2}
	\begin{aligned}
		0\rightarrow B_{\omepik,{\bm\lambda}_\bh}(V^{\BQ_p-\ana})[\fz^{\lrr}=\fo][\fm^\infty][
		\fm_{\chi}]\rightarrow &B_{\omepik,{\bm\lambda}_\bh}(\Pi)[\fz^{\lrr}=\fo][\fm^\infty][
		\fm_{\chi}]\\
		&\rightarrow B_{\omepik,{\bm\lambda}_\bh}(\Pi_1)[\fz^{\lrr}=\fo][\fm^\infty][
		\fm_{\chi}]\rightarrow 0.
	\end{aligned}
\end{equation}
Moreover,\;all the vector spaces in the last exact sequence are finite-dimensional.\;
\end{pro}
\begin{proof}The above proposition implies that taking $\overline{\fp}_{\Sigma_L}^{\lrr}$-invariant is exact on the short exact sequence (\ref{exactanaylytic}).\;Since taking $\bN^{\lrr}(L)_0$-invariant is exact on the category of smooth representations of $\bN^{\lrr}(L)_0$,\;we deduce the first exact sequence (\ref{techexactseq1}).\;To see the exactness of the second sequence of (\ref{techexactseq2}),\;we need to unwind the action of $\Delta_0$,\;$\iota_0(\bZ^{\lrr}(\cO_L))$ and $\iota_1(\bZ^{\lrr}(\cO_L))$ on $B_{\omepik,{\bm\lambda}_\bh}(W)$ for $W\in\{V^{\BQ_p-\ana},\Pi,\Pi_1\}$.\;As in the proof of \cite[Lemma 3.1.4]{Ding2021},\;we can choose a compact open subgroup of  $\bZ^{\lrr}(\cO_L)$ such that $B_{\omepik,{\bm\lambda}_\bh}(W)[\fz^{\lrr}=\fo]$ is a $\Delta_0$-equivalent direct summand of 
\begin{equation*}
	\begin{aligned}
		\big(J_{\bP^{\lrr}(L)}&(W)_{{\bm\lambda}_\bh}[\fz^{\lrr}=\fo\circ_{\bL^{\lrr}(L)}]\widehat{\otimes}_E\cC^{\BQ_p-\ana}(\bZ^{\lrr}(\cO_L),E)[\fz^{\lrr}=\fo]\otimes_E\sigma^\vee\big)^H\\
		\cong\;\bigoplus_{\delta'}&\big(\big(J_{\bP^{\lrr}(L)}(W)_{{\bm\lambda}_\bh}\otimes_E\sigma^\vee\big)^{H^D}[\fz^{\lrr}=\fo\circ\mathrm{det}_{\bL^{\lrr}(L)}][\fm_{\delta'}^\infty]\\&\widehat{\otimes}_E\cC^{\BQ_p-\ana}(\bZ^{\lrr}(\cO_L),E)[\fz^{\lrr}=\fo]\big)^{Z_H},		
	\end{aligned}
\end{equation*}
where $H^D:=H\cap \bD^{\lrr}(\cO_L)$ (resp.,\;$Z_H:=H\cap\bZ^{\lrr}(\cO_L)$),\;and $\delta'$ runs though the locally algebraic characters of $\iota_1(\bZ^{\lrr}(\cO_L))$ (and $[\fm_{\delta'}^\infty]$ is for the corresponding $\iota_1(\bZ^{\lrr}(\cO_L))$-action).\;Similar to the proof of \cite[Proposition 4.1,\;(4.3),\;Pages 10521-10522]{breuil2007first},\;we have a short exact sequence of finite-dimensional $E$-vector spaces:
\begin{equation*}
	\begin{aligned}
		0 \rightarrow \big(J_{\bP^{\lrr}(L)}(V^{\BQ_p-\ana})_{{\bm\lambda}_\bh}&\otimes_E\sigma^\vee\big)^{H^D}[\fz^{\lrr}=\fo\circ_{\bL^{\lrr}(L)}][\fm_{\delta'}^\infty] \\
		\rightarrow& \big(J_{\bP^{\lrr}(L)}(\Pi)_{{\bm\lambda}_\bh}\otimes_E\sigma^\vee\big)^{H^D}[\fz^{\lrr}=\fo\circ_{\bL^{\lrr}(L)}][\fm_{\delta'}^\infty]\\ & \rightarrow \big(J_{\bP^{\lrr}(L)}(\Pi_1)_{{\bm\lambda}_\bh}\otimes_E\sigma^\vee\big)^{H^D}[\fz^{\lrr}=\fo\circ_{\bL^{\lrr}(L)}][\fm_{\delta'}^\infty] \rightarrow 0.\;
	\end{aligned}
\end{equation*}
Then the result follows by taking $[\fm_\chi]$-eigenspaces and replacing character $\delta'$ of $\Delta_0$ by maximal ideals $\fm$ of $\mathfrak{Z}_{\omepik}$,\;as in the proof of \cite[Lemma 3.1.4]{Ding2021}.\;
\end{proof}

%Using the same strategy as in the proof of \cite[Prop.\;4.1]{breuil2020towards} and the fact that $\Supp(\cM_{\infty,\omepik,\bm{\lambda}_\bh}^0)\rightarrow\widehat{\fz^{\lrr}}$ has discrete fibers (see Proposition \ref{eigendiscretefibre}),\;we see that the three term in the second sequence are finite-dimensional vector space.\;A slight modification of \cite[Prop.\;4.1]{breuil2020towards} can show the exactness of the second sequence.\;

Let $\FI_y\subseteq \fm_y$ be a closed ideal of $R_{\infty}[1/p]$ such that $\dim_E (R_{\infty}[1/p]/\FI_y) < +\infty$ and that $\fm_y$ is the unique closed maximal ideal containing $\FI_y$ (e.g. $\FI_y=\fm_y^k$).\;Similar to \cite[Corollary 4.9]{2015Ding},\;we can obtain
\begin{cor}\label{cor: lgln-bij3}
Let $\st_{(r,k)}^{\infty}(\alpha,\pi_0,{\bm\lambda}_\bh)\subseteq W\subseteq \st_{(r,k)}^{\ana}(\alpha,\pi_0, {\bm\lambda}_\bh)$.\;If a morphism $$f: W\ra \Pi_{\infty}^{R_{\infty}-\ana}[\FI_y]$$ satisfies that $f|_{\st_{(r,k)}^{\infty}(\alpha,\pi_0,{\bm\lambda}_\bh)}\subseteq \Pi_{\infty}^{R_{\infty}-\ana}[\fm_y]$,\;then $f$ has image in $\Pi_{\infty}^{R_{\infty}-\ana}[\fm_y]$.
\end{cor}




%\begin{proof}
%  (1) follows from (\ref{soc2}).
%
%  \noindent
%
%\end{proof}
%\noindent Suppose $N^{n-1}\neq 0$ on $D_{\st}(D)$. For $i\in \Delta$, $\sigma\in \Sigma_L$, there exists $\bL_{i,\sigma}\in E$ such that $\bL(D)_i\cap \homo_{\sigma}(L^{\times}, E)$ is generated by $\psi_{i,\sigma}:=\log_{p, \sigma}-\bL_{i,\sigma} \valua_L$. The characters $\{\psi_{i,\sigma}\}_{\sigma\in \Sigma_L}$ form a basis of $\bL(D)_i$. Since $\psi_{i,\sigma}$ is locally $\sigma$-analytic, by (\ref{equ: BrLi}) Proposition \ref{prop: sigmaan}, we have (recall $\Sigma_{i,\sigma}(\alpha,\pi_0, {\bm\lambda}_\bh, \psi_{i,\sigma})$ is $\text{U}(\ug_{\Sigma_L\setminus\{\sigma\}})$-finite)
%\begin{equation*}\Sigma_i^{\lrr}(\alpha,\pi_0, {\bm\lambda}_\bh, \psi_{i,\sigma})\cong \Sigma_{i,\sigma}(\alpha,\pi_0, {\bm\lambda}_\bh, \psi_{i,\sigma}) \oplus_{\st_n^{\infty}(\alpha,\pi_0, {\bm\lambda}_\bh)} \big(\bigoplus_{\st_n^{\infty}(\alpha,\pi_0, {\bm\lambda}_\bh)}^{\tau\in \Sigma_L, \tau\neq \sigma} \Sigma_{i,\tau}(\alpha,\pi_0, {\bm\lambda}_\bh)\big).\end{equation*}
%In this case, using the basis $\{\psi_{i,\sigma}\}_{i\in \Delta, \sigma\in \Sigma_L}$ of $\bL(D)$, we see that $\Sigma^{\lrr}(\alpha,\pi_0, {\bm\lambda}_\bh, \bL(D))$ has the following form (where a line denotes a \emph{non-split} extension with the left object the sub, and the right object the quotient, and we enumerate the embeddings in $\Sigma_L$ for convenience):
%\begin{equation}\label{stru}
%  \begindc{\commdiag}[200]
%    \obj(0,3)[a]{$\st_n^{\infty}(\alpha,\pi_0, {\bm\lambda}_\bh)$}
%    \obj(4,6)[b]{$C_{1,\sigma_1}$}
%    \obj(8,6)[c]{$v_{\overline{P}_1}^{\infty}(\alpha,\pi_0,{\bm\lambda}_\bh)$}
%    \obj(4,5)[d]{$\vdots$}
%    \obj(8,5)[e]{$\vdots$}
%    \obj(4,4)[f]{$C_{1, \sigma_{d_L}}$}
%    \obj(8,4)[\bG]{$v_{\overline{P}_1}^{\infty}(\alpha,\pi_0, {\bm\lambda}_\bh)$}
%    \obj(4,3)[h]{$\vdots$}
%    \obj(8,3)[i]{$\vdots$}
%    \obj(4,2)[j]{$C_{n-1,\sigma_1}$}
%    \obj(8,2)[k]{$v_{\overline{P}_{n-1}}^{\infty}(\alpha,\pi_0, {\bm\lambda}_\bh)$}
%    \obj(4,1)[l]{$\vdots$}
%    \obj(8,1)[m]{$\vdots$}
%    \obj(4,0)[n]{$C_{n-1,\sigma_{d_L}}$}
%    \obj(8,0)[o]{$v_{\overline{P}_{n-1}}^{\infty}(\alpha,\pi_0, {\bm\lambda}_\bh)$}
%    \mor{a}{b}{}[+1,\solidline]
%    \mor{b}{c}{}[+1,\solidline]
%    \mor{a}{f}{}[+2,\solidline]
%    \mor{f}{g}{}[+1,\solidline]
%       \mor{a}{j}{}[+1,\solidline]
%    \mor{j}{k}{}[+1,\solidline]
%       \mor{a}{n}{}[+1,\solidline]
%    \mor{n}{o}{}[+1,\solidline]
%  \enddc,
%\end{equation}
%where the $(i,\sigma)$-th branch carries the exact information on $\bL_{i,\sigma}$, and is $\text{U}(\ug_{\Sigma_L\setminus \{\sigma\}})$-finite.
%\\
The following theorem is the main result of this paper.\;As soon as  we finish the previous preparation,\;the proof is almost completely parallel to the proof of \cite[Theorem 4.10]{2019DINGSimple}.\;We include a proof for the reader's convenience.\;

\begin{thm}\label{thm: lgln-main}
(1) The following restriction map is bijective
\begin{equation}\label{equ: lgln-main}
	\begin{aligned}
		\homo_{G}\big(\Sigma^{\lrr}(\alpha,\pi_0,{\bm\lambda}_\bh, \sL(\rho_L)),& \Pi_{\infty}^{R_{\infty}-\ana}[\fm_y]\big) \\
		&\lra \homo_{G}\big(\st_{(r,k)}^{\infty}(\alpha,\pi_0,{\bm\lambda}_\bh), \Pi_{\infty}^{R_{\infty}-\ana}[\fm_y]\big).
	\end{aligned}
\end{equation}


\noindent
(2) Let $0\neq \psi\in \homo(L^{\times},E)$ and $ir\in \Delta_n(k)$, an injection
\begin{equation*}
	f: \st_{(r,k)}^{\infty}(\alpha,\pi_0,{\bm\lambda}_\bh)\hooklongrightarrow \Pi_{\infty}^{R_{\infty}-\ana}[\fm_y]   \end{equation*}
can extend to an injection $\Sigma_{i}(\alpha, \pi,{\bm\lambda}_\bh, \psi)\hookrightarrow \Pi_{\infty}^{R_{\infty}-\ana}[\fm_y]$ if and only if $\psi\in \sL(\rho_L)_{ir}$.
\end{thm}

\noindent The rest of the section is to prove Theorem \ref{thm: lgln-main}, and we use the strategy of (\cite{2019DINGSimple}).\;\\
\begin{proof}
\textbf{(a)} The ``if " part  of Theorem \ref{thm: lgln-main} (2) is a consequence of (1).\;It suffices to prove the ``only if " part.\;Otherwise,\;we assume that there exists $\Sigma_{i}^{\lrr}(\alpha,\pi_0,{\bm\lambda}_\bh, \psi)\hookrightarrow \Pi_{\infty}^{R_{\infty}-\ana}[\fm_y]$ with $\psi\notin \sL(\rho_L)_{ir}$. Recall that $\sL(\rho_L)_{ir}$ is of codimension $1$ in $\homo(L^{\times}, E)$,\;we see that  $\sL(\rho_L)_{ir}+ E\psi =\homo(L^{\times}, E)$. Then the injection
\begin{equation*}
	\Sigma_i^{\lrr}(\alpha, \pi,{\bm\lambda}_\bh, \sL(\rho_L)_{ir})\oplus_{\Sigma_i^{\lrr}(\alpha,\pi_0,{\bm\lambda}_\bh)} \Sigma_{i}^{\lrr}(\alpha,\pi_0,{\bm\lambda}_\bh, \psi) \hooklongrightarrow \Pi_{\infty}^{R_{\infty}-\ana}[\fm_y]
\end{equation*}
induces
\begin{equation*}
	\Sigma_i^{\lrr}(\alpha, \pi,{\bm\lambda}_\bh, \psi_{\infty}) \hooklongrightarrow\Sigma_i^{\lrr}(\alpha, \pi,{\bm\lambda}_\bh, \sL(\rho_L)_{ir})\oplus_{\Sigma_i^{\lrr}(\alpha, \pi,{\bm\lambda}_\bh)} \Sigma_{i}^{\lrr}(\alpha, \pi,{\bm\lambda}_\bh, \psi) \hooklongrightarrow \Pi_{\infty}^{R_{\infty}-\ana}[\fm_y].
\end{equation*}
for any $0\neq \psi_{\infty}\in \homo_{\infty}(L^{\times},E)$.\;But by \cite[Remark 5.21]{2022ext1hyq}, $\Sigma_i^{\lrr}(\alpha,\pi_0,{\bm\lambda}_\bh, \psi_{\infty})$ contains $V\otimes_E L({\bm\lambda}_\bh)$ where $V$ is a smooth extension of $v_{\op^{\lrr}_I}^{\infty}(\pi)$ by $\st_{(r,k)}^{\infty}(\pi)$.\;By applying the (left exact) Jacquet-Emerton functor to $V\otimes_E L({\bm\lambda}_\bh)\hookrightarrow \Pi_{\infty}^{R_{\infty}-\ana}[\fm_y]$,\;we get a contradiction with Lemma \ref{detlameda}.\;This completes the proof of (2).\;
\\
\noindent \textbf{(b)} Since $\soc_{G} \Sigma^{\lrr}(\alpha,\pi_0,{\bm\lambda}_\bh=\st_{(r,k)}^{\infty}(\alpha,\pi_0,{\bm\lambda}_\bh)$,\;it follows from Lemma \ref{detlameda} and the step (a) of the proof of Proposition \ref{proadjuncitonformularight}) that the injectivity of (\ref{equ: lgln-main}) holds.\;It remains  to  show that (\ref{equ: lgln-main}) is surjective.\;By definition, it suffices to show that for any $ir\in \Delta_n(k)$, $\psi\in \sL(\rho_L)_{ir}$, the following restriction map is surjective
\begin{equation}\label{equ: lgln-rest}
	\homo_{\GLN_n(L)}\big(\Sigma_i^{\lrr}(\alpha ,\pi,{\bm\lambda}_\bh, \psi), \Pi_{\infty}^{R_{\infty}-\ana}[\fm_y]\big) \lra \homo_{\GLN_n(L)}\big(\st_{(r,k)}^{\infty}(\alpha,\pi_0,{\bm\lambda}_\bh), \Pi_{\infty}^{R_{\infty}-\ana}[\fm_y]\big).
\end{equation}
The key ingredient is the consequence of the surjectivity of (\ref{eigentangentmap2}).\;
By Corollary \ref{coreigentangentmap2},\;there exists a  $t: \Spec E[\epsilon]/\epsilon^2 \ra \bersteineigenvarpik$ (as an element in $T_{\bersteineigenvarpik,x}$)\;such that the $i$-th factor of the image of $t$ under (\ref{eigentangentmap2}) equals $\psi$, and the $j$-th factors for all $j\neq i$ are zero.\;Let $\Psi\in \homo(\bZ^{\lrr}(L),E)$ be the image of $t$ via the first map in (\ref{tangentpatchvar1}).\;By Proposition \ref{basicpropertyeigen} (3),\;we see that the coherent sheaf $\cM^{\infty}_{\omepik,{\bm\lambda}_\bh}$ is Cohen-Macaulay over $\mathcal{E}_{\omepik,\fp,{\bm\lambda}_\bh}^\infty(\overline{\rho})$.\;Since 
non-critical special point $x$ is a smooth point on $\bersteineigenvarpik$,\;we see that $\cM^{\infty}_{\omepik,{\bm\lambda}_\bh}$ is locally free in a certain neighborhood of $x$.\;let $\FI_t$ denotes the kernel of the morphism $R_{\infty}[1/p] \ra E[\epsilon]/\epsilon^2$ induced by $t$.\;We put \[\pi_0^{\lrr}(\alpha,{\bm\lambda}_\bh,\Psi):=\pi_0^{\lrr}(\alpha,{\bm\lambda}_\bh)\te({1+\Psi \epsilon})\circ\mathrm{det}_{\bL^{\lrr}(L)}.\;\]By the construction of $\cM^{\infty}_{\omepik,{\bm\lambda}_\bh}$,\;we deduce the following facts:
\begin{itemize}
	\item[(a)] $(x^*\cM^{\infty}_{\omepik,{\bm\lambda}_\bh})^{\vee}\cong \homo_{\bL^{\lrr}(L)}\big(\pi_0^{\lrr}(\alpha,{\bm\lambda}_\bh)\te \delta_{\bP^{\lrr}}, {J}_{\bP^{\lrr}(L)}\big(\Pi_{\infty}^{R_{\infty}-\ana}[\fm_y]\big)\big)$,\;and \\
	$(t^* \cM^{\infty}_{\omepik,{\bm\lambda}_\bh})^{\vee}\cong \homo_{\bL^{\lrr}(L)}\big(\pi_0^{\lrr}(\alpha,{\bm\lambda}_\bh,\Psi)\te \delta_{\bP^{\lrr}}, J_{\bP^{\lrr}(L)}\big(\Pi_{\infty}^{R_{\infty}-\ana}[\FI_t]\big)\big)$,\;\\which are closed subrepresentations of $J_{\bP^{\lrr}(L)}\big(\Pi_{\infty}^{R_{\infty}-\ana}[\fm_y])$ and $J_{\bP^{\lrr}(L)}(\Pi_{\infty}^{R_{\infty}-\ana}[\FI_t])$ respectively;
	% (since Jacquet Emerton functor is left exact and $\pi_0^{\lrr}(\alpha,{\bm\lambda}_\bh)\te \delta_{\bP^{\lrr}}\cong J_{\bP^{\lrr}(L)}\big(\st^{\infty}_{(r,k)}(\alpha,\pi_0,{\bm\lambda}_\bh)\big)$);
	\item[(b)] there are natural $\bL^{\lrr}(L)$-equivariant injections
	\begin{equation}\label{equ: lgln-injf}(x^*\cM^{\infty}_{\omepik,{\bm\lambda}_\bh})^{\vee}\hooklongrightarrow (t^* \cM^{\infty}_{\omepik,{\bm\lambda}_\bh})^{\vee} \hooklongrightarrow J_{\bP^{\lrr}(L)}(\Pi_{\infty}^{R_{\infty}-\ana}[\FI_t]).\end{equation}
\end{itemize}
By Lemma \ref{detlameda} and an easy variation of the proof of
\cite[Lemma 4.11]{2019DINGSimple},\;we can show that
\begin{lem} The morphisms of $\bL^{\lrr}(L)$-representations:
	\begin{equation}
		\begin{aligned}
			&(x^*\cM_{\infty,\omepik,{\bm\lambda}_\bh})^{\vee} \hookrightarrow J_{\bP^{\lrr}(L)}(\Pi_{\infty}^{R_{\infty}-\ana}[\fm_y]),\\
			\text{and\;}&(t^* \cM_{\infty,\omepik,{\bm\lambda}_\bh})^{\vee} \hookrightarrow
			J_{\bP^{\lrr}(L)}(\Pi_{\infty}^{R_{\infty}-\ana}[\FI_t])
		\end{aligned}
	\end{equation}
	are balanced (see \cite[Definition 0.8]{emerton2007jacquet},\;\cite[Definition 5.17]{Emerton2007summary}).
\end{lem}
\begin{proof}This lemma follows by an easy variation of the proof of \cite[Lemma 4.11]{2019DINGSimple}.\;We briefly indicate below the changes.\;The notation ''$U(\bB(L))$'' (resp.,\;$\cC^{\sm}_c\big(\bN(L),-)$,\;resp.,\;$\cC^{\lp}_c\big(\bN(L),-)$,\;resp., $\cC^{\BQ_p-\pol}_c\big(\bN(L),-)$) has to be replaced by $\text{U}(\fp^{\lrr}_{\Sigma_L})$ (resp.,\;$\cC^{\sm}_c\big(\bN^{\lrr}(L),-)$,\;resp.,$\cC^{\lp}_c\big(\bN^{\lrr}(L),-)$,\;resp., $\cC^{\BQ_p-\pol}_c\big(\bN^{\lrr}(L),-)$).\;The short exact sequence ``$0\rightarrow\chi\rightarrow\widetilde{\chi}\rightarrow\chi\rightarrow0$"  has to be replaced by the short exact sequence
	$$ ``0 \ra\pi_0^{\lrr}(\alpha,{\bm\lambda}_\bh)\te \delta_{\bP^{\lrr}}  \ra \pi_0^{\lrr}(\alpha,{\bm\lambda}_\bh,\Psi)\te \delta_{\bP^{\lrr}}\ra \pi_0^{\lrr}(\alpha,{\bm\lambda}_\bh)\te \delta_{\bP^{\lrr}} \ra 0".\;$$The "\cite[Lemma 4.6]{2019DINGSimple}" has to be replaced by Lemmas \ref{detlameda}.\;
\end{proof}

This lemma shows that the adjunction formula in \cite{emerton2007jacquet} is suitable for us.\;Let $V''$ be a  locally $\BQ_p$-analytic representation of $\bL^{\lrr}(L)$.\;We refer to \cite{emerton2007jacquet} for the detail of locally $\BQ_p$-analytic representation $I_{\op^{\lrr}}^G(V'')$ of $G$,\;which is a closed $G$-subrepresentation of $(\ind_{\op^{\lrr}(L)}^{G} V'')^{\BQ_p-\ana}$.\;By \cite[Theorem 0.13]{emerton2007jacquet},\;there exists an integer $r$,\;such that the injections in (\ref{equ: lgln-injf}) induce
\begin{equation}\label{equ: lgln-keymap}
	I_{\op^{\lrr}}^G \big(\pi_0^{\lrr}(\alpha,{\bm\lambda}_\bh)\big)^{\oplus r} \hooklongrightarrow I_{\op^{\lrr}}^G \big(\pi_0^{\lrr}(\alpha,{\bm\lambda}_\bh,\Psi)\big)^{\oplus r} \longrightarrow \pi_{\infty}^{R_{\infty}-\ana}[\FI_t].
\end{equation}
By \cite[Proposition 2.8.10]{emerton2007jacquet},\;we have $I_{\op^{\lrr}}^G (\pi_0^{\lrr}(\alpha,{\bm\lambda}_\bh))\cong i_{\op^{\lrr}}^G(\alpha,\pi_0, {\bm\lambda}_\bh)$.\;The natural exact sequence $$0\ra \;\pi_0^{\lrr}(\alpha,{\bm\lambda}_\bh)\; \ra \;\pi_0^{\lrr}(\alpha,{\bm\lambda}_\bh,\Psi)\; \ra \;\pi_0^{\lrr}(\alpha,{\bm\lambda}_\bh)\; \ra 0$$ induces a sequence (not necessary exact)
\begin{equation*}
	i_{\op^{\lrr}}^G(\alpha,\pi_0, {\bm\lambda}_\bh) \hooklongrightarrow I_{\op^{\lrr}}^G(\pi_0^{\lrr}(\alpha,{\bm\lambda}_\bh,\Psi)) \twoheadlongrightarrow  i_{\op^{\lrr}}^G(\alpha,\pi_0, {\bm\lambda}_\bh).
\end{equation*}
Then $v_{\op^{\lrr}_I}^{\infty}(\alpha,\pi_0, {\bm\lambda}_\bh)$ is an irreducible constituent in $I_{\op^{\lrr}}^G(\pi_0^{\lrr}(\alpha,{\bm\lambda}_\bh,\Psi)) $ of multiplicity $2$.\;

Using the natural embedding $$I_{\op^{\lrr}}^G(\pi_0^{\lrr}(\alpha,{\bm\lambda}_\bh,\Psi))  \hookrightarrow \Big(\ind_{\op^{\lrr}(L)}^{G} \pi_0^{\lrr}(\alpha,{\bm\lambda}_\bh,\Psi)\Big)^{\BQ_p-\ana},\;$$we can take inside  $\Big(\ind_{\op^{\lrr}(L)}^{G} \pi_0^{\lrr}(\alpha,{\bm\lambda}_\bh,\Psi)\Big)^{\BQ_p-\ana}$ the following intersections:
\begin{eqnarray*}U&:=&I_{\op^{\lrr}}^G(\pi_0^{\lrr}(\alpha,{\bm\lambda}_\bh,\Psi)) \cap \Big(\sum_{\emptyset\neq I\subseteq \Delta_n(k)} \BI_{\op^{\lrr}_{I}}^G(\alpha,\pi_0, {\bm\lambda}_\bh)\Big), \\
	W&:=&I_{\op^{\lrr}}^G(\pi_0^{\lrr}(\alpha,{\bm\lambda}_\bh,\Psi))\cap \BI_{\op^{\lrr}}^G(\alpha,\pi_0, {\bm\lambda}_\bh).
\end{eqnarray*}
Thus $\widetilde{\Sigma}^{\lrr}(\alpha,\pi_0, {\bm\lambda}_\bh)' :=W/U$ is a subrepresentation of $\st_{(r,k)}^{\ana}(\alpha,\pi_0, {\bm\lambda}_\bh)$, and  $I_{\op^{\lrr}}^G(\pi_0^{\lrr}(\alpha,{\bm\lambda}_\bh,\Psi))/U$ is an extension of $i_{\op^{\lrr}}^G(\alpha,\pi_0, {\bm\lambda}_\bh)$ by $\widetilde{\Sigma}^{\lrr}(\alpha,\pi_0, {\bm\lambda}_\bh)'$.\;As in Step (a) of the  proof of Proposition \ref{proadjuncitonformularight},\;we deduce from Lemma \ref{detlameda} that any irreducible constituent of $U$ can not appear in the socle of $\pi_{\infty}^{R_{\infty}-\ana}[\FI_t]$.\;Then (\ref{equ: lgln-keymap}) induces
\begin{equation}\label{equ: lgln-keymap2}
	\st_{(r,k)}^{\infty}(\alpha,\pi_0, {\bm\lambda}_\bh)^{\oplus r}\hooklongrightarrow \big(I_{\op^{\lrr}}^G (\pi_0^{\lrr}(\alpha,{\bm\lambda}_\bh,\Psi))/U\big)^{\oplus r} \lra \Pi_{\infty}^{R_{\infty}-\ana}[\FI_t].
\end{equation}
This composition is injective.\;Since $(x^* \cM_{\infty})^{\vee}$ has image in $J_{\bP^{\lrr}(L)}(\pi_{\infty}^{R_{\infty}-\ana}[\fm_y])$ via (\ref{equ: lgln-injf}),\;we see that it factors through an injection $\Pi_{\infty}^{R_{\infty}-\ana}[\fm_y]\hookrightarrow \Pi_{\infty}^{R_{\infty}-\ana}[\FI_t]$ by Corollary \ref{cor: lgln-bij3}.\;

Denote by $\Sigma_i^{\lrr}(\alpha,\pi_0, {\bm\lambda}_\bh)':=\widetilde{\Sigma}^{\lrr}(\alpha,\pi_0, {\bm\lambda}_\bh)' \cap \Sigma_i^{\lrr}(\alpha,\pi_0, {\bm\lambda}_\bh)$.\;Then \cite[Proposition 5.28]{2022ext1hyq} deduce $\soc_{G} \Sigma_i^{\lrr}(\alpha,\pi_0, {\bm\lambda}_\bh)'\cong \st_{(r,k)}^{\infty}(\alpha,\pi_0, {\bm\lambda}_\bh)$.\;Thus by Corollary \ref{cor: lgln-bij3},\;the  composition
\begin{equation}\label{equ: compW}
	(\Sigma_i^{\lrr}(\alpha,\pi_0, {\bm\lambda}_\bh)')^{\oplus r} \hooklongrightarrow \big(I_{\op^{\lrr}}^G (\pi_0^{\lrr}(\alpha,{\bm\lambda}_\bh,\Psi))/U\big)^{\oplus r} \lra \pi_{\infty}^{R_{\infty}-\ana}[\FI_t]
\end{equation}
also has image in $\pi_{\infty}^{R_{\infty}-\ana}[\fm_y]$.\;By Lemma \ref{detlameda},\;it is also injective.\;By Proposition \ref{proadjuncitonformularight}, we see that (\ref{equ: compW}) extends uniquely to   an injection
\begin{equation}\label{equ: compnoL}\Sigma_i^{\lrr}(\alpha,\pi_0, {\bm\lambda}_\bh)^{\oplus r} \lra \pi_{\infty}^{R_{\infty}-\ana}[\fm_y] \hooklongrightarrow \pi_{\infty}^{R_{\infty}-\ana}[\FI_t].\end{equation}
Combining (\ref{equ: compW}) with  (\ref{equ: compnoL}),\;we put $$V^+:=\big(I_{\op^{\lrr}}^G (\pi_0^{\lrr}(\alpha,{\bm\lambda}_\bh,\Psi))/U\big)\oplus_{\Sigma_i^{\lrr}(\alpha,\pi_0, {\bm\lambda}_\bh)'} \Sigma_i^{\lrr}(\alpha,\pi_0, {\bm\lambda}_\bh).$$
Therefore,\;(\ref{equ: compW}) and  (\ref{equ: compnoL}) give
\begin{equation}\label{equ: lgL1}
	\begin{aligned}
		\Sigma_i^{\lrr}(\alpha,\pi_0, {\bm\lambda}_\bh)^{\oplus r}\hooklongrightarrow \big(V^+\big)^{\oplus r}\lra \pi_{\infty}^{R_{\infty}-\ana}[\FI_t].
	\end{aligned}
\end{equation}
It has image in $\pi_{\infty}^{R_{\infty}-\ana}[\fm_y]$.\;It suffices to prove the following assertion.\;\\
\textbf{Claim.\;}$\Sigma_i^{\lrr}(\alpha,\pi_0, {\bm\lambda}_\bh, \psi)$ is a subrepresentation of $V^+$.\;\\
\textbf{Proof of the claim.\;}Let $V$ be the pull-back of $I_{\op^{\lrr}}^G (\pi_0^{\lrr}(\alpha,{\bm\lambda}_\bh,\Psi))/U$ via the injection \[i_{\op^{\lrr}_{ir}}^G(\alpha,\pi_0, {\bm\lambda}_\bh) \hookrightarrow i_{\op^{\lrr}}^G(\alpha,\pi_0, {\bm\lambda}_\bh).\;\]Then $V$ lies in a commutative diagram
\begin{equation}\label{commrevi}
	\xymatrix{0 \ar[r] & \widetilde{\Sigma}^{\lrr}(\alpha,\pi_0, {\bm\lambda}_\bh)' \ar[d]^{\subseteq}  \ar[r] & \hspace{20pt}V\hspace{20pt} \ar[d]^{\subseteq} \ar[r] & i_{\op^{\lrr}_{ir}}^{G}(\alpha,\pi_0,{\bm\lambda}_\bh) \ar@{=}[d] \ar[r] & 0  \\
		%0 \ar[r] & \widetilde{\Sigma}^{\lrr}(\alpha,\pi_0, {\bm\lambda}_\bh)'  \ar[r] \ar[d]^{\subseteq} & \BI_{\op^{\lrr}}^G (\pi_0^{\lrr}(\alpha,{\bm\lambda}_\bh,\Psi))/U \ar[r] \ar[d]^{\subseteq} & \BI_{\op^{lrr}_I}^{G}(\alpha,\pi_0,{\bm\lambda}_\bh)\ar[r] \ar[d]^{=} & 0\\
		0 \ar[r] &\st_{(r,k)}^{\ana}(\alpha,\pi_0, {\bm\lambda}_\bh) \ar[r] & \sE_{\{i\}}^\emptyset(\alpha,\pi_0,{\bm\lambda}_\bh,\Psi)^0/\widetilde{U} \ar[r] & i_{\op^{\lrr}_{ir}}^{G}(\alpha,\pi_0,{\bm\lambda}_\bh)\ar[r]&0,}
\end{equation}
where $\widetilde{U}:=\sum_{\emptyset\neq I\subseteq \Delta_n(k)} \BI_{\op^{\lrr}_I}^G(\alpha,\pi_0, {\bm\lambda}_\bh)$.\;Recall that the representation
$\sE_{\{ir\}}^\emptyset(\alpha,\pi_0,{\bm\lambda}_\bh,\Psi)^0$ is defined in the argument below \cite[Theorem 5.19]{2022ext1hyq}.\;By \cite[Theorem 5.19]{2022ext1hyq} and \cite[Remark 5.21]{2022ext1hyq},\;the pull-back of the bottom exact sequence via the injection
\begin{equation}\label{injabc}
	w_{\op^{\lrr}_{ir}}^{\infty}(\alpha,\pi_0,{\bm\lambda}_\bh):=w_{\op^{\lrr}_{ir}}^{\infty}(\alpha,\pi_0,{\bm\lambda}_\bh)\otimes_E \unr(\alpha)\circ \dett\hooklongrightarrow i_{\op^{\lrr}_{ir}}^G(\alpha,\pi_0, {\bm\lambda}_\bh)\end{equation}
is split.\;Since $\homo_G(w_{\op^{\lrr}_I}^{\infty}(\alpha,{\bm\lambda}_\bh), \st_{(r,k)}^{\ana}(\alpha,{\bm\lambda}_\bh)/\widetilde{\Sigma}^{\lrr}(\alpha,\pi_0, {\bm\lambda}_\bh)')=0$,\;we deduce
\begin{equation*}\ext^1_G(w_{\op^{\lrr}_{ir}}^{\infty}(\alpha,\pi_0,{\bm\lambda}_\bh), \widetilde{\Sigma}^{\lrr}(\alpha,\pi_0, {\bm\lambda}_\bh)')\hooklongrightarrow \ext^1_G(w_{\op^{\lrr}_{ir}}^{\infty}(\alpha,\pi_0,{\bm\lambda}_\bh),\st_{(r,k)}^{\ana}(\alpha,\pi_0,{\bm\lambda}_\bh))
\end{equation*}
Therefore,\;the pull-back of the top exact sequence of (\ref{commrevi}) via (\ref{injabc}) is also split.\;This implies that
$I_{\op^{\lrr}}^G (\pi_0^{\lrr}(\alpha,{\bm\lambda}_\bh,\Psi))/U$ contains a subrepresentation $\widetilde{\Sigma}_i^{\lrr}(\alpha,\pi_0, {\bm\lambda}_\bh, \psi)'$,\;which is isomorphic to an extension of $v_{\op^{\lrr}_{ir}}^{\infty}(\alpha,{\bm\lambda}_\bh)$ by $\widetilde{\Sigma}^{\lrr}(\alpha,\pi_0, {\bm\lambda}_\bh)'$.\;We have an isomorphism
\begin{equation*}
	\ext^1_G(v_{\op^{\lrr}_I}^{\infty}(\alpha,\pi_0,{\bm\lambda}_\bh), \Sigma_i^{\lrr}(\alpha,\pi_0, {\bm\lambda}_\bh)') \xrightarrow{\sim}  \ext^1_G(v_{\op^{\lrr}_I}^{\infty}(\alpha,\pi_0.{\bm\lambda}_\bh), \widetilde{\Sigma}^{\lrr}(\alpha,\pi_0, {\bm\lambda}_\bh)'),
\end{equation*}
by similar strategy in the proof of \cite[Proposition 5.29]{2022ext1hyq} and \cite[Proposition 5.34(2)]{2022ext1hyq}.\;This asserts that
$\widetilde{\Sigma}_i^{\lrr}(\alpha,\pi_0,{\bm\lambda}_\bh, \psi)'$ comes from some $\Sigma_i^{\lrr}(\alpha,\pi_0,{\bm\lambda}_\bh, \psi)'$,\;which is
an extension of $v_{\op^{\lrr}_{ir}}^{\infty}(\alpha,\pi_0, {\bm\lambda}_\bh)$ by the representation $\Sigma_i^{\lrr}(\alpha,\pi_0,{\bm\lambda}_\bh)'$. But the push-forward of $\Sigma_i^{\lrr}(\alpha,\pi_0, {\bm\lambda}_\bh, \psi)'$ via the injection $\Sigma_i^{\lrr}(\alpha,\pi_0,{\bm\lambda}_\bh)'\hookrightarrow \Sigma_i^{\lrr}(\alpha,\pi_0,{\bm\lambda}_\bh)$ is isomorphic to $\Sigma_i^{\lrr}(\alpha,\pi_0,{\bm\lambda}_\bh, \psi)$.\;The claim follows.\;

The composition in (\ref{equ: lgL1}) induces then
\begin{equation*}
	\Sigma_i^{\lrr}(\alpha,\pi_0, {\bm\lambda}_\bh)^{\oplus r}\hooklongrightarrow \Sigma_i^{\lrr}(\alpha,\pi_0,{\bm\lambda}_\bh, \psi)^{\oplus r} \lra \Pi_{\infty}^{R_{\infty}-\ana}[\FI_t].
\end{equation*}
By using the same argument as in the \cite[Page.\;8040]{2015Ding},\;we see that the image of the second morphism is also contained in $ \pi_{\infty}^{R_{\infty}-\ana}[\fm_y] $.\;The surjectivity of (\ref{equ: lgln-rest}) now follows.\;We complete the proof of Theorem \ref{thm: lgln-main}.
\end{proof}

\acknowl{This is part of the author's PhD thesis.\;I would like to express my sincere gratitude to my advisor Yiwen Ding for introducing me to this subject,\;for many helpful discussion and suggestions,\;and for his comments on earlier draft of this paper.\;}


%-------------------------------------------------------------------------%
% GENERAL INSTRUCTIONS                                                    %
% Please keep formatting macros to a minimum.                             %
% Avoid redundant source code such as unused definitions or longer        %
% passages of comments (% or {comment}).                                  %
% Do not worry about bad page breaks, etc. and avoid adding extra space   %
% and using glue to improve the appearance of the manuscript.             %
% Please ensure that the material has been carefully proofread and that   %
% all files used are made available to the editors/publisher.             %
%-------------------------------------------------------------------------%

%-------------------------------------------------------------------------%
% PACKAGES                                                                %
% Missing packages can be downloaded at www.tug.org/ctan.html.            %
%-------------------------------------------------------------------------%

%-------------------------------------------------------------------------%
% MULTILINE EQUATIONS                                                     %
% For numbered/unnumbered multiline equations use align/align* instead of %
% eqnarray. (Other available environments for multiline displays are      %
% gather, multline, aligned, split, etc.)                                 %
%-------------------------------------------------------------------------%

%-------------------------------------------------------------------------%
% QED SYMBOL                                                              %
% The \qedhere command forces the QED symbol to go on the present line,   %
% which might be useful if the last part of a proof environment is, e.g., %
% a displayed equation or list environment.                               %
%-------------------------------------------------------------------------%

%-------------------------------------------------------------------------%
% SOME MORE MATH                                                          %
% -- Avoid blank lines before or after a display, unless you really want  %
%    to start a new paragraph.                                            %
% -- Use equation*, gather* or the bracket pair \[ and \] instead of      %
%    double dollar signs $$.                                              %
% -- Use \operatorname{...} for math operators that are not predefined    %
%    or define new operators with \DeclareMathOperator. lim-like          %
%    operators can be defined with \DeclareOperator*                      %
% -- Use \quad for spacings in displayed formulas.                        %
% -- Use \substack{... \\ ...} for multiline subscripts on, e.g., sums.   %
%-------------------------------------------------------------------------%

%-------------------------------------------------------------------------%
% ENUMERATIONS                                                            %
% By default, the labels of first level enumerations are (i), (ii), etc.  %
% You may change them to become, e.g., A., B. etc., using an optional     %
% argument: \begin{enumerate}[A.]                                         %
%-------------------------------------------------------------------------%

%-------------------------------------------------------------------------%
% FIGURES AND TABLES                                                      %
% You may include a figure writing                                        %
% \begin{figure}[t]                                                       % 
%   \includegraphics[width=.8\textwidth]{xxx.yyy}                         %
%   \caption{Caption.}\label{fig:zzz}                                     %
% \end{figure}                                                            %
% A table is included similar using the table environment.                %
% Sub-captions are created as follows:                                    %
% \begin{figure}[t]                                                       %
%   \subfloat[Caption a]{\includegraphics[width=.45\textwidth]{...}}      %
%   \hfill                                                                %
%   \subfloat[Caption b]{\includegraphics[width=.45\textwidth]{...}}      %
%   \caption{Caption.}                                                    %
% \end{figure}                                                            %
%-------------------------------------------------------------------------%

%-------------------------------------------------------------------------%
% QUOTATION MARKS                                                         %
% Double quotation marks are produced using `` and ''. The opening quote  %
% then looks like a superscript 66 and the closing quote like a           %
% superscript 99.                                                         %
%-------------------------------------------------------------------------%

%-------------------------------------------------------------------------%
% HYPHEN VS. DASH                                                         %
% The hyphen (-) is used for compound words like $n$-dimensional.         %
% (Don't write things like $n-$dimensional!)                              %
% The en-dash (--) is used for number ranges and it can stand for `and',  %
% e.g., between two names: ``Various generalizations of the               %
% Cauchy--Schwarz inequality are presented on pages 63--67.''             %
% To partition a sentence one may use the em-dash (---) with no space     %
% on either side; we prefer the en-dash with a blank on both sides.       %
% The mathematical minus sign is produced by a single - within            %
% mathematics mode.                                                       %
%-------------------------------------------------------------------------%

%-------------------------------------------------------------------------%
% REFERENCES                                                              %
% Use \label and \ref for all cross-references to equations, figures,     %
% tables, sections, subsections, etc.                                     %
%-------------------------------------------------------------------------%

%-------------------------------------------------------------------------%
% BIBLIOGRAPHY                                                            %
% References should be collected at the end of the paper and numbered in  %
% alphabetical order of the authors' names. Titles of journals should be  %
% abbreviated as in Mathematical Reviews. Please do not use small caps    %
% for the names of the authors. The preferred style is shown in the       %
% examples below.                                                         %
% Refer to the bibliography entries using \cite commands.                 %
% Using BibTeX, you simply have to add the names of your databases as     %
%   \bibliography{mydatabase}                                             %
% The file crelle.bst then ensures the correct style.                     %
%-------------------------------------------------------------------------%

\begin{thebibliography}{9}
\bibitem{AST2009324R10}
\textit{J.~Bella\"{i}che} and \textit{G.~Chenevier},
 Families of Galois representations and Selmer groups (English, French summary), Ast\'{e}risque No. \textbf{324} (2009), xii+314 pp. 


\bibitem{bergdall_2017}
\textit{J.~Bergdall},
Paraboline variation over $p$-adic families of $(\phi,\Gamma)$-modules (English summary),
Compos. Math. \textbf{153} (2017), no. 1, 132-174.


\bibitem{bergdall2018adjunction}
\textit{J.~Bergdall} and \textit{P.~Chojecki},
An adjunction formula for the Emerton-Jacquet functor (English summary),
Israel J. Math. \textbf{223} (2018), no. 1, 1-52.


\bibitem{berger2008equations}
\textit{L.~Berger},
\'{E}quations diff\'{e}rentielles $p$-adiques et $(\varphi,N)$-modules filtr\'{e}s (French. English, French summary),
Ast\'{e}risque No. \textbf{319} (2008), 13-38. 

\bibitem{bernstein1984centre}
\textit{J.~N.~Bernstein}   and \textit{P.~Deligne},
{Le ``centre'' de Bernstein (French, Edited by P. Deligne) [The Bernstein ``center'']}, Travaux en Cours, Representations of reductive groups over a local field, 1-32, Hermann, Paris, 1984.


\bibitem{breuil2019ext1}
\textit{C.~Breuil},
{$\ext^1$} localement analytique et compatibilit\'{e} local-global, Amer. J. Math. \textbf{141} (2019), no. 3, 611-703.


\bibitem{Breuil03a}
\textit{C.~Breuil},
Sur quelques repr\'{e}sentations modulaires et {$p$}-adiques de {${\GLN}_2(\bQ_p)$}. {I}, Compositio Math. \textbf{138} (2003), no. 2, 165-188.


\bibitem{Breuil03b}
\textit{C.~Breuil},
Sur quelques repr\'{e}sentations modulaires et {$p$}-adiques de
	{${\GLN}_2(\bQ_p)$}. {II},
J. Inst. Math. Jussieu \textbf{2} (2003), no. 1, 23-58.


\bibitem{BCLINV2}
\textit{C.~Breuil}, 
S\'{e}rie sp\'{e}ciale p-adique et cohomologie \'{e}tale compl\'{e}t\'{e}e (French. English, French summary), Ast\'{e}risque  No. \textbf{331} (2010), 65-115. 


\bibitem{breuil2016socle}
\textit{C.~Breuil}, 
Socle localement analytique I. (French. English, French summary) [Locally analytic socle I], Ann. Inst. Fourier (Grenoble) \textbf{66} (2016), no. 2, 633-685.



\bibitem{breuil2015II}
\textit{C.~Breuil}, 
Vers le socle localement analytique pour GLn II (French. French summary) [Towards the locally analytic socle for GLn II],
Math. Ann. \textbf{361} (2015), no. 3-4, 741-785.


\bibitem{Ding2021}
\textit{C.~Breuil} and  \textit{Y.~Ding}, Bernstein eigenvarieties, 
preprint 2021, \url{http://arxiv.org/abs/2109.06696}.


\bibitem{HigherLinvariantsGL3Qp}
\textit{C.~Breuil} and  \textit{Y.~Ding}, 
Higher $L$-invariants for $\GLN_3(\BQ_p)$ and local-global
	compatibility (English summary),
Camb. J. Math. \textbf{8} (2020), no. 4, 775-951.


\bibitem{breuil2020probleme}
\textit{C.~Breuil} and \textit{Y.~Ding},
SUR UN PROBL{\`E}ME DE COMPATIBILIT{\'E} LOCAL-GLOBAL LOCALEMENT ANALYTIQUE, 
 Memoirs of the American Mathematical Society, 2020.

\bibitem{breuil2017smoothness}
\textit{C.~Breuil},\;\textit{E.~Hellmann} and \textit{B.~Schraen},
Smoothness and classicality on eigenvarieties (English summary),
Invent. Math. \textbf{209} (2017), no. 1, 197-274.


\bibitem{breuil2017interpretation}
\textit{C.~Breuil}, \textit{E.~ellmann} and \textit{B.~Schraen},
Une interpr{\'e}tation modulaire de la vari{\'e}t{\'e} trianguline.
Math. Ann. \textbf{367} (2017), no. 3-4, 1587-1645.


\bibitem{breuil2019local}
\textit{C.~Breuil}, \textit{E.~Hellmann} and \textit{B.~Schraen},
 A local model for the trianguline variety and applications (English summary), 
 Publ. Math. Inst. Hautes \'{E}tudes Sci. \textbf{130} (2019), 299-412.


\bibitem{breuil2020towards}
\textit{C.~Breuil} and \textit{F.~Herzig},
Towards the Finite Slope Part for $\GLN_n$ (English summary),
Int. Math. Res. Not. IMRN 2020, no. 24, 10495-10552.


\bibitem{breuil2007first}
\textit{C.~Breuil}  and \textit{P.~Schneider},
First steps towards $p$-adic langlands functoriality (English summary), 
J. Reine Angew. Math. \textbf{610} (2007), 149-180.



\bibitem{bushnell2016admissible}
\textit{C.~J.~Bushnell} and  \textit{P.~C.~Kutzko},
The Admissible Dual of $\GLN_n$ via Compact Open Subgroups, 
Annals of Mathematics Studies, \textbf{129.} Princeton University Press, Princeton, NJ, 1993. xii+313 pp. 

\bibitem{PATCHING2016}
\textit{A.~Caraiani}, \textit{M.~Emerton}, \textit{T.~Gee}, \textit{D.~Geraghty}, \textit{V.~Pa\v{s}k\={u}nas} and \textit{S.~W.~Shin},
Patching and the p-adic local Langlands correspondence (English summary),\;
Camb. J. Math. \textbf{4} (2016), no. 2, 197-287.


\bibitem{chenevier2011infinite}
\textit{G.~Chenevier}.
On the infinite fern of Galois representations of unitary type (English, French summary),
Ann. Sci. \'{E}c. Norm. Sup\'{e}r. \textbf{(4) 44} (2011), no. 6, 963-1019.


\bibitem{clozel2008automorphy}
\textit{L.~Clozel}, \textit{M.~Harris} and  \textit{R.~Taylor},
Automorphy for some $l$-adic lifts of automorphic mod l Galois representations (English summary, with Appendix $A$, summarizing unpublished work of Russ Mann, and Appendix $B$ by Marie-France Vign{\'e}ras),
Publ. Math. Inst. Hautes \'{E}tudes Sci. No. \textbf{108} (2008), 1-181.



\bibitem{PADINGL2COL}
\textit{P.~Colmez},
Repr\'{e}sentations de {${\GLN}_2(\bQ_p)$} et {$(\phi,\Gamma)$}-modules,
Ast\'{e}risque No. \textbf{330} (2010), 281-509. 


\bibitem{2015Ding}
\textit{Y.~Ding},
$\mathcal{L}$-invariants, partially de Rham families, and local-global compatibility (English, French summary),
 Ann. Inst. Fourier (Grenoble) \textbf{67} (2017), no. 4, 1457-1519.


\bibitem{2019DINGSimple}
\textit{Y.~Ding}, 
Simple $\sL$-invariants for $\GL_{n}$ (English summary),
Trans. Amer. Math. Soc. \textbf{372} (2019), no. 11, 7993-8042.


\bibitem{emerton2006jacquet}
\textit{M.~Emerton}.
Jacquet modules of locally analytic representations of p-adic reductive groups. I. Construction and first properties (English, French summary), Ann. Sci. \'{E}cole Norm. Sup. \textbf{(4) 39} (2006), no. 5, 775-839.

\bibitem{Em1}
\textit{M.~Emerton}, 
On the interpolation of systems of eigenvalues attached to automorphic Hecke eigenforms, Invent. Math. \textbf{164} (2006), no. 1, 1-84.


\bibitem{emerton2007jacquet}
\textit{M.~Emerton}.
Jacquet modules of locally analytic representations of $p$-adic reductive groups II. The relation to parabolic induction, J. Institut Math. Jussieu, 2007.

\bibitem{Emerton2007summary}
\textit{M.~Emerton}, 
Locally analytic representation theory of $p$-adic reductive groups: a summary of some recent developments, $L$-functions and Galois representations, 407-437,
London Math. Soc. Lecture Note Ser., \textbf{320}, Cambridge Univ. Press, Cambridge, 2007.



\bibitem{emerton2014geometric}
\textit{M.~Emerton} and \textit{T.~Gee}.
A geometric perspective on the breuil--m{\'e}zard conjecture (English summary), 
J. Inst. Math. Jussieu \textbf{13} (2014), no. 1, 183-223.


\bibitem{localglobalgl2}
\textit{M.~Emerton}, 
Local-global compatibility in the p-adic Langlands programme for $\mathrm{GL}_2/\mathbb{Q}$,
2011, available electronically at \url{"http://math.uchicago.edu/~emerton/preprints.html"} 


\bibitem{emerton2017locally}
\textit{M.~Emerton}, 
 Locally analytic vectors in representations of locally $p$-adic analytic groups,
Mem. Amer. Math. Soc. \textbf{248} (2017), no. 1175, iv+158 pp.  


\bibitem{2022ext1hyq}
\textit{Y.~He},
Extensions of locally analytic generalized parabolic Steinberg representations,
preprint 2022, \url{http://arxiv.org/abs/2211.00476}.


\bibitem{humphreysBGG}
\textit{J.~E.~Humphreys}, 
Representations of Semisimple Lie Algebras in the BGG Category $\cO$, volume~94.
Graduate Studies in Mathematics, 94. American Mathematical Society, Providence, RI, 2008. xvi+289 pp. 

\bibitem{kisin2008potentially}
\textit{M.~Kisin},
Potentially semi-stable deformation rings, J. Amer. Math. Soc. \textbf{21} (2008), no. 2, 513-546.


\bibitem{XLROBBA}
\textit{K.~S.~Kedlaya}, \textit{J.~Pottharst} and  \textit{L.~Xiao},
Cohomology of arithmetic families of {$(\varphi,\Gamma)$}-modules, J. Amer. Math. Soc. \textbf{27} (2014), no. 4, 1043-1115.


\bibitem{liu2007cohomology}
\textit{R.~Liu},
Cohomology and duality for ($\varphi$, $\Gamma$)-modules over the
robba ring, Int. Math. Res. Not. IMRN 2008, no. 3, Art. ID rnm150, 32 pp. 11S25

\bibitem{nakamura2009classification}
\textit{K.~Nakamura},
Classification of two-dimensional split trianguline representations
of $p$-adic fields (English summary), Compos. Math. \textbf{145} (2009), no. 4, 865-914.
11S25 (11F80 11F85)

\bibitem{orlik2015jordan}
\textit{S.~Orlik},
On Jordan-H\"{o}lder series of some locally analytic representations,
J. Amer. Math. Soc. \textbf{28} (2015), no. 1, 99-157.


\bibitem{orlik2014jordan}
\textit{S.~Orlik}, and \textit{B.~Schraen},
The {J}ordan-h{\"o}lder series of the locally analytic steinberg
representation. (English summary), Doc. Math. \textbf{19} (2014), 647-671.


\bibitem{schraen2011GL3}
\textit{B.~Schraen},
Repr\'{e}sentations localement analytiques de {${\GLN}_3(\BQ_p)$}, Ann. Sci. \'{E}c. Norm. Sup\'{e}r. \textbf{(4) 44} (2011), no. 1, 43-145.



\bibitem{Dilogarithm}
\textit{Z.~Qian}, 
Dilogarithm and higher {$\mathcal{L}$}-invariants for {${\GLN}_3({\BQ}_p)$}, Represent. Theory \textbf{25} (2021), 344-411.

\bibitem{schneider2002banach}
\textit{P.~Schneider} and \textit{J.~Teitelbaum},  
Banach space representations and Iwasawa theory (English summary), Israel J. Math. \textbf{127} (2002), 359-380.


\bibitem{scholze2013local}
\textit{P.~Scholze},
The local Langlands correspondence for $\GLN_{n}$ over $p$-adic fields, and the cohomology of compact unitary Shimura varieties. (English summary) Shimura varieties, 251-265, London Math. Soc. Lecture Note Ser., \textbf{457}, Cambridge Univ. Press, Cambridge, [2020], 2020. 

\bibitem{av1980induced2}
\textit{A.~V.~Zelevinsky},
Induced representations of reductive $p$-adic groups. II. On irreducible representations of $\GLN(n)$,  Ann. Sci. \'{E}cole Norm. Sup. \textbf{(4) 13} (1980), no. 2, 165-210




\end{thebibliography}

\end{document}
\bibitem{article}
\textit{T.~Angel} and \textit{I.~E.~Shparlinski},
Article in a journal, J. reine angew. Math. \textbf{647} (2010), 123--156.

\bibitem{preprint}
\textit{C.~Bonnaf\'e},
Preprint, preprint 2008, \url{http://arxiv.org/abs/0805.4100}.

\bibitem{proceedings}
\textit{P. Brooksbank},
Article in a conference proceedings,
in: Finite geometries (Pingree Park 2004), Oxford University Press, Oxford (2006), 1--16.

\bibitem{phdthesis}
\textit{J.~D.~King},
Unpublished dissertation,
Ph.D. thesis, University of Cambridge, 1995.

\bibitem{book}
\textit{F. Sukochev}, \textit{S. Lord} and \textit{D. Zanin},
Book, 2nd ed., De Gruyter, Berlin 2012.